% Paper source Paper7C-ContinuousQFT-v2_207.tex: v3.040 S3 literature mathematical/conditional comparison; physical EM packet unresolved.
% Framework v3.025: physical action-owner separation and conditional source equivalence.
% Current integrated paper source: Paper7C-ContinuousQFT-v2_204.tex (v3_017).
% Prior source: Paper7C-ContinuousQFT-v2_203.tex. New cross-DAG mathematics is scoped, not physical-occurrence certified.
% Paper-series v2.382: audit-identified standalone repair integration; mathematical scopes and open physical owners preserved.
% Paper-series v2.381: complete algebraic exceptional torsion atlas and all-cocycle twist-stability; original claims retained.
% Paper-series v2.380: unified bibliography, exact exceptional double/length/gauge mathematics in I-B; previous theorem labels retained.
% Paper-series v2.379: full equivariant Dirac/Hodge/eta/torsion atlas in Paper I-B; paper-root references synchronized.
% Paper-series v2.378: 2026 exceptional literature consolidation and reference synchronization. The mathematical additions are located in Paper I-B.
% Paper-series v2.377: reference-only rebind to current Paper I-B v1.1; original mathematics retained except new Paper I-B comparison chapter.
% Framework v2.370: derived exceptional-root/Hurwitz comparison in I-A/I-C; other papers reference-only synchronized.
% Framework v2.368: preserve prior II-B convention/unital fixes and verified II-D corrections; current references synchronized.
% Framework v2.358: current 30-root source reference synchronization; selectively updated III-E, III-F, VI-B, VI-C.
% Framework v2.356: new source-owned conditional results and source-registry synchronization.
% Framework v2.350: repaired Paper VII mathematics; exact finite identities separated from analytic/physical certificates.
% Framework v2.323: Petersen finite strong-source coherence, exact multi-core mirror spectra, QCD action-incidence and continuum firewalls. Existing source theorems preserved.
% Framework v2.275: exact positive dense-twin word cubature, finite repeatability and eighth-moment classification, nested sampling frames, joint action and POVM-jet exports; preserves v2.274.
% Framework v2.274: two-scalar positive convolution certificates, first-leaked exceptional coset detection, branch-symmetric twin mixing, and finite endpoint moments; all v2.273 branches preserved.
% Framework v2.273: restored missing Problem 3 theorems from v2.268; exact exceptional frames, twin geometry, convolution repeatability and projector/occurrence firewalls; merged on v2.272 without overwriting independent branches.
% Framework v2.272: merged v2.270 covariant native-readout mathematics with v2.271 Problem 4 finite-action Hopf hypercharge/holonomy corrections; no universal bulk-gauge uniqueness claim.
% Independent branches: v2.270 readout/guard exports and v2.270 K16 action-consumed charge closure; all prior scoped assumptions retained.
% Framework v2.267: exact Paper II-A converse/observable certificates, III-F action tests, V-B right orbit and V-C native gauge scope; prior v2.266 preserved.
% Framework v2.266: exact Paper II-C outcome-count/readout converse results and scoped downstream imports; 30-root references synchronized.
% Framework v2.265: II-C coordinate-stage constitution and exact first-symbol host inheritance; prior mathematical domains retained.
% Framework v2.264: II-B reverse and action-bridge integration; preserves prior mathematical domains; active 30-root references synchronized.
% Framework v2.263: Paper I-D adds bridge-controlled root-carry parity and source-relative old-cap cancellation; other root mathematical bodies retained; active 30-root references synchronized.
% Framework v2.262: merges Paper II-F three-gap closure from v2.260 into v2.261, preserving newer II-B/II-D repairs; 30-root references synchronized.
% Framework v2.261: merged Paper II-B GP2-level/geometric repairs; preserves v2.260 II-D repairs; active references synchronized.
% Framework v2.260: Paper II-D signed RS, finite-symbol typing, and conditional curvature/Jordan repairs; references synchronized.
% Framework v2.259: conditional native SU(2) atlas and explicit later discharge; v2.258 mathematical bodies retained elsewhere.
% Framework v2.258: preserves parallel II-A/dynamics update and adds complete readout/frequency closures.
% Framework v2.257: exact minimal complete position readout and arbitrary-preparation frequency law.
% Framework v2.256: exact native inertia and sharp bounded ordered-source closure; preserves concurrent v2.255 work.
% Framework v2.255: preserves parallel detector integration and adds exact finite readout closures.
% Framework v2.254: shared bibliography/reference sync; substantive Pasterski detector-Ward additions only IV-B/IV-C/VII-C.
% Framework v2.253: typed raw coordinate defect in I-A and synchronized VI-A import; active references synchronized.
% Framework v2.252: optional shell-frame source-return calculation in III-G; active references synchronized.
% Framework v2.251: forward common-action execution, joint source jets and scoped coupled control.
% Framework v2.251: exact full-heat variations, fermion source jets and explicit physical completion frontier.
% Framework v2.251: certified interacting source jets, pole stability and scoped completion proofs.
% Framework v2.251: interacting finite-source evaluation and coupled local vacuum classification.
% Framework v2.251: exact finite source execution and native causal common-parent bridge.
% Framework v2.251: constructive parent and scoped source-incidence audit; see release report.
% Framework v2.251: common-action two-jet, corrected source/response closure and synchronized references.
% Framework v2.244: scoped rooted source-order Schur diagnostics and synchronized paper references.
% Framework v2.242: signed rooted phase lift, complete transverse return and exact family cubic.
% Framework v2.241: proved rooted conormal geometry, conditional cap realization and scoped action tests.
% Framework v2.239: coherent 30-paper reference rebinding; selective new proofs on established owners.
% Framework v2.236: source registry synchronized; existing mathematical body retained except declared insertions.
\documentclass[12pt,a4paper]{article}
\usepackage{caption}
\usepackage{microtype}
\usepackage{makecell}
\PassOptionsToPackage{hidelinks}{hyperref}

% Paper-series v2.381: new I-B torsion/modular results, bibliography and current source roots.
% Paper-series v2.380: new exceptional bibliography and synchronized paper references.
% Paper-series v2.379 current 31-paper macro/aux registry.
% Paper-series v2.378: current 31-paper AUX/bibliography registry updated for Paper I-B v1.2 literature additions.
% Framework v2.371: current-source reference registry refreshed; old theorem labels retained.
% Framework v2.370: current 30-paper source/AUX root rebinding; native mathematical macros unchanged.
% Framework v2.368: synchronized merged 30-root registry; mathematical notation unchanged.
% Framework v2.358: active 30-paper source registry synchronized; existing mathematical macros unchanged.
% Framework v2.357: selective new source-owned certificates; active reference roots updated for five revised papers.
% Framework v2.356: active 30-root registry synchronized to the 30 papers in this release.
% Framework v2.275: exact positive dense-twin word cubature, finite repeatability and eighth-moment classification, nested sampling frames, joint action and POVM-jet exports; preserves v2.274.
% Framework v2.274: filename and root-source synchronization.
% Framework v2.273: restored missing Problem 3 theorems from v2.268; exact exceptional frames, twin geometry, convolution repeatability and projector/occurrence firewalls; merged on v2.272 without overwriting independent branches.
% Framework v2.272: merged v2.270 covariant native-readout mathematics with v2.271 Problem 4 finite-action Hopf hypercharge/holonomy corrections; no universal bulk-gauge uniqueness claim.
% Independent branches: v2.270 readout/guard exports and v2.270 K16 action-consumed charge closure; all prior scoped assumptions retained.
% Framework v2.269: synchronized 30 active AUX roots; mathematics and bibliography preserved from Macros-v2_161.
% Framework v2.266: Paper II-C converse-derived outcome resources; II-E, III-C, IV-A/B, V-A and VI-B downstream mathematical closures; 30-root registry synchronized.
% Framework v2.265: II-C coordinate-stage constitution and exact first-symbol host inheritance; prior mathematical domains retained.
% Framework v2.264: paper II-B reverse/audit and IV-A observable-minimal exchange; synchronized external references to 30 active roots.
% Framework v2.263: Paper I-C adds bridge-controlled root-carry parity and source-relative old-cap cancellation; other root mathematical bodies retained; active 30-root references synchronized.
% Framework v2.262: integrates Paper II-F radius, Casimir and readout results; active 30-root references synchronized.
% Framework v2.261: merged Paper II-B GP2-level/geometric repairs; preserves v2.260 II-D repairs; active references synchronized.
% Framework v2.260: Paper II-D signed RS, finite-symbol typing, and conditional curvature/Jordan repairs; references synchronized.
% Framework v2.259: conditional native SU(2) atlas and explicit later discharge; v2.258 mathematical bodies retained elsewhere.
% Framework v2.258: preserves parallel II-A/dynamics update and adds complete readout/frequency closures.
% Framework v2.257: exact minimal complete position readout and arbitrary-preparation frequency law.
% Framework v2.256: exact native inertia and sharp bounded ordered-source closure; preserves concurrent v2.255 work.
% Framework v2.255: preserves parallel detector integration and adds exact finite readout closures.
% Framework v2.254: Pasterski bibliography and synchronized 30-paper external-reference registry.
% Framework v2.253: typed raw coordinate defect in I-A and synchronized VI-A import; active references synchronized.
% Framework v2.252: synchronized source registry; optional III-G shell-frame calculation.
% Framework v2.251: forward common-action execution, joint source jets and scoped coupled control.
% Framework v2.251: exact full-heat variations, fermion source jets and explicit physical completion frontier.
% Framework v2.251: certified interacting source jets, pole stability and scoped completion proofs.
% Framework v2.251: synchronized active source registry; mathematical notation unchanged.
% Framework v2.251: synchronized active source registry; mathematical notation unchanged.
% Framework v2.251: synchronized active source registry; mathematical notation unchanged.
% Framework v2.251: synchronized source registry; no mathematical notation changes.
% Framework v2.244: synchronized 30-paper AUX registry and rooted phase-source action-order tests.
% Framework v2.242: signed rooted phase lift, complete transverse return and exact family cubic.
% Framework v2.241: proved rooted conormal geometry, conditional cap realization and scoped action tests.
% Framework v2.239: all 30 active source roots synchronized; no mathematical macro changes.
% Framework v2.238: Paper IV A/B/C split and 30-root reference registry.
% Framework v2.237: Paper VII-D technical split and 28-root reference registry.
% Framework v2.236: Paper VI-C technical-volume reference added; all current root bindings advanced.
% Framework v2.235: 26-root AUX and bibliography bindings synchronized; mathematical macros unchanged.
% BEGIN INLINE SOURCE: Paper00-TeXBuildCompatibility-v1_0.tex
% Rendering-only repairs for inherited long appendices and bitmap text fonts.
% Integer counters through 26 retain the original labels; later values use AA,
% AB, ... . This file changes no action coefficient, state or theorem premise.
\RequirePackage{alphalph}
\makeatletter
\let\BHPrimitiveAlph\@Alph
\newalphalph{\BHExtendedAlph}{\BHPrimitiveAlph}{26}
\let\@Alph\BHExtendedAlph
\AtBeginDocument{%
  \@ifpackageloaded{microtype}{\microtypesetup{expansion=false}}{}%
}
\makeatother

% END INLINE SOURCE: Paper00-TeXBuildCompatibility-v1_0.tex

% Framework v2.164: current reference/bibliography profile; mathematical definitions retained.
% Framework v2.162: finite metric, functional Hilbert and scoped Ward integration.
% Framework v2.161: inherits Macros-v2_57.tex exactly except the current external-reference binding.
% Macros v2.57 / framework v2.156: current reference registry only; mathematical definitions unchanged.
% Framework v2.143 / Macros v2.55: current external-reference bindings; mathematical notation unchanged.
% Framework v2.110 / Macros v2.54: current external-reference registry and bibliography; mathematical notation inherited.
% Framework v2.67 / Macros-v2_53: adds Born-action-cell and joint-EH readout symbols; absolute unit realization is metrological; MATCH_{J/scal}^{EH} remains the sole G24 normalization bridge.
% Framework v2.35 / Macros-v2_52: grouped paper numbering, active AUX registry, and Born-v6_20 bibliography binding synchronized; notation unchanged.
% Framework v2.26: active-source AUX and bibliography bindings synchronized; notation unchanged.
% Framework v2.09: reviewed early-reciprocity integration; scoped physical imports.
% Framework v2.08: Born--Euler mathematics and conditional complete-action potential criterion.
% Framework v2.07: active AUX registry includes Paper I-B v1.2; shared bibliography Born-v6_18.
% Macros-v2_46 / framework v2.04: mathematical atlas and downstream realization correction.
% Macros-v2_46 / framework v2.03: scoped structure integration and current dependency bindings.
% Macros-v2_46 / framework v2.01: shared source bindings; prior revision comments are historical.
% v2.42 / framework v1.99: Synchronized RG/AS owner and bibliography bindings.
% v2.41 / framework v1.99: active manuscript AUX bindings and geometric-closure bibliography synchronized.
% Release-v1.80 centralized source synchronization for scalar-circle/comparison-band integration.
% Release v1.79: Hurwitz hierarchy and kinetic-normalized action response integration; physical mass matching remains conditional.
% Release v1.78: approved M8 closure-batch integration; physical and analytic domains retained.
% Current v1.78 registry: all active TeX roots are synchronized below.
% Earlier release comments are historical scope records, not current status.
% Release-v1.77 synchronization: native same-owner Noether closure, generator--flow/history typing, analytic-certification firewalls, scalaron/K16, LF/QCD, quantum-Ward, and RAR boundary-owner transport are registered; shared mathematical notation is unchanged.
% Release-v1.76 synchronization: registry points to the v1.76 active paper roots.
% Release-v1.76 synchronization: registers the local A+Phi+g Noether/constraint closure and dynamics-frontier propagation; shared mathematical notation is unchanged.
% Release-v1.75 synchronization: registry points to the v1.75 active paper roots.
% Release-v1.75 synchronization: integrates the joint-Dirac Einstein bulk real form and the local odd-scalar heat-coefficient action on explicit analytic first-jet/interior domains; v1.74 Yang--Mills closure is retained. Fermionic reality, pure-metric curvature-squared/exact-nonlocal terms, noncanonical boundaries, physical vacuum selection, and actual solution-history tokens remain open.
% Release-v1.74 synchronization: registry points to the v1.74 active paper roots and retains the v1.73 theorem labels.
% Release-v1.74 synchronization: canonical Macros-v2_34.tex; generated local nonlinear Yang--Mills action, constrained connection-field phase space, source-free Gauss momentum map, and Legendre closure are integrated on their proved domains. The remaining non-gauge ACT_cl sectors and actual solution-history selection remain open.
% Release-v1.71 synchronization: closes the former JOINT_x/JOINT_5
% connection-occurrence gap by canonical structure-group extension of the
% nonflat Born connection, faithful associated representation, and the local
% cotangent--Clifford operator constructor.  Actual fluctuation members,
% actions, vacua, solutions, calibrations, and analytic QFT domains remain
% separately typed.  The active-AUX registry is advanced once.
% Release-v1.71 canonical macro and active-AUX registry.
% ======================================================================
% Macros-v2_39.tex -- canonical publication macro file
% Release v1.71 retains the v1.70 null-ray/QCD boundaries and additionally
% removes the duplicate structural connection actuality premise.  It does not
% claim global bundle topology, select a gauge/Higgs member, or promote the
% finite construction to an interacting continuum theory.
% Release v1.68 constructs the finite light-front total-fraction operators and
% parton number/momentum operator measures, proves their simplex and Hausdorff-limit
% theorems, and separates positive cut measures from scheme-renormalized
% light-ray distributions.  It also separates finite real Schur poles, stable
% norm-resolvent/Riesz pole-residue limits, and continued resonance
% pole/residue limits.  The Paper-VII-C scattering package is claim-indexed:
% local QFT and spectral claims import no S-matrix or infrared premise, while
% asymptotic amplitudes, LSZ, S-matrix poles, completeness, and cross sections
% import only their named massive or massless/long-range clauses.
% Release v1.67 factorizes BND48_H into a same-action Hessian/Schur
% descendant and still-open typed boundary-trace, critical-lift, and
% refinement-limit conditions.  It proves the finite-collar Dirichlet-to-
% Neumann law sqrt(y)tanh(L sqrt(y)), its nonuniform critical limit, and the
% finite-return ultra-infrared turnover.  The Paper-I-A Hardy constant mode is
% an exact geometric near-hit, not yet the physical metric boundary mode.
% The current local f(R) Hessian remains massive and therefore still refutes
% a RAR prediction on the minimal action branch.
% Release v1.66 derives the transverse depth-six normalizer bridge from the
% already generated C6 Reynolds image and Paper-I-A binary-dihedral normalizer.
% The bridge is the unique least 2D3-stable scalar-containing realization of
% the faithful quotient 2D3/C6=C2, has trace defect six and hence K=48, but is
% transverse to the nested active carrier.  Exact K=44 and K=52 contraction,
% semigroup, and energy-midpoint identities are registered as structural
% readouts rather than projector events.  The active E48 carrier gap remains
% unchanged.  A same-action boundary-symbol/Calderon intertwiner is isolated
% as the necessary and still unproved owner bridge to the static DtN sort;
% representation equivalence alone neither derives RAR nor creates a K=48
% action junction, pole, occupation, or cosmological transition.
% Release v1.65 derives the native E_48-centered carrier mass gap from the
% GP1--GP3-generated active spectrum, distinguishes its carrier-free interval
% from a particle located at E_48, and withdraws the unsupported promotion to
% a positive continuum Yang--Mills mass gap.  The latter remains a valid
% conditional transfer theorem on an enlarged continuum branch, but it is not
% an active framework prediction or an unfinished native-gap proof.
% Release v1.64 synchronizes Derivation Batches 14--22.  It registers the
% GP1--GP3 heat-profile orbit, same-Hessian electroweak response and pole
% ownership, K16/K48 action-response corrections, perturbation state and
% non-Gaussianity typing, finite-vacuum/gap/QCD limit repairs, factorized
% position-area ownership, the actual static metric response, and the ordered
% Wilson/string-tension limits.  Strict continuum convergence, positive
% physical gap, surviving higher cumulants, and positive tension remain open.
% Release v1.61 closes TIMEWIDTH16 as an independent direct-production datum.
% JUNC_{2|7} continues to own the finite intrinsic microscopic profile/width,
% while the actual same-action background/production history determines its
% physical-time pullback, support, and duration.  The current Paper-I-A--VII-C AUX
% registry is advanced once; shared mathematical notation is unchanged.
% Release v1.60 derives the final-octet primitive-remnant depletion, canonical
% Hopf-current representative, exact Hopf descent to dOmega/(4pi), and
% closed-FLRW a^{-3} neutral-current shape.  It retypes K=48 as the onset
% of the final one-unit geometric transfer and records cDM/RAR as distinct
% physical projections of that common neutral parent while preserving
% occurrence, Hilbert-owner, stress-matching, and static-kernel premises.
% Release v1.59 transports the v1.58 exceptional-shell/K16 closure through
% Papers III-F and VI-B--VII-C.  It replaces Paper VI-B's obsolete mixed-Hessian
% output typing by the third-variation pair vertex and Bogoliubov occupation
% support, derives the post-seed Higgs-pair feedback vertex and its exact
% alpha--delta kernel factors, propagates the unique K16/K48 action-junction
% split, identifies the QCD delta count with the upstream 2O shell atom only
% as a common-scalar numerical compatibility, and narrows THERM16 to dynamical
% approach/beta selection because the finite Gibbs/KMS target is already
% unique on a supplied simple observable algebra once its Hamiltonian and
% inverse temperature are supplied; direct-sum central weights remain typed.
% Release v1.58 registers the exceptional nearest-shell closure, the derived
% transverse K=16 action event, the unique interior action junction, the
% type-correct ordered pair vertex, the thin-layer squeeze law, and the
% radiation-like/matter-like kinematic split.  REG_0, transfer magnitude,
% physical layer width, and thermalization remain explicitly conditional.
% Release v1.57 registers the active-depth-six obstruction, the native
% K=48 carrier desert, the typed 7|9 owner reduction, and the dynamical
% reheating handoff across Papers III-B--VI-B.
% Release v1.56 registers the least-support n_2 inflation dressing of the
% persistent n_7 Laurent response, the analytic global-convexity and positive
% endpoint-measure proofs, the derived inflation core and K=16 coefficient
% freeze, and the no-particle dark-sector split into static RAR plus the
% conditional homogeneous effective-curvature response at reflected K=48.
% There is no n_6 action, scalaron-DM mass, or second Einstein--Hilbert term.
% Release v1.55 registers the protected-octet resolved/residual reflection
% q_sp+Z_W=8, its j<->8-j and K<->64-K scale covariance, and the derived
% rank-one-exit quadratic grade C_1^2/U=C_2 in Papers III-C, III-E, and VIII.
% The reflection is a typed trace/scale theorem, not a time reversal or a
% physical event-occurrence theorem; reheating transfer remains conditional.
% Release v1.54 registers the exact positive-part Laurent indices, the
% active-domain/terminal-endpoint action split, the no-second-EH firewall, and
% the q_DM=6/K_DM=48 scalaron readout typing across Papers III-D/E/F, VI,
% VIII, and IX.
% Release v1.53 registers the common nonnegative Laurent obstacle projection
% for n_1^+ and n_7^+, discharges the three former exit/freeze assumptions,
% and advances the Paper-III-D/E/F, VIII, and IX consumers.  Shared notation
% and bibliography ownership are unchanged.
% Release v1.51 registers the PQTI derivation batch and its updated
% Paper-III-A/C/D/E/F, VIII, and IX consumers.  Shared notation and
% bibliography ownership are unchanged.
% Release v1.50 advances the Paper-III-B AUX target to v1.40 after the
% rigorous audit repair; shared notation and bibliography ownership are unchanged.
% Release v1.49 advances the Paper-II-E/F AUX target to v3.14 after the
% minimal audit repair; shared notation and bibliography ownership are unchanged.
% Release v1.46 synchronizes Paper II-A's minimal frozen constructor-origin
% theorem and its five dependent AUX targets.  Shared notation is unchanged.
% Release v1.45 updates the Paper-II-B AUX target to v2.23 after the
% GP3-rigidity and appendix-placement repair; shared notation is unchanged.
% Release v1.44 advances the centralized Paper-I-A--VII-C AUX registry after the
% Paper-I-A v5.16 protected-completion scope repair.  Every current
% paper imports this one file; no paper-level opt-in or \externaldocument list
% is permitted.  The bibliography remains centralized here through Born-v6_13.
% Shared bibliography, theorem infrastructure, colors, notation alignment,
% and current Paper-I-A--VII-C external-reference registration.
% This file centralizes common imports, bibliography sorting, colors, theorem
% infrastructure, and the notation shared by the complete series.
% v2.9 (release v1.43) made the Paper-I-A--VII-C AUX registry automatic and
% centralized the bibliography and all external-reference ownership.
% v2.8 release-registry refresh (v1.33) registers the rooted generated-parent
% cascade and the POL/REG/free-law/11+4/Born-group closure revisions.  No
% shared mathematical notation or Paper-II-A fundamental principle is changed.
% v2.8 release-registry refresh (v1.32) registers the universal generated-
% continuation revisions in Papers II-A, III-A, III-C, IV, and VIII.  No
% shared mathematical notation or Paper-II-A fundamental principle is changed.
% v2.8 release-registry refresh (v1.30) registers the CRP4/projective-state
% source revisions and their direct registry-only consumers.  No shared
% mathematical notation or Paper-II-A axiom is changed.
% v2.7 release-registry refresh (v1.27) updates only the current external
% Paper-I-A--VII-C basenames after the placement/assumption synchronization; no
% shared mathematical notation or Paper-II-A axiom is changed.
% v2.7 synchronizes the geometric defect-repair revision: Paper-II-A v3.42,
% Paper-III-C v1.26, Paper-III-D v1.27, Paper-III-E v1.32, and Paper-III-F
% v1.34.  It registers the common trace-defect/spectral-filtration theorem,
% its owner-typed response pullback, and the reduced GRED and TERM interfaces.
% No shared mathematical notation is changed.
% v2.6 synchronizes Paper-III-E v1.31's cone-type repair, minimal-response
% classification, and explicit classical action-role boundary, together with
% Paper-III-F v1.32's downstream response-atlas synchronization.  No shared
% mathematical notation is changed.
% v2.5 synchronizes the Paper-III-D v1.26 faithful-trace-duality proof,
% exact ordered-owner joint locus, and typed alternative-presentation
% bookkeeping repair, together with its Paper-III-E and closure-frontier
% references.  No shared mathematical notation is changed.
% v2.4 synchronizes the Paper-III-C v1.25 orbit/process/owner repair, its
% downstream Papers III-B--F, V, VI, VIII, and IX revisions, and Born-v6_13.
% v2.3 synchronizes the Paper-II-A v3.41 witness/support/principal inheritance
% layer, the Paper-III-A v1.24 joint-realizability/readout firewall, all dependent
% host-rate revisions, and the current Paper-I-A--VII-B source registry.
% v2.2 synchronizes the generated active-realization, determinant-depth,
% dimension-weighted hierarchy, and typed cosmological-event interfaces.
% v2.1 synchronizes the exceptional-support interface revisions: Paper I-A
% v5.9, Paper II-A v3.30, Paper II-E/F v3.7, and Paper IV v4.16.
% v2.0 introduces the stage-safe G0 occurrence notation, synchronizes the
% active Paper-I-A--IV source registry (current rigorous Paper-III closure:
% III-A v1.17, III-B v1.22, III-C--E v1.17, III-F v1.14, IV v4.14), and makes
% external-reference self-skip depend on logical paper role rather than a
% versioned job name.
% ======================================================================

% ======================================================================
% macros.tex — Shared macros for the Born Geometry series
% ======================================================================

% ======================================================================
% 1. PACKAGES, BIBLIOGRAPHY, COLORS, AND UNIT DEFINITIONS
% ======================================================================
% Core packages shared across the complete Paper-I-A--VII-C series.  Loading them here makes
% the macro file a genuine series preamble while remaining compatible with existing
% paper preambles that load the same packages before \input.
\usepackage{amsmath}
\usepackage{amssymb}
\usepackage{amsthm}
\usepackage{mathrsfs}
\usepackage{array}
\usepackage{booktabs}
\usepackage{xcolor}
\makeatletter
\@ifpackageloaded{geometry}{}{\usepackage{geometry}\geometry{margin=1in}}
\makeatother
\usepackage{mdframed}

% Shared support packages required by the macros below and by most papers.
\usepackage{ragged2e}
\usepackage{longtable}
\usepackage{tabularx}
\usepackage{enumitem}
\usepackage[most]{tcolorbox}

% Series-wide numeric bibliography convention using the portable BibTeX path.
% All Papers I-A--VII-C use the ordinary \cite interface.
\usepackage[numbers,sort&compress]{natbib}
\providecommand{\printbibliography}{%
  \bibliographystyle{plainnat}%
  \bibliography{Born-v6_59}%
}

% Shared colors.  \providecolor avoids clobbering a paper-specific override.
\providecolor{postulatecolor}{RGB}{230,242,255}
\providecolor{compare}{RGB}{230,245,230}
\providecolor{definitioncolor}{RGB}{230,242,255}

\providecolor{openriskcolor}{RGB}{255,235,235}
\providecolor{bridgecolor}{RGB}{240,240,255}
\providecolor{gold}{RGB}{255,215,0}

% ======================================================================
% 2. CUSTOM ENVIRONMENTS
% ======================================================================
% Common table column types used in all papers.  Guarded so that a paper may
% define its own versions before loading this file.
\makeatletter
\@ifundefined{NC@find@L}{%
  \newcolumntype{L}[1]{>{\RaggedRight\arraybackslash}p{#1}}%
}{}
\@ifundefined{NC@find@C}{%
  \newcolumntype{C}[1]{>{\Centering\arraybackslash}p{#1}}%
}{}
\@ifundefined{NC@find@Y}{%
  \newcolumntype{Y}{>{\RaggedRight\arraybackslash}X}%
}{}
\makeatother

% ======================================================================
% 3. THEOREM STYLES AND ENVIRONMENTS
% ======================================================================
% Custom underlined theorem style (useful for remarks or notes)
\newtheoremstyle{underlined}
  {3pt}   % Space above
  {3pt}   % Space below
  {}      % Body font (empty means normal/roman text, typical for remarks)
  {}      % Indent amount
  {}      % Theorem head font (can be \itshape, \bfseries, etc.)
  {.}     % Punctuation after theorem head
  { }     % Space after theorem head
  {\underline{\thmname{#1}\thmnumber{ #2}}\thmnote{ (#3)}} % Theorem head formatting

% Standard plain theorem style environments.  Each environment is guarded so
% the shared macro file can be loaded after a paper-specific theorem setup.
\theoremstyle{plain}
\makeatletter
\@ifundefined{theorem}{\newtheorem{theorem}{Theorem}[section]}{}
\@ifundefined{lemma}{\newtheorem{lemma}{Lemma}[section]}{}
\@ifundefined{proposition}{\newtheorem{proposition}{Proposition}[section]}{}
\@ifundefined{corollary}{\newtheorem{corollary}{Corollary}[section]}{}
\makeatother

% Definition style environments (typically non-italicized body)
\theoremstyle{definition}
\makeatletter
\@ifundefined{definition}{\newtheorem{definition}{Definition}[section]}{}
\@ifundefined{postulate}{\newtheorem{postulate}{Postulate}[section]}{}
\@ifundefined{closure}{\newtheorem{closure}{Closure}[section]}{}
\makeatother

% Applying the custom underlined style
\theoremstyle{underlined}
\makeatletter
\@ifundefined{remark}{\newtheorem{remark}{Remark}[section]}{}
\@ifundefined{convention}{\newtheorem{convention}{Convention}[section]}{}
\@ifundefined{principle}{\newtheorem{principle}{Principle}[section]}{}
\makeatother

% Framed proof environment.
% Every \begin{proof} ... \end{proof} is automatically framed/shaded.
% This avoids mdframed/trivlist nesting, which can drop the top border after
% theorem-like environments.
\makeatletter
\renewenvironment{proof}[1][\proofname]{%
  \par\addvspace{6pt}%
  \noindent
  \begin{tcolorbox}[%
    enhanced,
    breakable,
    colback=black!4,
    colframe=black!35,
    boxrule=0.35pt,
    arc=1pt,
    left=6pt,
    right=6pt,
    top=4pt,
    bottom=4pt,
    boxsep=0pt,
    before skip=0pt,
    after skip=6pt,
    parbox=false
  ]%
  \pushQED{\qed}%
  \normalfont\itshape #1\@addpunct{.}\enspace
  \normalfont\ignorespaces
}{%
  \unskip\nobreak\hfill\popQED
  \end{tcolorbox}%
  \par\addvspace{2pt}%
}
\makeatother

\mdfdefinestyle{postulatebox}{
  backgroundcolor=yellow!10,
  linecolor=black,
  linewidth=0.4pt,
  skipabove=\baselineskip,
  skipbelow=\baselineskip,
  nobreak=true
}

% ======================================================================
% PROOF-READABILITY HELPERS
% ======================================================================

\newcommand{\ProofPlan}[1]{%
  \par\smallskip
  \noindent\textbf{Proof plan. }#1%
  \par\smallskip
}

\newlist{proofsteps}{enumerate}{2}
\setlist[proofsteps,1]{
  label=\textbf{(\alph*)},
  ref=\alph*,
  leftmargin=2.2em,
  itemsep=0.45\baselineskip,
  topsep=0.35\baselineskip
}
\setlist[proofsteps,2]{
  label=\textbf{(\roman*)},
  ref=\roman*,
  leftmargin=2.0em,
  itemsep=0.25\baselineskip,
  topsep=0.2\baselineskip
}


% ======================================================================
% 4. NUMBER SYSTEMS AND ALPHABETS
% ======================================================================
\providecommand{\RR}{\mathbb{R}} % Real numbers
\providecommand{\PP}{\mathbb{P}} % Projective space / Probability
\providecommand{\CC}{\mathbb{C}} % Complex numbers
\providecommand{\ZZ}{\mathbb{Z}} % Integers
\providecommand{\NN}{\mathbb{N}} % Natural numbers
\providecommand{\HH}{\mathbb{H}} % Quaternions / Hyperbolic space
\providecommand{\LFRayGeom}{\mathsf{LF\!\mbox{-}RAY\!\mbox{-}GEOM}}
\providecommand{\LFRep}{\mathsf{LF\!\mbox{-}REP}}
\providecommand{\LFState}{\mathsf{LF\!\mbox{-}STATE}}
\providecommand{\LFLim}{\mathsf{LF\!\mbox{-}LIM}}
\providecommand{\LFScat}{\mathsf{LF\!\mbox{-}SCAT}}
% Number systems / projective spaces
\providecommand{\CP}{\mathbb{CP}}
\providecommand{\RP}{\mathbb{RP}}

% ======================================================================
% 5. GROUPS AND LIE THEORY
% ======================================================================
\providecommand{\SO}{\ensuremath{\mathrm{SO}}}     % Special Orthogonal group
\providecommand{\SU}{\ensuremath{\mathrm{SU}}}     % Special Unitary group
\providecommand{\Spin}{\ensuremath{\mathrm{Spin}}} % Spin group
\providecommand{\CS}{\ensuremath{\mathrm{CS}}}     % Chern-Simons (capitalized)
\providecommand{\cs}{\ensuremath{\mathrm{cs}}}     % Chern-Simons (lowercase)

% ======================================================================
% 6. COMMON MANIFOLDS AND TOPOLOGY
% ======================================================================
\providecommand{\Sone}{S^1}                 % Standard circle

% ======================================================================
% 7. MATH OPERATORS (Declared via \DeclareMathOperator)
% ======================================================================
\DeclareMathOperator{\End}{End}   % Endomorphisms
\DeclareMathOperator{\Aut}{Aut}   % Automorphisms
\DeclareMathOperator{\spec}{spec} % Lowercase spectrum
\DeclareMathOperator{\tr}{tr}     % Lowercase trace (if you already have \Tr)
\DeclareMathOperator{\supp}{supp} % Support
\DeclareMathOperator{\TrA}{Tr_{A}} % Partial trace over subsystem A
\DeclareMathOperator{\TrB}{Tr_{B}} % Partial trace over subsystem B
% Lie/algebra operators
\DeclareMathOperator{\Ad}{Ad}
\DeclareMathOperator{\Sym}{Sym}
\DeclareMathOperator{\su}{su}
% ======================================================================
% 8. STANDARD OPERATOR-LIKE SYMBOLS (Safe definitions via \providecommand)
% ======================================================================
\providecommand{\vol}{\operatorname{vol}}   % Volume form
\providecommand{\Vol}{\operatorname{Vol}}   % Total volume
\providecommand{\Ric}{\operatorname{Ric}}   % Ricci tensor
\providecommand{\Rm}{\operatorname{Rm}}     % Riemann tensor
\providecommand{\Tr}{\operatorname{Tr}}     % Uppercase trace
\providecommand{\Spec}{\operatorname{Spec}} % Uppercase spectrum
\providecommand{\diag}{\operatorname{diag}} % Diagonal matrix
\providecommand{\rank}{\operatorname{rank}} % Rank
\providecommand{\scal}{\mathcal R} % Scalar curvature
\providecommand{\Cl}{\operatorname{Cl}}     % Clifford algebra
\providecommand{\diam}{\operatorname{diam}} % Diameter

% ======================================================================
% 9. FUNCTIONS, TRIGONOMETRY, AND CALCULUS
% ======================================================================

% Codifferential naming conventions
% (Avoid \ddag: commonly defined by physics.sty and symbol packages)
\providecommand{\codiff}{\mathrm{d}^{\dagger}} % Codifferential operator

% Quantum mechanics bra-ket notation (needed for basis tables)
\providecommand{\ket}[1]{\lvert #1 \rangle}
\providecommand{\bra}[1]{\langle #1 \rvert}
\providecommand{\braket}[2]{\langle #1 \mid #2 \rangle}

% ======================================================================
% 10. PHYSICAL SCALES AND CONSTANTS (Paper 1c / Born Geometry)
% ======================================================================
\providecommand{\Rp}{\ensuremath{R_{\mathrm P}}}     % matched common EH curvature scale (not defined from observed G)
\providecommand{\hBorn}{\ensuremath{h_{\mathrm B}}} % physical representative 2\pi\hbar of one primitive Born scalar action cell
\providecommand{\UEHjoint}{\ensuremath{U_{J}}} % joint-Dirac EH curvature-coordinate scale 8 kappa_J
\providecommand{\Gread}{\ensuremath{G_{\rm scal}}} % scalar-action coordinate c_op^3/(h_B U); physical matching explicit
\providecommand{\GreadJoint}{\ensuremath{G_{J,\mathrm{raw}}}} % raw ordinary-trace coordinate = G_scal/96
\providecommand{\Rds}{R_{\mathrm{dS}}}    % de Sitter Curvature
\providecommand{\Mpl}{M_{\mathrm{Pl}}}    % Planck mass
\providecommand{\lp}{\ell_{\mathrm{Pl}}}  % Planck length
\providecommand{\Pmax}{P_{\mathrm{max}}}  % Maximum momentum
\providecommand{\Mplbar}{\bar{M}_{\mathrm{Pl}}} % Reduced Planck mass
\providecommand{\eff}{\mathrm{eff}}       % Subscript for effective quantities
\providecommand{\Cosmo}{\ensuremath{\Lambda}}                % Cosmological constant (use \Cosmo in source)
\providecommand{\LCDM}{\ensuremath{\Cosmo\mathrm{CDM}}}  % \LambdaCDM (typesets with \Cosmo)
\providecommand{\Lcut}{\ensuremath{\Lambda_{\mathrm{cut}}}}  % Cutoff scale (brace when used as script)
\providecommand{\Rcs}{\ensuremath{R_{\mathrm{cs}}}}      % Configuration-space radius
\providecommand{\CosmoEff}{\ensuremath{\Cosmo_{\eff}}}   % Effective cosmological constant

% Disambiguation: \Rmom is the momentum-space (S^3) radius. 
% Reserve plain R for the 4D curvature scalar in f(R).
\providecommand{\Rmom}{\ensuremath{R_{\mathrm{mom}}}} 
\providecommand{\Punit}{P_{0}} % Momentum unit set by \Rmom (= \hbar/\Rmom)
\providecommand{\Rvec}{\ensuremath{R_{\mathrm{vec}}}} % invariant vector-space radius in the reduced round sector
% Specific numeric values and expressions
\providecommand{\Vlogcap}{L_{\mathrm{Hopf}}^{(1)}}
\providecommand{\Hopf}{\operatorname{Hopf}}
% Scalar phase (S^1) sector + universal cover notation (Paper I-A)
\providecommand{\tper}{T_\star}      % Time period/scale
\providecommand{\Eper}{E_\star}      % Energy period/scale
\providecommand{\ttilde}{\tilde t}   % Universal cover time
\providecommand{\Etilde}{\tilde E}   % Universal cover energy

% ======================================================================
% 10A. SERIES-WIDE LOGICAL / NOTATIONAL DISAMBIGUATIONS
% ======================================================================
% Use these symbols to avoid collisions among: (i) the Paper--III
% no-insertion/geometric-closure class, (ii) Paper--I-A trimer diagonals,
% (iii) Paper--V-A/V-C family-space diagonals, (iv) scalar-index running,
% and (v) protected shell-clock labels.
\providecommand{\Rprim}{\ensuremath{\mathsf R_{\mathrm{prim}}^{\Sigma}}} % primitive finite structural readouts
\providecommand{\Rloop}{\ensuremath{\mathsf R_{\circ}^{\Sigma}}} % closed-loop readout role
\providecommand{\Rbdry}{\ensuremath{\mathsf R_{\partial}^{\Sigma}}} % boundary-conditioned readout role
\providecommand{\Rgen}{\ensuremath{\mathsf R_{\mathrm{fin}}^{\Sigma}}} % finite generated closure of primitive roles
\providecommand{\PhysOcc}{\operatorname{PhysOcc}}
\providecommand{\Actual}{\operatorname{Actual}}
\providecommand{\GeomLic}{\operatorname{GeomLic}}
\providecommand{\InvDist}{\operatorname{InvDist}}
\providecommand{\Dcl}{\ensuremath{\mathsf{D}_1}}      % unmarked positive Dirac-even closed spectral-readout form class
\providecommand{\Dheat}{\ensuremath{\mathsf D_{\mathrm{heat}}}} % one-parameter heat-semigroup subclass
\providecommand{\DtriI}{\ensuremath{D^{\mathrm{tri}}_1}} % Paper--I-A cyclic-trimer diagonal
\providecommand{\DtriII}{\ensuremath{D^{\mathrm{tri}}_2}}
\providecommand{\DfamI}{\ensuremath{D^{\mathrm{fam}}_1}} % Family-space diagonal selector
\providecommand{\DfamII}{\ensuremath{D^{\mathrm{fam}}_2}}
\providecommand{\alpharun}{\ensuremath{\alpha_{\mathrm{run}}}} % scalar spectral-index running
\providecommand{\Ksh}{\ensuremath{K_{\mathrm{sh}}}} % protected shell-clock coordinate
\providecommand{\ddef}{\ensuremath{\delta_{\mathrm{def}}}} % universal defect scalar after Paper--III realization

% ======================================================================
% 10B. SERIES-WIDE LOGICAL STATUS TAGS
% ======================================================================
% Upright small-caps wrapper for audit/status labels.  It prevents LaTeX from
% trying to use unavailable small-caps italic font shapes inside theorem bodies
% such as TU/lmr/m/scit.
\providecommand{\AuditLabel}[1]{}
\renewcommand{\AuditLabel}[1]{{\normalfont\scshape #1}}

% These tags are bookkeeping labels for consistency audits. They do not
% create new mathematical hypotheses; they record how a datum is being used
% in the proof chain.
\providecommand{\StatusPrimitive}{}
\renewcommand{\StatusPrimitive}{\AuditLabel{Primitive}}
\providecommand{\StatusDerived}{}
\renewcommand{\StatusDerived}{\AuditLabel{Derived}}
\providecommand{\StatusImported}{}
\renewcommand{\StatusImported}{\AuditLabel{Imported}}
\providecommand{\StatusConditional}{}
\renewcommand{\StatusConditional}{\AuditLabel{Conditional}}
\providecommand{\StatusBridge}{}
\renewcommand{\StatusBridge}{\AuditLabel{Bridge-derived}}
\providecommand{\StatusComparison}{}
\renewcommand{\StatusComparison}{\AuditLabel{Comparison-only}}

% ======================================================================
% 10C. SERIES-WIDE CLAIM-WORDING AUDIT LABELS
% ======================================================================
% These are textual audit labels used to keep modal language consistent
% across the series.  They introduce no new mathematical assumptions.
\providecommand{\ClaimForced}{}
\renewcommand{\ClaimForced}{\AuditLabel{forced}}
\providecommand{\ClaimUnique}{}
\renewcommand{\ClaimUnique}{\AuditLabel{unique}}
\providecommand{\ClaimCanonical}{}
\renewcommand{\ClaimCanonical}{\AuditLabel{canonical}}
\providecommand{\ClaimMinimal}{}
\renewcommand{\ClaimMinimal}{\AuditLabel{minimal}}
\providecommand{\ClaimClosed}{}
\renewcommand{\ClaimClosed}{\AuditLabel{closed}}
\providecommand{\ClaimPrediction}{}
\renewcommand{\ClaimPrediction}{\AuditLabel{prediction}}

% ======================================================================
% 10D. SERIES-WIDE THEOREM-CHAIN AUDIT LABELS
% ======================================================================
% These labels mark the direction of load-bearing theorem chains. They are
% audit annotations only; they introduce no new mathematical assumptions.
\providecommand{\ChainInput}{}
\renewcommand{\ChainInput}{\AuditLabel{Chain input}}
\providecommand{\ChainOutput}{}
\renewcommand{\ChainOutput}{\AuditLabel{Chain output}}
\providecommand{\ChainGuard}{}
\renewcommand{\ChainGuard}{\AuditLabel{Chain guard}}
\providecommand{\ChainNoBackflow}{}
\renewcommand{\ChainNoBackflow}{\AuditLabel{No backflow}}

% ======================================================================
% 10E. SERIES-WIDE LOCAL THEOREM-AUDIT LABELS
% ======================================================================
% These labels mark a statement-level dependency audit around individual
% theorem/definition nodes. They are bookkeeping labels only.
\providecommand{\LocalAudit}{}
\renewcommand{\LocalAudit}{\AuditLabel{Local audit}}
\providecommand{\UsesOnly}{}
\renewcommand{\UsesOnly}{\AuditLabel{Uses only}}
\providecommand{\Exports}{}
\renewcommand{\Exports}{\AuditLabel{Exports}}
\providecommand{\DoesNotUse}{}
\renewcommand{\DoesNotUse}{\AuditLabel{Does not use}}

% ======================================================================
% 10F. SERIES-WIDE NUMERICAL-ROW AUDIT LABELS
% ======================================================================
% These labels classify individual rows in numerical, benchmark, status,
% and comparison tables.  They are audit annotations only and do not change
% any theorem, value, formula, or comparison number.
\providecommand{\RowAudit}{}
\renewcommand{\RowAudit}{\AuditLabel{Numerical-row audit}}
\providecommand{\RowTheorem}{}
\renewcommand{\RowTheorem}{\AuditLabel{Theorem-fixed row}}
\providecommand{\RowBridge}{}
\renewcommand{\RowBridge}{\AuditLabel{Bridge-fixed row}}
\providecommand{\RowRepresentative}{}
\renewcommand{\RowRepresentative}{\AuditLabel{Representative row}}
\providecommand{\RowComparison}{}
\renewcommand{\RowComparison}{\AuditLabel{Comparison row}}
\providecommand{\RowCalibrated}{}
\renewcommand{\RowCalibrated}{\AuditLabel{Calibrated row}}
\providecommand{\RowEstimate}{}
\renewcommand{\RowEstimate}{\AuditLabel{Estimate row}}
\providecommand{\RowOpen}{}
\renewcommand{\RowOpen}{\AuditLabel{Open row}}
\providecommand{\RowNoSelector}{}
\renewcommand{\RowNoSelector}{\AuditLabel{No selector}}

% ======================================================================
% 10G. TABLE HELPERS
% ======================================================================
% Table helpers used by Paper V-C.  The unused readout-label preset has been
% removed; only active table helpers remain.
\providecommand{\gpsection}[1]{%
  \addlinespace[0.55em]
  \multicolumn{3}{@{}l}{\textsc{#1}}\\[-0.25em]
  \midrule
}
\providecommand{\mathcell}[1]{%
  \(\begin{gathered}#1\end{gathered}\)%
}

\newcommand{\DerBy}[1]{%
  \begin{tabular}[t]{@{}l@{}}#1\end{tabular}%
}

\newcommand{\TableNotes}[1]{%
  \par\medskip
  \begin{minipage}{0.94\linewidth}
  \footnotesize\emph{Notes.} #1
  \end{minipage}\par
}
% ======================================================================
% 10H. PAPER VI-A / QCD NOTATION
% ======================================================================
% Shared symbols used by Paper VI-A and downstream QCD/Yang--Mills handoffs.
\providecommand{\Pqcd}{\mathfrak P_{\rm VIII}^{\rm QCD}}
\providecommand{\Iqcd}{\mathfrak I_{\rm QCD}^{(7)}}
\providecommand{\muQCD}{\mu_{\rm QCD}^{\rm int}}
\providecommand{\LambdaQCD}{\Lambda_{\rm QCD}}
\providecommand{\LambdaBH}{\Lambda_{\rm BH}^{\rm QCD}}
\providecommand{\HQCD}{\mathcal H_{\rm QCD}^{\rm BH}}
\providecommand{\HQCDfin}{H_{\rm QCD}^{\rm fin}}
\providecommand{\Hraw}{\mathcal H_{\rm raw}}
\providecommand{\HphysQCD}{\mathcal H_{\rm phys}^{\rm QCD}}
\providecommand{\Vcol}{V_C}
\providecommand{\Ps}{P_{\rm sing}}
\providecommand{\PiOct}{\Pi_8}
\providecommand{\MSbar}{\overline{\rm MS}}
\providecommand{\YM}{\mathrm{YM}}
\providecommand{\Ktri}{K_{\triangle}}
\providecommand{\KYM}{K_{\YM}^{\rm BH}}
\providecommand{\aBH}{a_{\rm BH}}
\providecommand{\epsBH}{\varepsilon_{\rm BH}}
\newcommand{\RowConstructive}{\textnormal{\textsc{Constructive}}}
\newcommand{\RowObligation}{\textnormal{\textsc{Bridge obligation}}}
\newcommand{\RowTarget}{\textnormal{\textsc{Target row}}}
\newcommand{\RowValidation}{\textnormal{\textsc{Validation row}}}
\providecommand{\slashed}[1]{\not{#1}}

% ======================================================================
% 10I. PAPER VI-B / CONTINUOUS-QFT NOTATION
% ======================================================================
% Shared symbols used by Paper VI-B and any downstream continuous-QFT or
% Yang--Mills handoff paper.  All definitions are guarded so individual
% papers may override them locally before loading this file if needed.
\providecommand{\A}{\mathfrak A}
\providecommand{\Hphys}{\mathcal H_{\rm phys}}
\providecommand{\QCD}{\mathrm{QCD}}
\providecommand{\BH}{\mathrm{BH}}
\providecommand{\loc}{\mathrm{loc}}

\providecommand{\RBH}{\mathcal R_{\rm BH}}
\providecommand{\DG}{\mathfrak D_G}
\providecommand{\PW}{\mathcal{P\!W}} % Peter--Weyl spectral-symbol subspace
\providecommand{\Toe}{\operatorname{Toe}} % Toeplitz compression/representative
\providecommand{\Hol}{\operatorname{Hol}} % Holonomy

% ======================================================================
% 11. LATE-LOADED PACKAGES (Must be at the end of the preamble)
% ======================================================================
% Load late packages only if a paper has not already loaded them, avoiding
% option clashes in paper-specific preambles.
\makeatletter
\@ifpackageloaded{hyperref}{}{\usepackage{hyperref}}
\makeatother
% --- Series-wide disambiguations (Pass 3) ---
\providecommand{\diracdelta}{\delta_{\rm D}} % Dirac delta distribution (delta reserved for hierarchy)
\providecommand{\Hg}{H_{\mathrm g}} % Hodge harmonic projector (keep H for Hubble)
\providecommand{\Gr}{G_{\mathrm r}} % Green operator (keep G for Newton)
\providecommand{\FD}{\mathrm{D}} % functional/variational derivative symbol (delta reserved)
\providecommand{\Id}{\mathbf{1}} % identity/Kronecker tensor (avoid \delta_{ij})


% BEGIN INLINE SOURCE: Paper00-ActiveExternalReferences-v2_130.tex
% Framework v2.191: current reference binding; existing definitions retained.
% Framework v2.190: current22-root registry with native clock-evolution proofs.
% Framework v2.189: exact current22-root external registry, including retained v2.188 hardening owners.
% Framework v2.187: current registry synchronization only; source release v2.186 preserved.
% Framework v2.186: authoritative current 22-paper registry.
% Framework v2.182: all current 22 logical paper AUX bindings.
% Framework v2.170: current I-B v1.43 external AUX binding.
% Framework v2.167: exact current 22-root registry after CP--BL2 integration.
% Framework v2.166: exact current 22-root compact-release AUX registry.
% Framework v2.165: authoritative current root bindings.
% Framework v2.164: exact current 22 manuscript AUX bindings.
% Framework v2.162: finite metric, functional Hilbert and scoped Ward integration.
% Framework v2.161: current 22 manuscript AUX bindings; source only.
% Framework v2.156 / registry v2.113: synchronized current 22 manuscript roots; source-only registry.
% Framework v2.143 / registry v2.111: active manuscript AUX basenames; no generated AUX or PDF included.
% ======================================================================
% Central external-reference imports, numbers only, no PDF jumping
% ======================================================================

% xr-hyper is loaded centrally so that papers may retain their existing
% \usepackage{xr-hyper} lines without redefining \externaldocument after this
% registry is installed.  If plain xr was loaded earlier, retain it.
\makeatletter
\@ifpackageloaded{xr-hyper}{}{%
  \@ifpackageloaded{xr}{}{\usepackage{xr-hyper}}%
}
\makeatother

% Make external-document registration idempotent.  The registry identifies a
% logical paper by its Paper<number><optional letter> prefix.  Prefix-qualified
% imports remain independent.  Current paper roots contain no local
% \externaldocument calls; this guard also makes accidental duplicates harmless.
\ExplSyntaxOn
\prop_new:N \g_bh_external_refs_prop
\cs_new_eq:NN \BHOriginalExternalDocument \externaldocument

\cs_new_protected:Npn \bh_external_document:nnnn #1#2#3#4
 {
  \tl_set:Nn \l_tmpa_tl {#3}
  \regex_replace_once:nnN
   { \A (?: xref/ )? (Paper[0-9]+[A-F]?) .* \Z }
   { \1 }
   \l_tmpa_tl
  \tl_set:Nx \l_tmpb_tl { #1 | \tl_use:N \l_tmpa_tl }
  \prop_if_in:NVTF \g_bh_external_refs_prop \l_tmpb_tl
   {
    \PackageInfo{BornHopf}
     {Skipping~duplicate~external~refs~for~#3}
   }
   {
    \prop_gput:NVn \g_bh_external_refs_prop \l_tmpb_tl {#3}
    \IfNoValueTF{#4}
     { \BHOriginalExternalDocument[#1][#2]{#3} }
     { \BHOriginalExternalDocument[#1][#2]{#3}[#4] }
   }
 }

\RenewDocumentCommand{\externaldocument}{O{} O{} m o}
 { \bh_external_document:nnnn {#1}{#2}{#3}{#4} }

\NewDocumentCommand{\BHExternalRefs}{m}
 {
  % Do not import the aux file of the logical paper currently being compiled,
  % even when the job and registry use different version suffixes.
  \tl_set:Nx \l_tmpa_tl { \c_sys_jobname_str }
  \regex_replace_once:nnN
   { \A (Paper[0-9]+[A-F]?) .* \Z }
   { \1 }
   \l_tmpa_tl
  \tl_set:Nn \l_tmpb_tl { #1 }
  \regex_replace_once:nnN
   { \A (Paper[0-9]+[A-F]?) .* \Z }
   { \1 }
   \l_tmpb_tl
  \str_if_eq:eeTF { \tl_use:N \l_tmpa_tl } { \tl_use:N \l_tmpb_tl }
   {
    \PackageInfo{BornHopf}
     {Skipping~external~refs~for~#1~because~this~is~the~current~document}
   }
   {
    \IfFileExists{xref/#1.aux}
     { \externaldocument[][nocite]{xref/#1}[#1.pdf] }
     {
      \IfFileExists{#1.aux}
       { \externaldocument[][nocite]{#1} }
       {
        \PackageInfo{BornHopf}
         {External~aux~file~#1.aux~not~found;~no~labels~imported}
       }
     }
   }
 }
\ExplSyntaxOff

% External references are automatic.  This file is the sole owner of both the
% bibliography and the current Paper-I-A--VII-D AUX registry (30 roots including I-C).  Missing auxiliaries
% are informational on the first pass and are populated by the source-only
% two-pass series build.  Basenames deliberately omit temporary download/copy
% suffixes such as ``(1)'' or ``-1''.
  % Paper I-A: intrinsic S3 and exceptional group introduction.
  \BHExternalRefs{Paper1A-Geometry-v5_175}

  % Paper I-B: complete exceptional geometry and finite spectral structures.
  \BHExternalRefs{Paper1B-ExceptionalGeometry-v1_5}
  % Paper I-C: reciprocal product M8 geometry.
  \BHExternalRefs{Paper1D-TechnicalMaterial-v1_65}
  % Paper I-D: terminal technical proofs and calculations.
  \BHExternalRefs{Paper1C-M8Geometry-v1_137}

\ifdefined\BHMathOnly
  % Papers I-A and I-B import only the mathematical atlas labels.
\else
  % Paper II: five physical-mathematical main papers plus one technical proof volume.
  \BHExternalRefs{Paper2A-GeometricConstitution-v3_221}
  \BHExternalRefs{Paper2B-MomentumCarrier-v2_144}
  \BHExternalRefs{Paper2C-ReciprocalLocalization-v2_169}
  \BHExternalRefs{Paper2D-SymmetrySpectra-v2_141}
  % The former II-F export labels remain in II-E; current II-F is the technical volume.
  \BHExternalRefs{Paper2E-KinematicExports-v3_138}
  % Paper II-F: complete technical proof and calculation owner for II-A--II-E.
  \BHExternalRefs{Paper2F-TechnicalMaterial-v1_67}

  % Paper III: Dynamics and Intrinsic Geometric Constitution (six components).
  \BHExternalRefs{Paper3A-GeneratedDynamics-v1_176}
  \BHExternalRefs{Paper3B-SpectralAction-v1_187}
  \BHExternalRefs{Paper3C-ProtectedHierarchy-v1_187}
  \BHExternalRefs{Paper3D-RunningResponse-v1_165}
  \BHExternalRefs{Paper3E-GravityAction-v1_197}
  \BHExternalRefs{Paper3F-CompleteDynamics-v1_263}
  % Paper III-G: complete preserved technical proof owner.
  \BHExternalRefs{Paper3G-TechnicalMaterial-v1_53}

  % Paper IV-A: generated finite quantum mechanics.
  \BHExternalRefs{Paper4A-FiniteQuantumMechanics-v4_214}
  % Paper IV-B: complete finite quantum processes and physical response.
  \BHExternalRefs{Paper4B-QuantumResponse-v1_49}
  % Paper IV-C: extended proofs and finite technical calculations.
  \BHExternalRefs{Paper4C-TechnicalMaterial-v1_49}

  % Paper V-A: The Standard Model.
  \BHExternalRefs{Paper5A-FiniteMatterGeometry-v3_209}

  % Paper V-B: Finite Spectrum and Flavor Closure.
  \BHExternalRefs{Paper5B-TerminalSpectrum-v5_199}

  % Paper V-C: Finite Interaction Boundaries of the Standard Model.
  \BHExternalRefs{Paper5C-InteractionResponse-v2_177}
  % Paper V-D: preserved technical proof and calculation owner.
  \BHExternalRefs{Paper5D-TechnicalMaterial-v1_53}

  % Paper VI-A: Gravity and Closed Cosmology.
  \BHExternalRefs{Paper6A-GravityCosmology-v2_260}

  % Paper VI-B: Geometric Dark Sector and Cosmological Readouts.
  \BHExternalRefs{Paper6B-CosmologyReadouts-v1_206}

  % Paper VI-C: complete technical proof and comparison owner.
  \BHExternalRefs{Paper6C-TechnicalMaterial-v1_50}

  % Paper VII-A: Finite QCD Kernel and Nonperturbative Closure.
\BHExternalRefs{Paper7A-QCDKernel-v1_172}

  % Paper VII-B: Hadronic, Thermal, and Nuclear QCD Readouts.
  \BHExternalRefs{Paper7B-QCDReadouts-v1_157}

  % Paper VII-C: Continuous Quantum Field Theory.
  \BHExternalRefs{Paper7C-ContinuousQFT-v2_202}

  % Paper VII-D: complete technical material for VII-A through VII-C.
  \BHExternalRefs{Paper7D-TechnicalMaterial-v1_53}

\fi

% END INLINE SOURCE: Paper00-ActiveExternalReferences-v2_130.tex



\setlength{\emergencystretch}{3em}


\title{\textbf{Born--Hopf Geometry in a Closed Universe VII-C: Bounded-Momentum Continuous QFT}\\[0.4em]
\large Fixed-UV, phase-selected reconstruction, quantum source identities and field certificates}
\author{Peter Kugelmann\thanks{\texttt{pkugelmann@proton.me}}}
\date{Version 2.208 --- 19 November 2026 \\ Framework source integration v3.025}
\setcounter{tocdepth}{2}
% Cross-DAG current-series import for the new source-owned mathematical bridge.
\BHExternalRefs{Paper7B-QCDReadouts-v1_159}

\BHExternalRefs{Paper7A-QCDKernel-v1_175}
\begin{document}
% Publication typesetting: controlled prose line-breaking, without scaling mathematics.
\emergencystretch=2em
\maketitle
\begin{abstract}
Paper VII-C gives the claim-indexed continuum-QFT certificate of the
 finite Born--Hopf prefix. Its native limit keeps the momentum ceiling
 $p_{\max}$ fixed while reciprocal-position radius increases and momentum
 spacing tends to zero. A selected bounded-local cluster phase and
 Osterwalder--Schrader/locality/regularity hypotheses yield a conditional QFT
 export; uniqueness, uniform gap, composite fields, asymptotic scattering and
 massless infrared completion require stronger separately named certificates.
 Source, action and Ward claims retain their actual finite state and operator
 owners. The v2.254 analysis isolates an additional uniform zero-frequency detector-limit certificate; this is not discharged by local QFT identity preservation alone. Detailed RG, subsystem and counterfamily calculations are in VII-D.
\end{abstract}
\tableofcontents
\listoftables
\clearpage

\section{Fixed-UV Born continuum, inputs and phase-selected status}
\label{p12:sec:v237-c1}
The controlling native limit is $R_x\to\infty$, $\Delta p\to0$ and $p_{\max}$ fixed. A selected interacting phase and an actual state/action ray are not created by a divergent ultraviolet momentum bound or by finite ancestry alone.

% Source Paper VII-C, original line 72; inherited proof bodies unchanged.
% Former logical part: Framework Prefix and Continuum-Certificate Boundary; integrated into the new thematic argument.

\noindent\textbf{Orientation.}
This part fixes the inherited prefix \(\mathfrak F_{\le7B}\), the one-way handoff from Paper~VII-B, and the status classes used throughout Paper~VII-C.  It is the dependency ledger for the continuous-QFT certificate; no new finite carrier, gauge sector, Lorentzian branch, or Standard--Model input is selected here.


% Source Paper VII-C, original line 77; inherited proof bodies unchanged.
\subsection{Framework status of Paper VII-C}
\label{p12:sec:scope}


% Source Paper VII-C, original line 80; inherited proof bodies unchanged.
\subsubsection{Framework prefix and native question}
\label{p12:sec:introduction-purpose}

Paper~VII-C is not an independent selection layer.  It is the continuum-QFT layer over the completed finite-core/QCD prefix
\begin{equation}
\label{p12:eq:framework-prefix}
\begin{aligned}
\mathfrak F_{\le 7B}
={}&\mathfrak F_{\le6B}
+\text{Paper-VII-A finite QCD kernel}\\
&+\text{Paper-VII-B finite QCD readout/export layer}.
\end{aligned}
\end{equation}
Here \(\mathfrak F_{\le6B}\) denotes the reference geometry, compact kinematics, dynamics, finite quantum closure, Standard--Model carrier, spectrum and interaction layers, gravitational realization, and cosmological readouts supplied by Papers~I-A--VI-B.  Paper~VII-A closes the finite \(SU(3)_c\) kernel, while Paper~VII-B constructs the physical readout layer and exports the pure-glue Wilson/curvature subdatum \(\mathfrak D_{7B,\YM}^{SU(3)}\) used here.

The native question of Paper~VII-C is therefore:
\[
\begin{gathered}
\text{when does a finite Born--Hopf regulator tower over }\mathfrak F_{\le7B}\\
\text{define an ordinary local-chart continuum QFT?}
\end{gathered}
\]
The output is a reconstruction theorem of the schematic form
\begin{equation}
\label{p12:eq:main-pipeline}
\left\{\A_L,\Omega_L,K_L,\chi_L\right\}_{L\in\RBH}
\Longrightarrow
\left(\Hphys,\Omega,H,\A(\mathcal O),U(\mathcal P_+^\uparrow)\right),
\end{equation}
where the left side is finite and Born--Hopf-native, while the right side is the reconstructed Lorentzian OS/Wightman local-chart QFT presentation, compatible on compact reciprocal-position windows.

\textbf{Global theorem-scope convention.}
Every theorem of Paper~VII-C is read with its displayed hypotheses.  In particular, an analytic continuum hypothesis is not promoted to a consequence merely because its finite carrier is inherited from \(\mathfrak F_{\le7B}\).  The logical form is
\begin{equation}
\label{p12:eq:global-theorem-scope}
\mathfrak F_{\le7B}+\{\text{displayed theorem hypotheses}\}
\vdash\text{the stated Paper-VII-C conclusion}.
\end{equation}
A theorem called ``native'' therefore means that no additional continuum \emph{structure} is inserted beyond the stated finite tower and comparison data; it does not mean that unresolved convergence, covariance, gap, or field-limit hypotheses are automatically discharged by ancestry.

Following the theorem-status split now explicit in Papers~II-A and III-A, Paper~VII-C distinguishes two uses of ordinary mathematics.  An \emph{internal theorem closure} is a consequence whose physical owners and symmetry/action antecedents are already generated on the stated domain; a separate theorem may then only unpack an unavoidable algebraic or variational consequence.  An \emph{analytic certification} still requires claim-specific hypotheses such as refinement Cauchy control, invariant-state or measure control, reflection positivity, local-flat covariance, regularity, clustering, residue bounds, or continuum convergence.  G0 actualization and GP ancestry do not prove those analytic hypotheses and do not actualize a continuum history.


% Source Paper VII-C, original line 123; inherited proof bodies unchanged.
\subsubsection{Layer ledger and relation to Papers VII-A--VII-B}
\label{p12:sec:why-separate}

Papers~VII-A--VII-B form the finite QCD layer.  Paper~VII-A owns the finite nonperturbative compiler, transfer family, generated algebra/commutant, exact Hodge geometry, and typed chiral/topological targets; physical confinement and continuum claims remain separate obligations, while a uniform gap is retained only for an optional branch that explicitly claims one.  Paper~VII-B owns the finite hadronic, thermal, nuclear, and validation readouts and exports the finite pure-glue subdatum used by Paper~VII-C.  Paper~VII-C begins only after that export is fixed.  Its role is to prove continuum state, positivity, locality, covariance, local-field, nontriviality, and sector-compatibility certificates, together with a conditional mass-gap transfer implication when its extra antecedent is supplied.

\begin{table}[!htbp]
\centering
\small
\renewcommand{\arraystretch}{1.12}
\begin{tabularx}{\textwidth}{@{}p{0.18\textwidth}p{0.25\textwidth}X@{}}
\toprule
\textbf{Layer} & \textbf{Status in Paper VII-C} & \textbf{Role} \\
\midrule
Paper~I-A & Reference import & Hopf-fibered \(S^3\), contact normalization, complete round-sphere spectra, the mathematical \(\mathcal H_m\) truncation family, and finite symmetric-coordinate boundary ranks. \\
Paper~II & Kinematic import & GP1--GP2, physical constitution of the compact momentum carrier, finite kinematic dictionaries, and symbol-density recovery. \\
Paper~III & \makecell[l]{Dynamical and\\gravitational import} & Positive Hamiltonian clock, compact momentum \(m_p=1\), reciprocal-position unfolding, Lorentzian sign selection, gravitational action, metric field equation, and scalar-response constitution. \\
Paper~IV & Finite-quantum import & Finite quantum algebra, Fock/interface sectors, Euclidean kernels, transfer positivity, local-net bridges, and quasi-free reconstruction. \\
Papers~V-A--V-C & Sector-carrier import & Standard--Model carrier, electromagnetic remnant, finite masses/mixings/Yukawa representatives, and QCD interface data. \\
Paper~VI-A / Paper~VI-B & Gravity/readout import & Paper~VI-A closed-cosmology and nontrivial gravitational realization, followed by Paper~VI-B cosmological observational readouts. \\
Paper~VII-A & Finite QCD kernel import & \(SU(3)_c\) compiler, singlet projection, finite Hamiltonian/transfer family, generated algebra/commutant, exact Hodge geometry, and typed chiral/topological targets.  The complete finite source/action/transfer datum and its licensed Haar/heat representative are imported on their stated domains; actual member/source matching, a physical position-space area law and continuum fields retain their own certificates. A uniform continuum gap is not an active claim. \\
Paper~VII-B & Finite QCD export import & Finite readout package and the pure-glue subdatum \(\mathfrak D_{7B,\YM}^{SU(3)}\): singlet Hilbert blocks, finite Hamiltonians/transfers, Wilson algebra, Hodge operator, vacuum states, constructed LF fraction measures, and isolated pole/residue limit theorems.  Continuum reconstruction, area law, and local-field limits remain Paper-VII-C obligations; uniform-gap transfer is optional and claim-specific. \\
Paper~VII-C & Native layer & OS/Wightman continuum certificate, sector-import compatibility, QFT1--QFT10 ledger, spectral UCS, Yang--Mills handoff, and the claim-indexed scattering/infrared firewall. \\
Downstream use & Export/comparison & Scattering, asymptotic completeness, collider observables, and external problem-statement comparisons only on their named claim domains. \\
\bottomrule
\end{tabularx}
\caption[Paper-VII-C framework layer ledger]{Paper-VII-C framework layer ledger.}
\label{p12:tab:framework-layer-ledger}
\TableNotes{The table fixes the architecture once.  Later theorem statements are read over \(\mathfrak F_{\le7B}\) without repeating the full prefix.}
\end{table}


% Source Paper VII-C, original line 153; inherited proof bodies unchanged.
\subsubsection{Imports from the previous papers}
\label{p12:sec:imports}

The present paper imports the following data from the prefix.
\begin{enumerate}[label=(I\arabic*), leftmargin=2.8em]

\item From Paper~I-A \cite{Kugelmann2026a}: the Hopf-fibered \(S^3\) geometry, contact normalization,
harmonic and higher-spin sphere spectra, the mathematical \(\mathcal H_m\) truncation bank, and the symmetric-coordinate boundary-rank geometry.

\item From Paper~II: GP1--GP2, the physical constitution of the compact momentum carrier, finite kinematic dictionaries, retained Peter--Weyl symbols, and
continuum symbol-density recovery.

\item From Paper~III \cite{Kugelmann2026c}: the one-dimensional dynamical reduction, positive
Hamiltonian clock, compact momentum \(m_p=1\) closure, reciprocal-position
unfolding, and Lorentzian sign-selection theorem.

\item From Paper~IV: the maximal finite envelope algebra and selector-subalgebra
criterion; the finite state, effect, Born, POVM, instrument, channel, tensor,
partial-trace, unitary, open-system, and conservation calculus; finite Fock
sectors, exact CAR and the finite bosonic CCR boundary defect; finite gauge
projection and Gibbs/KMS structure; finite Euclidean kernels, transfer
positivity, local-net and quasi-free bridge statements; the positive-frequency
finite dispersion theorem; and the \(m_p\) versus \(m_x\) regime dictionary.
These are Paper-IV-owned imports and are not independently rederived in
Paper~VII-C.

\item From Papers~V-A--V-C: the Standard--Model carrier, chiral gauge--matter
package, scalar stabilizer, electromagnetic remnant \(U(1)_{\rm em}\),
transverse photon/field-strength branch, finite masses, mixings, and Yukawa representatives,
the \(D_4\)-vertex-figure/\(A_2\)-projector origin of the strong color
support, and QCD interface data.

\item From Papers~VI-A--VI-B: the cosmological realization layer and the
branch-level data used by downstream large-scale comparison and validation
statements.  In particular, Paper~VI-B exports the corrected ordered pair vertex
\(V_{16}\), the ordinary occupation support
\(\overline{\operatorname{Ran}\Gamma_{16}^{\rm ord}}\), and the separation
between coherent finite Bogoliubov production and later kinetic
thermalization.  It does not export a selected inverse temperature, collision
semigroup, or convergence theorem.

\item From Paper~VII-A: the repaired finite \(SU(3)_c\) compiler, exact singlet projector, finite pure-glue Hamiltonian/transfer family, generated operator algebra and commutant, exact Hodge geometry, and the typed chiral/topological targets.  Paper~VII-A does not export a count-derived shell amplitude, fixed wall scale, positive continuum gap, physical string tension, or completed strong-CP invariant.

\item From Paper~VII-B: the repaired finite hadronic, thermal, nuclear, and validation readout layer together with the finite export datum \(\mathfrak D_{7B}^{\rm QCD}\) and its direct pure-glue subdatum \(\mathfrak D_{7B,\YM}^{SU(3)}\).  Paper~VII-B leaves continuum state convergence, OS reconstruction, position-space area law, and local-field limits to Paper~VII-C.  It exports a uniform-gap sequence only as an optional conditional input, not as an active mass-gap prediction.

\end{enumerate}

Throughout the repaired Papers~VII-A--VII-B handoff the physical finite pure-glue data are the self-adjoint singlet Hamiltonians \(H_{L,\YM}\), their vacuum energies \(E_{0,L}\), and the transfers \(e^{-t(H_{L,\YM}-E_{0,L})}\).  No universal \(Q_{\rm wall}\), primitive shell loss, or count-derived gap normalization is imported.  If \(\gamma_L:=E_{1,L}-E_{0,L}\) denotes the finite pure-glue gap, a continuum positive-gap claim is controlled by the single sharp quantity \(\gamma_*:=\inf_L\gamma_L\), after the relevant state/semigroup limit has been established.

\textbf{Paper-VII-C continuum-certificate convention.}
A continuum statement is closed when the finite Born--Hopf tower supplies the relevant certificate data: finite comparison maps, local-state convergence, reflection positivity, locality/covariance, a gap or controlled gapless branch, and local-field reconstruction.  The reconstructed Hilbert space, Hamiltonian, local net, and continuum fields are OS/Wightman outputs.  The native application is the Paper~VII-B export \(\mathfrak D_{7B,\YM}^{SU(3)}\); the one-way export dependency is Theorem~\ref{p12:thm:paper11-paper12-acyclic}.

\textbf{Momentum-first presentation.}
All position-space language in this paper is a reciprocal-position presentation of a spectral/momentum Born--Hopf regulator.  Paper~III fixes the microscopic compact momentum carrier at \(m_p=1\), and Paper~IV imports that realized carrier into the finite quantum calculus.  Large position resolution belongs to the reciprocal-position label \(m_x\gg1\).  The symbols \(\Gamma_L\) and \(\chi_L\), when they occur below, denote auxiliary large-\(m_x\) quadrature/window data used to present local observables in position charts.  Lorentzian signature and positive-frequency real-time propagation are imported from the sequential Papers~III--IV closure; Paper~VII-C reconstructs the continuum state/local-net side over that orientation.

\begin{theorem}[Momentum-rooted finite structural provenance]
\label{p12:thm:momentum-rooted-finite-structural-provenance}
Within the active Papers~I-A--VII-C constructor chain, let
\(S^3_{p,1}\) denote the \emph{constituted} level-one momentum-space
finite-geometry package of Paper~I-A, including the retained markings,
completion data, and constructors actually licensed there; it does not mean
the bare vector space \(\mathcal H_1\) in isolation.  Then the following
provenance statements hold.
\begin{enumerate}
\item The physical compact momentum carrier is the generated specialization
\[
   \mathcal H_p\cong\mathcal H_1,\qquad m_p=1.
\]
\item The reciprocal-position role \(S_x^3\) is the smooth Born/GP2-generated
reciprocal role.  It is not a second microscopic finite carrier, site set, or
position lattice.
\item Every finite \(x\)-side object used in the active chain is typed as a
harmonic-resolution window or shell, a frontier and its derived
least-resolution statistic, a transported representation, a finite
chart/quadrature presentation, or a downstream readout on that smooth role.
\item The finite gauge, charge, color, family, anomaly, and analogous internal
structural carriers or selectors used physically downstream are inherited
through the licensed finite constructor chain rooted in the constituted
momentum package; no active theorem introduces a second microscopic finite
structural root on the position side.
\item Actual frontier occurrences and histories, physical states,
dimensionful calibrations, action members, numerical flavor/Yukawa or
Majorana members, and analytic continuum certificates remain separately typed
where their owner papers keep them separate.  They are not asserted to follow
from the finite structural carrier alone.
\end{enumerate}
Consequently the active series has one microscopic finite structural root on
the momentum side together with a generated smooth reciprocal-position role
and its finite-resolution presentations.
\end{theorem}

\begin{proof}
Paper~I-A constitutes the level-one momentum package and explicitly keeps its
retained completion/marking data distinct from the bare carrier.  Paper~II-B
then fixes the physical compact specialization
\(\mathcal H_p\cong\mathcal H_1\), \(m_p=1\).  Papers~II-A--II-C derive the
reciprocal B-role by intrinsic evaluation and GP2 and identify the physical
\(x\)-side finite datum with harmonic resolution rather than a site set or
second carrier; in particular the primitive frontier has
\(m_x^{\rm res}=1\) as a derived resolution statement.  Papers~III--IV retain
that typed split while constructing the reciprocal product, dynamics, and
finite quantum calculus.  Papers~V-A--V-C generate their finite internal
structural sectors from the inherited finite chain while leaving the stated
physical member/fiber data separate.  Papers~VI-A--VII-C use the smooth
Born-derived reciprocal-position geometry for cosmological, QCD, and
continuum presentations and do not insert a primitive position lattice.
The separately displayed occurrence, calibration, member, and analytic
certificate boundaries in those papers establish the exclusions in item~5.
\end{proof}

\begin{corollary}[No equal-level ontology]
\label{p12:cor:no-equal-level-ontology}
The primitive bookkeeping coincidence
\[
   (m_x^{\rm res},m_p)=(1,1)
\]
does not imply
\(\mathcal P_x(1)\cong\mathcal H_p\), equal carrier dimensions, equal site
counts, or a second position carrier.  The two entries label differently
typed primitive stages: reciprocal-position resolution and compact momentum
carrier level, respectively.
\end{corollary}

\begin{corollary}[Reciprocal-position resolution growth does not generate internal species]
\label{p12:cor:v259-resolution-growth-no-internal-species}
Fix the active momentum-rooted finite structural chain of
Theorem~\ref{p12:thm:momentum-rooted-finite-structural-provenance} and vary
only the reciprocal-position frontier through complete retained bands
\(\mathcal P_x(m_x)\).  Then changing \(m_x\) alone cannot generate or change
any finite internal structural carrier or selector already inherited from the
momentum-rooted chain, including the number of color directions, family slots,
gauge-factor carriers, or anomaly-sector carriers.

If a downstream sector is explicitly represented on a product or tensor
presentation of the form
\[
 \mathcal P_x(m_x)\otimes V_{\rm int}
\]
with fixed inherited internal factor \(V_{\rm int}\), increasing \(m_x\) may
increase only the reciprocal-position/external mode multiplicity; it does not
alter \(V_{\rm int}\) or create a new internal species.  Any such internal
change requires an additional licensed finite witness, selector, or
parent-completion constructor and cannot be inferred from spatial refinement.

This is a structural provenance statement.  It does not freeze physical
states, couplings, masses, occupation numbers, or effective parameters that
may run or change under their separately owned dynamics.
\end{corollary}

\begin{proof}
Items 2--4 of
Theorem~\ref{p12:thm:momentum-rooted-finite-structural-provenance} type every
finite x-side frontier as resolution/presentation data and place the physical
finite internal structures in the licensed constructor chain rooted in the
constituted momentum package.  Varying only \(m_x\) therefore changes the
x-side factor but supplies no new internal constructor.  In an explicit tensor
presentation the dimension factorizes as
\(\dim\mathcal P_x(m_x)\,\dim V_{\rm int}\), so only the first factor varies.
The final exclusions follow from item 5, which keeps states, members and
analytic certificates separately typed.
\end{proof}

\begin{corollary}[Mode-density transport is provenance, not field-sector constitution]
\label{p12:cor:v2180-mode-density-not-sector-constructor}
The reciprocal-position trace and no-new-internal-species theorems determine
the external mode-density and fixed-internal-factor consequences of a
downstream sector already represented on
\[
 \mathcal P_x(m_x)\otimes V_{\rm int}.
\]
They do not by themselves generate an interacting CAR/quark occurrence, a
local color representation on that carrier, or gauge-covariant first-order
Dirac shifts.  Those are separately typed physical constructors.
\end{corollary}

\begin{proof}
Corollary~\ref{p12:cor:v259-resolution-growth-no-internal-species} is explicitly
conditional on a downstream sector being represented on the product, while
Theorem~\ref{p12:thm:momentum-rooted-finite-structural-provenance} types
$m_x$ as reciprocal-position resolution/presentation data.  Neither statement
contains a CAR-field occurrence or a covariant quark-shift constructor.
\end{proof}

\begin{definition}[Paper-VII-C certificate status classes]
\label{p12:def:status-classes}
A statement in this paper has one of the following statuses.
\begin{description}[leftmargin=3.4em]
\item[Imported.] A result proved in Papers~I-A--VII-B and used with its stated theorem domain.
\item[Native certificate.] A finite-to-continuum implication proved in Paper~VII-C from the Born--Hopf tower and its certificate obligations.
\item[Sector extension.] A rigorous extension once the displayed optional finite sector package is supplied.
\item[Comparison/export.] A translation into standard QFT, Clay, scattering, or continuum terminology, downstream of the native certificate.
\end{description}
\end{definition}

\begin{proposition}[Native regulator and reconstruction boundary]
\label{p12:prop:momentum-first-no-lattice}
\label{p12:prop:no-continuum-primitive-inserted}
\label{p12:prop:gp1gp2-generation-principle}
Paper~VII-C uses the compact momentum carrier \(m_p=1\), reciprocal-position resolution \(m_x\), and the Lorentzian orientation supplied by Papers~III--IV.  Retained spectral symbols, finite Toeplitz/Wilson representatives, positive finite quadratures, and the geometric local-flat chart exhaustion are finite spectral/operator presentations or derived finite-dimensional constructions.  By contrast, interacting cross-level state convergence, any continuum chiral determinant-line/global-integrability certificate, a uniform gap, all-order residue cumulant bounds, and unique local-field limits are not generated by GP2 and must be separately proved when used.  Continuum Hilbert space, vacuum, Hamiltonian, local net, and any reconstructed local fields are outputs of the corresponding verified limit hypotheses.
\end{proposition}
\begin{proof}
The fixed compact-momentum/reciprocal-position regime split is an upstream import.  Peter--Weyl density and finite compression produce the bounded symbol sector, while Proposition~\ref{p12:prop:tier2-quadrature} derives positive exact quadrature from finite-dimensional convexity and Lemma~\ref{p12:lem:tier2-flat-chart-covariance} derives the geometric chart exhaustion.  The counterexamples in Propositions/Lemmas~\ref{p12:prop:tier2-qft4-estimates}, \ref{p12:lem:native-residue-bound}, \ref{p12:thm:euclidean-corridor-transfer}, and \ref{p12:lem:finite-cumulant-connectivity} show that GP2 and finite-carrier ancestry do not imply the missing interacting convergence estimates.  Therefore the only continuum objects declared derived are those reconstructed after their actual analytic hypotheses have been verified.
\end{proof}


% Source Paper VII-C, original line 357; inherited proof bodies unchanged.
\subsection{Handoff boundary and dependency ledger}
\label{p12:sec:handoff-dependency-ledger}


% Source Paper VII-C, original line 360; inherited proof bodies unchanged.
\subsubsection{Acyclic handoff from Paper VII-B}
\label{p12:sec:acyclic-handoff}

\begin{theorem}[Acyclic repaired Paper-VII-B--Paper-VII-C handoff]
\label{p12:thm:paper11-paper12-acyclic}
The current dependency graph is acyclic in the following strict sense:
\[
\begin{gathered}
\text{Papers I-A--VII-A finite structures}\\[-2pt]
\downarrow\\[-2pt]
\mathfrak D_{7B,\YM}^{SU(3)}\\[-2pt]
\downarrow\\[-2pt]
\text{Paper-VII-C continuum obligations and reconstruction}.
\end{gathered}
\]
Paper~VII-B exports only finite pure-glue singlet Hilbert blocks, vacuum states, self-adjoint Hamiltonians/transfers, the bounded Wilson algebra, and exact Hodge data.  Paper~VII-C may use those finite objects, but it may not import a fixed wall scale, count-derived shell loss, uniform continuum gap, position-space area law, continuum local field, or physical \(\bar\theta\) unless that object is separately constructed and proved.
\end{theorem}
\begin{proof}
The repaired Paper~VII-B export definition contains the finite family
\[
\{\mathcal H_{L,\YM}^{\rm sing},\Omega_{L,\YM},H_{L,\YM},e^{-t(H_{L,\YM}-E_{0,L})},\mathfrak A_{W,L}^{SU(3)},\Delta_{2,H}\}_{L\in\mathcal R_{\rm adm}}
\]
and explicitly states that continuum reconstruction, OS positivity, a position-space area law, and local-field limits remain Paper-VII-C obligations.  A uniform mass gap appears only in a claim-specific conditional branch.  Release v1.64 factorizes Paper~VII-A's former position-area item: local symbol transport, metric area, holonomy, and chart exhaustion are generated, whereas strict positive surface free energy remains unresolved.  Hence every continuum arrow starts only after the finite datum is fixed, while no Paper-VII-C conclusion can be used retroactively to define an entry of Papers~VII-A--VII-B.  This proves acyclicity and the stated source boundary.
\end{proof}



% Source Paper VII-C, original line 8176; inherited proof bodies unchanged.
\subsection{Source-domain tests and projective state comparison}
\label{p12:app:v202-source-comparison}
For the positive two-point observable $O=\pm1$ with equal probabilities,
$Z(J)=\cosh J$ vanishes at $J=i\pi/2$. Hence the universal commuting
source disk in Theorem~\ref{p12:thm:v202-bounded-source-domain} cannot
cross that boundary for the whole stated class. Positivity is essential:
the signed weight example $Z(J)=101-100e^J$ has $Z(0)=1$ but a zero
at $\log(101/100)<0.01$.

For positive weights on a common fixed configuration space put
$d_H(F,G)=\operatorname{ess\,osc}\log(F/G)$. Positive linear integration
contracts this distance, since $mG\le F\le MG$ implies
$m\mathcal BG\le\mathcal BF\le M\mathcal BG$.
Normalized probabilities satisfy the sharp comparison
\[
 \|P_F-P_G\|_{\rm TV}\le\tanh[d_H(F,G)/4],
 \qquad \|P-Q\|_{\rm TV}:=\sup_A|P(A)-Q(A)|.
\]
Indeed, for $X=dP_F/dP_G\in[m,M]$ with mean one, the chord bound on
$(X-1)_+$ gives $(M-1)(1-m)/(M-m)$. Maximizing at fixed $M/m=e^d$
gives $\tanh(d/4)$; proportional weights give zero directly. See
\cite{CohenFaustiHilbert2024} for the general comparison.

On a countable actual tower of standard Borel configuration spaces with
consistent measurable projections, the sufficient condition
\[
 \|\pi_{n,n-1*}P_n-P_{n-1}\|_{\rm TV}\le\delta_n,
 \qquad \sum_n\delta_n<\infty
\]
makes every fixed coarse marginal Cauchy: pushforward contracts total
variation and the telescoping remainder is bounded by the tail sum.
The limits give the unique compatible projective state for that
sequence, and Theorem~\ref{p12:thm:v202-bounded-source-inheritance}
applies to its bounded cylinders. The native summability estimate is not
proved here; general quantum comparison maps require their own typed
state-space formulation. This normalization comparison concerns fixed
states. Source- or metric-dependent constants remain in complete-action
derivatives and gravitational vacuum bookkeeping.
% END V202 SOURCE DOMAIN DIAGNOSTICS

\begingroup
\small
\RaggedRight
\setlength{\emergencystretch}{3em}



% BEGIN V283 PHASE-SELECTED FIXED-UV QFT EXPORT

% Source Paper VII-C, original line 8224; inherited proof bodies unchanged.
\subsection{fixed-UV phase-selected QFT export and continuum-status retyping}
\label{p12:sec:v283-phase-selected-fixed-uv-export}

\begin{theorem}[Fixed-UV Born continuum]
\label{p12:thm:v283-fixed-uv-born-continuum}
On the native Born refinement branch,
\[
 R_x\to\infty,\qquad \frac{m_x}{R_x}\to\Lambda_F,
 \qquad p_{\max}=\hbar\Lambda_F\ \text{fixed}.
\]
Thus the continuum presentation densifies the momentum spectrum inside a fixed momentum ball and enlarges reciprocal-position extent.  It does not remove the ultraviolet cutoff and does not introduce a primitive position-space lattice spacing tending to zero.
\end{theorem}
\begin{proof}
This is the typed content of the existing fixed-\(m_p=1\) sharp-frontier/momentum-ball limit: the frontier label grows proportionally to reciprocal-position radius while the physical momentum ceiling remains the same protected compact-momentum owner.  Therefore \(\Delta p\sim R_x^{-1}\to0\) but \(p_{\max}\) is fixed.
\end{proof}

\begin{definition}[Cluster-state existence and canonicality]
\label{p12:def:v283-qft4-existence-canonicity}
Let \(\mathfrak C_\alpha\) be the weak-* cluster set of the bounded local state net restricted to one supplied/selected physical phase \(\alpha\).  Define
\[
 \mathrm{QFT4}_E:\quad \mathfrak C_\alpha\ne\varnothing,
 \qquad
 \mathrm{QFT4}_C:\quad |\mathfrak C_\alpha|=1.
\]
QFT4$_E$ is continuum-state existence.  QFT4$_C$ is the stronger canonical same-phase refinement theorem.
\end{definition}

\begin{theorem}[QFT4 existence on a constructed bounded-word algebra]
\label{p12:thm:v283-qft4-cluster-split}
For a directed family of finite unital $C^*$ algebras $\mathcal A_L$
with states $\omega_L$ and uniformly bounded representatives of each
fixed cylinder word, there exists a common asymptotic bounded-word
$C^*$ algebra and at least one positive normalized cluster state on it.
This proves QFT4$_E$ for that \emph{constructed algebra} without first
assuming inter-level $*$-homomorphisms.
Identification with a separately prescribed physical local cylinder algebra
requires that its relations descend in norm (or in the selected GNS
representation), together with the relevant completion and locality tests.
QFT4$_C$ on the constructed algebra holds exactly when the full net of
fixed cylinder-word expectations converges to a unique
refinement-path-independent state.
\end{theorem}
\begin{proof}
Let $\mathfrak P=\prod_L\mathcal A_L$ be the bounded direct product
and $\mathfrak J=\{(A_L):\lim_L\|A_L\|=0\}$ its closed two-sided
$*$-ideal. Then $\mathfrak C_\infty=\mathfrak P/\mathfrak J$ is a unital
$C^*$ algebra. Evaluation of $\omega_L$ on the $L$-th product entry
is a state $\Phi_L$ on the \emph{same} $\mathfrak P$.
Banach--Alaoglu gives a weak-* convergent subnet $\Phi_{L_i}\to\Phi$.
For $A\in\mathfrak J$, $|\Phi_{L_i}(A)|\le\|A_{L_i}\|\to0$;
thus $\Phi$ descends to a positive normalized state on
$\mathfrak C_\infty$. Fixed bounded cylinder word families generate
a unital $C^*$ subalgebra. A singleton cluster set is equivalent to
convergence of the entire net by compactness and the usual subnet
argument; that is a stronger condition than cluster existence.
\end{proof}
\begin{proposition}[Exact cylinder-relation descent test]
\label{p12:prop:v382-cylinder-relation-descent}
For the free $*$-algebra on the declared bounded cylinder generators,
let $\pi_L$ evaluate words in the actual finite operators and put
$p(a)=\limsup_L\|\pi_L(a)\|$. An intended relation $r=0$ descends to
$\mathfrak C_\infty$ if and only if $p(r)=0$.
In a selected cluster GNS representation the weaker exact condition is
$\omega_\alpha(b^*r^*rb)=0$ for every bounded cylinder word $b$.
Identification with a preassigned completed algebra additionally
requires norm-continuity of this quotient representation.
\end{proposition}
\begin{proof}
The direct-product quotient norm is the tail $\limsup$ norm;
its kernel is precisely the asymptotic null ideal. A relation has
zero GNS operator exactly when $\|\pi_\omega(r)\pi_\omega(b)\Omega\|^2
=\omega(b^*r^*rb)=0$ on a dense set of word vectors.
\end{proof}


\begin{theorem}[Countable local-test criterion for a unique bounded continuum state]
\label{p12:thm:v292-dense-local-test-convergence}
Let $\mathfrak A$ be a unital separable quasi-local $C^*$-algebra and
$\{\omega_L\}_{L\in\mathcal D}$ a directed net of states on the same
algebra, or a norm-one compatible local state net for which each fixed
local test is eventually defined.  Fix a countable norm-dense collection
$\{A_k:k\geq1\}\subset\mathfrak A$ of bounded tests.  Then
\[
\boxed{\quad
\forall k:\ \lim_L\omega_L(A_k)\ \text{exists}
\quad\Longleftrightarrow\quad
\omega_L\xrightarrow{\text{weak-*}}\omega\in S(\mathfrak A).
\quad}
\]
In particular, on the bounded same-phase net of
Theorem~\ref{p12:thm:v283-qft4-cluster-split}, the left condition is an
exact sufficient-and-necessary certificate for $\mathrm{QFT4}_C$, rather
than merely $\mathrm{QFT4}_E$.  The criterion consumes the full directed
net and a norm-dense test family; no finite fixed list of observables or
single cofinal subnet certifies it in general.  It does not establish
unbounded-field, scattering, or source-phase occurrence conclusions.
\end{theorem}
\begin{proof}
For $A\in\mathfrak A$ and $\varepsilon>0$ choose $A_k$ with
$\|A-A_k\|<\varepsilon$.  Since each $\omega_L$ has norm one, for
$L,M$ sufficiently late,
\[
 |\omega_L(A)-\omega_M(A)|
 \leq 2\varepsilon+|\omega_L(A_k)-\omega_M(A_k)|.
\]
The rightmost term tends to zero on the full directed net.  Thus
$\omega_L(A)$ is Cauchy for every $A$; its limit defines a bounded
normalized positive linear functional $\omega$, hence a state, and
$\omega_L\to\omega$ weak-*.  For compatible local nets the same argument
first defines the state on the norm-dense local union and then extends it
uniquely by continuity.  Conversely weak-* convergence gives convergence
on each $A_k$.  Uniqueness of the limit and the singleton-cluster
criterion of Theorem~\ref{p12:thm:v283-qft4-cluster-split} establish the
claimed bounded $\mathrm{QFT4}_C$ interpretation.
\end{proof}

\begin{proposition}[No universal fixed finite-test certificate for phase uniqueness]
\label{p12:prop:v292-finite-tests-phase-countermodel}
For $\mathfrak A=C(X)$ with $X$ any infinite compact metric space and
any finite family of real continuous tests $f_1,\ldots,f_m$, there are
distinct states $\omega_+,\omega_-$ that agree on every $f_j$.  Hence a
net alternating between them has constant values on all these tests,
but does not converge weak-*.  A finite list of scalar readouts therefore
cannot replace the norm-dense-family hypothesis of
Theorem~\ref{p12:thm:v292-dense-local-test-convergence} without an
additional state-separation or finite-dimensionality certificate.
\end{proposition}
\begin{proof}
Choose $m+2$ distinct points $x_1,\ldots,x_{m+2}$ in $X$.
Their vectors
\[
(1,f_1(x_i),\ldots,f_m(x_i))\in\mathbb R^{m+1}
\]
lie in a vector space of dimension $m+1$ and hence are linearly
dependent, so there are real coefficients
$c_i$, not all zero, with $\sum_i c_i=0$ and
$\sum_i c_if_j(x_i)=0$ for every $j$.  The nonzero positive and negative
parts of this signed measure have equal mass.  Normalize those parts
to obtain different probability measures $\mu_+,\mu_-$ with identical
$f_j$ integrals.  Their states on $C(X)$ are distinct, so some other
continuous function separates them.  Alternating the two states along
the natural-number refinement net proves the claim.
\end{proof}

\begin{corollary}[Reflection positivity and bounded regularity for every phase cluster]
\label{p12:cor:v283-cluster-rp-regularity}
Finite reflection positivity and refinement-uniform bounded test-function estimates pass to every weak-* cluster state.  Hence neither QFT5 nor bounded QFT7 consumes QFT4$_C$.
\end{corollary}
\begin{proof}
Reflection positivity is a finite matrix of scalar inequalities in bounded local expectations and is closed under scalar limits.  The bounded regularity estimates have refinement-independent seminorm bounds, which survive the same limit.
\end{proof}

\begin{theorem}[Local positive-measure 1PI transfer on a fixed source chart]
\label{p12:thm:v283-local-1pi-transfer}
Let a fixed finite classical source tuple have positive finite-level measures and locally uniform bounded-source convergence \(W_L(J)\to W(J)\) on a common complex neighborhood.  Let \(a+E\) be the minimal affine support of the limiting source law.  Then
\[
 \ker D^2W(J)=E^\perp,\qquad D^2W(J)|_E>0.
\]
On every compact noncritical source neighborhood the source-field maps \(\phi_L=\nabla W_L\) are eventually uniformly invertible on \(E\); their local Legendre transforms satisfy
\[
 \Gamma_L\to\Gamma,\qquad D^n\Gamma_L\to D^n\Gamma
\]
for every fixed derivative order on smaller compact neighborhoods.
\end{theorem}
\begin{proof}
The Hessian is the tilted covariance.  A covariance direction vanishes exactly when that linear source coordinate is constant on the minimal affine support, giving kernel \(E^\perp\) and strict positivity on \(E\).  On a compact noncritical neighborhood the smallest positive eigenvalue is bounded away from zero.  Local uniform holomorphic convergence gives uniform convergence of the Hessians on smaller neighborhoods, so the finite inverse-function theorem applies uniformly.  Convergence of the inverses and of the Legendre transforms, together with all fixed derivatives, follows by composition.  Critical points, growing source dimension, unbounded QFT10 fields and global phase-crossing Legendre transforms are outside this theorem.
\end{proof}

\begin{theorem}[Phase-selected bounded prediction transfer]
\label{p12:thm:v283-phase-prediction-transfer}
Let \(\omega_{L_\beta}\to\omega_\alpha\) be a phase-selected weak-* subnet and let \(A_{1,L},\ldots,A_{n,L}\) be compatible bounded local representatives.  Then
\[
 \omega_\alpha(A_1\cdots A_n)
 =\lim_\beta\omega_{L_\beta}(A_{1,L_\beta}\cdots A_{n,L_\beta}),
\]
and every fixed connected cumulant and bounded effect/POVM probability has the corresponding limit.  Any prediction that is a continuous function of finitely many such controlled data therefore agrees with the phase-selected continuum QFT prediction.
\end{theorem}
\begin{proof}
The product is again a bounded local observable on a common sufficiently large region.  Weak-* convergence gives the moment limit.  Connected cumulants are finite polynomials in moments, and effects are bounded local observables.  Continuity of the final finite-dimensional prediction functional completes the result.
\end{proof}

\begin{corollary}[QFT8 and QFT10 status]
\label{p12:cor:v283-qft8-qft10-status}
QFT8 is a dependent observable non-Gaussianity certificate and, when claimed, a stronger same-source dynamical interaction certificate, not a prerequisite for QFT existence.  QFT10 is an extension layer: source-owned local fields descend under their common source/UCS certificate, while independently normalized external composites retain claim-scoped mixed-moment obligations.  Neither QFT8 nor QFT10 is required for the bounded local OS net.
\end{corollary}

\begin{remark}[Scattering status]
The massive finite-window/M\o ller/cross-section theorem remains owned by Paper~IV and is claim-scoped to channels satisfying its locality, leakage, return, tail, detector and flux hypotheses.  Massless/long-range channels retain their separate inclusive/dressed/algebraic-detector infrared certificate.  Neither asymptotic layer is part of the base local-QFT existence theorem.
\end{remark}

\begin{corollary}[Continuum frontier retyping after v2.83]
\label{p12:cor:v283-continuum-frontier-retyping}
The finite QCD regulator/action/measure owner and exact coarse-map square are not ultraviolet-cutoff-removal problems.  The bounded phase-cluster state is derived by compactness. Exact ordinary
local-QFT export remains the implication of Theorem~\ref{p12:thm:main-qft}
on its actual locality, full covariance and reconstruction certificate;
those analytic tests are not new finite physical owners.  Remaining canonical same-phase uniqueness, external-composite field limits, ordered chiral/topological comparison, and infrared/scattering claims are analytic comparison or claim-scoped certificates and must not be counted as new native physical owners.
\end{corollary}
% END V283 PHASE-SELECTED FIXED-UV QFT EXPORT

% BEGIN V282 G20 CONTINUUM REROUTE

% Source Paper VII-C, original line 8315; inherited proof bodies unchanged.
\subsection{chiral/topological obligations after native G20 retirement}
\label{p12:sec:v282-g20-continuum-reroute}
Paper~VII-A now treats the overlap/GW index as the native finite topological
charge on its gapped domain and proves populated primitive \(\nu=\pm1\)
sectors of positive native magnitude.  Paper~IV already owns the canonical
fixed-carrier paired Berezin measure.  Consequently neither native topology
nor the finite vectorlike color measure waits for a global chiral determinant
line or an unrestricted clutching identification.

\begin{definition}[Continuum chiral/topological comparison certificate]
\label{p12:def:v282-chiral-topological-comparison-certificate}
Whenever a continuum claim uses a gauge-field-dependent chiral/Weyl subspace,
or identifies the native overlap index with a continuum/geometric clutching
class on a larger admitted family, its certificate must establish the
claim-relevant subset of:
\begin{enumerate}
\item source-uniform gapped/projector control along the refinement family;
\item determinant-line curvature cancellation and trivial holonomy, or the
      precise anomaly obstruction, for the actual varying chiral owner;
\item unrestricted overlap/clutching comparison on the domain where that
      identification is asserted;
\item ordered chiral/topological refinement limits and anomaly/Ward transport.
\end{enumerate}
These are analytic/reconstruction obligations under G39 and the ordered-limit
layer where consumed.  They are not a second headline physical-owner gap.
\end{definition}

\begin{corollary}[No G20 duplication in the continuum frontier]
\label{p12:cor:v282-no-g20-duplication}
The global chiral determinant-line and unrestricted overlap/clutching tasks
are counted once, under the continuum certificate of
Definition~\ref{p12:def:v282-chiral-topological-comparison-certificate}.
They must not be duplicated as an independent native G20 residual.
\end{corollary}
% END V282 G20 CONTINUUM REROUTE

% BEGIN G33-SYNC v2.84
\paragraph{Current G33 scope.}
Native finite geometric--quantum reconstruction is closed and is not a
headline residual.  A stronger \(\mathrm{G33\text{-}COMP}\) certificate is
retained only on claims that require complete-action/source-uniform multiscale
control or continuum-universality estimates.  Interacting continuum state and
action existence remain under their separately typed QFT4/\texttt{QFT4\_C} and G39
owners; they are not renamed as native G33.
% END G33-SYNC v2.84

% Bibliography consolidated at the end of this volume
\endgroup

% BEGIN SAFEGUARD-CONTINUUM-SELECTOR-SYNC v2.86


\section{Refinement, projective state/action rays and phase existence}
\label{p12:sec:v237-c2}
Weak-* cluster existence on a supplied phase is separate from same-phase canonical uniqueness. Reflection positivity does not by itself select an interacting state/action ray.

% Source Paper VII-C, original line 384; inherited proof bodies unchanged.
% Former logical part: Admissible Regulator Towers and Sector Compatibility; integrated into the new thematic argument.

\noindent\textbf{Orientation.}
This part defines the finite regulator towers and the certificate obligations that make an ordinary continuum-QFT presentation admissible over the Born--Hopf prefix.  It also records how Standard--Model, QCD, controlled gapless, and scattering-boundary data enter as imported or optional sector packages rather than new primitive choices.


% Source Paper VII-C, original line 389; inherited proof bodies unchanged.
\subsection{Admissible Born--Hopf refinement systems}
\label{p12:sec:admissible-refinements}


% Source Paper VII-C, original line 392; inherited proof bodies unchanged.
\subsubsection{Continuum-certificate data ledger}
\label{p12:sec:refinement-ledger}

\begin{table}[htbp]
\centering
\small
\begin{tabularx}{\textwidth}{@{}
>{\raggedright\arraybackslash}p{0.55\textwidth}
>{\raggedright\arraybackslash}X
@{}}
\toprule
\textbf{Continuum certificate datum} & \textbf{Discharged by} \\
\midrule
Admissible Born--Hopf refinement set
& Definition~\ref{p12:def:admissible-refinement-set}. \\

Finite Born--Hopf QFT datum
& Definition~\ref{p12:def:finite-bh-qft-datum}. \\

Local observable comparison system
& Definition~\ref{p12:def:local-comparison-system}. \\

Asymptotic local commutativity
& Definition~\ref{p12:def:asymptotic-local-commutativity}. \\

Born--Hopf continuum-QFT obligations
& Definition~\ref{p12:def:qft-obligations}. \\
\bottomrule
\end{tabularx}
\caption[Born--Hopf refinement and QFT obligations]{Born--Hopf refinement and QFT obligations.}
\label{p12:tab:refinement-qft-obligations}
\TableNotes{The following sections turn the finite regulator tower into a precise
continuum certificate problem.  The table records the objects whose rows are
later closed by compactness, positivity, OS reconstruction, gap transfer, and
gauge-specialized Wilson-loop arguments.}
\end{table}


% Source Paper VII-C, original line 429; inherited proof bodies unchanged.
\subsubsection{Admissible refinement sets}
\label{p12:sec:admissible-refinement-sets}

The finite theory at a fixed truncation level is already part of the earlier papers.  The new problem is the refinement system: a directed family of finite systems whose local observables can be compared on compact regions of a limiting Euclidean chart.

\begin{definition}[Admissible Born--Hopf refinement set]
\label{p12:def:admissible-refinement-set}
An admissible Born--Hopf refinement set is a directed set \(\RBH\) with a scale map
\[
L\longmapsto a_L>0,
\qquad
L\preceq L'\Rightarrow a_{L'}\le a_L,
\qquad
\lim_La_L=0.
\]
For the Born--Hopf regulator the native refinement data are spectral-symbol data, not primitive point lattices.  At level \(L\), the spatial retained symbol space is a finite Peter--Weyl space
\[
\PW_L\subset C^\infty(S^3),
\]
with finite representative map of Toeplitz type
\[
\Toe_L(f)=\Pi_LM_f\Pi_L\in\End(\mathcal H_L).
\]
The refinement is admissible when
\[
\PW_L\subseteq\PW_{L'},\qquad L\preceq L',
\qquad
\overline{\bigcup_L\PW_L}^{\|\cdot\|_\infty}=C(S^3),
\]
and when, after adjoining the scalar Euclidean-time cover from Papers~III--IV,
restrictions of retained symbols are uniformly dense on compact local product
charts \(I\times U\subset\RR\times S^3\).  Optional discrete windows
\(\Gamma_L\) may be used for quadrature or Wilson-loop bookkeeping, but they
are not continuum inputs and no metric convergence of finite point sets is
assumed as a starting axiom.  In particular, \(\Gamma_L\) is a presentation
of retained spectral data and is not the definition of the reciprocal-position
window itself; changing an exact quadrature while preserving the retained
functionals does not change the underlying Born-derived position role.
\end{definition}


% Source Paper VII-C, original line 469; inherited proof bodies unchanged.
\subsubsection{Finite QFT data and local comparison}
\label{p12:sec:finite-qft-local-comparison}

\begin{definition}[Finite Born--Hopf QFT datum]
\label{p12:def:finite-bh-qft-datum}
A finite Born--Hopf QFT datum over \(\RBH\) is a family
\begin{equation}
\label{p12:eq:finite-bh-qft-datum}
\mathfrak D_{\rm QFT}
:=
\left\{
(\Gamma_L,\chi_L,\mathcal H_L,\Omega_L,H_L,K_L^{\mathrm{ref}},\A_L,\tau_{L,t})
\right\}_{L\in\RBH},
\end{equation}
where
\begin{enumerate}[label=(\roman*), leftmargin=2em]
\item \(\mathcal H_L\) is a finite-dimensional Hilbert space;
\item \(\Omega_L\in\mathcal H_L\) is a unit vacuum vector;
\item \(H_L\ge0\) is a finite Hamiltonian and \(K_L^{\mathrm{ref}}=e^{-a_LH_L}\) is the one-step Euclidean refinement transfer;
\item \(\A_L\subset\mathcal B(\mathcal H_L)\) is a finite observable algebra;
\item \(\tau_{L,t}(A)=e^{tH_L}Ae^{-tH_L}\) is the finite Euclidean time evolution on its domain of definition.
\end{enumerate}
The finite state is
\begin{equation}
\label{p12:eq:finite-state}
\omega_L(A):=\langle\Omega_L,A\Omega_L\rangle.
\end{equation}
\end{definition}

\begin{definition}[Refinement transfer versus fixed-time sector transfer]
\label{p12:def:transfer-scale-convention}
For the repaired handoff no external wall scale is required.  Two uses of the same finite Hamiltonian are kept distinct.
\begin{enumerate}[label=\textnormal{(\roman*)}, leftmargin=2.2em]
\item The \emph{refinement transfer}
\begin{equation}
\label{p12:eq:refinement-transfer}
K_L^{\mathrm{ref}}:=e^{-a_L(H_L-E_{0,L})}
\end{equation}
uses the reciprocal-position comparison step \(a_L\) and is a bookkeeping representative for reflection-positive finite Euclidean time increments.
\item For any fixed physical Euclidean time \(t_0>0\), the \emph{sector transfer block}
\begin{equation}
\label{p12:eq:physical-sector-transfer}
\widehat K_{L,\mathcal S}(t_0):=e^{-t_0(H_{L,\mathcal S}-E_{0,L})}
\end{equation}
is determined by the supplied finite Hamiltonian; \(t_0\) is a comparison time, not a new QCD normalization parameter.
\end{enumerate}
Let \(P_{0,L,\mathcal S}:=\mathbf 1_{\{E_{0,L}\}}(H_{L,\mathcal S})\)
be the orthogonal spectral projector onto the \emph{entire} sector ground
eigenspace (possibly of rank greater than one), and assume the excited
complement \(\mathcal H_{L,\mathcal S}^{\rm exc}:=
\operatorname{ran}(I-P_{0,L,\mathcal S})\) is nonzero.
Define \(\gamma_L>0\) to be the least eigenvalue of
\(H_{L,\mathcal S}-E_{0,L}\) on that excited complement. Finite spectral
functional calculus then gives exactly
\[
\left\|\widehat K_{L,\mathcal S}(t_0)
 \big|_{\mathcal H_{L,\mathcal S}^{\rm exc}}\right\|
=e^{-t_0\gamma_L}.
\]
Therefore a refinement-uniform Hamiltonian gap
\(\gamma_*:=\inf_L\gamma_L>0\) is equivalent to the fixed-time
transfer contraction on the full-ground-complement,
\(\sup_L\|\widehat K_{L,\mathcal S}(t_0)
|_{\mathcal H_{L,\mathcal S}^{\rm exc}}\|
\le e^{-t_0\gamma_*}<1\).
Using only \(\Omega_L^\perp\) requires a separately proved rank-one
ground eigenspace; it is not implicit here. No \(Q_{\rm wall}\) or
shell-loss coefficient is needed.
\end{definition}

\begin{definition}[Local observable comparison system]
\label{p12:def:local-comparison-system}
For each bounded open \(\mathcal O\subset\RR^4\), a local observable comparison system assigns finite-dimensional subalgebras
\[
\A_L(\mathcal O)\subset\A_L
\]
for all sufficiently refined \(L\), together with comparison maps
\[
\iota_{L}^{\mathcal O}:\A_0(\mathcal O)\longrightarrow \A_L(\mathcal O)
\]
from a common algebraic cylinder algebra \(\A_0(\mathcal O)\).  If \(\mathcal O_1\subset\mathcal O_2\), the inclusions are compatible up to errors \(O(a_L^\eta)\) in all finite vacuum expectations, for some \(\eta>0\).
\end{definition}

\begin{remark}[Common cylinder-algebra comparison]
The finite algebras \(\A_L\) need not literally embed into each other.  The continuum limit only requires that fixed local questions have compatible representatives at every sufficiently fine level.  This is the same logical form as the finite-symbol and local-net bridges already used in Paper~IV, but here it is made the primary object.
\end{remark}


% Source Paper VII-C, original line 539; inherited proof bodies unchanged.
\subsubsection{Asymptotic locality}
\label{p12:sec:asymptotic-locality}

\begin{definition}[Asymptotic local commutativity]
\label{p12:def:asymptotic-local-commutativity}
The refinement datum is asymptotically local if, whenever \(\mathcal O_1,\mathcal O_2\subset\RR^4\) are spacelike-separated after OS reconstruction, representatives \(A_L\in\A_L(\mathcal O_1)\) and \(B_L\in\A_L(\mathcal O_2)\) satisfy
\begin{equation}
\label{p12:eq:asymptotic-commutator}
\|[A_L,B_L]\|_{\omega_L}
:=
\omega_L([A_L,B_L]^*[A_L,B_L])^{1/2}
\le C_{A,B}a_L^\eta .
\end{equation}
\end{definition}


% Source Paper VII-C, original line 554; inherited proof bodies unchanged.
\subsection{Continuum certificate obligations for a general interacting QFT}
\label{p12:sec:qft-obligations}


% Source Paper VII-C, original line 557; inherited proof bodies unchanged.
\subsubsection{Certificate-obligation ledger}
\label{p12:sec:qft-obligation-list}

The following list is the general-QFT analogue of the Paper~VII-B BH1--BH10 Yang--Mills ledger.  We use four blocks throughout:
\[
\begin{array}{ccl}
\text{QFT1--QFT3} &:& \text{symbols, local cylinder algebras, boundedness},\\
\text{QFT4--QFT7} &:& \text{state convergence, OS positivity, covariance, regularity/cluster},\\
\text{QFT8--QFT9} &:& \text{nontriviality and gap or controlled gapless sector},\\
\text{QFT10} &:& \text{regulator-normalized local-field moment convergence}.
\end{array}
\]

\begin{definition}[Assumption-minimized Born--Hopf continuum-QFT obligations]
\label{p12:def:qft-obligations}
A finite Born--Hopf QFT datum is organized in three logical tiers.
\begin{enumerate}[label=\textnormal{(QFT\arabic*)}, leftmargin=3.4em]
\item \textbf{Spectral-symbol chart recovery.} Retained Peter--Weyl symbols are uniformly dense on the compact Born-derived reciprocal-position geometry and admit an exhausting large-\(m_x\) local flat-chart presentation.
\item \textbf{Common local cylinder algebra.} Each fixed bounded local question has representatives at all sufficiently fine levels, compatible under inclusion of regions. Literal embeddings of the whole finite matrix algebras are not required.
\item \textbf{Uniform boundedness.} Bounded local cylinder/Wilson representatives obey a refinement-uniform operator bound.
\item \textbf{Local-state limit.} QFT4 has two parts.  QFT4\(_E\) is existence of at least one weak-* bounded-local cluster state and is derived from QFT2--QFT3 by compactness.  QFT4\(_C\) is the stronger assertion that the cluster set in a supplied/selected physical phase is a singleton, equivalently refinement-path-independent canonicity within that phase.
\item \textbf{Reflection positivity.} Every weak-* cluster state of finite reflection-positive presentations is reflection positive; no QFT4\(_C\) uniqueness premise is consumed.
\item \textbf{Native covariance and local flat covariance.} Exact finite covariance follows from an equivariant finite Hamiltonian and its normalized ground projector.  Flat-chart state covariance is derived whenever the representative transport converges in operator norm on the compact comparison window.  Any remaining representative/chart comparison is an analytic export obligation, not an independent vacuum-state selector.
\item \textbf{Regularity.} Uniform bounded cylinder/Wilson test-function estimates pass to every cluster state.  Unbounded distributional fields are handled only by the QFT10 extensions below.  Clustering or purity is required only for claims that need a unique pure vacuum/phase.
\item \textbf{Observable non-Gaussianity and interaction classification when claimed.} A surviving higher connected cumulant certifies non-Gaussianity of the stated source family. For nonlinear Wilson/curvature/composite readouts, a stronger underlying-dynamics interaction claim additionally consumes the same-source free-comparator or actual collision certificate of Corollary~\ref{p12:cor:v2210-lit-interaction-comparator}. QFT8 is not a prerequisite for QFT existence.
\item \textbf{Gap status.} A positive-gap claim consumes a refinement-uniform physical gap certificate; otherwise the phase is typed controlled-gapless or unresolved.  The electromagnetic branch is the controlled-gapless alternative.
\item \textbf{Local-field extension.} QFT10\(_{\rm source}\) denotes source-owned regulator-normalized fields whose mixed moments descend from a uniformly controlled source family.  QFT10\(_{\rm ext}\) denotes independently normalized external composites and is claim-scoped.  Neither is required to reconstruct the bounded local observable net.
\end{enumerate}
Thus bounded cluster-state existence consumes QFT1--QFT3. An exact ordinary
local-QFT export additionally consumes a selected cluster phase, finite
reflection positivity, and the actual locality, full Euclidean covariance,
time-zero and regularity reconstruction certificate of
Definition~\ref{p12:def:v2210-lit-bounded-os-certificate} (or its verified
corrected field-OS replacement).  QFT4\(_C\), QFT8, positive-gap QFT9, QFT10 fields and scattering are stronger or claim-scoped extensions.
\end{definition}




% Source Paper VII-C, original line 4778; inherited proof bodies unchanged.
\subsection{Projective spectral refinement and the remaining interacting boundary}
\label{p12:sec:g7-g8-update}

Let
\[
 V_m=\mathcal P_x(m)=\bigoplus_{k\le m}E_k^x,
 \qquad
 J_{nm}:V_m\hookrightarrow V_n
\]
for $m\le n$.  Define
\[
 \mathcal E_{n\to m}(A)=J_{nm}^*AJ_{nm}.
\]

\begin{theorem}[Canonical projective coarse-graining]
\label{p12:thm:canonical-ucp-coarse-graining}
The maps $\mathcal E_{n\to m}$ are unital completely positive and satisfy
\[
 \mathcal E_{n\to m}\mathcal E_{\ell\to n}=\mathcal E_{\ell\to m}.
\]
For every bounded parent symbol $f$ and
$T_m(f)=P_mM_fP_m|_{V_m}$,
\[
 \boxed{\mathcal E_{n\to m}(T_n(f))=T_m(f).}
\]
\end{theorem}

\begin{definition}[First-order distinguishability quotient]
\label{p12:def:minimal-source-quotient}
At finite level \(L\), let \(\mathfrak P_L\) be a finite-dimensional real
parameter manifold of owned action or Hamiltonian coefficients and let
\[
 \mathscr O_L:\mathfrak P_L\longrightarrow
 \operatorname{Op}^{1}(\mathcal E_L)
\]
be the differentiable map to the real vector space of first-order owned
operator insertions on the finite carrier.  At \(\varphi\in\mathfrak P_L\),
define the distinguishable tangent and its source dual by
\[
 V_{L,\varphi}
 :=T_\varphi\mathfrak P_L/\ker(d\mathscr O_L)_\varphi,
 \qquad
 \mathcal J_{L,\varphi}:=V_{L,\varphi}^*.
\]
Only the quotient classes are first-order observable; a vector in the kernel
is source-invisible at \(\varphi\).
\end{definition}

\begin{theorem}[Canonical minimal first-order source space]
\label{p12:thm:minimal-source-quotient}
There are canonical finite-dimensional identifications
\[
 V_{L,\varphi}\cong\operatorname{im}(d\mathscr O_L)_\varphi,
 \qquad
 \mathcal J_{L,\varphi}
 \cong\bigl(\operatorname{im}(d\mathscr O_L)_\varphi\bigr)^*
 \cong\operatorname{Ann}\!\bigl(\ker(d\mathscr O_L)_\varphi\bigr).
\]
Any linear source parametrization that separates every distinguishable
first-order class has dimension at least
\(\operatorname{rank}(d\mathscr O_L)_\varphi\); the displayed source dual
attains that bound.  If the rank is constant on a parameter stratum, these
spaces assemble into vector bundles on that stratum.  No global
constant-rank assertion is made.
\end{theorem}

\begin{proof}
The first isomorphism theorem gives
\(T_\varphi\mathfrak P_L/\ker d\mathscr O_L\cong\operatorname{im}d\mathscr O_L\).
Dualizing identifies the quotient dual with the annihilator of the kernel.
If a source space \(Q\) separates all quotient classes, the induced map
\(V_{L,\varphi}\to Q^*\) is injective, hence
\(\dim Q\ge\dim V_{L,\varphi}=\operatorname{rank}d\mathscr O_L\).  The
constant-rank theorem supplies the final bundle statement.
\end{proof}

\begin{definition}[Positive projective Schwinger functional]
\label{p12:def:positive-projective-schwinger-functional}
Let \(\mathcal A_L\) be a finite unital \(*\)-algebra, let
\(0\ne\rho_L\in\mathcal A_L^{*+}\), and retain only its positive ray
\([\rho_L]\).  For a real source
\(J\in\mathcal J_{L,\varphi}\), let
\(\Psi_L(J)\in\mathcal A_L^{\rm sa}\) be the associated linear source
insertion and set
\[
 Z_L(J;\rho_L):=\rho_L\!\left(e^{\Psi_L(J)}\right)>0,
 \qquad
 \widehat Z_L(J;[\rho_L])
 :=\frac{Z_L(J;\rho_L)}{\rho_L(\mathbf1)}.
\]
Because \(Z_L(J;c\rho_L)=cZ_L(J;\rho_L)\) for \(c>0\), the unnormalized
functional is a section of the dual tautological line over the projective
positive cone, while \(\widehat Z_L\) is a well-defined scalar on the ray.
The connected generator
\[
 W_L(J;[\rho_L]):=\log\widehat Z_L(J;[\rho_L])
\]
is well defined for every real source, and analytically near \(J=0\) after
complexification.  Its derivatives at zero are the connected source
cumulants.  This definition uses \(\rho_L(A)\), not an unspecified trace.
\end{definition}

\begin{theorem}[Generated finite Schwinger and effective-action closure]
\label{p12:thm:generated-schwinger-effective-action-closure}
Assume the finite level $L$ carries the generated positive state/action ray of
Paper~IV together with the minimal source quotient of
Theorem~\ref{p12:thm:minimal-source-quotient}.  Then
$\widehat Z_L$ and $W_L=\log\widehat Z_L$ of
Definition~\ref{p12:def:positive-projective-schwinger-functional} are
descendants of those data and require no additional normalization selector.
If $W_L$ is proper lower-semicontinuous and convex on the real source quotient,
define its convex conjugate
\[
 \boxed{
 \Gamma_L(\varphi)
 =\sup_J\{\langle J,\varphi\rangle-W_L(J)\}.}
\]
This $\Gamma_L$ is the generated finite effective action on the expectation
coordinate dual to the owned source space.  On a differentiable strictly
convex branch, $\varphi=dW_L(J)$ is locally invertible and
\[
 d\Gamma_L(\varphi)=J;
\]
hence the zero-source stationary condition is $d\Gamma_L=0$.
\end{theorem}

\begin{proof}
Paper~IV fixes the finite positive ray and the exponential weight on its stated
branch.  Projective normalization makes $\widehat Z_L$ independent of the ray
representative, and the logarithm is the continuous map converting products of
independent normalized generating functionals into sums of connected
generators.  Convex conjugation is canonical for the generated source/expectation
pairing.  The derivative identity is finite-dimensional Legendre duality on
the strict-convexity domain.
\end{proof}

\begin{remark}[Effective-action constructor firewall]
\label{p12:rem:effective-action-constructor-firewall}
The theorem does not assert that every sector already owns the required
positive reference measure, integrability domain, chiral phase, gravity
quotient/contour, or refinement limit.  It says that once those upstream data
are actually generated, $Z$, $W$ and $\Gamma$ are descendants rather than new
physical members.  In particular a generated bare action is not by itself a
proof of a broken vacuum.
\end{remark}

\begin{theorem}[Canonical directed-limit descendant]
\label{p12:thm:canonical-limit-descendant}
Let $(I,\preceq)$ be the generated directed comparison/refinement system for a
claim, and let $O_i$ be the generated readout in a Hausdorff target $Y$ at
level $i\in I$.  Suppose there is $O_*\in Y$ such that every licensed cofinal
refinement route (equivalently the full directed net) converges to the same
limit $O_*$.  Then $O_*$ is the unique generated endpoint descendant of the
finite/readout family and requires no additional endpoint selector.

The theorem does not prove convergence, path independence, or the hypotheses
that make a candidate route licensed; those remain analytic/domain
obligations.
\end{theorem}

\begin{proof}
A convergent net in a Hausdorff target has a unique limit.  Because the net
and its licensed comparison maps are generated descendants, that unique limit
is the only compatible endpoint value once convergence and route independence
are established.  An extra selector would therefore duplicate an already
unique descendant.
\end{proof}

\begin{corollary}[Certified finite-precision closure]
\label{p12:cor:certified-finite-precision-closure}
Let $O$ be an exact target quantity and $O_i$ a generated finite/readout
approximation in a metric target.  If the framework proves an error certificate
\[
 d(O,O_i)\le\varepsilon_i,
\]
then any claim whose declared tolerance is $\delta>0$ is certified at level
$i$ whenever $\varepsilon_i<\delta$.  This is closure of the stated
finite-precision claim only; it does not assert existence of an exact
continuum object when that stronger object is itself part of the claim.
\end{corollary}

\begin{definition}[Asymptotically natural state/action ray]
\label{p12:def:asymptotically-natural-state-action-ray}
Let \(\mathcal E_{M\to L}:\mathcal A_M\to\mathcal A_L\) be the canonical UCP
coarse maps and put \(\omega_L=\rho_L/\rho_L(\mathbf1)\).  A family of
positive rays \([\rho_L]\) is \emph{asymptotically natural} on a declared
bounded local cylinder family if, for every compatible bounded representative
\(A_M\),
\begin{equation}
\label{p12:eq:asymptotically-natural-ray}
 \sup_{M\succeq L}
 \left|\omega_M(A_M)
 -\omega_L\!\left(\mathcal E_{M\to L}(A_M)\right)\right|
 \longrightarrow0
\end{equation}
cofinally in \(L\), with the same condition for the declared bounded
finite-time cylinder/semigroup insertions.  On an action-first presentation,
\[
 \rho_L(A)=\int A e^{-S_L}\,d\mu_{0,L},
\]
the replacement \(S_L\mapsto S_L+c_L\) rescales \(\rho_L\) and hence
represents the same state/action ray.  Thus additive vacuum constants are
projective gauge, whereas condition
\eqref{p12:eq:asymptotically-natural-ray} is an analytic property of the
actual interacting sequence.
\end{definition}

\begin{theorem}[Projective form of bounded QFT4]
\label{p12:thm:projective-state-action-ray-qft4}
An asymptotically natural state/action ray defines a unique normalized,
positive, refinement-path-independent state on the bounded local cylinder
algebra.  For every bounded finite source family whose exponentials remain in
that cylinder class, the projective Schwinger functionals \(\widehat Z_L\)
are Cauchy and have the same path-independent limit.  Conversely, bounded
QFT4 Cauchy compatibility on the declared cylinder family implies
\eqref{p12:eq:asymptotically-natural-ray} after passage to a common
refinement.  Hence asymptotic naturality is the projective statement of
bounded QFT4, not an automatic consequence of the existence of UCP maps.
\end{theorem}

\begin{proof}
For a compatible representative \(A_L\), define
\(\omega(A)=\lim_L\omega_L(A_L)\).  Equation
\eqref{p12:eq:asymptotically-natural-ray} and a common-refinement comparison
make the limit Cauchy and independent of the representative and refinement
path.  Positivity and \(\omega(\mathbf1)=1\) pass to the limit.  Applying the
same argument to \(e^{\Psi_L(J)}\) gives convergence of
\(\widehat Z_L\).  Conversely, compare levels \(L\preceq M\) through a
common finer level; the bounded QFT4 Cauchy property, applied once to each
representative, gives the displayed naturality defect.  Rescaling \(\rho_L\)
or adding a constant to \(S_L\) cancels in every normalized expression.
\end{proof}

\begin{remark}[Irreducible occurrence boundary]
\label{p12:rem:projective-ray-occurrence-boundary}
Projectivization removes only normalization.  It does not construct an
asymptotically natural ray, prove its occurrence, or force the covariance,
gap, locality, or residue estimates needed elsewhere in QFT1--QFT10.  For the
actual interacting Born--Hopf sequence, QFT4 remains unverified until
\eqref{p12:eq:asymptotically-natural-ray} or an equivalent defect estimate is
proved.
\end{remark}

This does not contradict the terminal-shell no-go for the stronger fine
intertwiner $J_{nm}T_m(f)-T_n(f)J_{nm}$: projective coarse-graining and
inductive fine intertwining are different requirements.  Consequently the
bounded parent-symbol operator comparison is not an independent QFT4
assumption in the projective formulation.  What remains is compatibility of
the \emph{interacting state and semigroup} with these already generated maps.

For the gap branch, the finite statement
\[
 \gamma_L\ge\gamma_*>0
\]
is exactly equivalent to the Poincare/form inequality
\[
 H_L-E_{0,L}\ge\gamma_*(1-P_{0,L}).
\]
The exact geometric Hodge gap does not supply this physical comparison without
an additional uniform form bound.  Dynamical locality is independent and must
be proved through the Paper-IV DLOC/weighted-conjugation or equivalent
finite-propagation criterion.  Gap plus locality yields pair exponential
mixing, but the existing Bernoulli counterexample still proves that all-order
connected-tree domination is a separate theorem.

Uniform residue normalization is likewise independent of the gap/cluster
package and must come from canonical owner/source normalization plus a uniform
elliptic/Sobolev estimate on the residue compiler.  Mixed-moment Cauchy
uniqueness remains an interacting cross-level state theorem.

Finally, if topological blocks of the regulated physical measure are mutually
orthogonal and nonnegative and the populated charge lattice has gcd $d$, then
nonnegative reweighting by $e^{i\nu\bar\theta}$ forces
\[
 \bar\theta\in\frac{2\pi}{d}\ZZ.
\]
The primitive case $d=1$ gives $\bar\theta=0$.  This is a conditional theorem
about an already generated positive sector measure; it does not generate the
sector weights or primitive-sector occurrence.



% BEGIN SC v1.80 SPATIAL-TRACE-LIMIT

% Source Paper VII-C, original line 5061; inherited proof bodies unchanged.
\subsubsection{Large-radius spatial-transport trace and fourth-power free readout}
\label{p12:sec:v180-spatial-transport-trace-limit}

\begin{proposition}[Scaled spatial transport and vanishing commutators]
\label{p12:prop:scaled-spatial-transport-commutators}
On a complete sharp position band define
\[
 \Pi_i^{(a)}=\frac{\hbar}{a}Q_i.
\]
Then exactly at finite level
\[
 [\Pi_i^{(a)},\Pi_j^{(a)}]
 =\frac{2i\hbar}{a}\epsilon_{ijk}\Pi_k^{(a)},\qquad
 \sum_i(\Pi_i^{(a)})^2=\hbar^2\Delta_a,
 \qquad \|\Pi_i^{(a)}\|=\frac{\hbar m_x}{a}.
\]
Along the existing complete sharp-frontier comparison with fixed
$\Lambda_F$ and $m_x/a\to\Lambda_F$, the norms remain bounded by
$p_F=\hbar\Lambda_F$ while the commutator norms tend to zero.
\end{proposition}

\begin{remark}[Cotangent meaning of the scaled spatial-transport triple]
\label{p12:rem:cotangent-meaning-spatial-transport}
Paper~II-D,
Corollary~\ref{p2d:cor:finite-spatial-momentum-principal-symbol}, identifies
these same operators by
\[
 \Pi_i^{(a)}=\frac{\hbar}{a}Q_i=-i\hbar e_i,
 \qquad
 \sigma_{\rm phys}(\Pi_i^{(a)})(x,p)=p(e_i(x)).
\]
Paper~III-A and Paper~III-F identify $p\in T_x^*U$ as the localized physical
reciprocal generator role.  Hence the commuting variables $p_i$ below are the
continuum coordinates of that cotangent role.  They are not obtained by
raising the microscopic compact momentum label: the compact carrier remains
$m_p=1$, and there is no finite-level operator equality
$\Pi_i^{(a)}=\widehat P_i$ on a common Hilbert space.
\end{remark}

\begin{theorem}[Momentum-ball trace limit]
\label{p12:thm:spatial-transport-momentum-ball-trace}
For each fixed noncommutative polynomial $F$,
\[
 \boxed{
 \lim_{a\to\infty}\frac1{V_a}\operatorname{Tr}
 F(\Pi_1^{(a)},\Pi_2^{(a)},\Pi_3^{(a)})
 =\int_{|p|\le\hbar\Lambda_F}
 \frac{d^3p}{(2\pi\hbar)^3}F_{\rm comm}(p).}
\]
The result is a trace-moment comparison of finite represented operators.  The
commuting integration variables $p_i$ are therefore the continuum coordinates
of the cotangent momentum role described in
Remark~\ref{p12:rem:cotangent-meaning-spatial-transport}.  The result is not a
finite-level joint projection-valued measure, a continuum CCR construction,
or a common-Hilbert-space norm limit, and it does not define an
$m_p\to\infty$ flow.
\end{theorem}

\begin{corollary}[Free spectral density and fourth-power counting functional]
\label{p12:cor:free-spectral-density-at-fixed-cutoff}
On a compatible family of the already admitted free scalar blocks,
\[
 \lim_{a\to\infty}\frac1{V_a}\operatorname{Tr}f(H_{\rm free}^{(a)})
 =\int_{|p|\le p_F}\frac{d^3p}{(2\pi\hbar)^3}
 f\!\left(\sqrt{c^2|p|^2+\mu^2c^4}\right).
\]
In the massless sector the energy-valued counting trace is therefore
\[
 \boxed{
 \mathscr E_{\rm count}
 =\frac{\hbar c\,\Lambda_F^4}{8\pi^2}}
\]
(or $\Lambda_F^4/(8\pi^2)$ in natural units).  The same position geometry
therefore supplies a native fourth-degree spectral readout, matching the
fourth-power degree of the local scalar action without identifying the two
functionals.  In particular no theorem here sets
$\mathscr E_{\rm count}=\Delta V_H$ or
$\hbar c\Lambda_F=E_{32}$.
\end{corollary}

\begin{remark}[Finite comparison band versus continuum frontier]
The value $m_{\rm cmp}=2$ belongs to the exact finite scalar action comparison.
The large-radius limit instead uses the physical complete frontier
$m_x(a)\to\infty$ with $m_x(a)/a\to\Lambda_F$.  These roles are compatible
precisely because the comparison band is not a truncation of position growth.
By Corollary~\ref{p12:cor:v259-resolution-growth-no-internal-species}, this
large-$m_x$ refinement also cannot by itself generate additional color,
family, gauge-factor, anomaly, or particle species; it refines the
reciprocal-position presentation of the already inherited finite internal
structure.  The interacting QFT, local-state, gap, and measure certificates
retain their previous statuses.
\end{remark}

% BEGIN v2.60 FIXED-CUTOFF / TRACE-STABILITY CONSEQUENCES
\begin{theorem}[Large $m_x$ at fixed frontier scale is volume refinement, not UV-cutoff removal]
\label{p12:thm:v260-fixed-cutoff-not-rg}
Along the complete sharp-frontier large-radius family with fixed $\Lambda_F$,
\[
 \Pi_i^{(a)}=\frac{\hbar}{a}Q_i,
 \qquad
 \|\Pi_i^{(a)}\|=\frac{\hbar m_x}{a},
 \qquad
 \frac{m_x}{a}\longrightarrow\Lambda_F.
\]
Hence the physical momentum ceiling tends to the fixed value
$p_F=\hbar\Lambda_F$ while the number of reciprocal-position modes grows with
spatial volume.  Therefore
\[
 \boxed{
 m_x\to\infty\ \text{at fixed }\Lambda_F
 \quad\text{is large-volume/spatial refinement at fixed UV ceiling,}}
\]
not $\Lambda_F\to\infty$, not cutoff removal, and not by itself an RG flow.
No beta function, asymptotic-safety trajectory, or G39 certificate may be
inferred from $m_x$ growth alone.
\end{theorem}

\begin{corollary}[Resolution stability of purely internal trace identities]
\label{p12:cor:v260-internal-trace-resolution-stability}
On a downstream sector of the explicitly factorized form
\[
 \mathcal P_x(m)\otimes V_{\rm int}
\]
with fixed inherited finite internal carrier $V_{\rm int}$, every purely
internal operator $A$ satisfies
\[
 \boxed{
 \operatorname{Tr}_{\mathcal P_x(m)\otimes V_{\rm int}}(I\otimes A)
 =D_m\operatorname{Tr}_{V_{\rm int}}A,
 \qquad
 D_m=\dim\mathcal P_x(m).}
\]
Consequently any exact internal trace/anomaly coefficient that is zero remains
zero at every reciprocal-position resolution, and a nonzero one cannot be
``repaired'' merely by increasing $m_x$.  More strongly, with the canonical
normalized trace
\[
 \tau_{\mathcal P_x(m)\otimes V_{\rm int}}(B)
 :=\frac{\operatorname{Tr}(B)}{D_m\dim V_{\rm int}},
\]
every purely internal factorized observable satisfies the exact identity
\[
 \boxed{
 \tau_{\mathcal P_x(m)\otimes V_{\rm int}}(I\otimes A)
 =\frac{\operatorname{Tr}_{V_{\rm int}}A}{\dim V_{\rm int}}
 =\tau_{V_{\rm int}}(A).}
\]
Thus normalized purely-internal observables are exactly $m_x$-independent;
reciprocal-position refinement alone cannot generate their running.  Mixed
$x$--internal operators, nonfactorized states/measures, extensive quantities,
or genuine RG flow are outside this statement.  Global chiral-line and unrestricted topology-comparison obligations, when consumed by a continuum claim, remain under the G39/chiral-continuum certificate; native G20 is retired as a separate physical-owner gap.
\end{corollary}

\begin{theorem}[Exact finite frontier count converges to the momentum-ball Weyl density]
\label{p12:thm:v260-frontier-count-weyl-density}
At the first-entry event $\tau_k$,
\[
 D_k=\frac{(k+1)(k+2)(2k+3)}6,
 \qquad
 V_k=2\pi^2a_k^3
 =\frac{2\pi^2[k(k+2)]^{3/2}}{\Lambda_F^3}.
\]
Therefore
\[
 \boxed{
 \frac{D_k}{V_k}
 =\frac{\Lambda_F^3}{12\pi^2}
 \frac{(k+1)(k+2)(2k+3)}{[k(k+2)]^{3/2}},}
\]
and
\[
 \boxed{
 \lim_{k\to\infty}\frac{D_k}{V_k}
 =\frac{\Lambda_F^3}{6\pi^2}
 =\frac{\operatorname{Vol}(B^3_{\Lambda_F})}{(2\pi)^3}.}
\]
Thus the leading scalar density-of-states normalization in
Theorem~\ref{p12:thm:spatial-transport-momentum-ball-trace} is already forced
by exact finite sharp-frontier harmonic counting.  The continuum theorem adds
the operator/moment limit; it does not introduce this density coefficient as
an independent normalization.
\end{theorem}

\begin{theorem}[Directional Hopf marginal of the finite frontier]
\label{p12:thm:v261-hopf-directional-momentum-ball-marginal}
Let \(d_k(q)\) be the exact Hopf-charge multiplicity on the complete scalar
frontier \(\mathcal P_x(k)\).  Then for \(q/k\to x\in[-1,1]\),
\[
 d_k(q)=\frac{k^2-q^2}{4}+O(k),
 \qquad
 D_k=\frac{k^3}{3}+O(k^2),
\]
so the normalized charge measure converges to
\[
 \boxed{f(x)\,dx=\frac34(1-x^2)\,dx,\qquad -1\le x\le1.}
\]
Its moments are
\[
 \boxed{
 \int_{-1}^{1}x^{2n}f(x)\,dx
 =\frac{3}{(2n+1)(2n+3)},
 \qquad
 \int_{-1}^{1}x^{2n+1}f(x)\,dx=0.}
\]
This is exactly the one-coordinate marginal of uniform volume on the unit
three-ball.  At the first-entry event \(\tau_k\), with
\(p_F=\hbar\Lambda_F\),
\[
 \boxed{
 \tau_k\!\left[
 \left(\frac{\Pi_i^{(a_k)}}{p_F}\right)^2\right]
 =\frac{k+3}{5(k+2)}\longrightarrow\frac15.}
\]
Thus the finite representation multiplicities already encode a nontrivial
directional statistic of the continuum momentum ball, in addition to the
total Weyl-density coefficient of
Theorem~\ref{p12:thm:v260-frontier-count-weyl-density}.
\end{theorem}

\begin{proof}
Paper~II-D, Proposition~\ref{p2d:prop:complete-position-hopf-multiplicities},
gives the exact finite multiplicity polynomial.  Expanding it with
\(q/k\to x\) and normalizing by \(D_k\) gives the density.  The moment formula
follows by elementary integration.  The finite second-moment identity is the
chiral-isotropy specialization
\(\tau_k(Q_i^2)=k(k+3)/5\) together with
\(a_k^2\Lambda_F^2=k(k+2)\).
\end{proof}

\begin{remark}[Two inequivalent classicalization modes]
\label{p12:rem:v261-coordinate-versus-transport-classicalization}
The scaled transport generators of
Proposition~\ref{p12:prop:scaled-spatial-transport-commutators} satisfy
\[
 \|[\Pi_i^{(a)},\Pi_j^{(a)}]\|\longrightarrow0
\]
along the fixed-\(\Lambda_F\) large-radius family.  The compressed scalar
coordinate multipliers of Paper~I-C behave differently.  For
\(C_a=P_kM_{x_a}P_k\),
\[
 \sum_aC_a^2
 =I-\frac{k+2}{2(k+1)}\Pi_k,
 \qquad
 [C_a,C_b]=\frac1{2(k+1)}\mathcal L_{ab}\Pi_k,
\]
and their strong/normalized-trace defects vanish, but
\[
 \boxed{
 \left\|I-\sum_aC_a^2\right\|\longrightarrow\frac12,
 \qquad
 \|[C_a,C_b]\|\longrightarrow\frac12.}
\]
Hence transport/momentum classicalization occurs in operator norm.
By contrast, compressed-coordinate classicalization is only strong and
tracial at the UV edge. The momentum-ball theorem is therefore not a
common-Hilbert-space operator-norm commuting-coordinate theorem.
\end{remark}

% END v2.60 FIXED-CUTOFF / TRACE-STABILITY CONSEQUENCES

% END SC v1.80 SPATIAL-TRACE-LIMIT



% Source Paper VII-C, original line 6806; inherited proof bodies unchanged.
\subsection{Two-site reflection-positive state/action counterfamily}
\label{p12:app:two-site-reflection-positive-no-go}

This appendix isolates the negative result used by the projective-state
boundary.  Let \(x_+\) and \(x_-\) be two real sites exchanged by Euclidean
reflection, fix \(m>0\), and for \(g\ge0\) put
\[
 h_{m,g}(x)
 :=\exp\!\left[-\frac{1+m^2}{2}x^2-gx^4\right],
 \qquad
 d\mu_{m,g}
 :=Z_{m,g}^{-1}h_{m,g}(x_+)h_{m,g}(x_-)
   e^{x_+x_-}\,dx_+dx_-.
\]
The strict mass term makes \(Z_{m,g}<\infty\), including at \(g=0\).

\begin{proposition}[Reflection positivity does not select an action ray]
\label{p12:prop:two-site-reflection-positive-no-go}
Every \(\mu_{m,g}\), \(g\ge0\), is reflection positive.  Nevertheless the
normalized state and action ray vary nontrivially with \(g\).  Hence finite
reflection positivity, positivity of the state, and removal of the partition-
function normalization do not select a unique interacting state/action ray.
\end{proposition}

\begin{proof}
For a bounded positive-time test \(F=F(x_+)\), the series is dominated by
the integrable kernel with \(e^{|x_+x_-|}\) because \(m>0\).  Dominated
convergence and
\(e^{x_+x_-}=\sum_{n\ge0}(x_+x_-)^n/n!\) therefore give
\[
 \int\overline{F(x_-)}F(x_+)\,d\mu_{m,g}
 =\frac1{Z_{m,g}}
  \sum_{n=0}^{\infty}\frac1{n!}
  \left|\int_{\RR}F(x)h_{m,g}(x)x^n\,dx\right|^2
 \ge0.
\]
Thus every member is reflection positive.  On the other hand, with
\(R=x_+^4+x_-^4\), differentiation under the integrable density yields
\[
 \frac d{dg}\,\mathbb E_{m,g}[R]
 =-\operatorname{Var}_{m,g}(R)<0.
\]
The variance is strict because the density has full support and \(R\) is not
constant.  Therefore the normalized fourth moment, and hence the normalized
positive functional, changes with \(g\).  Since
\(S_{m,g}-S_{m,g'}=(g-g')R\) is not a constant for \(g\ne g'\), the action
rays are distinct as well.
\end{proof}

\begin{remark}[Placement and scope]
The counterfamily is appendix material because it is a no-go/minimality
witness.  It does not alter the positive main-chain result that an
asymptotically natural ray, once verified, is exactly a bounded projective
QFT4 datum.  It proves only that finite positivity cannot perform that
selection.
\end{remark}



\section{Bounded local states, reflection positivity and OS reconstruction}
\label{p12:sec:v237-c3}
Conditional local QFT reconstruction consumes consistent bounded cylinders, reflection positivity, locality, regularity and the stated analytic certificates. Finite covariance and source maps must be verified in the same state.

% Source Paper VII-C, original line 595; inherited proof bodies unchanged.
\subsection{Full-framework sector imports and QFT compatibility}
\label{p12:sec:full-framework-sector-compatibility}


% Source Paper VII-C, original line 598; inherited proof bodies unchanged.
\subsubsection{Sector-import interface}
\label{p12:sec:sector-import-interface}

Standard--Model and QCD sector physics enter Paper~VII-C as finite sector packages constructed in Papers~V-A--V-C and Papers~VII-A--VII-B.  The role of the present paper is to certify a phase-selected bounded-local QFT presentation first and to attach stronger field/gap/scattering certificates only when claimed.  Thus the base full-framework statement is sector-by-sector:
\[
\begin{gathered}
\text{finite sector package from Papers~IV--V-C and VII-A--VII-B}
+\text{QFT1--QFT10}
\\
\Longrightarrow
\text{continuum-QFT-compatible sector}.
\end{gathered}
\]

\begin{definition}[Imported finite sector package]
\label{p12:def:imported-sector-package}

\begin{center}
\begingroup
\small
\begin{tabularx}{\linewidth}{@{}
>{\raggedright\arraybackslash}p{0.55\linewidth}
>{\raggedright\arraybackslash}X
@{}}
\toprule
\textbf{Sector-interface claim} & \textbf{Discharged by} \\
\midrule
Imported finite sector package
& This definition. \\

Sector-import QFT compatibility
& Theorem~\ref{p12:thm:sector-import-qft-compatibility}. \\

Controlled gapless electromagnetic branch
& Proposition~\ref{p12:prop:controlled-gapless-em-branch}. \\

Graded sector compatibility
& Proposition~\ref{p12:prop:graded-chiral-sector-compatibility}. \\

Finite identities survive reconstruction
& Theorem~\ref{p12:thm:finite-identity-preservation}. \\

Renormalized data are comparison outputs
& Proposition~\ref{p12:prop:renormalized-sector-data-comparison}. \\

Scattering is outside the local certificate
& Proposition~\ref{p12:prop:scattering-nonclaim} and
Theorem~\ref{p12:thm:scattering-hypothesis-firewall}. \\
\bottomrule
\end{tabularx}

\captionsetup{hypcap=false}
\captionof{table}[Sector imports and QFT compatibility]{Sector imports and QFT compatibility.}
\captionsetup{hypcap=true}
\label{p12:tab:sector-import-compatibility}
\label{p12:tab:qft-sector-compatibility-ledger}

\TableNotes{The imported Standard-Model, electromagnetic, fermionic, and QCD
structures enter only as finite sector packages. Paper~VII-C certifies their
local-QFT compatibility and marks scattering, collider, and detector questions
as later comparison layers.}
\endgroup
\end{center}

An imported finite sector package is a family
\[
\mathfrak S=
(\mathcal H_L^{\mathfrak S},
 \mathcal A_L^{\mathfrak S},
 \omega_L^{\mathfrak S},
 T_L^{\mathfrak S},
 \Sigma_{\mathfrak S})_{L\in\RBH}
\]
constructed in Papers~IV--V-C and VII-A--VII-B.  Here \(\Sigma_{\mathfrak S}\) denotes the finite selector data of the sector: internal charges, gauge-invariant subalgebra, chirality or grading, mass/gap or controlled-gapless status, and the finite local residue observables used for field reconstruction.  The package is \emph{base-QFT-admissible} when a supplied/selected weak-*
cluster phase satisfies the bounded-local conditions together with the actual
ordinary-net reconstruction certificate of
Definition~\ref{p12:def:v2210-lit-bounded-os-certificate} or a verified
corrected field-OS replacement. Cluster compactness alone is not that
ordinary-net certificate.  QFT4$_C$, QFT8, positive-gap QFT9, QFT10 fields and scattering are attached only when their stronger claims are made.
\end{definition}


% Source Paper VII-C, original line 679; inherited proof bodies unchanged.
\subsubsection{Continuum compatibility theorem}
\label{p12:sec:sector-compatibility-theorem}

\begin{theorem}[Sector-import QFT compatibility theorem]
\label{p12:thm:sector-import-qft-compatibility}
Let \(\mathfrak S\) be an imported finite sector package from Papers~IV--V-C and VII-A--VII-B.  If \(\mathfrak S\) is base-QFT-admissible on a supplied/selected cluster phase \(\alpha\), then Paper~VII-C reconstructs the ordinary local-chart OS/Wightman presentation of that phase:
\[
\mathfrak S
\longmapsto
(\mathcal H_{\mathfrak S},
 \Omega_{\mathfrak S},
 H_{\mathfrak S},
 \mathcal A_{\mathfrak S}(\mathcal O),
 U_{\mathfrak S},
 \Phi_{\mathfrak S}).
\]
No Standard-Model or QCD-specific datum is selected, tuned, or rederived in Paper~VII-C.  Those data remain the finite sector data constructed in Papers~V-A--V-C and Papers~VII-A--VII-B.  Paper~VII-C supplies only the continuum-QFT compatibility certificate.
\end{theorem}

\begin{proof}
QFT1--QFT3 give the local cylinder algebra and uniformly bounded finite representatives.  Weak-* compactness supplies the selected cluster state.  Finite reflection positivity and bounded regularity pass to that state, while Theorem~\ref{p12:thm:native-bh-covariance} reduces flat-state covariance to the representative/chart comparison.  Theorem~\ref{p12:thm:main-qft} then gives the bounded local OS/Wightman export.  QFT4$_C$ only upgrades that export to same-phase path canonicity; QFT8, QFT9, QFT10 and scattering attach their respective stronger claims.  Since the internal carrier, selectors, charges, matter content, gauge algebra, electroweak stabilizer, anomaly status, color/singlet data and finite transfer data are already part of \(\mathfrak S\), Paper~VII-C introduces no additional sector physics.
\end{proof}


% Source Paper VII-C, original line 702; inherited proof bodies unchanged.
\subsubsection{Controlled gapless and graded sectors}
\label{p12:sec:controlled-gapless-graded-sectors}

\begin{proposition}[Controlled gapless electromagnetic branch]
\label{p12:prop:controlled-gapless-em-branch}
The electromagnetic sector imported from Papers~V-A--V-C is a controlled gapless sector in the sense of QFT9.  Its finite-sector content consists of the scalar-stabilizer line
\[
L_Q=\mathbb R Q,
\qquad
Q=T_L^3+Y,
\]
the residual gauge group \(U(1)_{\rm em}\), the abelian connection
\[
A_{\rm em}\in\Omega^1(M;L_Q),
\]
and the reciprocal-position interface equations
\[
F_{\rm em}=dA_{\rm em},
\qquad
 dF_{\rm em}=0,
\qquad
 dG_{\rm em}=J_Q,
\qquad
 G_{\rm em}=\chi(F_{\rm em}).
\]
In the source-free vacuum branch, gauge invariance removes the longitudinal connection component and leaves the transverse photon polarizations.  Therefore Paper~VII-C does not infer a mass gap, exponential clustering, or massive LSZ particle structure for this sector.  It certifies only the local continuum-QFT compatibility of the already-derived massless electromagnetic branch.
\end{proposition}

\begin{proof}
Papers~V-A--V-C supply the finite electromagnetic stabilizer, charge operator, connection line, transverse photon branch, and massless-shell status.  These data enter Paper~VII-C as the sector package \(\mathfrak S_{\rm em}\).  Since QFT9 allows either a uniform gap or an explicitly specified controlled gapless sector, the electromagnetic branch satisfies the gap-status requirement by the second alternative.

The OS/Wightman reconstruction of local fields, observables, and covariance
uses QFT1--QFT8 and QFT10; it does not require a positive transfer gap.
Exponential clustering requires the appropriate gap hypothesis, while any
massive single-particle scattering claim separately uses
Definition~\ref{p12:def:claim-indexed-smat-package}.  Hence the electromagnetic
sector is compatible with the Paper~VII-C continuum construction as a controlled
gapless branch, while all massive-gap conclusions are explicitly excluded for
it.
\end{proof}

\begin{proposition}[Graded chiral sector compatibility]
\label{p12:prop:graded-chiral-sector-compatibility}
Let \(\mathfrak S\) be an imported finite sector package whose finite carrier contains fermionic or chiral degrees of freedom constructed in Papers~V-A--V-C.  Assume that the finite Fock/CAR construction of Paper~IV applies to this carrier and that the corresponding even observable algebra and smeared sector-changing field residues satisfy QFT1--QFT10, including an independently verified graded spacetime-locality certificate for the selected physical localization and causal evolution.  Then the Paper~VII-C continuum reconstruction is a graded local QFT presentation:
\[
\mathcal A_{\mathfrak S}^{\rm even}(\mathcal O)
\subset
\mathcal A_{\mathfrak S}(\mathcal O),
\]
with physical observables in the even local algebra and fermionic fields realized as graded sector-changing operator-valued distributions.  Chirality, hypercharge and family carriers, together with the actual same-action assembled finite source/mass member and its invariant readouts, remain finite-sector data inherited from Papers~V-A--V-C.  Paper~VII-C reconstructs their local continuum field presentation over those finite sector data.
\end{proposition}

\begin{proof}
The finite Fock/CAR construction supplies a \(\mathbb Z_2\)-graded finite operator algebra.  Its even part is the physical observable algebra, while odd CAR generators are field or sector-changing operators.  The CAR relations alone do not imply regional graded locality: a finite one-particle overlap produces a nonzero cross anticommutator. The present graded locality conclusion instead uses the additional physical localizer/causal-evolution hypothesis in the proposition and the corresponding QFT certificates.

Applying QFT1--QFT10 to the even local representatives reconstructs the continuum observable net.  Applying QFT10 to the normalized odd smeared residues reconstructs the corresponding fermionic fields as distributions affiliated with the graded local net.  Since the chiral projectors, gauge charges and family carriers are generated finite data, while the actual source/mass member is supplied by G0 and transported by G16, the continuum reconstruction preserves these inputs and their invariant readouts rather than acting as a new selector.
\end{proof}


\begin{proposition}[Exact spectral CAR locality scope and obstruction]
\label{p12:prop:v382-fixed-radius-smooth-car-locality}
On a \emph{fixed-radius} round $S^3(R_x)$, let $P_m$ project onto scalar
harmonics of degree at most $m$, and let $f,g$ be fixed smooth one-particle
probes with disjoint spatial supports. Their projected CAR overlap obeys
\[
 \|\{a(P_mf),a^\dagger(P_mg)\}\|
 \le(1+\lambda_{m+1})^{-s}\|f\|_{H^s}\|g\|_{H^s},
 \qquad\lambda_{m+1}=\frac{(m+1)(m+3)}{R_x^2}
\]
for every $s\ge0$. Fixed finite smooth disjoint probe families
therefore become equal-time graded local faster than any fixed power of
$m^{-1}$ in the \emph{fixed-radius} $m\to\infty$ limit.
Nevertheless, at every finite cutoff and for every nonempty open
$O\subset S^3$, one has $P_mC_c^\infty(O)=\operatorname{ran}P_m$;
there is no uniform operator-norm locality claim for all cutoff-dependent
regional probes.
For the physical fixed-UV Born limit
$m/R_x\to\Lambda_F<\infty$, the omitted eigenvalue instead tends to
$\Lambda_F^2$, and the displayed estimate alone proves no vanishing
locality defect. Full spacetime locality retains the action-owned
localization, causal evolution, and QFT state/field-domain requirements.
\end{proposition}
\begin{proof}
The CAR anticommutator is the projected one-particle overlap.
For disjoint probes $\langle P_mf,P_mg\rangle
=-\langle(1-P_m)f,(1-P_m)g\rangle$. Spectral tail estimates give the
bound. For the obstruction, any retained spherical harmonic orthogonal
to every projection of a test supported in $O$ vanishes on $O$; real
analyticity forces it to vanish globally. Thus the regional projected
image spans the full finite mode space. Finally
$\lambda_{m+1}\to\Lambda_F^2$ in the fixed-UV scaling, so the
fixed-radius decay premise cannot be silently reused there.
\end{proof}

% Source Paper VII-C, original line 760; inherited proof bodies unchanged.
\subsubsection{Finite identities under reconstruction}
\label{p12:sec:finite-identity-reconstruction}

\begin{proposition}[Finite quantum Ward antecedent and anomaly firewall]
\label{p12:prop:v248-finite-quantum-ward-antecedent}
Let \(\alpha_s\) be a continuous one-parameter \(*\)-automorphism group of a
finite regulated observable algebra, generated by a constituted symmetry of
the complete finite owner, and let \(\omega_L\) be the finite state used in
the regulated correlation functions.  If the state is invariant,
\[
\omega_L\!\circ\alpha_s=\omega_L,
\]
then for every finite product of differentiable controlled representatives
\(A_{1,L},\ldots,A_{n,L}\), with
\(\delta_L A:=\left.\frac{d}{ds}\alpha_s(A)\right|_{s=0}\),
\begin{equation}
\label{p12:eq:v248-finite-ward}
\boxed{
\sum_{k=1}^n
\omega_L\!\left(
 A_{1,L}\cdots (\delta_LA_{k,L})\cdots A_{n,L}
\right)=0 .}
\end{equation}
Thus a finite Ward identity is internal theorem closure once both the native
symmetry and the invariant finite state/measure realization have been
generated or independently proved.  Classical same-owner action invariance
alone does not imply the second antecedent.  In a measure realization, a
non-invariant regulator or Jacobian contributes the corresponding anomaly or
contact term; such a term must remain explicitly in the finite identity.
\end{proposition}

\begin{proof}
Differentiate
\(
\omega_L(\alpha_s(A_{1,L})\cdots\alpha_s(A_{n,L}))
=\omega_L(A_{1,L}\cdots A_{n,L})
\)
at \(s=0\).  The finite product rule gives
Eq.~\eqref{p12:eq:v248-finite-ward}.  If the regulated realization is written
as a finite measure/change-of-variables formula, failure of exact invariance
appears precisely as the extra Jacobian/regulator/contact contribution, so it
cannot be deleted by the classical Noether identity.
\end{proof}


\begin{remark}[Two finite Ward domains and their reconstruction handoff]
\label{p12:rem:v2106-two-finite-ward-domains}
The preceding proposition concerns products in one fixed invariant state.
Paper~IV, Theorem~\ref{p4:thm:v4106-canonical-stress-ward}, instead gives
an on-equation source-history identity with the state transported by the
actual finite canonical symmetry. Its operator defect and first source
variation must vanish, all source/clock/carrier marks and external forces
must be retained, and the stress is the complete functional Hilbert insertion
of Theorem~\ref{p4:thm:v4106-functional-hilbert-response}.
No invariance of an arbitrary evolving density under every local
diffeomorphism is required for that different identity.
Neither statement retires the other's hypotheses. The guarded free
specialization of Paper~VI-A,
Theorem~\ref{p8:thm:v2149-guarded-free-metric-ward}, transfers stress
conservation through the consumed first-response jet order; it is not
a generic interacting window certificate. Continuum transport still
requires the controlled representatives and limits of the theorem below.
\end{remark}

% BEGIN M8 v1.78 XII_IMPORT
\begin{remark}[Finite/local imports and quantum continuation]
\label{p12:rem:m8-batch-quantum-inputs}
The finite paired Berezin identity and first-order
fermion clock theorem of IV close those finite
operations on their stated owners.  The finite local
bosonic real form of III-F extends local covariant
Euler/Noether identities through a declared finite
jet order.  Neither result alone proves invariance
of an interacting quantum measure, absence of
anomalies, an OS/refinement limit, a local field
reconstruction or a physical pole/scattering theorem.
The claim-complete action contract replaces a generic
missing-action wrapper while the quantum Ward,
boundary and analytic certificates remain explicit.
A line involution and a mathematical conjugate-line
block do not supply the physical global \(J_F\) or
a regulator-compatible Majorana sector.
\end{remark}
% END M8 v1.78 XII_IMPORT



\begin{theorem}[Finite identity preservation under continuum reconstruction]
\label{p12:thm:finite-identity-preservation}
Let \(\mathfrak S\) be a QFT-admissible imported sector package.  Suppose that for each refinement level \(L\) the finite representatives satisfy an exact linear or polynomial local identity
\[
\mathcal I_L(A_{1,L},\ldots,A_{n,L})=0
\]
inside finite correlation functions, where the involved observables or smeared fields have QFT4 limits and, when distributional fields occur, QFT10 residue limits.  Then the corresponding continuum identity
\[
\mathcal I(A_1,\ldots,A_n)=0
\]
holds in the reconstructed OS/Wightman theory.  In particular, finite Ward identities, anomaly-cancellation identities, gauge-invariance identities, and Schwinger--Dyson identities survive the Paper~VII-C continuum reconstruction whenever their terms belong to the controlled QFT1--QFT10 class.  No identity is passed through the limit unless all contact terms, current divergences, regulator residues, and composite-renormalization terms present in the finite identity are included in the controlled family.
\end{theorem}

\begin{proof}
Each term in the finite identity is a finite expectation value or distributional pairing built from representatives whose continuum limits exist by QFT4 and QFT10.  Addition, scalar multiplication, finite products of controlled local representatives, and the required distributional pairings commute with the limit by the boundedness, regularity, and residue-control hypotheses.  Taking the limit of
\[
\mathcal I_L(A_{1,L},\ldots,A_{n,L})=0
\]
therefore gives
\[
\mathcal I(A_1,\ldots,A_n)=0.
\]
If the finite identity is a Ward identity, the limiting identity is the continuum Ward identity including its controlled contact terms.  If the finite identity is anomaly cancellation, no anomalous obstruction appears in the continuum limit beyond the regulator residues already displayed and controlled in the finite identity.  If the finite identity is a Schwinger--Dyson relation, the limiting relation is the reconstructed continuum Schwinger--Dyson identity for the controlled field class, with composite terms interpreted through the QFT10 residue limits.
\end{proof}


% BEGIN V2.254 FINITE SOFT-NO-OMISSION CONTINUUM GATE
\begin{remark}[Static noise requires an additional uniform infrared-limit certificate]
\label{p12:rem:v254-static-noise-limit-gate}
Paper~IV-C, Theorem~\ref{p4c:thm:v254-zero-mode-no-omission},
shows that the complete finite connected in-in response retains a
positive stationary zero-frequency sector, while retarded commutators
alone do not determine it. This is mathematically distinct from the
asymptotic soft omissions discussed by
\cite{JorstadPasterski2026SoftOmissions}, the memory/detector identity of
\cite{MoultNarayananPasterski2026MemoryCorrelators}, and the
infrared-dressed detector construction of
\cite{OertelMoultPasterski2026Detectors}.
For an actual sequence of \emph{same-action} finite detector operators,
write $C_L(T)$ for the symmetric time-averaged connected covariance and
$C_L^{(0)}=\lim_{T\to\infty}C_L(T)$ from IV-C. A sufficient (not
currently established) uniform interchange certificate is
\[
\sup_L|C_L(T)-C_L^{(0)}|\le\eta(T),\quad
\eta(T)\xrightarrow[T\to\infty]{}0,\quad
C_L^{(0)}\xrightarrow[L\to\infty]{}C_\infty^{(0)}.
\]
It entails $\lim_{T\to\infty}\lim_{L\to\infty}C_L(T)=
\lim_{L\to\infty}\lim_{T\to\infty}C_L(T)=C_\infty^{(0)}$
provided the finite-time continuum limit exists. It does
\emph{not} itself construct a null-infinity detector, control the
simultaneous radius/soft-frequency limits, or generate a physical BMS
charge in the terminal curved background. QFT4/QFT10 identity transfer
cannot silently replace these missing infrared and detector certificates.
\end{remark}
% END V2.254 FINITE SOFT-NO-OMISSION CONTINUUM GATE

\begin{corollary}[Current and stress-tensor sum rules under QFT reconstruction]
\label{p12:cor:v264-current-stress-pdf-sum-rule-transport}
Assume the finite QCD source family satisfies the current and forward spin-two
owner identities of Paper~VII-A, Theorem~\ref{p10:thm:v264-twist-two-current-stress-owner},
and that all terms entering those identities satisfy QFT4 and, where composite
or distributional limits occur, QFT10.  Then the reconstructed continuum
current charges and forward stress-tensor identity obey the Paper~VII-B
flavor, baryon, electric-charge, and momentum sum rules of
Theorem~\ref{p11:thm:native-pdf-sum-rules} whenever the corresponding
renormalized light-ray limits exist.

This corollary does not prove support of the nonlocal PDF on \([0,1]\),
positivity after scheme subtraction, Wilson-line/light-ray renormalization,
rapidity control, or small-\(x\)/zero-mode bounds.  Those remain separately
consumed analytic certificates.
\end{corollary}
\begin{proof}
The finite \(n=1\) identity is a Ward/current identity and the forward \(n=2\)
identity is a finite local source/stress identity.  Theorem~\ref{p12:thm:finite-identity-preservation}
passes each identity to the continuum provided every contact, regulator,
anomaly and composite term belongs to the controlled family.  Paper~VII-B
then identifies the corresponding renormalized light-ray Mellin moments with
the current and twist-two matrix elements.  None of these steps supplies the
additional nonlocal spectral or positivity properties listed in the final
paragraph.
\end{proof}


\begin{remark}[Classical Noether versus quantum Ward firewall]
\label{p12:rem:v176-classical-noether-quantum-ward-firewall}
The classical gauge and diffeomorphism Noether identities derived from the
local Paper-III-F \(A+\Phi+g\) action are native same-owner closure on that
classical domain; they do not by themselves establish an exact quantum Ward
identity.  Proposition~\ref{p12:prop:v248-finite-quantum-ward-antecedent}
shows the additional finite quantum antecedent: the regulated state/measure
must realize the symmetry, with every Jacobian, contact term, regulator
residue, and composite-renormalization contribution retained.  The
reconstructed quantum identity is then available only when those finite terms
also satisfy the QFT4/QFT10 hypotheses of
Theorem~\ref{p12:thm:finite-identity-preservation}.  An anomaly is therefore a
quantum obstruction to the stronger Ward claim, not a failure of the
classical theorem on its own domain.
\end{remark}


\begin{corollary}[\texorpdfstring{Finite \(D_4/A_2\) strong selector survives QFT reconstruction}{Finite D4/A2 strong selector survives QFT reconstruction}]
\label{p12:cor:d4-a2-selector-survives-qft}
Let the imported strong-sector package contain the finite color support
constructed in Paper~V-A,
\[
C=\CC^{T_+},
\qquad
\End_0(C)\cong\mathfrak{sl}_3(\CC),
\qquad
\mathfrak{su}(C)\cong\mathfrak{su}(3)_c,
\]
where the trimer \(T_+\) is the axis-removed vertex-figure slice of the
Paper-I-A \(D_4\)/Hurwitz \(24\)-cell and the trace-zero projector algebra on
\(T_+\) carries the \(A_2\) Cartan/root decomposition.  If this imported
strong package satisfies QFT1--QFT10, then the reconstructed Paper-VII-C
continuum pure-glue sector has gauge algebra \(\mathfrak{su}(C)\), Wilson
algebra in finite-dimensional representations of \(SU(C)\), curvature
residue fields valued in \(\mathfrak{su}(C)\), and gauge-invariant local net
obtained from the same finite singlet/Haar selector.  No new continuum
internal group, no extra triality sector, and no position-lattice copy of the
\(D_4\) shell is introduced.
\end{corollary}
\begin{proof}
Paper~V-A proves the finite algebraic chain
\[
D_4\text{/Hurwitz vertex figure}
\longrightarrow
T_+
\longrightarrow
C=\CC^{T_+}
\longrightarrow
\End_0(C)=\mathfrak h_{A_2}^{\CC}\oplus
\bigoplus_{r\ne s}\CC E_{rs},
\]
with \(\dim\End_0(C)=2+6=8\), and then takes the compact real form
\(\mathfrak{su}(C)\).  Paper~VII-B imports this as the finite QCD color
support and separates the resulting \(SU(3)_c\) center grading from the
\(S_3\) diagram triality of \(D_4\).  These are exact finite sector identities
and selector data.  By Theorem~\ref{p12:thm:sector-import-qft-compatibility},
Paper~VII-C reconstructs the continuum OS/Wightman presentation of a
QFT-admissible imported package without reselecting its internal sector data.
By Theorem~\ref{p12:thm:finite-identity-preservation}, the finite polynomial
and representation identities defining \(C\), \(\End_0(C)\), the invariant
trace form, the center grading, and the singlet selector survive inside
controlled continuum correlation functions.  Hence the local continuum gauge
objects are the reconstructed fields and Wilson observables of the already
selected finite group \(SU(C)\).  The \(D_4\) geometry remains the upstream
finite selector origin of \(C\); it is not promoted to a continuum
spacetime lattice or to an additional gauge symmetry.
\end{proof}


% Source Paper VII-C, original line 968; inherited proof bodies unchanged.
\subsubsection{Renormalized comparison data and scattering boundary}
\label{p12:sec:renormalized-data-scattering-boundary}

\begin{proposition}[Renormalized data as finite-scale comparison outputs]
\label{p12:prop:renormalized-sector-data-comparison}
Let \(\mathfrak S\) be an imported finite sector package with finite couplings, masses, charge normalizations, threshold data, or local composite representatives constructed in Papers~V-A--V-C and Papers~VII-A--VII-B.  In the Paper~VII-C continuum presentation, their renormalized values are comparison coordinates of the same finite data along the Born--Hopf refinement system:
\[
\mathcal R_{L'\to L}^{\mathfrak S}:
\mathcal A_{L'}^{\mathfrak S}(\mathcal O_{L'})
\longrightarrow
\mathcal A_L^{\mathfrak S}(\mathcal O_L).
\]
Thus renormalization does not add independent sector parameters.  It records how finite sector representatives are compared between refinement scales and between continuum schemes.
\end{proposition}

\begin{proof}
The finite sector package fixes the structural carrier, selectors, and finite readout data before the continuum limit is taken.  Paper~VII-C introduces comparison maps only to relate representatives at different refinement levels or in different continuum coordinate schemes.  Therefore a renormalized coupling, mass, charge, or composite-operator normalization is not an additional primitive datum; it is the coordinate value assigned to the finite sector observable under a chosen comparison map.  Changing scheme changes the representative coordinate, not the underlying finite sector datum.
\end{proof}

\begin{proposition}[Scattering is not part of the local QFT certificate]
\label{p12:prop:scattering-nonclaim}
The Paper~VII-C reconstruction theorem certifies local OS/Wightman QFT data:
state, vacuum, Hamiltonian, covariance, local net, and reconstructed local
fields.  It does not, merely by reconstruction, assert an \(S\)-matrix, LSZ
formula, asymptotic completeness, scattering cross section, or detector rate.
Those assertions are available only on the claim-indexed domains defined below.
\end{proposition}

\begin{proof}
OS/Wightman reconstruction supplies the listed local objects but contains no
large-time comparison dynamics, asymptotic channel Hilbert space, wave
operator, detector limit, or infrared prescription.  Hence none of the omitted
objects follows from QFT1--QFT10 alone.
\end{proof}

\begin{definition}[Claim-indexed massive/short-range scattering package]
\label{p12:def:claim-indexed-smat-package}
For a declared channel class \(\mathcal C\), write
\(\mathsf{SMAT}_{\mathcal C}\) for the following cumulative package.  A claim
uses only the clauses named for it.
\begin{enumerate}[label=\textnormal{(SM\arabic*)},leftmargin=3.2em]
\item The reconstructed theory on the claimed sector is local or graded-local,
      Poincar\'e covariant, has a unique invariant vacuum, and satisfies the
      spectrum condition.  These are outputs of the local certificate when its
      hypotheses have been discharged, not new scattering axioms.
\item Every declared massive external particle has an isolated finite-multiplicity
      mass hyperboloid, stable under the exact dynamics, and a chosen local
      interpolator with nonzero one-particle residue.  Thresholds are excluded
      from the chosen wave-packet supports.
\item Haag--Ruelle estimates, or a stated Cook/finite-window substitute, hold
      for separated velocity supports and give isometric wave operators
      \begin{equation}
      \label{p12:eq:claim-indexed-wave-operators}
      \Omega_\pm:\mathcal H_{\mathcal C}^{\rm as}\longrightarrow\mathcal H,
      \qquad
      S_{\mathcal C}:=\Omega_+^*\Omega_- .
     \end{equation}
      This clause is required to assert an asymptotic channel
      operator~\cite{Haag1996}.
\item If an LSZ formula is asserted, the connected time-ordered distributions
      have the required boundary values, the external poles are simple with
      fixed nonzero residues, the amputated distributions possess on-shell
      limits, and boundary terms plus interchange of large-time, regulator,
      and on-shell limits are controlled.
\item If a resonance is asserted specifically as an \(S\)-matrix pole, the
      channel amplitudes have the stated threshold boundary values and a unique
      meromorphic continuation to the named sheet.  This is additional to the
      correlator/resolvent continuation used to locate a spectral resonance.
\item Asymptotic completeness---equality of the relevant interacting sector
      with the in/out scattering ranges---is required only when surjectivity,
      full-sector unitarity, or exhaustion by the stated particle channels is
      claimed.  Without it, \(S_{\mathcal C}\) is only the partial channel
      operator on the constructed scattering subspaces.
\item A differential or inclusive cross section additionally requires
      convergent detector-cell/instrument limits, a stable incident-flux
      normalization, and the relevant threshold, phase-space, and
      resolution limits.  An optical-theorem or channel-unitarity identity is
      asserted only on a common asymptotic domain on which the required sums
      over open channels converge.
\end{enumerate}
\end{definition}

\begin{definition}[Massless and long-range infrared alternatives]
\label{p12:def:claim-indexed-infrared-package}
If \(\mathcal C\) contains massless gauge quanta, charged sectors with
long-range fields, soft accumulation, or collinear degeneracy, the ordinary
massive Fock/LSZ clauses are not imported by default.  An \(S\)-matrix-level
claim must instead specify and prove at least one compatible route:
\begin{enumerate}[label=\textnormal{(IR\arabic*)},leftmargin=3em]
\item an inclusive, gauge-invariant, finite-resolution observable together
      with the required sum over degenerate soft/collinear states, uniform
      regulator cancellation, and factorization/renormalization control as the
      auxiliary regulator is removed at fixed nonzero detector resolution;
      any subsequent zero-resolution limit is a separate claim and is not
      presumed to exist;
\item a dressed or modified asymptotic dynamics, including its dressing,
      charged-sector/Gauss-law representation, existence and isometry of the
      modified wave operators, and independence of auxiliary infrared
      regulators for the claimed observable; or
\item an algebraic detector/infraparticle construction with explicit detector
      limits and a proof that the claimed inclusive transition measure exists.
\end{enumerate}
In particular, a positive local-QFT certificate or a massless photon field does
not imply a sharp charged one-particle pole.  Ordinary charged-particle LSZ may
be used only if that pole and its asymptotic construction are separately proved.
\end{definition}

\begin{theorem}[Scattering-hypothesis firewall]
\label{p12:thm:scattering-hypothesis-firewall}
The following local, spectral, and operator claims consume no clause of
\(\mathsf{SMAT}_{\mathcal C}\): OS/Wightman reconstruction, local correlators
and Ward identities, finite-window channel probabilities, Riesz convergence of
isolated stable poles and source residues, a continued correlator/resolvent
resonance pole, the Paper-VII-B light-front fraction operators and measures, and
DGLAP/ERBL evolution of renormalized light-ray operators.

The additional hypotheses attach only as follows.
\begin{enumerate}[label=\textnormal{(\alph*)},leftmargin=2.4em]
\item An asymptotic amplitude or channel \(S_{\mathcal C}\) uses
      \textnormal{(SM1)--(SM3)}; an LSZ representation additionally uses
      \textnormal{(SM4)}.
\item Calling a continued pole an \(S\)-matrix pole uses
      \textnormal{(SM1)--(SM3)} and \textnormal{(SM5)}; interpreting its
      imaginary part as a partial decay rate also uses the applicable
      open-channel normalization and phase-space parts of \textnormal{(SM7)}.
\item Full-sector unitarity or particle-channel exhaustion uses
      \textnormal{(SM6)} in addition to the constructed channel operator.
      Cross sections and detector rates use \textnormal{(SM1)--(SM3)} and
      \textnormal{(SM7)}, plus \textnormal{(SM4)} if an LSZ representation is
      used.
\item Every such claim in a massless/long-range channel additionally uses the
      selected route in
      Definition~\ref{p12:def:claim-indexed-infrared-package}.  No infrared
      package is imposed on a purely local or isolated-spectral claim.
\end{enumerate}
Thus scattering/infrared hypotheses are claim-local and do not enlarge the
premise set of the finite framework or the local continuum certificate.
\end{theorem}

\begin{proof}
Items in the first paragraph are defined by finite spectral calculus,
operator-valued distributions, compact-contour resolvents, or local
reconstruction; none contains a large-time comparison map.  By contrast,
Equation~\eqref{p12:eq:claim-indexed-wave-operators} is the definition of the
asymptotic channel comparison and directly requires \textnormal{(SM1)--(SM3)}.
LSZ is an identification of this limit with amputated on-shell correlators and
therefore requires \textnormal{(SM4)}.  A pole of a continued correlator need
not lie in any constructed scattering amplitude, so the \(S\)-matrix
interpretation needs both the channel construction and \textnormal{(SM5)}.
Surjectivity/completeness and detector normalization are exactly the extra
conclusions of \textnormal{(SM6)} and \textnormal{(SM7)}.  Finally, in a
long-range or massless charged sector the
sharp massive-particle premise can fail; each IR route replaces that premise by
the stated inclusive, dressed, or detector construction.  This proves the
claim-by-claim allocation.
\end{proof}

\begin{corollary}[Claim-local light-ray, optional LFRep, and scattering dependency split]
\label{p12:cor:lf-ray-limit-scattering-split}
The Paper-VII-B packages propagate without recombination:
\begin{enumerate}[label=\textnormal{(\alph*)},leftmargin=2.4em]
\item physical local current/twist-two readouts consume the derived
      \(\LFRayGeom\), the physical null/source owner, and the relevant
      \(\LFState_{h,L}\), but no exact finite \(\LFRep_L\), \(\LFLim\), or
      \(\mathsf{SMAT}_{\mathcal C}\) clause;
\item a renormalized collinear PDF consumes the physical Wilson/light-ray
      operator plus clauses~\textnormal{(L2)--(L4)}, while a TMD also consumes
      clause~\textnormal{(L5)};
\item the stronger finite commuting fraction simplex and its positive
      number/momentum POVMs additionally consume the exact \(\LFRep_L\)
      certificate; a weak limit of those optional positive measures further
      consumes clause~\textnormal{(L1)} and one of Paper~VII-B's Hausdorff or
      bounded-Lipschitz convergence criteria;
\item only a factorized scattering cross section, fragmentation/out-state
      sum, width, or asymptotic amplitude imports the applicable
      \(\mathsf{SMAT}_{\mathcal C}\) and long-range infrared clauses.
\end{enumerate}
Thus exact finite LFRep, continuum light-ray renormalization, and scattering
completion are three differently typed theorem layers.  The finite fraction
simplex is an optional stronger realization and is not the owner of the
physical renormalized PDF.
\end{corollary}

\begin{proof}
Item~(a) is the Paper-VII-A source-owner split and Paper-VII-B current/twist-two
construction.  Item~(b) names exactly the path, color-holonomy, ultraviolet and
Fourier/Mellin clauses required by the physical light-ray definition.  Item~(c)
is the optional exact occupation-space branch retained in the Paper-VII-B
appendix.  Item~(d) names exactly the claims whose definitions use the wave,
detector, or open-channel maps in
Definitions~\ref{p12:def:claim-indexed-smat-package}
and~\ref{p12:def:claim-indexed-infrared-package}.  Theorem~\ref{p12:thm:scattering-hypothesis-firewall}
therefore gives the displayed allocation.
\end{proof}

\begin{corollary}[Paper-IV scattering import]
\label{p12:cor:paper-iv-scattering-import}
Paper~IV's finite-window contraction and leakage identities require none of
\(\mathsf{SMAT}_{\mathcal C}\).  Its conditional large-window M\o ller theorem
supplies clause \textnormal{(SM3)} only on the displayed DLOC, compatibility,
leakage, return, and Cook-tail domain; its detector/flux limit supplies
\textnormal{(SM7)} when that limit is proved.  It does not by itself assert
\textnormal{(SM6)} or an infrared route.
\end{corollary}

\begin{proof}
This is a direct match between the hypotheses and conclusions of
Theorem~\ref{p4:thm:finite-channel-scattering-lsz-interface} and
Theorem~\ref{p4:thm:moller-bridge} and the claim-indexed clauses above.
\end{proof}

\begin{center}
\small
\renewcommand{\arraystretch}{1.15}
\begin{tabularx}{\textwidth}{>{\raggedright\arraybackslash}p{0.30\textwidth}
                            >{\raggedright\arraybackslash}p{0.18\textwidth}
                            >{\raggedright\arraybackslash}X}
\toprule
\textbf{QFT ingredient} & \textbf{Source} & \textbf{Paper VII-C role} \\
\midrule
Finite state/effect/operator calculus & Paper~IV & Imported as the finite quantum layer. \\
Finite Fock/CAR/interface construction & Paper~IV & Imported for bosonic and fermionic finite field representatives. \\
SM chiral gauge--matter carrier & Paper~V-A & Imported sector data; not rederived. \\
Hypercharge, scalar stabilizer, \(U(1)_{\rm em}\) & Paper~V-A & Imported electromagnetic sector; Paper~VII-C only certifies controlled gapless local continuum compatibility. \\
Photon transverse branch & Papers~V-A--V-C & Imported as a massless controlled gapless branch. \\
Finite masses, mixings, Yukawa representatives & Paper~V-B & Imported finite-core readouts; continuum presentation only. \\
Color carrier, singlet projection, strong-sector interface & Papers~V-A, V-C, and X--XI & Paper~V-A supplies structural color, Paper~V-C the interaction/QCD handoff, and Papers~VII-A--VII-B the compiled strong sector and readouts. \\
Pure-glue finite transfer, finite gap sequence, Wilson/Hodge data & Paper~VII-B & Finite import only; interacting nontriviality, area law, and continuum limits are Paper-VII-C obligations; uniform-gap transfer is optional if claimed. \\
Continuum OS/Wightman reconstruction & Paper~VII-C & Proved from QFT1--QFT10. \\
Renormalized comparison maps & Paper~VII-C & Proved as refinement/scheme comparison structure. \\
LSZ, all-order amplitudes, detector rates & Downstream sector analysis & Claimed only with the applicable clauses of Definitions~\ref{p12:def:claim-indexed-smat-package} and~\ref{p12:def:claim-indexed-infrared-package}. \\
\bottomrule
\end{tabularx}
\end{center}


\paragraph{Full-framework QFT compatibility statement.}
Standard--Model and QCD sector physics enter from Papers~V-A--V-C and
Papers~VII-A--VII-B.  Paper~VII-C proves that any such sector package satisfying
QFT1--QFT10 has an ordinary local-chart OS/Wightman QFT presentation.  Massive
sectors inherit their finite transfer gap when QFT9 is realized by a uniform
positive gap.  The electromagnetic sector is instead imported as a controlled
gapless branch: its \(U(1)_{\rm em}\) stabilizer, Maxwell/interface system, and
transverse photon modes are supplied by Papers~V-A--V-C, while Paper~VII-C supplies
only the local continuum compatibility certificate.  Scattering, detector
response, and all-order collider observables are available only through the
claim-indexed firewall of
Theorem~\ref{p12:thm:scattering-hypothesis-firewall}; they are not additional
axioms of the local framework.


% Source Paper VII-C, original line 1218; inherited proof bodies unchanged.
% Former logical part: Euclidean Positivity, OS Reconstruction, and Local States; integrated into the new thematic argument.

\noindent\textbf{Orientation.}
This part proves the finite-to-continuum reconstruction mechanism: finite states and transfer data produce OS-positive local states, a reconstructed Hilbert space, a Hamiltonian, and local algebras.  The Lorentzian and positive-frequency branches remain imports from Papers~III--IV; Paper~VII-C supplies the continuum state/local-net certificate.


% Source Paper VII-C, original line 1223; inherited proof bodies unchanged.
\subsection{Euclidean positivity and OS reconstruction}
\label{p12:sec:os-reconstruction}


% Source Paper VII-C, original line 1226; inherited proof bodies unchanged.
\subsubsection{OS outputs from finite states}
\label{p12:sec:os-reconstruction-ledger}

\begin{table}[htbp]
\centering
\small
\begin{tabularx}{\textwidth}{@{}
>{\raggedright\arraybackslash}p{0.55\textwidth}
>{\raggedright\arraybackslash}X
@{}}
\toprule
\textbf{OS/local-field step} & \textbf{Discharged by} \\
\midrule
Subsequential local states
& Proposition~\ref{p12:prop:subsequential-states}. \\

Spectral-symbol QFT1--QFT3 control
& Proposition~\ref{p12:prop:spectral-symbol-qft123}. \\

Canonical local state
& Theorem~\ref{p12:thm:canonical-local-state}. \\

OS-admissible limit
& Definition~\ref{p12:def:os-admissible-limit}. \\

Native Born--Hopf covariance
& Theorem~\ref{p12:thm:native-bh-covariance}. \\

Bounded-cylinder regularity
& Lemma~\ref{p12:lem:bounded-cylinder-regularity}. \\

Reflection positivity survives the limit
& Lemma~\ref{p12:lem:rp-survives-limit}. \\

Local flat chart recovery
& Lemma~\ref{p12:lem:euc-invariance-survives}. \\

OS reconstruction implication
& Theorem~\ref{p12:thm:bh-os-reconstruction}. \\

Continuum fields are outputs
& Corollary~\ref{p12:cor:fields-are-outputs}. \\
\bottomrule
\end{tabularx}
\caption[OS/local-field certificate ledger]{OS/local-field certificate ledger.}
\label{p12:tab:os-local-field-ledger}
\TableNotes{This ledger separates finite-state compactness, the canonical local-state limit, and OS reconstruction.  Continuum Hilbert space, Hamiltonian, local net, and local fields are outputs of the certificate, not continuum inputs.}
\end{table}

This section first constructs the limiting Euclidean state on local observables once the Cauchy obligation holds, and then reconstructs the Lorentzian quantum theory from the resulting functional.  The reconstruction step requires the Osterwalder--Schrader conditions, with flat covariance recovered only on local reciprocal-position charts.


% Source Paper VII-C, original line 1277; inherited proof bodies unchanged.
\subsubsection{Canonical local state and bounded cylinder control}
\label{p12:sec:canonical-local-state-bounded-control}

\begin{proposition}[Finite boundedness and subsequential states]
\label{p12:prop:subsequential-states}
Assume QFT2--QFT3 on uniformly bounded finite representatives of each
fixed cylinder word. Their evaluation states have weak-* cluster points
on the common asymptotic bounded-word algebra of
Theorem~\ref{p12:thm:v283-qft4-cluster-split}.
Such a state descends to the \emph{intended} local cylinder algebra
only after the exact relation and completion compatibility conditions
of Proposition~\ref{p12:prop:v382-cylinder-relation-descent} are checked.
\end{proposition}
\begin{proof}
Theorem~\ref{p12:thm:v283-qft4-cluster-split} constructs the common
$C^*$ algebra from the finite algebras and obtains cluster states by
weak-* compactness. No unproved multiplication-preserving symbol map is
needed for existence on that constructed algebra. The intended-algebra
identification is a separate quotient statement.
\end{proof}



\begin{proposition}[Finite spectral compression is not multiplicative in norm]
\label{p12:prop:v382-toeplitz-product-obstruction}
For the round-$S^3$ Peter--Weyl compression $T_m(f)=P_mM_fP_m$
and real fundamental character $t(g)=\frac12\operatorname{Tr}g$,
\[
 T_m(t^2)-T_m(t)^2=P_mM_t(1-P_m)M_tP_m\succeq0,
 \qquad\|T_m(t^2)-T_m(t)^2\|\ge\frac14
\]
for every $m\ge1$. Thus bounds on each symbol representative do not by
themselves descend the classical product relations in operator norm.
\end{proposition}
\begin{proof}
The identity is the decomposition of $P_mM_t^2P_m$ across
$P_m+(1-P_m)$. On the normalized degree-$m$ class character $\chi_m$, the exact
$SU(2)$ character product gives
$t\chi_m=(\chi_{m+1}+\chi_{m-1})/2$ (with $\chi_{-1}=0$).
The omitted $\chi_{m+1}$ component has norm $1/2$; its squared
norm supplies the lower bound $1/4$.
\end{proof}

\begin{proposition}[Spectral-symbol regulator and bounded cylinder sector]
\label{p12:prop:spectral-symbol-qft123}
For the Born--Hopf tower imported from Papers~I-A--II and IV, QFT1--QFT3 hold for the bounded local cylinder sector.
\end{proposition}
\begin{proof}
\ProofPlan{We verify QFT1--QFT3 in four steps: density of retained spectral symbols, the local Euclidean cylinder presentation, the finite Toeplitz/cylinder algebra, and uniform boundedness with symbol-level status.}
\begin{proofsteps}
\item \textbf{Peter--Weyl density of retained symbols.}
Paper~II supplies the retained Peter--Weyl spaces
\[
\PW_m=\operatorname{span}\{D^j_{ab}:2j\le m\}\subset C^\infty(S^3).
\]
For every matrix coefficient \(D^j_{ab}\), choosing \(m\ge2j\) places it in \(\PW_m\).  Hence
\[
\bigcup_m\PW_m
=
\operatorname{span}\{D^j_{ab}:j\in\tfrac12\ZZ_{\ge0}\},
\]
whose uniform closure is \(C(S^3)\) by Peter--Weyl.  Local restrictions to coordinate patches are therefore uniformly dense on compact subsets.

\item \textbf{Local Euclidean cylinder presentation.}
Adjoining the scalar Euclidean-time cover of Papers~III--IV gives local product charts \(I\times U\subset\RR\times S^3\), hence the Euclidean \(\RR^4\) cylinder presentation.  This proves QFT1 in the spectral-symbol sense.

\item \textbf{Finite Toeplitz/cylinder algebra.}
For retained symbols, Paper~IV uses finite Toeplitz representatives
\[
\Toe_m(f)=\Pi_mM_f\Pi_m\in\End(\mathcal H_m).
\]
Finite products, adjoints, linear combinations, and finite mode/Fock representatives give the local cylinder algebra, proving QFT2.

\item \textbf{Uniform boundedness and symbol-level status.}
If \(f\) is bounded, then
\[
\|\Toe_m(f)\|\le\|\Pi_m\|^2\|M_f\|=\|f\|_\infty.
\]
Finite cylinder polynomials are therefore uniformly bounded by a polynomial in the generator sup norms.  This proves QFT3 for the bounded cylinder sector.  The conclusion is deliberately symbol-level: the finite matrix algebras are not being identified with \(C(S^3)\) as commutative algebras.
\end{proofsteps}
\end{proof}

\begin{theorem}[Canonical local state from the Cauchy obligation]
\label{p12:thm:canonical-local-state}
Assume QFT2--QFT4.  Then there is a unique continuum local state
\[
\omega_{\mathcal O}:\A_0(\mathcal O)\to\CC
\]
for every bounded open \(\mathcal O\subset\RR^4\), given by
\begin{equation}
\label{p12:eq:omega-local-limit}
\omega_{\mathcal O}(A)=\lim_{L}\omega_L(A_L).
\end{equation}
These local states are compatible under inclusions and therefore define a state \(\omega\) on the algebraic inductive local algebra
\[
\A_{\loc}:=\bigcup_{\mathcal O\Subset\RR^4}\A_0(\mathcal O)/\mathcal N_{\mathcal O}.
\]
\end{theorem}
\begin{proof}
Existence and refinement-path independence are exactly QFT4.  Positivity follows because each finite \(\omega_L\) is positive and the representative of \(A^*A\) is positive up to the null quotient.  Normalization follows from \(\omega_L(\mathbf1)=1\).  Compatibility under inclusion follows from QFT2.  The inductive union state is therefore well-defined by assigning to each local element the value computed in any region containing its support.
\end{proof}



\begin{longtable}{>{\raggedright\arraybackslash}p{0.24\textwidth} >{\raggedright\arraybackslash}p{0.36\textwidth} >{\raggedright\arraybackslash}p{0.28\textwidth}}
\caption[Local-state compactness/Cauchy discharge]{Local-state compactness/Cauchy discharge.}\label{p12:tab:local-state-compactness-cauchy-discharge}\\
\toprule
Step & Functional-analytic input & Paper-VII-C consequence \\
\midrule
\endfirsthead
\toprule
Step & Functional-analytic input & Paper-VII-C consequence \\
\midrule
\endhead
Bounded local algebra & QFT2--QFT3 give uniformly bounded finite representatives and a local \(C^*\)-seminorm quotient & Each finite vacuum functional is a state on the local completion. \\
Subsequential existence & Banach--Alaoglu weak-\(*\) compactness of the state space & Every bounded local region has cluster states along refinement subnets. \\
Uniqueness mechanism & QFT4 Cauchy/path-independence of all local matrix elements & All cluster states agree on the local cylinder algebra. \\
Inclusion compatibility & QFT2 local representatives restrict consistently under \(\mathcal O_1\subset\mathcal O_2\) & The local states glue to the inductive local algebra. \\
Positivity and normalization & Positivity and \(\omega_L(\mathbf1)=1\) pass through the limit & The glued functional is a continuum state, not only a pointwise limit. \\
\bottomrule
\end{longtable}
\TableNotes{The ledger separates weak-\(*\) subsequential compactness from the stronger QFT4 Cauchy statement that gives the canonical local state.}


\begin{corollary}[Canonical local-state quotient discharge]
\label{p12:cor:canonical-local-state-quotient-discharge}
Under the hypotheses of Theorem~\ref{p12:thm:canonical-local-state}, the limiting functional is independent of the representative chosen in the finite local cylinder quotient.  Equivalently, if two local representatives differ by a QFT4-null family, then they have the same continuum expectation value, and the inclusion-compatible state descends to the inductive local quotient algebra.
\end{corollary}
\begin{proof}
Let \(A_L\) and \(A'_L\) be two representatives of the same local observable in a bounded region \(\mathcal O\), and write \(N_L=A_L-A'_L\).  By definition of the local null quotient and the QFT4 Cauchy obligation, the scalar net \(\omega_L(N_L)\) converges to zero, and the same statement holds after multiplying by any fixed bounded local cylinder representative needed to test compatibility.  Hence
\[
\lim_L\omega_L(A_L)=\lim_L\omega_L(A'_L).
\]
Positivity is inherited because the representative family of \(A^*A\) is positive modulo the same null ideal, and normalization is inherited from \(\omega_L(\mathbf 1)=1\).  If \(\mathcal O_1\subset\mathcal O_2\), QFT2 gives common finite representatives on refinements of the larger region, and QFT4 makes the resulting scalar limits path-independent.  Thus the local states glue and descend to the inductive quotient algebra claimed in Table~\ref{p12:tab:local-state-compactness-cauchy-discharge}.
\end{proof}


% Source Paper VII-C, original line 1382; inherited proof bodies unchanged.
\subsubsection{OS admissibility and native covariance}
\label{p12:sec:os-admissibility-native-covariance}

\begin{remark}[Analytic reconstruction firewall]
\label{p12:rem:v248-analytic-reconstruction-firewall}
The same-owner equivariance theorem of Paper~II-A can constrain an already
licensed finite Euclidean or Lorentzian owner, but it cannot create the
cross-owner continuum reconstruction map.  In particular it does not imply
QFT4 refinement Cauchy convergence, an invariant arbitrary interacting state,
reflection positivity of a limit, QFT6 local-flat covariance, temperedness,
clustering, or the Osterwalder--Schrader reconstruction hypotheses.  The
local-chart OS theorem below is therefore an \emph{analytic certification}:
once its displayed hypotheses are proved, the reconstructed Hilbert space,
vacuum, Hamiltonian, semigroup, and local net are theorem outputs rather than
primitive continuum inputs.  This is the continuum counterpart of the
Euclidean/Lorentzian constructor firewall in Paper~III-E.
For the actual normalized full-ground/Gibbs branch, however, an equivariant
finite Hamiltonian already gives an invariant density by
Theorem~\ref{p12:thm:native-bh-covariance} and
Corollary~\ref{p12:cor:v227-covariance-provider}. On that branch no additional
finite vacuum selector is an analytic reconstruction hypothesis.
\end{remark}

\begin{definition}[OS-admissible Born--Hopf limit]
\label{p12:def:os-admissible-limit}
A continuum local state \(\omega\) obtained from a Born--Hopf refinement system is OS-admissible if its Euclidean Schwinger functionals satisfy:
\begin{enumerate}[label=(\roman*), leftmargin=2em]
\item exact native \(\RR_\tau\rtimes\Spin(4)\) covariance and ordinary flat Euclidean covariance on every compact local-chart exhaustion window;
\item reflection positivity;
\item symmetry/permutation covariance for the relevant fields or loop variables;
\item regularity/temperedness sufficient for reconstruction;
\item clustering or a specified vacuum-sector uniqueness condition when a unique-vacuum or pure-phase claim consumes it; the base bounded algebraic net certificate does not require that stronger claim.
\end{enumerate}
\end{definition}
For a bounded ordinary local-net claim this admissibility is implemented by
Definition~\ref{p12:def:v2210-lit-bounded-os-certificate} or an independently
verified equivalent reconstruction certificate. Vacuum commutator norm alone
is weaker than representation-wide locality, as in
Lemma~\ref{p12:lem:v2210-lit-locality-sandwich}.

\begin{definition}[Local-flat-chart OS exhaustion]
\label{p12:def:local-flat-chart-os-exhaustion}
A Born--Hopf refinement tower admits a local-flat-chart OS exhaustion if there are product charts
\[
\varphi_L:I_L\times U_L\longrightarrow \RR\times S^3_x(R_{x,L}),
\qquad R_{x,L}\asymp m_{x,L}\ell_{\BH},
\]
with chart radii \(\rho_L\) satisfying
\begin{equation}
\label{p12:eq:local-chart-exhaustion-radii}
\rho_L\longrightarrow\infty,
\qquad
\frac{\rho_L}{R_{x,L}}\longrightarrow0,
\end{equation}
such that every fixed compact set \(K\Subset\RR^4\) is contained in \(\varphi_L(I_L\times U_L)\) for all sufficiently large \(L\).  Pulling finite Schwinger functionals back through \(\varphi_L^{-1}\), the chart-curvature and spectral-window comparison errors vanish uniformly on each such \(K\).  The exhaustion is only a local reciprocal-position presentation of the native compact momentum regulator; it is not a new global \(E(4)\)-action on \(S^3\).
\end{definition}

\begin{theorem}[Native Born--Hopf covariance under an equivariant finite Hamiltonian]
\label{p12:thm:native-bh-covariance}
Let \(U_L\) be the finite \(\Spin(4)\) representation and assume \([H_L,U_L(g)]=0\) for all \(g\in\Spin(4)\).  Let \(P_{0,L}\) be the ground spectral projector, of arbitrary finite rank, and define the symmetric vacuum density
\[
 \rho_{0,L}=P_{0,L}/\operatorname{rank}P_{0,L}.
\]
Then
\[
 U_L(g)\rho_{0,L}U_L(g)^{-1}=\rho_{0,L}
\]
and the corresponding finite vacuum state is exactly invariant under the native compact spatial symmetry and shifted Euclidean time transfer.  Moreover, if a flat comparison transport \(g\) has representatives \(g_L\) satisfying
\[
 \|(gA)_L-U_L(g_L)A_LU_L(g_L)^{-1}\|\to0,
\]
then
\[
 |\omega_L((gA)_L)-\omega_L(A_L)|\to0.
\]
Thus no independent interacting state-covariance defect is required on this symmetric finite-vacuum branch beyond representative/chart convergence.
\end{theorem}
\begin{proof}
Commutation of \(H_L\) with \(U_L(g)\) implies commutation of every spectral projector, in particular \(P_{0,L}\), by finite spectral functional calculus.  Hence \(\rho_{0,L}\) and its state are invariant.  For the local flat representative, invariance and the norm-one property of a state give
\[
 |\omega_L((gA)_L)-\omega_L(A_L)|
 \le \|(gA)_L-U_L(g_L)A_LU_L(g_L)^{-1}\|,
\]
which tends to zero by hypothesis.
\end{proof}


% Source Paper VII-C, original line 1468; inherited proof bodies unchanged.
\subsubsection{Regularity, reflection positivity, and local-chart reconstruction}
\label{p12:sec:regularity-positivity-local-chart}

\begin{lemma}[Bounded-cylinder regularity]
\label{p12:lem:bounded-cylinder-regularity}
Let
\[
S_{n,L}(f_1,\ldots,f_n)=\omega_L(A_L(f_1)\cdots A_L(f_n))
\]
be Schwinger functionals built from bounded local cylinder or bounded Wilson-loop observables.  If
\[
\|A_L(f_i)\|\le C_i p_i(f_i)
\]
for test-function seminorms \(p_i\), uniformly in \(L\), then the QFT4 limit \(S_n\) is a continuous multilinear functional, hence a distribution.  Thus QFT7 regularity is closed for the bounded cylinder/Wilson algebra.
\end{lemma}
\begin{proof}
The finite operator bound gives
\[
|S_{n,L}(f_1,\ldots,f_n)|
\le
\prod_{i=1}^n \|A_L(f_i)\|
\le
\prod_{i=1}^n C_i p_i(f_i).
\]
The right-hand side is independent of \(L\).  Passing to the QFT4 limit preserves the same seminorm bound, which is precisely continuity in the test-function topology.  This argument does not apply to the normalized curvature residues \(X_{p,L}=12a_L^{-4}(1-\operatorname{Re}W_F(p))\); those are handled by QFT10/UCS.
\end{proof}

\begin{lemma}[Reflection positivity survives the limit]
\label{p12:lem:rp-survives-limit}
Assume QFT4.  Suppose that each finite transfer presentation is reflection positive on positive-time cylinder observables, so that for every finite list \(A_{i,L}\)
\[
\sum_{i,j}\overline{c_i}c_j\,\omega_L(\Theta_LA_{i,L}^*A_{j,L})\ge0.
\]
Then the limiting state is OS reflection positive.
\end{lemma}
\begin{proof}
\ProofPlan{We pass finite reflection positivity to the continuum limit by keeping reflected cylinder products local, applying QFT4 to the matrix elements, and taking the limit of finite nonnegative quadratic forms.}
\begin{proofsteps}
\item \textbf{Reflected cylinder products remain local.}
For a fixed finite list of continuum cylinder observables, the reflected products \(\Theta A_i^*A_j\) are again bounded local cylinder observables after enlarging the time-reflected support.

\item \textbf{QFT4 passage to the limiting state.}
QFT4 gives
\[
\omega_L(\Theta_LA_{i,L}^*A_{j,L})\longrightarrow
\omega(\Theta A_i^*A_j).
\]

\item \textbf{Limit of finite positive quadratic forms.}
Passing the finite nonnegative quadratic form to the limit gives
\[
\sum_{i,j}\overline{c_i}c_j\,\omega(\Theta A_i^*A_j)
=
\lim_L
\sum_{i,j}\overline{c_i}c_j\,\omega_L(\Theta_LA_{i,L}^*A_{j,L})
\ge0.
\]

\item \textbf{No extra positivity constant.}
Finite reflection positivity is supplied on the licensed Hamiltonian/time-
ordered or reflected-action construction with its actual positive-time
algebra and same-owner identification. Pointwise positive face weight alone
is insufficient, as in Proposition~\ref{p10:prop:v2211-central-transfer-positive-type}.
On that certified finite domain the remaining step here is passage of the
reflected products to the limit.
\end{proofsteps}
\end{proof}

% Framework v2.217 helper integration.
\begin{remark}[Finite GW source positivity has a different premise]
\label{p12:rem:v2217-gw-source-os-scope}
The configuration-wise and finite-ensemble source positivity of
Paper~VII-A, Theorem~\ref{p10:thm:v2217-temporal-gw-source}, is a
Hilbert--Schmidt-square statement for a cyclic Euclidean regulator.
It is not the reflected positive-time quadratic form assumed by the
preceding lemma. A determinant-positive gauge measure and a nonzero
off-clock response do not supply that missing transfer identification or
OS certificate; the existing physical clock and continuum hypotheses
remain in force.
\end{remark}


\begin{lemma}[Native covariance and local flat chart recovery]
\label{p12:lem:euc-invariance-survives}
Assume QFT1, QFT4, and QFT6.  Then the canonical continuum state \(\omega\) is exactly invariant under the native group \(\RR_\tau\rtimes\Spin(4)\).  In local product charts \(I\times U\subset\RR\times S^3\), this gives the usual Euclidean covariance on compact comparison windows up to the vanishing chart-curvature errors specified by QFT1/QFT6.
\end{lemma}
\begin{proof}
Exact native covariance is Theorem~\ref{p12:thm:native-bh-covariance} passed through QFT4.  For local flat covariance, restrict to a compact comparison window in a product chart.  The spectral-symbol recovery of QFT1 supplies representatives of the transported local observable, and QFT6 requires the chart transport error in expectation to vanish under refinement.  Therefore the limiting functional is native-covariant globally and locally Euclidean in the flat chart limit.  No primitive global \(E(4)\)-action on the compact Born-derived reciprocal-position geometry is being inserted.
\end{proof}

\begin{theorem}[Local-chart OS reconstruction]
\label{p12:thm:local-chart-os-reconstruction}
Assume QFT1--QFT7, reflection positivity, the vacuum uniqueness/clustering condition of Definition~\ref{p12:def:os-admissible-limit}, and a local-flat-chart OS exhaustion in the sense of Definition~\ref{p12:def:local-flat-chart-os-exhaustion}.  Then the pulled-back bounded Schwinger functionals on every fixed compact \(K\Subset\RR^4\) satisfy the ordinary Osterwalder--Schrader hypotheses in the inductive local limit.  Consequently the Born--Hopf local-chart limit reconstructs a Hilbert space, cyclic vacuum, nonnegative Hamiltonian, Euclidean-time semigroup, and local Wightman/observable net on the reconstructed flat local theory.  If QFT10 is also supplied, the same local reconstruction includes the small-observable field distributions.
\end{theorem}
\begin{proof}
\ProofPlan{We reduce the native compact-spatial reconstruction to the standard
OS theorem by using an exhaustion by flat charts whose curvature error tends to
zero on every fixed compact window.}

\begin{proofsteps}

\item \textbf{Local chart covariance.}\par\noindent
Let \(K\Subset\RR^4\) and choose \(L\) so large that \(K\) and all
reflected/permuted supports under consideration lie inside the chart of radius
\(\rho_L\). Lemma~\ref{p12:lem:tier2-flat-chart-covariance} derives an exhausting chart sequence with geometric defect \(O(\rho_L^2/R_{x,L}^2)\to0\).  QFT6 also consumes the actual transported Schwinger representative comparison. On the normalized full-ground branch with the actual equivariant comparison unitary, Theorem~\ref{p12:thm:native-bh-covariance} removes a separate state term; Corollary~\ref{p12:cor:v227-covariance-provider} gives the quantitative action-defect bound on approximate canonical branches. Arbitrary states or time--space comparisons not implemented by such a unitary retain their actual correlation comparison. Combining these tests with the geometric exhaustion gives local Euclidean covariance on the fixed compact window.

\item \textbf{Reflection positivity and permutation covariance.}\par\noindent
The reflected products remain inside a compact window, after enlarging the
window slightly. Lemma~\ref{p12:lem:rp-survives-limit} passes the positive
finite OS form to the limit there. The finite Wilson/\allowbreak cylinder
products have the required adjoint and permutation covariance, and QFT4 passes
those finite identities to the same local limit.

\item \textbf{Regularity and clustering.}\par\noindent
For bounded cylinder/\allowbreak Wilson observables,
Lemma~\ref{p12:lem:bounded-cylinder-regularity} gives the test-function seminorm
bounds needed for the local observable net. Curvature-residue fields are added
only in sectors where QFT10 supplies the corresponding distributional moment
bounds. The assumed vacuum uniqueness/\allowbreak clustering condition fixes
the vacuum sector on the local inductive limit.

\item \textbf{Inductive local consistency.}\par\noindent
If \(K_1\subset K_2\), the two chart limits agree on observables supported in
\(K_1\) by QFT4 path independence and the common-refinement argument of
Theorem~\ref{p12:thm:qft4hij-local-state-convergence}. Hence the compact-window
OS reconstructions form an inductive local net.

\item \textbf{Application of standard OS reconstruction.}\par\noindent
The preceding steps are exactly the standard OS hypotheses on each compact flat
window, with compatible restrictions. Applying the Osterwalder--Schrader
reconstruction theorem%
~\cite{OsterwalderSchrader1973,OsterwalderSchrader1975,OsterwalderSeiler1978}
window by window and passing through the inductive local net gives the stated
Hilbert space, vacuum, positive Hamiltonian, Euclidean-time semigroup, and local
Wightman/\allowbreak observable net. No global flat \(E(4)\)-action on the
compact carrier has been inserted; flat covariance is recovered only on the
local exhaustion.

\end{proofsteps}
\end{proof}

\begin{theorem}[Born--Hopf OS reconstruction implication]
\label{p12:thm:bh-os-reconstruction}
Assume QFT1--QFT8, the local-flat-chart OS exhaustion of Definition~\ref{p12:def:local-flat-chart-os-exhaustion}, and the vacuum uniqueness/clustering part of Definition~\ref{p12:def:os-admissible-limit}.  Then the Born--Hopf continuum functional reconstructs a Hilbert space \(\mathcal H\), vacuum \(\Omega\), nonnegative Hamiltonian \(H\), Euclidean-time semigroup \(e^{-tH}\), and a local observable net \(\mathcal A(\mathcal O)\) in the local-chart OS/Wightman sense.  Combining this reconstructed state/net with the Papers~III--IV Lorentzian sign and positive-frequency closure gives the Lorentzian real-time presentation used by the rest of the series.
\end{theorem}
\begin{proof}
Theorem~\ref{p12:thm:canonical-local-state} gives a canonical Euclidean functional.  Lemma~\ref{p12:lem:rp-survives-limit} gives reflection positivity.  Lemma~\ref{p12:lem:euc-invariance-survives} gives exact native covariance and the qualitative local flat chart limit, while Lemma~\ref{p12:lem:tier2-flat-chart-covariance} and Definition~\ref{p12:def:local-flat-chart-os-exhaustion} supply the quantitative compact-window exhaustion needed for ordinary flat OS covariance.  QFT7 supplies regularity, and QFT8 supplies the declared sector nontriviality witness when an interacting sector is claimed.  Therefore Theorem~\ref{p12:thm:local-chart-os-reconstruction} applies.  The Lorentzian sign is not selected here: it is imported from Paper~III, with the finite positive-frequency dispersion input supplied by Paper~IV.  Hence Paper~VII-C reconstructs the continuum state/local-net layer and then attaches it to the already derived Lorentzian orientation rather than inserting a new branch.
\end{proof}



\begin{longtable}{>{\raggedright\arraybackslash}p{0.23\textwidth} >{\raggedright\arraybackslash}p{0.37\textwidth} >{\raggedright\arraybackslash}p{0.28\textwidth}}
\caption[OS reconstruction hypothesis-matching ledger]{OS reconstruction hypothesis-matching ledger.}\label{p12:tab:os-hypothesis-matching-ledger}\\
\toprule
OS requirement & Born--Hopf discharge & Reconstruction output \\
\midrule
\endfirsthead
\toprule
OS requirement & Born--Hopf discharge & Reconstruction output \\
\midrule
\endhead
Euclidean local algebra & QFT1--QFT3 and Proposition~\ref{p12:prop:spectral-symbol-qft123} & Local cylinder/Wilson representatives on compact chart windows. \\
Canonical state & Theorem~\ref{p12:thm:canonical-local-state} & A positive normalized local functional compatible under inclusions. \\
Reflection positivity & Lemma~\ref{p12:lem:rp-survives-limit} and finite positive transfers & Positive OS form on positive-time local observables. \\
Regularity & Lemma~\ref{p12:lem:bounded-cylinder-regularity} for bounded observables; QFT10/UCS for curvature residues & Schwinger functionals are distributions in the relevant sector. \\
Covariance & Native covariance plus local-flat-chart exhaustion and Lemma~\ref{p12:lem:tier2-flat-chart-covariance} & Flat Euclidean covariance on every compact reciprocal-position chart window. \\
Clustering/vacuum sector & The vacuum uniqueness/clustering clause of Definition~\ref{p12:def:os-admissible-limit} & The reconstructed cyclic vacuum is the intended sector vacuum. \\
Lorentzian orientation & Papers~III--IV import Lorentzian sign and positive-frequency closure & The OS Hilbert space/net is attached to the already-fixed real-time presentation. \\
\bottomrule
\end{longtable}
\TableNotes{This table is a hypothesis match for the standard OS theorem used in Theorems~\ref{p12:thm:local-chart-os-reconstruction} and \ref{p12:thm:bh-os-reconstruction}; it does not add a second continuum axiom system.}


\begin{corollary}[OS hypothesis ledger discharge]
\label{p12:cor:os-hypothesis-ledger-discharge}
The hypotheses used in Theorem~\ref{p12:thm:bh-os-reconstruction} are exactly the ordinary OS reconstruction hypotheses after pullback to each compact local chart window and passage through the QFT4 local limit.
\end{corollary}
\begin{proof}
Table~\ref{p12:tab:os-hypothesis-matching-ledger} matches each ordinary OS input to a Born--Hopf source.  QFT1--QFT3 give the local Euclidean cylinder algebra; Theorem~\ref{p12:thm:canonical-local-state} and Corollary~\ref{p12:cor:canonical-local-state-quotient-discharge} give the positive normalized local functional on the quotient; Lemma~\ref{p12:lem:rp-survives-limit} gives reflection positivity; Lemma~\ref{p12:lem:euc-invariance-survives} gives native covariance once the finite Hamiltonian is equivariant, while Lemma~\ref{p12:lem:tier2-flat-chart-covariance} derives the geometric flat exhaustion and the canonical covariance provider reduces its state term to the actual representative/action-defect comparison; QFT7 and QFT10/UCS give regularity for bounded observables and curvature residues; and the OS-admissible clustering clause fixes the vacuum sector.  Therefore no independent continuum OS datum is inserted at Theorem~\ref{p12:thm:bh-os-reconstruction}; the theorem applies the standard OS reconstruction theorem to the already matched local-chart data.
\end{proof}

\begin{corollary}[Continuum reconstructed structures are conditional descendants]
\label{p12:cor:fields-are-outputs}
Let a generated finite Born--Hopf refinement family satisfy the analytic
QFT/reconstruction hypotheses used in
Theorem~\ref{p12:thm:bh-os-reconstruction}.  Then the reconstructed Hilbert
space, vacuum ray, Hamiltonian, local observable net, symmetry
representation, and QFT10 local fields are theorem outputs of that certified
finite tower.  They require no second premise asserting the actuality of
those continuum structures.  The QFT4/QFT6/QFT8--QFT10 convergence,
covariance, nontriviality, gap, residue, and reconstruction hypotheses are
nevertheless genuine analytic obligations on the generated tower; they are
not selectors and are not proved by the finite Paper-II-A structural
induction alone.
\end{corollary}
\begin{proof}
QFT1--QFT4 construct the limiting local Euclidean functional; reflection
positivity, covariance/regularity, and the remaining stated OS hypotheses
then invoke the reconstruction theorem, whose outputs include the Hilbert
space, vacuum, Hamiltonian and local net.  QFT10 supplies the controlled
operator-valued distributional field limits.  Thus, conditional on the
analytic certificate, the continuum objects are descendants rather than new
roots.  Conversely the finite structural-geometrization theorem does not
establish convergence of an infinite refinement family, OS regularity,
non-Gaussianity, a uniform gap, or residue convergence, so those hypotheses
cannot be deleted by this corollary.
\end{proof}


% Source Paper VII-C, original line 9320; inherited proof bodies unchanged.
\subsection{Quantitative propagation of the canonical covariance provider}
\label{p12:sec:v227-covariance-provider}
\begin{corollary}[The canonical state term has an action-defect bound]
\label{p12:cor:v227-covariance-provider}
On each actual finite normalized full-ground branch let $U_L$ be the
unitary comparison, $\epsilon_L=\|[H_L,U_L]\|$, and $\Delta_L>0$ its
gap above the full ground eigenspace. For a bounded local representative put
$\delta_L(A)=\|(gA)_L-U_LA_LU_L^*\|$. Then
\[
 |\omega_L((gA)_L)-\omega_L(A_L)|
 \le\delta_L(A)+2\|A_L\|
             \min\{1,\epsilon_L/\Delta_L\}.
\]
On a Gibbs branch replace $\epsilon_L/\Delta_L$ by
$\beta_L\epsilon_L$. Fixed canonical block mixtures use the weighted
occupied-block bounds. If these actual residuals vanish, there is no
additional state-covariance selector to verify. Exact equivariance sets
the state term to zero at every level.
\end{corollary}
\begin{proof}
Apply Corollary~\ref{p1b:cor:v227-chart-covariance} on the actual
finite carriers, and its occupied-block extension when applicable.
For bounded products used as Schwinger representatives the same bound
applies to the complete product/transfer word, with its actual norm and
transport residual. If an equivariant reference exists, the action
commutator defect is at most twice the actual operator residual.
Proposition~\ref{p12:prop:v226-refinement-coefficients} can bound that
residual from operator-valued native coefficients without a growing outer
dimension factor.
\end{proof}
The ratios, rather than $\epsilon_L\to0$ alone, are essential if a gap
closes or temperature changes. No uniform positive ground gap is needed
when $\epsilon_L/\Delta_L\to0$. A proper pure substate of a degenerate
ground space is not covered. A time--space Euclidean motion not implemented
by the actual comparison unitary still requires its full correlation
comparison, as in Proposition~\ref{p12:prop:v2210-lit-covariance-scope}.

This completes the missed provider-to-ledger handoff in the earlier BH5
and OS proof: on the canonical equivariant branch the remaining local-flat
obligation is the actual representative/chart comparison, not a second
interacting vacuum selector. QFT4 same-phase canonicity, reflection
positivity, local-net identification, clustering, local fields and physical
scattering retain their separate stated hypotheses. Finite response sampling
is levelwise and cannot replace any uniform refinement estimate.

% END FRAMEWORK v2.227: DOWNWARD CLOSURES


% BEGIN FRAMEWORK v2.230: SHARP GIBBS PROVIDER PROPAGATION
\begin{corollary}[Sharpened finite Gibbs-state covariance provider]
\label{p12:cor:v230-sharp-gibbs-provider}
On the finite canonical Gibbs branch in
Corollary~\ref{p12:cor:v227-covariance-provider}, with the same actual
unitary comparison, $\epsilon_L=\|[H_L,U_L]\|$ and
$\delta_L(A)=\|(gA)_L-U_LA_LU_L^*\|$, the thermal inequality sharpens to
\[
 \boxed{
 |\omega_L((gA)_L)-\omega_L(A_L)|
 \le\delta_L(A)+\|A_L\|\min\{2,\beta_L\epsilon_L\}.}
\]
It follows from Theorem~\ref{p1b:thm:v230-sharp-gibbs-covariance}
and Corollary~\ref{p1b:cor:v230-sharp-gibbs-chart}; the earlier
full-ground bound is unchanged.  Finite exact covariance at
$\epsilon_L=0$ creates no independent canonical-vacuum selector,
but does not by itself prove representative convergence,
reflection positivity, clustering, boosts, local field reconstruction
or an interacting continuum scattering matrix.
\end{corollary}
% END FRAMEWORK v2.230: SHARP GIBBS PROVIDER PROPAGATION


% End of predecessor document.

\section{Gauge, Wilson, gap and spectral continuum obligations}
\label{p12:sec:v237-c4}
Finite Wilson/Schur/Hodge data are not unconditional physical confinement or a uniform continuum mass gap. Those claims consume their own limits, area inequalities and state realization.

% Source Paper VII-C, original line 1676; inherited proof bodies unchanged.
% Former logical part: Gap Transfer and Local-Field Reconstruction; integrated into the new thematic argument.

\noindent\textbf{Orientation.}
This part transfers finite spectral gap information to the reconstructed Hamiltonian and records the local-field reconstruction step.  The pure-glue positive-gap branch, controlled gapless sectors, and distributional field limits are separated so that massive-gap consequences are claimed only where the finite prefix supplies the required gap data.


% Source Paper VII-C, original line 1681; inherited proof bodies unchanged.
\subsection{Mass gap and transfer estimates}
\label{p12:sec:mass-gap-transfer}


% Source Paper VII-C, original line 1684; inherited proof bodies unchanged.
\subsubsection{Gap and small-field ledger}
\label{p12:sec:gap-small-field-ledger}

\begin{table}[htbp]
\centering
\small
\begin{tabularx}{\textwidth}{@{}
>{\raggedright\arraybackslash}p{0.55\textwidth}
>{\raggedright\arraybackslash}X
@{}}
\toprule
\textbf{Gap or small-field step} & \textbf{Discharged by} \\
\midrule
Uniform transfer gap
& Definition~\ref{p12:def:uniform-transfer-gap}. \\

Mass-gap transfer to the OS Hamiltonian
& Theorem~\ref{p12:thm:mass-gap-transfer}. \\

proved uniform pure-glue gap
& Corollary~\ref{p12:cor:paper11-pure-glue-gap-floor}. \\

Gapless sectors separated from the pure-glue claim
& Remark~\ref{p12:rem:gapless-sectors}. \\

Renormalized small-observable family
& Definition~\ref{p12:def:small-observable-family}. \\

Distributional field limit
& Definition~\ref{p12:def:distributional-field-limit}. \\

Small-observable field reconstruction
& Theorem~\ref{p12:thm:small-observable-field-reconstruction}. \\
\bottomrule
\end{tabularx}
\caption[Mass gap and small-observable reconstruction]{Mass gap and small-observable reconstruction.}
\label{p12:tab:mass-gap-small-observable-ledger}
\label{p12:tab:gap-small-field-ledger}
\TableNotes{The table records the passage from finite transfer gaps to an
OS-Hamiltonian gap and then to distribution-valued local fields.  Gapless
photon/chiral sectors remain explicitly separated from the pure-glue
Yang--Mills handoff.}
\end{table}


% Source Paper VII-C, original line 1728; inherited proof bodies unchanged.
\subsubsection{Finite-gap transfer theorem}
\label{p12:sec:uniform-transfer-gap}

A finite transfer gap at one level is not enough for a continuum mass gap.  The continuum statement requires a uniform lower bound along refinement and a reconstruction theorem ensuring that the bound survives the limit.

\begin{definition}[Uniform transfer gap]
\label{p12:def:uniform-transfer-gap}
A Born--Hopf refinement system has a uniform transfer gap in a sector \(\mathcal S\) if there is \(\Delta_\ast>0\) such that for every sufficiently refined \(L\),
\begin{equation}
\label{p12:eq:uniform-transfer-gap}
\Spec(H_L|_{\mathcal S_L})\cap(0,\Delta_\ast)=\varnothing .
\end{equation}
Equivalently, for vacuum-subtracted local observables in the sector,
\begin{equation}
\label{p12:eq:uniform-correlation-decay}
|\omega_L((\tau_{L,t}A_{0,L})B_{0,L})|
\le C_{A,B}e^{-\Delta_\ast t}+O(a_L^\eta).
\end{equation}
\end{definition}



\begin{lemma}[Sector spectral-measure gap discharge]
\label{p12:lem:sector-spectral-measure-gap-discharge}
Assume the hypotheses of Theorem~\ref{p12:thm:mass-gap-transfer}.  For every vacuum-subtracted reconstructed local vector \(\psi=A_0\Omega\), its positive Hamiltonian spectral measure \(\mu_\psi\) satisfies
\begin{equation}
\label{p12:eq:sector-spectral-measure-gap-discharge}
\mu_\psi((0,\Delta_*))=0.
\end{equation}
By polarization the same holds for cross spectral measures.  If such local vectors are cyclic in the sector, then \(\operatorname{Spec}(H|_{\mathcal S})\cap(0,\Delta_*)=\varnothing\).
\end{lemma}
\begin{proof}
For \(\psi=A_0\Omega\), the limiting Euclidean autocorrelation satisfies
\[
0\le\langle\psi,e^{-tH}\psi\rangle=\int_0^\infty e^{-t\lambda}\,d\mu_\psi(\lambda)\le C_\psi e^{-\Delta_*t}.
\]
If \(\mu_\psi((0,\Delta_* -\varepsilon])=m>0\), then the left side is at least \(m e^{-(\Delta_* -\varepsilon)t}\), contradicting the upper bound for sufficiently large \(t\).  Hence there is no mass in \((0,\Delta_* -\varepsilon]\) for every \(\varepsilon>0\), proving the claim.  Cross measures follow by the polarization identity applied to \(A_0\Omega\pm B_0\Omega\) and \(A_0\Omega\pm iB_0\Omega\).  Cyclicity then excludes a nonzero subgap spectral subspace orthogonal to all local vectors.
\end{proof}


% Source Paper VII-C, original line 1767; inherited proof bodies unchanged.
\subsubsection{Transfer to the reconstructed Hamiltonian}
\label{p12:sec:gap-transfer-reconstructed-hamiltonian}

\begin{theorem}[Mass-gap transfer]
\label{p12:thm:mass-gap-transfer}
Assume QFT1--QFT7 and QFT9 in the positive-gap branch.  If local vectors generated by the reconstructed local algebra are cyclic in the sector under consideration, then the reconstructed Hamiltonian satisfies
\begin{equation}
\label{p12:eq:continuum-gap}
\Spec(H)\cap(0,\Delta_\ast)=\varnothing .
\end{equation}
\end{theorem}
\begin{proof}
For finite \(L\), QFT9 gives the exponential Euclidean correlation bound \eqref{p12:eq:uniform-correlation-decay}.  QFT4 and QFT7 allow the bound to pass to the reconstructed Euclidean correlator.  Lemma~\ref{p12:lem:sector-spectral-measure-gap-discharge} converts this surviving decay bound into the spectral-measure statement \(\mu_{A,B}((0,\Delta_\ast))=0\) for all vacuum-subtracted local vectors.  Since those local vectors are cyclic in the sector by hypothesis, no subgap spectral subspace remains in the reconstructed sector, giving \eqref{p12:eq:continuum-gap}.
\end{proof}



\begin{longtable}{>{\raggedright\arraybackslash}p{0.24\textwidth} >{\raggedright\arraybackslash}p{0.36\textwidth} >{\raggedright\arraybackslash}p{0.28\textwidth}}
\caption[Mass-gap transfer proof ledger]{Mass-gap transfer proof ledger.}\label{p12:tab:mass-gap-transfer-proof-ledger}\\
\toprule
Input & Use in the transfer proof & Consequence \\
\midrule
\endfirsthead
\toprule
Input & Use in the transfer proof & Consequence \\
\midrule
\endhead
Uniform finite transfer gap & Supplies the bound \(\|R_L\|\le e^{-\Delta_\ast t}\) or the equivalent Euclidean correlation decay & Low-energy spectral weight below \(\Delta_\ast\) is absent at finite level uniformly in refinement. \\
QFT4 local convergence & Passes finite matrix elements to the reconstructed local correlators & The exponential decay bound survives on local vectors. \\
QFT7 regularity & Makes the limiting correlators distributions/OS inputs rather than only formal pointwise limits & The OS Hamiltonian spectral measure is well-defined on the reconstructed local sector. \\
Cyclicity of local vectors & Local vectors generate the sector under the reconstructed algebra & Absence of subgap spectral support on local vectors implies absence in the whole sector. \\
Paper-VII-B finite gap sequence & Provides \(\gamma_L=E_{1,L}-E_{0,L}\) & A continuum positive gap follows only if \(\gamma_*=\inf_L\gamma_L>0\) and the local semigroup limit is established. \\
\bottomrule
\end{longtable}
\TableNotes{The ledger expands Theorem~\ref{p12:thm:mass-gap-transfer}: finite positivity alone gives a floor at each level; QFT4/QFT7 and cyclicity are the analytic mechanism that transfers it to the reconstructed Hamiltonian.}


% Source Paper VII-C, original line 1803; inherited proof bodies unchanged.
\subsubsection{Pure-glue branch and gapless-sector boundary}
\label{p12:sec:pure-glue-gapless-boundary}

\begin{corollary}[Paper-VII-B finite pure-glue gap sequence and sharp continuum criterion]
\label{p12:cor:paper11-pure-glue-gap-floor}
For the repaired Paper-VII-B \(SU(3)_c\) pure-glue export, define the exact finite gaps
\begin{equation}
\label{p12:eq:paper11-finite-gap-sequence}
\gamma_L:=E_{1,L}-E_{0,L}\ge0,
\qquad
\gamma_*:=\inf_{L\in\mathcal R_{\rm adm}}\gamma_L .
\end{equation}
No count-derived positive lower bound is part of the current import.  If \(\gamma_*>0\) and the QFT4/OS semigroup matrix elements of the
\emph{full-ground-subtracted} finite local-sector vectors converge to a
cyclic family of reconstructed local vectors, then the reconstructed
pure-glue Hamiltonian satisfies
\begin{equation}
\label{p12:eq:paper11-continuum-gap-criterion}
\operatorname{Spec}(H_{\YM})\cap(0,\gamma_*)=\varnothing .
\end{equation}
If only \(\gamma_L>0\) level by level is known, no positive continuum lower bound follows.
\end{corollary}
\begin{proof}
The definition of \(\gamma_L\) is the spectral theorem applied to the supplied finite self-adjoint Hamiltonian.  If \(\gamma_*>0\), let \(Q_{0,L}:=I-P_{0,L}\) project away the entire
finite pure-glue ground space and put
\(\psi_{A,L}:=Q_{0,L}A_L\Omega_L\),
\(\psi_{B,L}:=Q_{0,L}B_L\Omega_L\). For \(t\ge0\), the spectral theorem gives
\[
|\langle \psi_{A,L},e^{-t(H_L-E_{0,L})}\psi_{B,L}\rangle|
\le \|\psi_{A,L}\|\,\|\psi_{B,L}\|e^{-\gamma_*t}.
\]
This is a \emph{full-ground} subtraction; subtracting only the chosen
vacuum expectation is equivalent only when the sector ground state is unique
or the stated local probes have no overlap with the other zero modes.
The explicitly assumed QFT4/OS semigroup limit of these projected
matrix elements passes this inequality to the reconstructed local vectors.  Theorem~\ref{p12:thm:mass-gap-transfer} then excludes spectral support in \((0,\gamma_*)\).  Conversely, finite positivity alone cannot provide a uniform lower bound: the explicit family \(\gamma_L=1/L\) has \(\gamma_L>0\) for every finite \(L\) but \(\inf_L\gamma_L=0\).  Therefore \(\gamma_*>0\) is the sharp missing quantitative statement; no \(Q_{\rm wall}\) or \(24/15\) normalization enters.
\end{proof}

\begin{remark}[Gapless sectors]
\label{p12:rem:gapless-sectors}
The theorem is intentionally sector-sensitive.  A framework sector may reconstruct a controlled gapless continuum theory, but Yang--Mills in the Clay sense requires the positive-gap branch.
\end{remark}


% Source Paper VII-C, original line 1877; inherited proof bodies unchanged.
% Former logical part: Gauge Specialization, Wilson Data, and Native Spectral UCS; integrated into the new thematic argument.

\noindent\textbf{Orientation.}
This part specializes the general certificate to compact gauge data, Wilson-loop observables, plaquette/residue representatives, and the spectral uniform-cluster-summability criteria needed for local fields.  The gauge-theoretic statements are written broadly where they are representation-theoretic, while the repaired Paper~VII-B \(SU(3)_c\) pure-glue handoff remains the principal conditional application.


% Source Paper VII-C, original line 1882; inherited proof bodies unchanged.
\subsection{Gauge specialization and Wilson-loop continuum data}
\label{p12:sec:gauge-specialization}


\begin{table}[htbp]
\centering
\small
\begin{tabularx}{\textwidth}{@{}
>{\raggedright\arraybackslash}p{0.55\textwidth}
>{\raggedright\arraybackslash}X
@{}}
\toprule
\textbf{Gauge-continuum step} & \textbf{Discharged by} \\
\midrule
Pure \(G\)-Wilson Born--Hopf datum
& Definition~\ref{p12:def:pure-g-wilson-datum}. \\

Contact-normalized plaquette residue
& Lemma~\ref{p12:lem:contact-plaquette-residue}. \\

No hidden relevant pure-gauge branch
& Lemma~\ref{p12:lem:no-hidden-relevant-gauge-branch}. \\

OS-positive \(\theta=0\) sector
& Lemma~\ref{p12:lem:theta-zero-os-selection}. \\

Curvature-residue modes
& Definition~\ref{p12:def:normalized-wilson-residue}. \\

QFT10a--QFT10e structural closure
& Proposition~\ref{p12:prop:qft10a-e-closed}. \\

Uniform cumulant susceptibility and subladder
& Definitions~\ref{p12:def:uniform-small-loop-cumulant-susceptibility}
and~\ref{p12:def:ucs-subladder}. \\

Native spectral UCS
& Theorem~\ref{p12:thm:native-spectral-ucs}. \\

QFT10f--QFT10l from UCS
& Theorem~\ref{p12:thm:ucs-implies-qft10}. \\
\bottomrule
\end{tabularx}
\caption[Gauge and Yang--Mills continuum ledger]{Gauge and Yang--Mills continuum ledger.}
\label{p12:tab:gauge-ym-continuum-ledger}
\label{p12:tab:gauge-wilson-ledger}
\label{p12:tab:ym-continuum-ledger}
\TableNotes{This ledger is the gauge-specialized certificate.  The plaquette
chart is a large-\(m_x\) representative of the compact-momentum spectral
residue; it is not inserted as a fundamental rough lattice datum.}
\end{table}


% Source Paper VII-C, original line 1934; inherited proof bodies unchanged.
\subsubsection{Wilson datum and plaquette residue}
\label{p12:subsec:wilson-datum-plaquette-residue}

We now specialize the general QFT construction to compact simple gauge theory.  The algebraic Wilson-loop, Haar-positivity, heat-kernel, and plaquette-residue statements are written for a compact simple group \(G\), because those parts are representation-theoretic.  The main \(SU(3)_c\) result is the sharp conditional handoff of Theorem~\ref{p12:thm:native-bh-ym-continuum}; the repaired finite prefix does not pre-discharge its continuum hypotheses.  An arbitrary compact-simple \(G\) becomes a theorem only after its own finite gap and nontriviality witness are supplied, as recorded in Tier~4 below.

For the native branch, however, the compact simple group is not arbitrary
input.  By Corollary~\ref{p12:cor:d4-a2-selector-survives-qft}, it is
\[
G_c=SU(C),
\qquad
C=\CC^{T_+},
\qquad
\mathfrak g_c=\mathfrak{su}(C),
\]
where \(T_+\) is the finite trimer selected inside the \(D_4\)/Hurwitz
vertex-figure geometry.  Thus the fundamental Wilson representation used
below is the trimer representation \(C\), while the adjoint curvature index
runs over the eight trace-zero directions of \(\End_0(C)\).  The small-loop
normalization for the native QCD handoff is therefore the contact-normalized
Wilson residue of this finite \(A_2\)-projector color closure, not an inserted
continuum normalization.

\begin{definition}[Pure \(G\)-Wilson Born--Hopf datum]
\label{p12:def:pure-g-wilson-datum}
Let \(G\) be a compact simple Lie group with Lie algebra \(\mathfrak g\) and positive invariant form \(\langle\cdot,\cdot\rangle_{\mathfrak g}\).  A pure \(G\)-Wilson Born--Hopf datum is a spectral/momentum regulator family together with auxiliary large-\(m_x\) reciprocal-position quadrature windows
\begin{equation}
\label{p12:eq:pure-g-datum}
\DG
:=
\left\{
(\Gamma_L,\mathcal H_{L,G},\Omega_{L,G},H_{L,G},K^{\mathrm{ref}}_{L,G},\mathfrak A_{W,L}^G,\chi_L)
\right\}_{L\in\RBH},
\end{equation}
where \(K^{\mathrm{ref}}_{L,G}=e^{-a_LH_{L,G}}\) is the finite Euclidean refinement transfer, physical gap statements use the sector block \(\widehat K_{L,G}=e^{-H_{L,G}/\mu_G}\) when such a clock is supplied, and \(\mathfrak A_{W,L}^G\) is generated by normalized Wilson loops
\begin{equation}
\label{p12:eq:normalized-wilson-loop}
W_{\rho,C}^{(L)}
:=
\frac1{\dim\rho}\Tr_{\rho}\Hol_L(C),
\end{equation}
with \(\rho\) ranging over finite-dimensional unitary representations of \(G\).  Here \(\Gamma_L\) and \(\chi_L\) are not fundamental position-lattice data; they are quadrature/chart representatives of the large-\(m_x\) reciprocal-position presentation used to write local Wilson observables.  The finite state is
\begin{equation}
\label{p12:eq:finite-wilson-state}
\omega_{L,G}(A)=\langle\Omega_{L,G},A\Omega_{L,G}\rangle .
\end{equation}
\end{definition}


\begin{lemma}[Contact-normalized plaquette residue]
\label{p12:lem:contact-plaquette-residue}
Let \(G\) be compact simple and let \(\rho\) be a unitary representation with anti-Hermitian generators normalized by
\[
\operatorname{tr}_\rho(T_aT_b)=-T_\rho\delta_{ab}.
\]
For a small plaquette \(p\) of side \(a\),
\[
U_p=\exp(a^2F_{\mu\nu}+O(a^3)),
\]
and
\[
1-\frac1{d_\rho}\operatorname{Re}\operatorname{tr}_\rho U_p
=
\frac{T_\rho}{2d_\rho}a^4F_{\mu\nu}^aF_{\mu\nu}^a+O(a^5).
\]
For \(SU(3)\) in the fundamental representation this coefficient is \(1/12\).
\end{lemma}
\begin{proof}
Expanding the holonomy gives
\[
\operatorname{Re}\operatorname{tr}_\rho U_p
=d_\rho+\frac{a^4}{2}\operatorname{tr}_\rho(F_{\mu\nu}^2)+O(a^5).
\]
Since \(F_{\mu\nu}=F_{\mu\nu}^aT_a\), the trace normalization gives
\[
\operatorname{tr}_\rho(F_{\mu\nu}^2)=-T_\rho F_{\mu\nu}^aF_{\mu\nu}^a.
\]
Substitution yields the stated residue.  The remaining global coefficient is fixed by the Born--Hopf contact/cell normalization, not by fitting to a continuum observable.  For the fundamental representation of \(SU(3)\), \(T_F=1/2\) and \(d_F=3\), hence \(T_F/(2d_F)=1/12\).
\end{proof}

\begin{remark}[Trimer origin of the fundamental plaquette coefficient]
\label{p12:rem:trimer-origin-plaquette-coefficient}
In the native \(SU(3)_c\) branch, the fundamental representation in
Lemma~\ref{p12:lem:contact-plaquette-residue} is the color trimer support
\(C=\CC^{T_+}\).  Hence \(d_F=\dim C=3\), and the trace normalization
\(T_F=1/2\) is the usual fundamental normalization for the compact real form
of \(\End_0(C)\).  The factor
\[
\frac{T_F}{2d_F}=\frac1{12}
\]
therefore records the continuum chart normalization of the finite
\(D_4\)-selected \(A_2\) color closure.  It should not be read as evidence
that the \(D_4\) root shell itself is a four-dimensional spacetime lattice;
it fixes the internal color residue that is later displayed on large-\(m_x\)
local charts.
\end{remark}

\begin{lemma}[Dimension-\texorpdfstring{\(\le4\)}{<=4} pure-gauge local-term classification]
\label{p12:lem:no-hidden-relevant-gauge-branch}
Let \(G\) be compact simple and assume a smooth local four-dimensional Euclidean gauge-field presentation with the canonical engineering dimension \([F]=2\).  Up to a vacuum constant and total derivatives, the gauge-invariant Euclidean scalar bulk terms of engineering dimension at most four are generated by
\[
\operatorname{tr}(F_{\mu\nu}F_{\mu\nu}),
\qquad
\operatorname{tr}(F_{\mu\nu}\widetilde F_{\mu\nu}).
\]
Thus there is no additional CP-even relevant/marginal pure-gauge bulk operator.  The CP-odd coefficient is not removed by dimensional analysis or GP2; fixing the physical \(\bar\theta\) branch is a separate action/measure question.
\end{lemma}
\begin{proof}
Gauge invariance excludes a local mass term \(A_\mu A_\mu\) and non-invariant bare polynomials in the connection.  Simplicity/tracelessness exclude a gauge-invariant term linear in \(F\).  At dimension four, two curvature tensors can be contracted with either the Euclidean metric or the orientation tensor, giving exactly \(\operatorname{tr}(F^2)\) and \(\operatorname{tr}(F\widetilde F)\).  Any further local gauge-invariant term contains additional curvatures or covariant derivatives and therefore has engineering dimension greater than four.  The second dimension-four invariant is independent of the first and its coefficient is not determined by this classification.  This lemma therefore classifies the generic continuum presentation only; it neither creates nor reopens an action coefficient absent from the native generated strong owner.
\end{proof}

\begin{lemma}[Positive resolved-sector weight criterion for a \(\theta\)-factor]
\label{p12:lem:theta-zero-os-selection}
Suppose, as an additional finite action/measure realization, that integer topological sectors \(Q\in\mathbb Z\) are resolved as mutually orthogonal positive blocks with nonnegative base weights and are reweighted only by \(e^{i\theta Q}\).  If a charge-one block has nonzero weight and each reweighted block is required separately to retain a real nonnegative transfer weight, then
\[
\theta=0\pmod{2\pi}.
\]
This is a conditional positivity criterion for that alternative resolved-block realization.  Native strong-CP closure does not depend on it: the v2.78 branch already has trivial sector character and rooted positive quark determinant orientation upstream.
\end{lemma}
\begin{proof}
On a nonzero \(Q=1\) positive block the multiplicative factor is \(e^{i\theta}\).  Requiring that block weight itself to remain real and nonnegative gives \(e^{i\theta}=1\), hence \(\theta=0\bmod2\pi\).  The conclusion uses the stated blockwise-positive representation; without such an action/measure realization, reflection positivity cannot be inferred from the phase factor alone.  In particular, this conditional continuum argument should not be substituted for the native v2.78 action/source theorem.
\end{proof}



% Source Paper VII-C, original line 2057; inherited proof bodies unchanged.
\subsubsection{Native residue UCS: the technical keystone}
\label{p12:sec:qft10-subladder}

QFT10 is the small-field part of the certificate.  Bounded spectral/Wilson data are controlled by Tier-2 estimates; renormalized curvature-density fields are controlled by native residue UCS:
\[
\begin{aligned}
&\text{finite gap}+\text{Tier-2 locality}
 +\text{native residue bound}+\text{summable kernel}\\
&\hspace{3em}\Longrightarrow \text{tree cumulant bound}
\Longrightarrow \text{UCS}.
\end{aligned}
\]
\begin{definition}[Native curvature-residue modes and chart plaquette presentation]
\label{p12:def:normalized-wilson-residue}
Let \(\mathcal I_L\) denote the finite index set of native curvature-residue modes in the Born--Hopf spectral/momentum presentation, with positive quadrature weights \(w_{\alpha,L}\) when the modes are displayed in a reciprocal-position chart.  Let
\[
\mathcal X_{\alpha,L}
\]
be the normalized finite scalar curvature-residue observable associated with \(\alpha\in\mathcal I_L\).  In a large-\(m_x\) local chart, an elementary fundamental \(SU(3)\) plaquette representative \(p\) gives the presentation
\begin{equation}
\label{p12:eq:normalized-wilson-residue}
X_{p,L}:=12\,a_L^{-4}\bigl(1-\operatorname{Re}W_F(p)\bigr).
\end{equation}
This plaquette expression is not the definition of an arbitrary rough lattice variable.  It is only the large-\(m_x\) chart representative of the already defined native spectral residue \(\mathcal X_{\alpha,L}\).  By Lemma~\ref{p12:lem:contact-plaquette-residue}, this chart representative has leading symbol \(F^a_{\mu\nu}F^a_{\mu\nu}\).  The native smeared centered finite field is
\begin{equation}
\label{p12:eq:centered-smeared-wilson-residue}
:\!\mathcal F_L\!:(f)
:=
\sum_{\alpha\in\mathcal I_L}w_{\alpha,L}\,\widehat f_{\alpha,L}\,\mathcal X_{\alpha,L}
-
\omega_L\!\left(
\sum_{\alpha\in\mathcal I_L}w_{\alpha,L}\,\widehat f_{\alpha,L}\,\mathcal X_{\alpha,L}
\right)\mathbf1 ,
\end{equation}
where \(\widehat f_{\alpha,L}\) is the retained spectral/quadrature coefficient of the test function in the chosen large-\(m_x\) presentation.  In a plaquette chart this reads
\[
:\!\mathcal F_L\!:(f)
=
12\sum_{p} f(x_p)\bigl(1-\operatorname{Re}W_F(p)\bigr)
-
\omega_L\!\left(12\sum_p f(x_p)\bigl(1-\operatorname{Re}W_F(p)\bigr)\right)\mathbf1 ,
\]
up to the quadrature convention.  Equivalently, the native chart normalization is \(Z_L=\frac1{12}a_L^4\) for the scalar \(F^2\)-density.  This normalization is fixed by the residue and the Born--Hopf contact/cell normalization, not by data fitting.
\end{definition}

\begin{proposition}[QFT10a--QFT10e: exact finite pieces and the remaining normalization/locality boundary]
\label{p12:prop:qft10a-e-closed}
For the fundamental \(SU(3)\) Wilson residue family, the following statements are exact at finite level: (a) the small-loop coefficient \(1/12\) of Lemma~\ref{p12:lem:contact-plaquette-residue}; (b) the algebraic definition and centering of a chosen smeared residue family; and (c) support in whichever finite local algebra is actually assigned to that residue by the regulator compiler.  By contrast, the claims that \(12a_L^{-4}(1-\operatorname{Re}W_F)\) is uniformly bounded, that this normalization is the unique physical composite-field normalization, and that its local support is uniformly controlled are separate QFT10 hypotheses unless proved by the finite compiler.
\end{proposition}
\begin{proof}
The coefficient \(1/12\) is the representation-theoretic Taylor coefficient proved in Lemma~\ref{p12:lem:contact-plaquette-residue}.  Once a finite residue operator and smearing rule are specified, subtracting its expectation gives exact centering.  Local-algebra membership is exact only for the support assignment supplied by the finite regulator.  Lemma~\ref{p12:lem:native-residue-bound} proves that uniform boundedness requires a uniform symbol envelope, while Theorem~\ref{p12:thm:ucs-implies-qft10} proves that unique continuum-field reconstruction additionally requires mixed-moment Cauchy convergence.  Hence the Taylor coefficient does not by itself close the normalization or continuum-limit problem.
\end{proof}

\begin{definition}[Uniform spectral connected cumulant susceptibility]
\label{p12:def:uniform-small-loop-cumulant-susceptibility}
The native curvature-residue modes satisfy uniform spectral connected cumulant susceptibility if, for every \(n\ge2\),
\begin{equation}
\label{p12:eq:ucs-estimate}
\sup_{L,\alpha_1}
\sum_{\alpha_2,\ldots,\alpha_n\in\mathcal I_L}
 w_{\alpha_2,L}\cdots w_{\alpha_n,L}
\left|
\omega_L^{\mathrm{conn}}
\left(\mathcal X_{\alpha_1,L}\cdots \mathcal X_{\alpha_n,L}\right)
\right|
\le C_n<\infty .
\end{equation}
This is the final Paper-VII-C local-field estimate, abbreviated \emph{spectral UCS}.  In the large-\(m_x\) plaquette chart it has the presentation
\begin{equation}
\label{p12:eq:ucs-plaquette-chart}
\sup_{L,p_1}
\sum_{p_2,\ldots,p_n}
 a_L^{4(n-1)}
\left|
\omega_L^{\mathrm{conn}}
\left(X_{p_1,L}\cdots X_{p_n,L}\right)
\right|
\le C_n<\infty ,
\end{equation}
with \(X_{p,L}\) as in \eqref{p12:eq:normalized-wilson-residue}.  Equation~\eqref{p12:eq:ucs-plaquette-chart} is not a fundamental lattice axiom; it is the reciprocal-position quadrature form of \eqref{p12:eq:ucs-estimate}.
\end{definition}

\begin{remark}[Spectral-residue domain of the plaquette estimate]
\label{p12:rem:ucs-mp-mx-distinction}
Equation~\eqref{p12:eq:ucs-plaquette-chart} must not be read as a statement about an unconstrained rough lattice plaquette variable.  In the Born--Hopf construction the microscopic compact momentum carrier is fixed at
\[
 m_p=1,
\]
and the chart parameter
\[
 m_x=L\to\infty
\]
is a reciprocal-position presentation/window refinement.  Thus \(X_{p,L}=12a_L^{-4}(1-\operatorname{Re}W_F(p))\) is only the large-\(m_x\) chart display of a native spectral curvature-residue mode \(\mathcal X_{\alpha,L}\).  The admissible residue has a smooth compact-momentum symbol before it is displayed in a reciprocal-position chart.  Consequently the trivial rough bound \(|1-\operatorname{Re}W_F(p)|\le2\) is the wrong estimate: it forgets that admissible Born--Hopf plaquettes are bandlimited spectral residues with \(1-\operatorname{Re}W_F(p)=O(a_L^4)\) in the chart presentation.  The estimate therefore applies to the fixed-\(m_p\) spectral-residue class and its large-\(m_x\) chart presentation, not to arbitrary rough plaquette variables.
\end{remark}

\begin{definition}[UCS subladder]
\label{p12:def:ucs-subladder}
The spectral-UCS proof is decomposed into the following subtargets.
\begin{enumerate}[label=\textnormal{(UCS\arabic*)}, leftmargin=5.2em]
\setcounter{enumi}{-1}
\item native curvature-residue compiler and chart plaquette presentation;
\item centering of the scalar curvature-density field;
\item spectral/window equivalence with the large-\(m_x\) chart formula;
\item separated connected decay from the pure-glue gap and Tier-2 locality;
\item two-point contact susceptibility;
\item full scalar UCS;
\item cumulant tree bound;
\item tensor, oriented, and derivative residue extension;
\item irrelevant-residue suppression;
\item distributional reconstruction.
\end{enumerate}
The finite contact expansion and centering parts of UCS0--UCS1 are algebraic.  UCS2--UCS9 require genuine regulator estimates: a uniform residue envelope, actual local support/dynamical locality, a proved uniform gap if exponential mixing is used, a direct all-order tree/cumulant domination estimate, and mixed-moment Cauchy convergence for a unique field limit.  Lemma~\ref{p12:lem:finite-cumulant-connectivity} proves that the former split-mixing shortcut cannot discharge the tree step.
\end{definition}


\begin{longtable}{@{}>{\raggedright\arraybackslash}p{0.22\textwidth} >{\raggedright\arraybackslash}p{0.72\textwidth}@{}}
\caption[UCS acceptance ledger]{UCS acceptance ledger.}\label{p12:tab:ucs-acceptance-ledger}\\
\toprule
UCS ingredient & Source, uniformity, and proof role \\
\midrule
\endfirsthead
\toprule
UCS ingredient & Source, uniformity, and proof role \\
\midrule
\endhead
Native residue norm & \textbf{Source:} Lemma~\ref{p12:lem:native-residue-bound}\par\textbf{Uniform in $L$?:} conditional on a uniform symbol envelope\par\textbf{Proof role:} Compression is bounded by the symbol norm; a uniform envelope must be proved.  \\
Dynamical corridor/locality & \textbf{Source:} Theorem~\ref{p12:thm:euclidean-corridor-transfer}\par\textbf{Uniform in $L$?:} conditional\par\textbf{Proof role:} Requires a finite-propagation/Lieb--Robinson-type dynamical theorem; equal-time locality is insufficient.  \\
Pure-glue transfer gap & \textbf{Source:} Corollary~\ref{p12:cor:paper11-pure-glue-gap-floor}\par\textbf{Uniform in $L$?:} conditional\par\textbf{Proof role:} The exponential factor is available only after \(\gamma_*>0\) is proved for the actual finite sequence.  \\
Contact kernel & \textbf{Source:} Equation~\eqref{p12:eq:residue-cluster-kernel}\par\textbf{Uniform in $L$?:} yes\par\textbf{Proof role:} Handles near-contact plaquette/cell pairs without singular chart insertion.  \\
Kernel summability & \textbf{Source:} Equation~\eqref{p12:eq:residue-kernel-summable}\par\textbf{Uniform in $L$?:} yes\par\textbf{Proof role:} Makes the leaf-by-leaf tree sum finite.  \\
Cumulant tree domination & \textbf{Source:} Lemma~\ref{p12:lem:finite-cumulant-connectivity} and Lemma~\ref{p12:lem:ucs-tree-bound}\par\textbf{Uniform in $L$?:} conditional\par\textbf{Proof role:} Split mixing is insufficient; a direct all-order tree bound must be proved.  \\
Tree reduction & \textbf{Source:} Lemma~\ref{p12:lem:ucs-tree-bound} and Cayley tree count\par\textbf{Uniform in $L$?:} yes\par\textbf{Proof role:} Produces the uniform \(C_n\) in the UCS estimate.  \\
Distributional output & \textbf{Source:} Theorem~\ref{p12:thm:ucs-implies-qft10}\par\textbf{Uniform in $L$?:} yes for fixed test-function order\par\textbf{Proof role:} Converts uniform cumulant susceptibility into local field distributions.  \\
\bottomrule
\end{longtable}
\TableNotes{This ledger is the acceptance test for the native spectral-UCS proof.  Each entry is finite or refinement-uniform before the large-\(m_x\) chart is read as reciprocal position space.}

\begin{corollary}[UCS constants are uniform only after their analytic envelopes are proved]
\label{p12:cor:ucs-constants-finite-carrier}
If the residue symbol envelope \(C_r\), a tree-domination constant \(A_n\), and the kernel norm
\[
I_K:=\sup_{L,\alpha}\sum_\beta w_{\beta,L}K_L(\alpha,\beta)
\]
are finite and independent of \(L\), then the UCS constant may be chosen
\[
C_n=A_n n^{n-2}I_K^{n-1},
\]
which is refinement independent.  Finite-carrier ancestry alone does not prove the uniformity of \(C_r\), \(A_n\), or \(I_K\).
\end{corollary}
\begin{proof}
Lemma~\ref{p12:lem:ucs-tree-bound} gives the displayed formula by leafwise summation and Cayley's count.  Lemma~\ref{p12:lem:native-residue-bound} proves that finite-dimensional support does not imply a uniform coefficient envelope, and Lemma~\ref{p12:lem:finite-cumulant-connectivity} proves that pair/split mixing does not imply a uniform tree constant.  Hence the resulting \(C_n\) is a finite analytic constant only after those envelopes are separately established.
\end{proof}

\begin{lemma}[Compression bound and necessity of a uniform residue envelope]
\label{p12:lem:native-residue-bound}
For every orthogonal projection \(P_L\) and bounded multiplication symbol \(r_{\alpha,L}\), the compressed residue
\begin{equation}
\label{p12:eq:residue-compiler-map}
\mathcal X_{\alpha,L}=P_LM_{r_{\alpha,L}}P_L
\end{equation}
satisfies the exact operator estimate
\begin{equation}
\label{p12:eq:native-residue-uniform-bound}
\|\mathcal X_{\alpha,L}\|\le \|r_{\alpha,L}\|_\infty.
\end{equation}
Consequently a refinement-uniform residue bound follows from, and is reduced to, the explicit coefficient-envelope condition
\begin{equation}
\label{p12:eq:residue-symbol-envelope}
C_r:=\sup_{L,\alpha}\|r_{\alpha,L}\|_\infty<\infty.
\end{equation}
Membership of all \(r_{\alpha,L}\) in one fixed finite-dimensional symbol space does not imply \eqref{p12:eq:residue-symbol-envelope}.
\end{lemma}
\begin{proof}
Because \(P_L\) is an orthogonal projection,
\[
\|P_LM_{r_{\alpha,L}}P_L\|\le\|M_{r_{\alpha,L}}\|=\|r_{\alpha,L}\|_\infty,
\]
which proves the first two claims.  For the last claim, fix any nonzero symbol \(r_0\) in the alleged finite-dimensional residue space and set \(r_L=Lr_0\).  Every \(r_L\) remains in that same finite-dimensional space, while \(\|r_L\|_\infty=L\|r_0\|_\infty\to\infty\).  Thus finite dimension controls equivalence of norms for each vector but not the size of an unbounded sequence of coefficients.  A uniform compiler/source normalization such as \eqref{p12:eq:residue-symbol-envelope} is therefore the minimal missing statement.
\end{proof}

\begin{theorem}[Dynamical-locality no-go and sufficient corridor criterion]
\label{p12:thm:euclidean-corridor-transfer}
Equal-time locality of the finite observable algebras and positivity of the transfer do not imply a lower bound relating spatial chart distance to the number of non-vacuum transfer steps.  A sufficient dynamical hypothesis is the following: there exist \(r_{\rm step},r_0,c_0<\infty\), independent of \(L\), such that one transfer step can enlarge the support of every admissible local excitation by at most \(r_{\rm step}\) outside a fixed contact region.  Under this additional finite-propagation hypothesis, any connected transfer corridor joining supports at chart distance \(d_L(I,J)>r_0\) has length
\begin{equation}
\label{p12:eq:euclidean-corridor-transfer-length}
N_L(I,J)\ge \frac{d_L(I,J)-c_0}{r_{\rm step}}.
\end{equation}
\end{theorem}
\begin{proof}
The sufficient statement is elementary: after \(N\) steps the support can advance by at most \(Nr_{\rm step}\), plus the fixed endpoint/contact thickening \(c_0\); therefore bridging distance \(d_L\) requires \(Nr_{\rm step}+c_0\ge d_L\).  The hypothesis is genuinely dynamical and cannot be derived from equal-time locality alone.  For a counterexample, take a finite spin system whose local algebras on distant sites commute exactly, but whose positive Hamiltonian contains a direct endpoint coupling \(\frac{\lambda}{2}(I-X_1X_N)\).  One Hamiltonian term couples arbitrarily distant supports regardless of their chart distance.  Hence a corridor theorem requires a finite-propagation/Lieb--Robinson-type estimate or an equivalent transfer-kernel locality theorem; spectral support bookkeeping by itself is insufficient.
\end{proof}

\begin{lemma}[Uniform gap plus a proved dynamical corridor gives exponential residue mixing]
\label{p12:lem:residue-cluster-mixing}
Assume (i) the uniform residue envelope \eqref{p12:eq:residue-symbol-envelope}; (ii) a fixed-time transfer decomposition \(T_L=P_L+R_L\) with \(P_L=|\Omega_L\rangle\langle\Omega_L|\) and \(\|R_L\|\le q_*<1\); and (iii) for every separated pair of local residue products \(X_I,X_J\), the connected matrix element admits a transfer-corridor representation with \(N_L(I,J)\) satisfying Theorem~\ref{p12:thm:euclidean-corridor-transfer}.  Then
\begin{equation}
\label{p12:eq:residue-partition-mixing}
|\omega_L(X_IX_J)-\omega_L(X_I)\omega_L(X_J)|
\le B_n e^{-d_L(I,J)/\xi}
\end{equation}
outside the contact region, with \(\xi=r_{\rm step}/[-\log q_*]\) after absorbing fixed endpoint constants.  If the reciprocal-position quadrature has uniformly bounded local volume, the contact-plus-exponential kernel
\begin{equation}
\label{p12:eq:residue-cluster-kernel}
K_L(\alpha,\beta)=\mathbf1_{d_L\le r_0}+e^{-d_L(\alpha,\beta)/\xi}\mathbf1_{d_L>r_0}
\end{equation}
is uniformly summable:
\begin{equation}
\label{p12:eq:residue-kernel-summable}
\sup_{L,\alpha}\sum_\beta w_{\beta,L}K_L(\alpha,\beta)<\infty.
\end{equation}
\end{lemma}
\begin{proof}
Center the two products.  By the assumed corridor representation and \(P_LX_J^\circ\Omega_L=0\),
\[
|\operatorname{Cov}_L(X_I,X_J)|
\le \|X_I^\circ\|\,\|R_L\|^{N_L(I,J)}\,\|X_J^\circ\|.
\]
The residue envelope bounds the product norms uniformly at each fixed order \(n\).  The corridor criterion gives \(N_L\ge(d_L-c_0)/r_{\rm step}\), hence \(q_*^{N_L}\le C e^{-d_L/\xi}\) with \(\xi=r_{\rm step}/[-\log q_*]\).  In the contact region boundedness supplies a finite constant.  Uniform positive quadrature converts the exponential tail to a bounded comparison with \(\int_{\mathbb R^4}e^{-|x|/\xi}d^4x\), while the fixed-radius contact ball has bounded quadrature volume.  This proves the displayed bounds.  The theorem does not assert that the current Paper-VII-B tower satisfies hypotheses (ii)--(iii); those are the sharp gap/dynamical-locality obligations.
\end{proof}

\begin{lemma}[Split mixing is insufficient for a cumulant tree bound]
\label{p12:lem:finite-cumulant-connectivity}
The two-block estimate
\[
|\omega(X_IX_J)-\omega(X_I)\omega(X_J)|
\le B_n\sum_{i\in I,j\in J}K_{ij}
\]
for every nontrivial split does not, by itself, imply a connected-graph or tree-product bound for the full \(n\)-point cumulant with a constant depending only on \(n\), a uniform operator bound, \(B_n\), and \(\sup K\).
\begin{equation}
\label{p12:eq:finite-cumulant-connectivity}
\begin{gathered}
\text{Therefore a direct all-order cumulant/tree-domination}\\
\text{estimate is an independent analytic input.}
\end{gathered}
\end{equation}
\end{lemma}
\begin{proof}
It suffices to give a bounded four-variable counterfamily.  Let \(X_1=X_2=X_3=X_4=Y\), where \(Y\) is Bernoulli with \(\mathbb P(Y=1)=p\).  Every nontrivial split has defect \(k=p(1-p)\).  Define \(K_{12}=1\) and \(K_{ij}=k/3\) for the other five pairs.  Every cut has cross-edge sum at least \(k\), so the split estimate holds with \(B_4=1\), \(\|Y\|\le1\), and \(\sup K=1\).  However
\[
\kappa_4(Y)=p(1-p)(1-6p+6p^2)=O(p).
\]
Every connected graph on four vertices must contain at least two of the small \(O(p)\) edges, because the single order-one edge connects only vertices 1 and 2.  Hence the sum of all connected-graph products is \(O(p^2)\).  The ratio of the cumulant to that graph sum diverges as \(p\downarrow0\).  Thus no uniform graph/tree constant follows from split mixing alone.
\end{proof}

\begin{lemma}[Conditional two-point residue susceptibility]
\label{p12:lem:ucs-two-point-susceptibility}
Under the hypotheses of Lemma~\ref{p12:lem:residue-cluster-mixing}, the native scalar residues satisfy
\begin{equation}
\label{p12:eq:ucs-two-point-susceptibility}
\sup_{L,\alpha}\sum_{\beta\in\mathcal I_L}w_{\beta,L}
|\omega_L^{\rm conn}(\mathcal X_{\alpha,L}\mathcal X_{\beta,L})|
\le C_2<\infty.
\end{equation}
For the current Paper-VII-B export this is therefore conditional on the proved uniform gap/residue envelope/dynamical corridor hypotheses, not an unconditional finite-carrier result.
\end{lemma}
\begin{proof}
Set \(n=2\) in Lemma~\ref{p12:lem:residue-cluster-mixing}.  Its exponential/contact kernel gives
\[
|\omega_L^{\rm conn}(\mathcal X_{\alpha,L}\mathcal X_{\beta,L})|
\le B_2K_L(\alpha,\beta).
\]
Summing against \(w_{\beta,L}\) and applying \eqref{p12:eq:residue-kernel-summable} gives \(C_2\le B_2I_K\).  No additional inference is used.
\end{proof}

\begin{lemma}[Direct tree domination implies uniform connected-cumulant summability]
\label{p12:lem:ucs-tree-bound}
Assume that for each fixed \(n\ge2\) the residue cumulants satisfy the direct tree-domination estimate
\begin{equation}
\label{p12:eq:ucs-tree-bound}
|\omega_L^{\rm conn}(\mathcal X_{\alpha_1,L}\cdots\mathcal X_{\alpha_n,L})|
\le A_n\sum_{T\in\mathcal T_n}\prod_{(ij)\in T}K_L(\alpha_i,\alpha_j),
\end{equation}
where \(A_n\) is independent of \(L\) and the nonnegative kernel obeys
\[
I_K:=\sup_{L,\alpha}\sum_\beta w_{\beta,L}K_L(\alpha,\beta)<\infty.
\]
Then spectral UCS holds with
\[
C_n\le A_n n^{n-2}I_K^{n-1}.
\]
\end{lemma}
\begin{proof}
For each labeled tree, sum successively over leaves other than the fixed root \(\alpha_1\).  Every leaf sum contributes at most \(I_K\), so a tree with \(n-1\) edges contributes at most \(I_K^{n-1}\).  Cayley's formula gives \(|\mathcal T_n|=n^{n-2}\).  Summing the tree bound proves the stated UCS constant.  By Lemma~\ref{p12:lem:finite-cumulant-connectivity}, the direct tree-domination premise cannot be replaced by the weaker split-mixing estimate used in the previous version.
\end{proof}

\begin{theorem}[Sharp spectral-UCS criterion for the native residue family]
\label{p12:thm:native-spectral-ucs}
Let \(\mathcal X_{\alpha,L}\) be the native pure-glue residue family.  Suppose
\begin{enumerate}[label=\textnormal{(\roman*)}]
\item the normalized symbol envelope \eqref{p12:eq:residue-symbol-envelope} is finite;
\item the connected cumulants obey the direct tree-domination estimate \eqref{p12:eq:ucs-tree-bound}; and
\item the corresponding kernel satisfies \eqref{p12:eq:residue-kernel-summable}.
\end{enumerate}
Then uniform spectral connected cumulant susceptibility holds:
\begin{equation}
\label{p12:eq:ucs-closed-estimate}
\sup_{L,\alpha_1}\sum_{\alpha_2,\ldots,\alpha_n}
 w_{\alpha_2,L}\cdots w_{\alpha_n,L}
|\omega_L^{\rm conn}(\mathcal X_{\alpha_1,L}\cdots\mathcal X_{\alpha_n,L})|
\le C_n<\infty.
\end{equation}
For the current Paper-VII-B finite export, conditions (i)--(iii) are continuum analytic obligations; they are not consequences of fixed \(m_p=1\), finite dimensionality, or the existence of a positive finite transfer.
\end{theorem}
\begin{proof}
Condition (i) makes every fixed-order residue product uniformly bounded.  Conditions (ii)--(iii) are exactly the hypotheses of Lemma~\ref{p12:lem:ucs-tree-bound}, which yields \(C_n\le A_n n^{n-2}I_K^{n-1}\).  The last sentence follows from the counterexamples in Lemmas~\ref{p12:lem:native-residue-bound}, \ref{p12:thm:euclidean-corridor-transfer}, and \ref{p12:lem:finite-cumulant-connectivity}: each removes one former shortcut from the current finite data to the UCS hypotheses.
\end{proof}

\begin{corollary}[Finite residue-type stability of a proved UCS bound]
\label{p12:cor:tensor-oriented-ucs}
Suppose a finite family of scalar, oriented, tensorial, derivative, or irrelevant residue types is obtained from a UCS-controlled base family by linear maps/contractions whose operator norms are uniformly bounded in \(L\), and suppose the direct tree-domination/kernel bound of Theorem~\ref{p12:thm:native-spectral-ucs} holds for the enlarged finite family.  Then the same spectral-UCS conclusion holds for every residue type, with constants multiplied only by the finite uniform coefficient bounds.  The extension is not automatic when differentiation or regulator normalization introduces \(L\)-dependent coefficients.
\end{corollary}
\begin{proof}
Connected cumulants are multilinear.  Applying a uniformly bounded finite linear combination/contraction multiplies each \(n\)-point cumulant by at most the product of the corresponding coefficient/operator norms.  The finite number of residue labels can be absorbed into the constant \(A_n\), after which Lemma~\ref{p12:lem:ucs-tree-bound} gives the same summability.  If the coefficient norms grow with \(L\), no uniform conclusion follows, so that envelope is an explicit necessary premise.
\end{proof}


\textbf{UCS consequence.}
The finite carrier identifies the residue type, but spectral UCS follows only after the explicit uniform envelope, direct tree-domination, and kernel-summability criteria of Theorem~\ref{p12:thm:native-spectral-ucs} have been proved.  The large-\(m_x\) plaquette formula is a chart representative; its regulator-normalized boundedness and unique continuum limit are QFT10 statements, not consequences of finite dimensionality.


\begin{theorem}[UCS plus residue-moment Cauchy convergence implies QFT10]
\label{p12:thm:ucs-implies-qft10}
Assume uniform spectral UCS for a regulator-normalized residue family and assume, in addition, that every mixed moment
\[
\omega_L\!\left(A_{0,L}:\!\mathcal F_L\!:(f_1)\cdots:\!\mathcal F_L\!:(f_n)B_{0,L}\right)
\]
is Cauchy and refinement-path independent for bounded local \(A_0,B_0\).  Then the connected and ordinary moments have uniform distributional seminorm bounds, their limits are unique, and the limiting field is an operator-valued distribution affiliated with the reconstructed local net.  Without the additional residue-moment Cauchy property, UCS gives tightness/subsequential compactness but not uniqueness.
\end{theorem}
\begin{proof}
Expand each connected smeared moment in the spectral residue coefficients.  UCS bounds the sum uniformly by a finite product of test-function seminorms and the finite weighted support volume.  The moment--cumulant partition formula transfers these bounds to ordinary moments, giving distributional compactness.  The additional Cauchy assumption now fixes every mixed scalar limit independently of subsequence and refinement path; this is the step that QFT4 on the bounded cylinder algebra alone cannot supply for an \(L\)-dependent unbounded residue family.  Local support of the representatives gives affiliation with the reconstructed local net, and reflection positivity/covariance pass through the same mixed-moment limits.  For necessity of the extra convergence condition, take centered bounded residue variables with law \(\pm1\) at even \(L\) and \(\pm\sqrt2\) at odd \(L\): all fixed-order cumulants are uniformly bounded, while the second moments alternate between 1 and 2, so two subsequential field limits remain.
\end{proof}


\textbf{QFT10 consequence.}
The contact expansion and centering are finite algebraic facts.  Uniform residue normalization, all-order connected-moment bounds, and mixed-moment Cauchy convergence are separate analytic statements.  Theorem~\ref{p12:thm:ucs-implies-qft10} proves local-field reconstruction only when those statements are supplied.  Thus QFT10 remains an executable continuum obligation for the current Paper-VII-B pure-glue residue family rather than a consequence of finite support.




\begin{longtable}{>{\raggedright\arraybackslash}p{0.23\textwidth} >{\raggedright\arraybackslash}p{0.37\textwidth} >{\raggedright\arraybackslash}p{0.28\textwidth}}
\caption[UCS-to-local-field discharge ledger]{UCS-to-local-field discharge ledger.}\label{p12:tab:ucs-local-field-discharge}\\
\toprule
QFT10 block & Discharged by & Output \\
\midrule
\endfirsthead
\toprule
QFT10 block & Discharged by & Output \\
\midrule
\endhead
Contact/Taylor coefficient & Lemma~\ref{p12:lem:contact-plaquette-residue} & The coefficient \(1/12\) is exact for the smooth fundamental \(SU(3)\) small-loop expansion. \\
Uniform residue normalization & Lemma~\ref{p12:lem:native-residue-bound} & Requires a proved uniform symbol/coefficient envelope; finite dimension alone is insufficient. \\
Uniform connected moments & Theorem~\ref{p12:thm:native-spectral-ucs} & Conditional on direct tree domination and a summable kernel for the actual residue family. \\
Ordinary moment tightness & Moment--cumulant expansion & Follows from uniform connected-moment seminorm bounds. \\
Limit uniqueness & Theorem~\ref{p12:thm:ucs-implies-qft10} & Requires mixed residue moments to be Cauchy/path independent; QFT4 on bounded observables alone is insufficient. \\
Local-net affiliation & Local support plus the reconstructed OS net & Follows only after the regulator compiler supplies compatible local support and the mixed limits exist. \\
Tensor/oriented extensions & Corollary~\ref{p12:cor:tensor-oriented-ucs} & Conditional on uniformly bounded residue-type maps and the same tree/kernel control; stress-tensor normalization also needs the finite action/metric response. \\
\bottomrule
\end{longtable}
\TableNotes{This table separates exact finite contact/centering facts from the analytic QFT10 obligations.  Uniform residue normalization, all-order connected-moment control, mixed-moment Cauchy convergence, and tensor/action normalization must be proved for the actual finite sequence.}

\begin{definition}[Yang--Mills continuum certificate obligations]
\label{p12:def:ym-continuum-obligations}
The pure \(G\)-Wilson datum satisfies the Yang--Mills continuum package if the following obligations hold.
\begin{enumerate}[label=\textnormal{(BH\arabic*)}, leftmargin=3.1em]
\item \textbf{Compact-simple pure-gauge compiler.}  For every compact simple \(G\), construct the pure connection/curvature compiler \(\mathfrak C_{L,G}^{\YM}\) with no matter modules, producing a datum \(\DG\) of Definition~\ref{p12:def:pure-g-wilson-datum}.
\item \textbf{Uniform reflection positivity.}  Finite reflection positivity is supplied by the positive transfer presentation; in the continuum package it is the passage of this nonnegative form through BH6/QFT4.
\item \textbf{Contact-normalized plaquette residue.}  Use Lemma~\ref{p12:lem:contact-plaquette-residue} and the Born--Hopf contact/cell normalization to fix the quadratic curvature residue.  The coefficient is algebraic/geometric, not fitted.
\item \textbf{Dimension-\(\le4\) action classification and irrelevant remainder control.}  Lemma~\ref{p12:lem:no-hidden-relevant-gauge-branch} leaves exactly the CP-even \(F^2\) term and the independent CP-odd \(F\widetilde F\) term, besides the vacuum constant.  A physical CP branch therefore belongs to the finite action/measure package.  Higher-dimensional remainders vanish only under the coefficient and residue-moment bounds of Theorem~\ref{p12:thm:tier3-irrelevant-suppression}.
\item \textbf{Asymptotic local chart recovery.}  Proposition~\ref{p12:prop:tier2-quadrature} derives positive finite quadrature and Lemma~\ref{p12:lem:tier2-flat-chart-covariance} derives the geometric flat-chart exhaustion.  The interacting expectation-value covariance defect must still be shown to vanish for the actual tower.
\item \textbf{Canonical Wilson-state limit.}  Prove that \(\omega_{L,G}(A)\) is a Cauchy net for every bounded local Wilson-loop polynomial \(A\), with a refinement-independent limit.
\item \textbf{Uniform continuum mass-gap transfer.}  For the Paper~VII-B \(SU(3)_c\) pure-glue branch the finite gap sequence is known but its uniform positive infimum is not; for any compact simple \(G\), prove a refinement-uniform pure-glue singlet gap
\[
\inf_L(E_{1,L,G}-E_{0,L,G})\ge\Delta_{\ast,G}>0
\]
and show that it survives OS reconstruction.
\item \textbf{Uniform continuum nontriviality witness.}  Prove a witness excluding the empty and generalized-free limits.  A sufficient local criterion is a nonzero limiting connected/truncated correlator of order \(n\ge3\).  A positive Wilson area exponent alone is not sufficient.
\item \textbf{OS regularity and Wightman reconstruction.}  Prove OS axioms with sufficient regularity/temperedness to reconstruct a local-chart OS/Wightman theory with the imported Lorentzian/Poincar\'e presentation on compact windows.
\item \textbf{Regulator-normalized curvature-residue field reconstruction.}  Prove a uniform residue envelope, distributional moment bounds (for example via a valid UCS/tree estimate), and mixed-moment Cauchy/path independence.  Tensor/stress residues additionally require the finite action or metric-response normalization.  Only then do the distributional limits define the claimed \(F^2\), \(F\widetilde F\), and \(T_{\mu\nu}\) local fields.
\end{enumerate}
\end{definition}

\begin{proposition}[BH obligations imply QFT obligations]
\label{p12:prop:bh-implies-qft}
For a pure \(G\)-Wilson datum, BH1--BH10 imply QFT1--QFT10 for the bounded Wilson-loop local algebra and the small-loop gauge-invariant field family.
\end{proposition}
\begin{proof}
BH1 supplies the finite refinement datum and Wilson local cylinder algebra.  BH5 supplies chart and Euclidean recovery, hence QFT1 and QFT6.  Bounded normalized Wilson loops give QFT3.  BH6 gives the canonical Cauchy property QFT4.  BH2 gives QFT5.  BH8 gives QFT8.  BH7 gives the positive-gap branch of QFT9.  BH9 gives OS regularity, and BH10 gives local-field reconstruction.  BH3 supplies the small-loop quadratic coefficient, while BH4 requires the separately typed action/measure branch and irrelevant-term control needed to identify the Yang--Mills universality class.  Thus the implication is logical once the BH obligations are actually established; it is not a proof that the current Paper-VII-B export satisfies them.
\end{proof}


% Source Paper VII-C, original line 2443; inherited proof bodies unchanged.
% Former logical part: Yang--Mills Continuum Handoff and Refinement Compatibility; integrated into the new thematic argument.

\noindent\textbf{Orientation.}
This part converts the Wilson/residue and UCS machinery into the native \(SU(3)_c\) Yang--Mills continuum handoff, then records renormalization as compatibility of finite refinements.  It is the terminal continuum bridge from the Paper~VII-B finite pure-glue datum to the Paper~VII-C OS/Wightman presentation.


% Source Paper VII-C, original line 2448; inherited proof bodies unchanged.
\subsection{Conditional \(SU(3)_c\) Yang--Mills continuum handoff}
\label{p12:sec:ym-identification}


\begin{table}[htbp]
\centering
\small
\begin{tabularx}{\textwidth}{@{}
>{\raggedright\arraybackslash}p{0.55\textwidth}
>{\raggedright\arraybackslash}X
@{}}
\toprule
\textbf{Yang--Mills closure step} & \textbf{Discharged by} \\
\midrule
Yang--Mills continuum certificate obligations
& Definition~\ref{p12:def:ym-continuum-obligations}. \\

Born--Hopf obligations imply QFT obligations
& Proposition~\ref{p12:prop:bh-implies-qft}. \\

Wilson-state compactness and BH/QFT-consistency uniqueness under BH1--BH6
& Lemma~\ref{p12:lem:wilson-state-compactness}. \\

Yang--Mills Schwinger--Dyson identification
& Lemma~\ref{p12:lem:ym-sd-identification}. \\

Yang--Mills nontriviality
& Lemma~\ref{p12:lem:ym-nontriviality} and
Corollary~\ref{p12:cor:wilson-algebra-nontriviality}. \\

Tier-3 irrelevant suppression
& Theorem~\ref{p12:thm:tier3-irrelevant-suppression}. \\

Tier-3 Schwinger--Dyson and local fields
& Theorems~\ref{p12:thm:tier3-sd-from-ucs}
and~\ref{p12:thm:tier3-local-fields-from-ucs}. \\

Conditional \(SU(3)_c\) Yang--Mills continuum handoff
& Theorem~\ref{p12:thm:native-bh-ym-continuum}. \\

Compatibility with the full Paper~VII-B compiler
& Corollary~\ref{p12:cor:p11-final-export-compatibility}. \\
\bottomrule
\end{tabularx}
\caption[Yang--Mills identification and tier closure]{Yang--Mills identification and tier closure.}
\label{p12:tab:ym-identification-tier-closure}
\label{p12:tab:ym-identification-ledger}
\TableNotes{The row order separates general Wilson-state compactness from the
native \(SU(3)_c\) handoff.  Compactness gives existence of cluster states;
refinement-independent uniqueness is supplied by BH/QFT local Cauchy consistency,
not by compactness alone.  General-\(G\), gapless, and external adjudication
questions are scope extensions, not hidden premises of the Paper~VII-B finite
QCD closure.}
\end{table}


% Source Paper VII-C, original line 2503; inherited proof bodies unchanged.
\subsubsection{Yang--Mills obligations and Wilson-state compactness}
\label{p12:subsec:ym-obligations-wilson-compactness}

\begin{lemma}[Wilson-state compactness and BH/QFT-consistency uniqueness under BH1--BH6]
\label{p12:lem:wilson-state-compactness}
Assume BH1--BH6.  Then the finite Wilson states \(\omega_{L,G}\) have a weak-\(*\) continuum limit \(\omega_G\) on the bounded local Wilson-loop algebra, and this limit is unique because of the BH/QFT local Cauchy consistency, not because of compactness alone.
\end{lemma}
\begin{proof}
\ProofPlan{We separate existence from uniqueness.  Compactness gives cluster states; BH3--BH6 force all cluster states to agree on every bounded local Wilson observable.}

\begin{proofsteps}
\item \textbf{Compactness gives subsequential limits.}
For every bounded local Wilson-loop algebra \(\mathcal W(\mathcal O)\), normalized Wilson loops obey
\[
\|W_{\rho,C}^{(L)}\|\le1 .
\]
Hence each finite \(\omega_{L,G}\) is a state on the \(C^*\)-completion of \(\mathcal W(\mathcal O)\).  The state space is weak-\(*\) compact, so every refinement subnet has a weak-\(*\) cluster state on \(\mathcal W(\mathcal O)\).  This step proves existence of subsequential local limits only.

\item \textbf{BH consistency identifies the possible local limit.}
BH3--BH4 identify the only relevant local pure-gauge bulk term as the Yang--Mills quadratic curvature term, with the contact-normalized coefficient.  Thus two cluster states cannot differ by choosing a different relevant or marginal local pure-gauge action inside the same BH1--BH6 datum.  BH5 makes chart changes and Euclidean motions equivalent up to irrelevant errors on compact local windows, so cluster states obtained through different local chart presentations evaluate the same local Wilson observable after passing to the refinement limit.

\item \textbf{BH6 supplies the uniqueness mechanism.}
Let \(A\in\mathcal W(\mathcal O)\) be a bounded local Wilson-loop polynomial.  BH6 says that the net \(\omega_{L,G}(A)\) is Cauchy under admissible refinements.  Therefore \(\omega_{L,G}(A)\) has a refinement-independent scalar limit.  If \(\omega'\) and \(\omega''\) are two weak-\(*\) cluster states, both are obtained as subnet limits of the same Cauchy scalar net, hence
\[
\omega'(A)=\lim_L\omega_{L,G}(A)=\omega''(A).
\]
Since this holds for every \(A\) in every bounded local Wilson algebra, all local cluster states agree.  Compatibility under inclusions of local regions then gives a single continuum Wilson state \(\omega_G\) on the bounded local Wilson-loop algebra.
\end{proofsteps}
\end{proof}

\begin{lemma}[Finite action-measure realization criterion\\for Yang--Mills Schwinger--Dyson identities]
\label{p12:lem:ym-sd-identification}
Suppose the actual finite pure-glue vacuum/source member admits a
matched complete Haar/action-measure realization from its licensed native
datum
\looseness=1
\[
\omega_L(A)=Z_L^{-1}\int A(U)e^{-S_L(U)}\prod_e dU_e,
\]
with differentiable gauge-invariant \(S_L\), and suppose the differentiated action/observable insertions have QFT10 mixed-moment limits.  If the local action expansion has leading term
\[
S_{\YM,G}=\frac1{2g_G^2}\int\langle F_{\mu\nu},F_{\mu\nu}\rangle_{\mathfrak g}\,d^4x
\]
and all higher-dimensional remainders vanish in the controlled correlators, then the continuum limits satisfy the Yang--Mills Schwinger--Dyson identities.  The finite source/action/transfer lineage and readout descendants are already supplied by Definitions~\ref{p10:def:complete-finite-qcd-source-transfer-datum}, \ref{p11:def:imported-paper-x-core-datum}, Theorem~\ref{p10:thm:master-finite-nonperturbative-qcd-closure}, and Corollary~\ref{p11:cor:v147-generated-readout-descendants}. The licensed heat/Haar representative is Theorem~\ref{p10:thm:canonical-haar-heat-representative}. Matching it to the actual consumed vacuum/source member, its finite invariance, and differentiated insertion limits remain the displayed hypotheses; they are not a second universal action/measure owner.
\end{lemma}
\begin{proof}
Haar invariance at finite level gives the exact integration-by-parts identity
\[
0=\int \nabla_e(Ae^{-S_L})\prod_e dU_e
=\langle\nabla_eA-A\nabla_eS_L\rangle_L.
\]
By the stated QFT10 convergence, both differentiated terms have unique distributional limits.  The assumed local action expansion identifies the dimension-four contribution with the Yang--Mills functional, while controlled irrelevant suppression removes the remaining terms.  Taking the limit therefore gives \(\omega(\delta A)-\omega(A\,\delta S_{\YM,G})=0\).  Without matching the licensed complete action/measure to this actual transfer/vacuum source member, the transfer operator alone does not identify the probability density to differentiate. This is a claim-specific member/source matching and identity certificate. The historical `X-ACT` constructor wording is superseded on the licensed finite native branch by the preceding providers; it does not reopen a universal finite action/measure gap.
\end{proof}

\begin{lemma}[Non-Gaussian connected-cumulant criterion on the stated observable family]
\label{p12:lem:ym-nontriviality}
Let a reconstructed gauge-invariant local field or bounded Wilson observable family have a connected/truncated continuum correlator of some order \(n\ge3\) that is nonzero.  Then that continuum sector is neither empty nor generalized-free/Gaussian on that observable family.  A positive Wilson area exponent by itself does not imply this conclusion.
\end{lemma}
\begin{proof}
In an empty/trivial sector every centered local observable vanishes in correlation functions.  In a generalized-free/Gaussian family, Wick's rule gives zero connected/truncated correlators of every order \(n\ge3\).  Hence a nonzero higher connected correlator excludes both possibilities.  To show that an area exponent is insufficient, take independent centered Gaussian fluxes \(\phi_p\) of variance \(2\sigma\) and define \(W(C)=\exp(i\sum_{p\in A(C)}\phi_p)\).  Gaussian characteristic functions give exactly \(\mathbb E W(C)=e^{-\sigma A(C)}\), while all flux cumulants above second order vanish.  Therefore an area law may coexist with a Gaussian underlying field and cannot serve alone as the claimed non-Gaussian witness.
\end{proof}

\begin{corollary}[Continuum nontriviality from a surviving higher connected Wilson/curvature cumulant]
\label{p12:cor:wilson-algebra-nontriviality}
Assume QFT4/QFT10 convergence for a bounded Wilson or curvature-residue family and suppose that for some \(n\ge3\) and fixed local test data
\[
\lim_L\omega_L^{\rm conn}(A_{1,L}\cdots A_{n,L})=c\ne0.
\]
Then the reconstructed continuum family is nonempty and non-generalized-free on that observable family; for a nonlinear composite this does not alone classify the underlying dynamics as interacting, by Proposition~\ref{p12:prop:v2210-lit-composite-cumulant}. No position-space area law or string-tension pairing is required for this local nontriviality criterion.
\end{corollary}
\begin{proof}
The convergence hypothesis gives a continuum connected correlator equal to \(c\ne0\).  Lemma~\ref{p12:lem:ym-nontriviality} then excludes the empty and generalized-free alternatives.  Since the witness is purely local/correlational, the separate strict positive-area theorem is irrelevant to this conclusion; it is needed only for a physical area/string-tension claim.
\end{proof}



% Source Paper VII-C, original line 2575; inherited proof bodies unchanged.
\subsubsection{Tier-3 consequences of spectral UCS}
\label{p12:subsec:tier3-ucs-consequences}

Theorem~\ref{p12:thm:native-spectral-ucs} supplies spectral UCS for native compact-momentum residues.  The finite irrelevant-insertion estimate below does not use QFT4; QFT4 enters only when one passes these finite estimates to a unique continuum state.  The following results derive irrelevant-remainder suppression, the Yang--Mills Schwinger--Dyson identities, and local reconstruction of the standard gauge-invariant curvature composites from that estimate.

\begin{definition}[UCS-stable curvature-residue family]
\label{p12:def:ucs-stable-curvature-residue-family}
A finite family of native local residue modes
\[
\mathcal R_{a,\alpha,L}
\]
is UCS-stable if the spectral UCS estimate \eqref{p12:eq:ucs-estimate} holds with each \(\mathcal X_{\alpha,L}\) replaced by any member of the family, and by finite products/contractions generated by the family, with constants depending only on the finite list of labels.  The required family for pure gauge theory consists of the scalar curvature residue, the oriented pseudoscalar residue, tensorial anisotropic curvature residues, and the differentiated residue modes appearing in the finite Haar integration-by-parts identities.
\end{definition}

\begin{theorem}[Spectral UCS suppresses irrelevant remainders]
\label{p12:thm:tier3-irrelevant-suppression}
Assume the Tier-2 spectral-window estimates and UCS-stability for the local residue modes in the pure-gauge action expansion.  Write
\[
S_{L,G}=S_{\YM,L}+c_{0,L}\mathbf 1+R_L,
\qquad
R_L=\sum_{d>4}a_L^{d-4}\sum_{\alpha}w_{\alpha,L}c_{d,\alpha,L}\mathcal O_{d,\alpha,L}.
\]
If the irrelevant expansion is finite at each fixed approximation order, or if
\[
\sum_{d>4}|c_d|C_{A,B,d}<\infty
\]
for the UCS constants controlling bounded local insertions \(A_L,B_L\), then
\[
\omega_L(A_LR_LB_L)\longrightarrow 0
\]
for all bounded local spectral-window observables \(A,B\).
\end{theorem}
\begin{proof}
\ProofPlan{We show that each irrelevant local term is suppressed by its positive power of the refinement scale, then sum finite or absolutely summable families of such terms.}
\begin{proofsteps}
\item \textbf{Single irrelevant dimension-\(d\) term.}
For one dimension-\(d\) term,
\[
R_{d,L}=a_L^{d-4}\sum_\alpha w_{\alpha,L}c_{d,\alpha,L}\mathcal O_{d,\alpha,L}.
\]

\item \textbf{UCS insertion bound.}
UCS-stability gives a uniform local insertion bound
\[
\left|\omega_L\!\left(A_L\sum_\alpha w_{\alpha,L}c_{d,\alpha,L}\mathcal O_{d,\alpha,L}B_L\right)\right|
\le C_{A,B,d}.
\]

\item \textbf{Vanishing by \(a_L^{d-4}\).}
Therefore
\[
|\omega_L(A_LR_{d,L}B_L)|\le a_L^{d-4}C_{A,B,d}\to0
\]
because \(d>4\).

\item \textbf{Finite or summable expansion closure.}
A finite irrelevant expansion follows by summing finitely many terms.  The stated absolute summability condition gives dominated convergence for an infinite local expansion.
\end{proofsteps}
\end{proof}

\begin{corollary}[Uniform action interpolation removes finite regulator-member dependence]
\label{p12:cor:uniform-action-interpolation-universality}
Let $S_{L,s}=S_L+sR_L$, $0\le s\le1$, be a family of finite positive
action measures on the same canonical Haar reference space.  If the finite
irrelevant-insertion estimate of Theorem~\ref{p12:thm:tier3-irrelevant-suppression}
holds uniformly in $s$ for $A$ and for the unit insertion, then
\[
 \boxed{
 \sup_{0\le s\le1}|\operatorname{Cov}_{L,s}(A,R_L)|\longrightarrow0}
\]
and therefore
\[
 \boxed{|\omega_{L,1}(A)-\omega_{L,0}(A)|\longrightarrow0}
\]
for every fixed bounded local cylinder observable $A$.
\end{corollary}
\begin{proof}
The exact finite identity
$\frac d{ds}\omega_{L,s}(A)=-\operatorname{Cov}_{L,s}(A,R_L)$
combined with
$|\operatorname{Cov}(A,R)|\le|\omega(AR)|+\|A\|\,|\omega(R)|$
gives the result after integration in $s$.
\end{proof}


\begin{corollary}[QFT4 and source limits are inherited across a uniformly irrelevant regulator deformation]
\label{p12:cor:v185-universality-inherits-continuum-state}
Let $S_L^{(1)}=S_L^{(0)}+R_L$ be two positive regulator representatives on
the same canonical Haar reference space.  Suppose the uniform interpolation
hypotheses of Corollary~\ref{p12:cor:uniform-action-interpolation-universality}
hold for every fixed bounded local observable in the claimed cylinder class.
If the $S^{(0)}$ sequence satisfies QFT4 with continuum state $\omega_\infty$,
then the $S^{(1)}$ sequence satisfies QFT4 with the same continuum state.
If, in addition, the interpolation estimate is uniform for the source
exponentials of Theorem~\ref{p12:thm:v185-uniform-source-effective-action-stability}
on a common polydisc, then the two regulator sequences have the same limiting
projective Schwinger functional and the same source-owned connected cumulants.
\end{corollary}
\begin{proof}
For a fixed bounded $A$ and $M\succeq L$, insert the $S^{(0)}$ expectations:
\[
 \begin{aligned}
 |\omega_M^{(1)}(A_M)-\omega_L^{(1)}(A_L)|
 &\le |\omega_M^{(1)}(A_M)-\omega_M^{(0)}(A_M)|\\
 &\quad+|\omega_M^{(0)}(A_M)-\omega_L^{(0)}(A_L)|\\
 &\quad+|\omega_L^{(0)}(A_L)-\omega_L^{(1)}(A_L)|.
 \end{aligned}
\]
The first and third terms vanish cofinally by the uniform interpolation
corollary and the middle term by QFT4, proving the common state limit and path
independence.  Applying the same triangle inequality uniformly to the bounded
source exponentials gives equality of the limiting $\widehat Z$; the source
cumulant statement follows from
Theorem~\ref{p12:thm:v185-uniform-source-effective-action-stability}.
\end{proof}

\begin{theorem}[Action-measure plus residue convergence identifies the Yang--Mills Schwinger--Dyson limit]
\label{p12:thm:tier3-sd-from-ucs}
Assume the finite action-measure realization of Lemma~\ref{p12:lem:ym-sd-identification}, the contact-normalized quadratic curvature expansion, QFT4, and QFT10 convergence for all differentiated residue insertions.  If the higher-dimensional action remainders satisfy Theorem~\ref{p12:thm:tier3-irrelevant-suppression}, then the continuum finite identities converge to the Euclidean Yang--Mills Schwinger--Dyson identities.
\end{theorem}
\begin{proof}
This is Lemma~\ref{p12:lem:ym-sd-identification} with the irrelevant remainder controlled by Theorem~\ref{p12:thm:tier3-irrelevant-suppression}.  The crucial finite identity is Haar integration by parts in the independently supplied action measure; QFT10 supplies convergence of the differentiated insertions.  No Schwinger--Dyson conclusion is drawn from the transfer operator alone.
\end{proof}

\begin{theorem}[Regulator-normalized residue convergence reconstructs gauge-invariant local fields]
\label{p12:thm:tier3-local-fields-from-ucs}
Assume QFT10 in the strengthened form of Theorem~\ref{p12:thm:ucs-implies-qft10} for the scalar, oriented, and tensorial regulator-normalized curvature-residue families, with the tensor/stress representative fixed by the finite action or metric-response variation.  Then the corresponding limits reconstruct local operator-valued distributions representing
\[
:F_{\mu\nu}^aF_{\mu\nu}^a:,
\qquad
F_{\mu\nu}^a\widetilde F_{\mu\nu}^a,
\qquad
T_{\mu\nu},
\]
affiliated with the reconstructed gauge-invariant local net.
\end{theorem}
\begin{proof}
For each residue type, QFT10 gives unique mixed distributional moments with local support.  The contact expansion identifies the scalar leading symbol, orientation identifies the pseudoscalar contraction, and the independently fixed finite action/metric-response variation identifies the tensor representative.  OS reconstruction then produces the corresponding Wightman distributions affiliated with the local net.  Uniform cumulant bounds alone are not used to infer uniqueness; the residue-moment Cauchy condition is essential.
\end{proof}

\begin{corollary}[Tier-3 implication ledger]
\label{p12:cor:tier3-closure}
The Tier-3 conclusions separate into three independent implications: (i) irrelevant remainders vanish under the coefficient/UCS hypotheses of Theorem~\ref{p12:thm:tier3-irrelevant-suppression}; (ii) Schwinger--Dyson identities follow only with the finite action-measure realization and differentiated-residue convergence of Theorem~\ref{p12:thm:tier3-sd-from-ucs}; and (iii) local gauge-invariant composite fields follow only with the regulator-normalized residue-moment convergence of Theorem~\ref{p12:thm:tier3-local-fields-from-ucs}.  None of (ii)--(iii) is an automatic consequence of finite dimensionality or of a positive finite transfer.
\end{corollary}
\begin{proof}
Each item is exactly the theorem cited in the statement.  Their hypotheses are logically independent in the current finite export: an action measure is not determined by a transfer operator, and uniform distributional boundedness does not determine a unique residue limit.  Hence the three implications must remain separately typed.
\end{proof}



\begin{longtable}{>{\raggedright\arraybackslash}p{0.13\textwidth} >{\raggedright\arraybackslash}p{0.41\textwidth} >{\raggedright\arraybackslash}p{0.35\textwidth}}
\caption[\texorpdfstring{BH1--BH10 status census for the repaired Paper-VII-B \(SU(3)_c\) handoff}{BH1--BH10 status census for the repaired Paper-VII-B SU(3)c handoff}]{BH1--BH10 status census for the repaired Paper-VII-B \(SU(3)_c\) handoff.}\label{p12:tab:bh1-bh10-su3-discharge-census}\\
\toprule
BH item & Required object & Current status/source \\
\midrule
\endfirsthead
\toprule
BH item & Required object & Current status/source \\
\midrule
\endhead
BH1 & finite pure-glue singlet Hamiltonian/transfer and bounded Wilson algebra & \textbf{Finite input available} from repaired Paper~VII-B. \\
BH2 & finite reflection-positive form and passage to the limit & \textbf{Finite positivity available conditionally}; continuum positivity is derived once QFT4 reflected products converge. \\
BH3 & smooth small-loop/contact coefficient & \textbf{Derived coefficient}: Lemma~\ref{p12:lem:contact-plaquette-residue}; identification of an actual local connection/action chart remains domain dependent. \\
BH4 & dimension-\(\le4\) action branch and irrelevant control & \textbf{Structurally narrowed}: the only local marginal directions are \(F^2\) and \(F\widetilde F\); product Haar and the common-heat gauge normalization are derived on the source-native branch.  The remaining obligation is the quasilocal/irrelevant effective-action stability estimate. \\
BH5 & local chart geometry and interacting flat covariance & \textbf{Geometry and canonical symmetry derived}; on the equivariant full-ground/Gibbs branch only the actual representative/chart comparison remains. Other states and unimplemented time--space transports retain their correlation test. \\
BH6 & canonical Wilson-state Cauchy limit & \textbf{Open Class-F execution}: prove the interacting state/effective-action stability theorem; the refinement category/comparison maps are not missing structural data. \\
BH7 & positive pure-glue continuum gap & \textbf{Open inequality}: \(\gamma_*:=\inf_L(E_{1,L}-E_{0,L})>0\), plus semigroup convergence. \\
BH8 & observable non-Gaussianity; stronger dynamical interaction when claimed & \textbf{Open execution}: a higher cumulant tests the chosen observable family; nonlinear-source particle-interaction claims also consume \ref{p12:cor:v2210-lit-interaction-comparator}. \\
BH9 & OS/Wightman reconstruction & \textbf{Derived implication} once BH2, BH5, BH6 and the ordinary OS regularity/vacuum hypotheses are verified. \\
BH10 & regulator-normalized local composite fields & \textbf{Open execution}: uniform residue envelope/distribution bounds and mixed-moment Cauchy convergence; stress tensor also needs action/metric-response normalization. \\
\bottomrule
\end{longtable}
\TableNotes{The table distinguishes finite data already exported by Paper~VII-B from continuum implications and still-open executions.  No row treats the repaired finite export as an advance continuum certificate.}

\begin{corollary}[BH1--BH10 status for the repaired \(SU(3)_c\) handoff]
\label{p12:cor:bh1-bh10-su3-discharge}
For the repaired Paper-VII-B pure-glue datum, the finite algebraic inputs of BH1--BH3 are available at finite level, while the continuum parts of BH4--BH10 reduce to explicit obligations.  In particular: the CP-odd action coefficient is action/measure dependent; QFT4 same-phase local-state canonicity and the actual flat representative/chart comparison must be proved for the claimed tower; on the equivariant canonical branch the finite state term is already derived, with approximate bounds in Corollary~\ref{p12:cor:v227-covariance-provider}; a positive gap requires \(\gamma_*>0\); interacting nontriviality requires a surviving higher connected cumulant or equivalent witness; Schwinger--Dyson identities require an action-measure realization; and local fields require regulator-normalized mixed-moment convergence.  Thus the repaired finite export does not, by itself, discharge BH1--BH10.
\end{corollary}
\begin{proof}
BH1--BH3 are finite construction/classification statements and the contact-normalized small-holonomy coefficient is algebraic once the smooth local connection chart exists.  Corollary~\ref{p12:cor:paper11-pure-glue-gap-floor} shows that BH7 is exactly the open inequality \(\gamma_*>0\), not a count-derived floor.  Theorems~\ref{p12:thm:euclidean-corridor-transfer}, \ref{p12:thm:native-spectral-ucs}, and \ref{p12:thm:ucs-implies-qft10} isolate the dynamical locality/tree/residue-convergence requirements.  Lemma~\ref{p12:lem:ym-sd-identification} isolates the action-measure requirement, while Corollary~\ref{p12:cor:wilson-algebra-nontriviality} supplies a valid non-Gaussian criterion.  Therefore the finite export is a correct input datum but not an advance certificate of the continuum conclusions.
\end{proof}

\begin{theorem}[Sharp conditional \(SU(3)_c\) Yang--Mills continuum handoff]
\label{p12:thm:native-bh-ym-continuum}
Consider the repaired Paper-VII-B pure-glue finite datum.  Suppose the following are independently established for that actual sequence:
\begin{enumerate}[label=\textnormal{(\roman*)}]
\item QFT4 local-state/finite-time-semigroup Cauchy convergence and the actual QFT6 flat representative/correlation comparison; on the equivariant canonical branch no separate state-covariance selector is required, as in Corollary~\ref{p12:cor:v227-covariance-provider};
\item finite reflection positivity and the remaining OS regularity/vacuum-sector hypotheses;
\item if a positive mass gap is claimed, \(\gamma_*:=\inf_L(E_{1,L}-E_{0,L})>0\);
\item a nonzero limiting connected gauge-invariant cumulant of order \(n\ge3\) when non-Gaussianity of the stated family is claimed; a stronger underlying-dynamics interaction claim additionally requires the actual same-source comparator or collision witness of Corollary~\ref{p12:cor:v2210-lit-interaction-comparator};
\item for each claimed local composite field, regulator normalization, uniform distributional bounds, and mixed-moment Cauchy convergence as in QFT10;
\item if Yang--Mills Schwinger--Dyson identification is claimed, the finite action-measure realization and controlled irrelevant-action limit of Theorem~\ref{p12:thm:tier3-sd-from-ucs}.
\end{enumerate}
Then the ordinary local-chart OS/Wightman reconstruction yields the corresponding \(SU(3)_c\) gauge-invariant continuum QFT.  Under (iii) its reconstructed pure-glue Hamiltonian has gap at least \(\gamma_*\).  Under (v) the stated local composite fields exist, and under (vi) the continuum Schwinger--Dyson identities hold.  A physical positive area law/string tension is not part of this local continuum certificate and still requires the separate area-extensive same-action lower bound.
\end{theorem}
\begin{proof}
Items (i)--(ii) provide the hypotheses of Theorem~\ref{p12:thm:bh-os-reconstruction} on the bounded gauge-invariant local algebra, yielding the Hilbert space, vacuum, Hamiltonian, and local net.  Item (iii), together with local-vector cyclicity, gives the stated gap by Corollary~\ref{p12:cor:paper11-pure-glue-gap-floor} and Theorem~\ref{p12:thm:mass-gap-transfer}.  Item (iv) gives observable-family non-Gaussianity by Lemma~\ref{p12:lem:ym-nontriviality}; its additional same-source comparator clause supplies any stronger underlying-dynamics interaction classification actually claimed.  Item (v) invokes Theorems~\ref{p12:thm:ucs-implies-qft10} and \ref{p12:thm:tier3-local-fields-from-ucs}.  Item (vi) invokes Theorem~\ref{p12:thm:tier3-sd-from-ucs}.  None of these steps uses a count-derived wall scale, a universal shell loss, or a position-space area law.  The latter belongs to confinement observables and is logically downstream of local QFT reconstruction.
\end{proof}



\begin{corollary}[Compatibility boundary of the final Paper-VII-B export]
\label{p12:cor:p11-final-export-compatibility}
The repaired Paper-VII-B pure-glue subdatum is a valid finite input to Theorem~\ref{p12:thm:native-bh-ym-continuum}; it is not, by itself, a proof of that theorem's continuum hypotheses.  Additional fermionic, thermal, nuclear, heavy-source, source-deformed, or open-channel continuum claims require their own sector-import QFT obligations.
\end{corollary}
\begin{proof}
Paper~VII-B exports the finite pure-glue singlet Hilbert spaces, vacua, Hamiltonians/transfers, Wilson algebra, and Hodge data and assigns continuum state convergence, OS reconstruction, area law, and local-field limits to Paper~VII-C.  A uniform gap is consumed only if a separate positive-gap claim is added.  The conditional theorem above consumes exactly those finite objects plus independently proved limit hypotheses.  Hence the handoff is compatible and acyclic without treating the finite export as an advance continuum certificate.
\end{proof}



% Source Paper VII-C, original line 2785; inherited proof bodies unchanged.
\subsubsection{Scope separation}
\label{p12:subsec:tier4-scope-closure}

The main handoff is the Paper~VII-B \(SU(3)_c\) pure-glue branch.  Arbitrary compact-simple \(G\) and controlled gapless sectors are separate extensions with their own finite inputs.

\begin{theorem}[Main and extension scopes after assumption reduction]
\label{p12:thm:tier4-scope-closure}
Paper~VII-C has three distinct scopes.  The source-native gauge branch may use either the operator-first UCP comparison or the action-first effective-action theorem above; these are alternative sufficient implementations of the same QFT4 state requirement.
\begin{enumerate}[label=(\roman*)]
\item The main application is the repaired Paper-VII-B \(SU(3)_c\) pure-glue finite datum together with the explicit continuum hypotheses of Theorem~\ref{p12:thm:native-bh-ym-continuum}; the finite prefix supplies its licensed source/action/measure lineage but does not pre-discharge a stronger dynamical interaction witness, residue convergence, or actual member/source matching and regulated identity obligations.  Its uniform-gap clause is inactive unless a separate positive-gap branch is claimed.
\item For arbitrary compact simple \(G\), representation-theoretic Wilson/Haar identities and the contact coefficient \(T_\rho/(2d_\rho)\) are universal, while the actual compatible interacting tower, gap, nontriviality, residue convergence, and action measure are group-specific inputs unless proved.
\item Controlled gapless sectors may use the same local OS reconstruction but do not inherit positive-gap or exponential-clustering conclusions.
\end{enumerate}
The strict positive-area/string-tension theorem is outside the local QFT existence theorem in all three scopes.
\end{theorem}
\begin{proof}
Part (i) is Corollary~\ref{p12:cor:p11-final-export-compatibility}.  Part (ii) follows from compact-group representation theory and Lemma~\ref{p12:lem:contact-plaquette-residue}, while the remaining properties are statements about the selected dynamics and its continuum limit rather than group theory.  Part (iii) follows because OS reconstruction itself does not require a positive mass gap; only gap-specific spectral and clustering conclusions do.  Finally, local gauge-invariant observables and their OS reconstruction do not require a strict positive-area estimate, so that estimate is a confinement/string-tension obligation rather than a local-QFT existence premise.
\end{proof}



% Source Paper VII-C, original line 7939; inherited proof bodies unchanged.
\subsection{QCD state, gap, interaction, and Wilson limits}
\label{p12:sec:v164-qcd-state-gap-wilson}

\begin{theorem}[Finite vacuum and continuum-state obligation]
\label{p12:thm:v164-finite-vacuum-continuum-state}
For the finite compiled Hamiltonian, the declared symmetric vacuum is the
zero-temperature Gibbs descendant
\[
 \rho_{0,L}=P_{0,L}/\operatorname{rank}P_{0,L}.
\]
No independent finite vacuum selector remains on this branch.  Its continuum
state is nevertheless not automatic: for every bounded local question the
actual refinement family must prove the asymptotic naturality defect
\[
 \delta_{ML}(A_M)
 :=\omega_M(A_M)-\omega_L(\mathcal E_{M\to L}A_M),
 \qquad
 \sup_{M\succeq L}|\delta_{ML}(A_M)|\longrightarrow0.
\]
This is QFT4 for the generated vacuum ray and finite-time semigroup cylinders.
Finite-dimensionality and the existence of UCP maps do not prove the limit.
\end{theorem}

\begin{theorem}[Physical gap and interaction witnesses]
\label{p12:thm:v164-physical-gap-interaction-witness}
With physical transfer step \(a_L\) and largest vacuum-orthogonal eigenvalue
\(q_L\), the finite physical gap is
\[
 \Delta_L^{\rm phys}=-a_L^{-1}\log q_L.
\]
If a positive continuum gap is separately claimed, its required antecedent is
\(\inf_L\Delta_L^{\rm phys}>0\), followed by QFT4/QFT7 and cyclic-sector
transfer.  This is not an active framework target.  Positive finite gaps alone are insufficient, and
\(q_L\to1\) is compatible with fixed positive mass when \(a_L\to0\).

Likewise, finite nonlinear or non-Gaussian actions do not prove an interacting
continuum.  QFT8 requires at least one regulator-normalized connected
cumulant of order \(n\ge3\) to converge to a nonzero value.  QFT10 separately
requires regulator-normalized mixed residue moments to be Cauchy.  These are
same-action limit calculations, not new coefficients.
\end{theorem}

\begin{theorem}[Fixed Wilson limits and derived surface exhaustion]
\label{p12:thm:v164-fixed-wilson-surface-exhaustion}
For every fixed compact contour, the normalized Wilson loop is a bounded
local cylinder observable, so its continuum expectation is exactly the QFT4
limit.  No second `WCONV' axiom is required.  Moreover, with
\[
 \rho_L=\sqrt{R_{x,L}\ell_0},
 \qquad R_L=T_L=\rho_L/4,
\]
the existing local-flat chart theorem gives rectangles \(C_L\) and spanning
surfaces \(\Sigma_L\) satisfying
\[
 A_{g_x}(\Sigma_L)
 =\frac{\rho_L^2}{16}
 \left[1+O\!\left(\frac{\ell_0}{R_{x,L}}\right)\right]
 \longrightarrow\infty.
\]
Thus `SURF' is derived on the local-flat continuum branch.
\end{theorem}

\begin{theorem}[Ordered static-potential and tension limits]
\label{p12:thm:v164-ordered-static-tension-limits}
Take the regulator limit first at fixed \((R,T)\).  Cutting the rectangle on
its reflection plane gives an OS vector \(\Psi_R\) and positive spectral
measure \(\mu_R\) with
\[
 \omega(W(C_{R,T}))
 =\langle\Psi_R,e^{-T(H_R-E_0)}\Psi_R\rangle
 =\int_0^\infty e^{-TE}\,d\mu_R(E).
\]
Hence
\[
 V(R)=-\lim_{T\to\infty}T^{-1}\log\omega(W(C_{R,T}))
 =\inf\operatorname{supp}\mu_R.
\]
Spatial reflection positivity makes \(V\) concave after the fixed endpoint
self-energy subtraction.  With \(V(0)=0\), \(V(R)/R\) is nonincreasing, so
the stable static sector has
\[
 \boxed{
 \sigma=\lim_{R\to\infty}\frac{V(R)}R
 =\inf_{R>0}\frac{V(R)}R.}
\]
This proves existence of the slope but permits \(\sigma=0\).  Paper~VII-A
v1.71 now proves both the square specialization and the finite-type
triality/branch-core criterion.  For the positive asymptotic slope claimed on
the pure-glue or unscreened sector of Paper~VII-A,
Definition~\ref{p10:def:v2187-wilson-sector-scope}, APBI is a sufficient
extensive same-owner route; an equivalent direct surface bound on the actual
nonlinear owner is an alternative.  A screened fundamental sector can instead
have \(\sigma=0\).  An area exponent cannot substitute for the independent
QFT8 non-Gaussian witness.
\end{theorem}

\begin{theorem}[Same-owner presentation transport and the APBI confinement bridge]
\label{p12:thm:v2155-same-owner-apbc-bridge}
\label{p12:thm:v2158-same-owner-apbi-bridge}
Let the canonical finite heat action and measure be transported by an exact
sampled realization or local-flat presentation map in the sense of the
framework's action/measure transport theorems.  Then the transported action,
measure, Wilson observables and conditional blocking maps are the same physical
owner: exact presentation transport contributes zero action or blocking
discrepancy.  In particular, using a reciprocal-position face chart does not
insert a primitive position lattice or a new QCD regulator.

The positive-slope conclusion concerns the claimed pure-glue/unscreened
Wilson sector.  It does not impose a positive fundamental asymptotic tension
on full QCD with a screening channel.

This same-owner statement does not by itself imply that the complete nonlinear
heat measure admits an extensive bounded-incidence face realization.  Define
APBI (action-preserving bounded-incidence presentation) to mean that on a
cofinal fixed-UV local-flat family the same nonlinear heat action has finite
face presentations with: (i) a level-independent finite set of rooted local
incidence types; (ii) physical face areas bounded by $A_+<\infty$; and (iii) a
level-independent strict supersolution of Paper~VII-A
Theorem~\ref{p10:thm:v2158-finite-type-triality-branch-core} at some $z>1$,
including its charged/neutral corridor contraction hypotheses.  Then
Paper~VII-A Corollary~\ref{p10:cor:v2158-apbi-string-tension-reduction} gives
\[
 V(R)\ge\frac{\log z}{A_+}R-O(1),
 \qquad
 \boxed{\sigma\ge\frac{\log z}{A_+}>0.}
\]
The previous APBC condition is the square-face, valence-six sufficient
specialization with $z=1.1$.
\end{theorem}
\begin{proof}
Exact sampled action transport is conjugation/pullback of the already existing
action and preserves all retained derivatives and symmetries.  Exact measure
transport pushes the reference measure, action and coarse map together, so
conditional blocking commutes with the presentation change and normalized
integrals are unchanged.  Hence a chart is a proof representation, not a new
physical action.

The plaquette chart theorems, however, identify local curvature-residue
representatives; they do not create a new refinement level or prove that the
complete nonlinear measure factorizes into one-face heat factors on an
extensive complex.  Paper~VII-A Proposition~\ref{p10:prop:multiface-heat-blocking}
shows that exact four-dimensional elimination generates multiface
interactions, while Proposition~\ref{p10:prop:v2158-heat-subdivision-area-nogo}
shows that exact semigroup subdivision of one fixed heat face does not
manufacture an extensive area exponent.  Thus APBI is an independent
same-owner fine-action theorem.  Once APBI is supplied, the VII-A finite-type
surface majorant applies at the fine level; exact pushforward may change the
blocked action without changing the already bounded Wilson expectation.  QFT4
passes fixed-loop bounds to the continuum, and the ordered limits give the
positive asymptotic slope.
\end{proof}

\begin{remark}[Current confinement status]
\label{p12:rem:v2155-confinement-status}
\label{p12:rem:v2158-confinement-status}
The representation sum, triality branching, neutral decorations, cycles and
handles are no longer the open confinement problem on the canonical heat
branch.  The finite-type theorem also shows that square faces and valence six
are not fundamental: any finite rooted incidence family with one strict
triality supersolution is sufficient.  Neither are reciprocal-position area,
Wilson-loop convergence, existence of large surfaces, the static-potential
limit, or existence of the asymptotic slope independent residuals on the
certified domains.  For the claimed pure-glue/unscreened positive slope,
APBI supplies a sufficient fine heat-action route with the stated
bounded-incidence/area and uniform-recursion properties; a direct same-owner
surface theorem is an alternative.  Until one such actual certificate is
proved, that positive slope is not an unconditional physical theorem.
\end{remark}

\begin{remark}[Action-owner residue inside APBI]
\label{p12:rem:v2159-apbi-action-owner-residue}
The nonlinear law's origin is the complete action parent; the cofinal
extensive family and uniform surface estimate are later analytic content.
A licensed whole-action presentation change transports the former, not the
latter.  An equivalent direct estimate on the actual generated multiface
action can establish the stated confinement target without making a product
of one-face heat kernels its mandatory physical ancestor.

The v2.159 action-owner audit narrows APBI without closing it.  Paper~VII-A
Corollary~\ref{p10:cor:v2159-same-hessian-nonheat-family} gives an explicit
positive central one-face family with the same Yang--Mills Hessian as the heat
action but different nonlinear density.  Hence exact quadratic cell stiffness,
positive quadrature and local plane isotropy do not select the complete
magnetic face law.  Appendix
Theorem~\ref{p10:thm:v2159-conditional-electric-heat} shows that a supplied
same-owner finite holonomy--flux kinetic realization does force a compact-group
heat kernel for the electric/temporal kinetic transfer, while
Proposition~\ref{p10:prop:v2159-holonomy-cotangent-firewall} shows that
holonomy alone does not generate that forward cotangent realization.  Finally,
Proposition~\ref{p10:prop:v2159-fixed-hurwitz-area-nogo} excludes a fixed
homothetic Hurwitz copy as the required bounded-area extensive family.
Therefore the remaining APBI content is specifically a cofinal extensive
same-owner nonlinear fine-action construction (or an equivalent direct surface
activity theorem for the actual generated multiface action); it is not a
missing local-incidence, Wilson-limit, or quadratic-normalization result.
\end{remark}


% Source Paper VII-C, original line 8136; inherited proof bodies unchanged.
\subsection{native mass gap and continuous no-promotion boundary}
\label{p12:sec:v165-native-gap-continuum-boundary}

\begin{theorem}[Finite carrier mass gap is the active prediction]
\label{p12:thm:v165-native-gap-active-prediction}
The active framework imports from Papers~III-C and V-B the theorem
\begin{equation}
 \label{p12:eq:v165-imported-native-carrier-gap}
 \boxed{
 \mathbf1_{(E_{56},E_{40})}(H_{\rm car})=0,
 \qquad
 E_{48}=\sqrt{E_{40}E_{56}},
 \qquad
 P_{48}^{\rm act}=0.}
\end{equation}
This is the predicted native mass gap: a finite gap in generated carrier
support centred logarithmically at the empty \(E_{48}\) coordinate.

The finite-QCD quantities
\(Delta_L^{\rm phys}=-a_L^{-1}\log q_L\) remain computable diagnostics, and
Theorem~\ref{p12:thm:mass-gap-transfer} remains a valid implication if an
independent branch supplies a refinement-uniform lower bound.  However, the
carrier theorem admits both gapped and gapless auxiliary refinement
families, so it cannot prove that antecedent.  Consequently a positive
continuous Yang--Mills mass gap is not an active framework prediction and is
removed from the mandatory frontier.  This statement does not prove a
gapless continuum; it proves only non-promotion from the native carrier data.
\end{theorem}

\begin{corollary}[Status of QFT9]
\label{p12:cor:v165-qft9-claim-specific-status}
QFT9 has two uses.  Controlled gapless-sector typing remains part of the
ordinary continuum certificate.  Its positive-gap alternative is a
claim-specific extension: if asserted, it must independently prove the
uniform lower bound and then may invoke the transfer theorem.  The positive
alternative is not an unfinished proof of
Eq.~\eqref{p12:eq:v165-imported-native-carrier-gap}.
\end{corollary}

% BEGIN V202 SOURCE DOMAIN DIAGNOSTICS


\section{RG, action compatibility and fixed-band continuum transport}
\label{p12:sec:v237-c5}
Native refinement with fixed physical ultraviolet momentum is separated from ordinary ultraviolet cutoff removal. Higher-order RG universality and source-uniform convergence are distinct stronger claims.

% Source Paper VII-C, original line 2805; inherited proof bodies unchanged.
\subsection{Renormalization as finite refinement compatibility}
\label{p12:sec:renormalization}


\begin{table}[htbp]
\centering
\small
\begin{tabularx}{\textwidth}{@{}
>{\raggedright\arraybackslash}p{0.55\textwidth}
>{\raggedright\arraybackslash}X
@{}}
\toprule
\textbf{Refinement/RG step} & \textbf{Discharged by} \\
\midrule
Born--Hopf refinement-compatibility map
& Definition~\ref{p12:def:bh-renormalization-map}. \\

Two-scale gapped refinement comparison
& Definition~\ref{p12:def:two-scale-gapped-comparison}. \\

Two-scale comparison implies QFT4
& Theorem~\ref{p12:thm:two-scale-comparison-qft4}. \\

Cross-vacuum diagnostic and same-side residuals
& Definition~\ref{p12:def:local-vacuum-comparison-seminorm}. \\

Sufficient same-side compatibility
& Definition~\ref{p12:def:admissible-local-vacuum-compatibility}. \\

Vacuum-line alignment
& Lemma~\ref{p12:lem:qft4g-projection-stability}. \\

Local state convergence
& Theorem~\ref{p12:thm:qft4hij-local-state-convergence}. \\

Bounded spectral-window class
& Definition~\ref{p12:def:tier2-spectral-window-class}. \\

Tier-2 estimates and quadrature
& Proposition~\ref{p12:prop:tier2-qft4-estimates} and
Proposition~\ref{p12:prop:tier2-quadrature}. \\

Flat-chart covariance and asymptotic locality
& Lemmas~\ref{p12:lem:tier2-flat-chart-covariance}
and~\ref{p12:lem:tier2-asymptotic-locality}. \\

RG flow as refinement bookkeeping
& Proposition~\ref{p12:prop:rg-not-new-input}. \\
\bottomrule
\end{tabularx}
\caption[Renormalization and spectral-window refinement]{Renormalization and spectral-window refinement.}
\label{p12:tab:renormalization-spectral-window-ledger}
\TableNotes{Renormalization is treated as a comparison law between finite local
presentations.  It does not supply new continuum inputs, new fitted
couplings, or a second microscopic position lattice.}
\end{table}

The framework should not treat renormalization as an external patch.  In the Born--Hopf setting, renormalization is the comparison law between finite local representatives at different refinement scales.


% Source Paper VII-C, original line 5324; inherited proof bodies unchanged.
\subsection{Exact RG generators and inheritance of the geometric UV fixed point}
\label{p12:sec:v199-rgas}

The starting datum is the framework's supplied UV geometry \(X_*\), with
the full protected carrier and its given geometric rescaling law.
The following finite RG identities and quantum transfer theorem do not
require a second geometric UV occurrence.  Derivatives are taken on the
displayed smooth domains.

\begin{proposition}[Heat representative and Gaussian interaction flow]
\label{p12:prop:v199-rgas-heat-gaussian}
On a compact connected Riemannian configuration manifold without boundary,
let \(P=-\Delta\ge0\), \(w_0\) be smooth and strictly positive, and
\(w_\tau=e^{-\tau P/2}w_0\).  Then \(S_\tau=-\log w_\tau\) satisfies
\begin{equation}
\label{p12:eq:v199-rgas-heat}
 \partial_\tau S_\tau=\tfrac12\Delta S_\tau
                             -\tfrac12|\nabla S_\tau|^2 .
\end{equation}
Separately, on an admitted finite real linear sector with fixed
nonnegative self-adjoint kinetic operator \(P\), \(0<k<\Lambda\), define
\[
 C_{\Lambda,k}=\int_{\Lambda^{-2}}^{k^{-2}}e^{-sP}\,ds,\qquad
 e^{-V_k(\phi)}
  =\int d\gamma_{C_{\Lambda,k}}(\chi)e^{-V_\Lambda(\phi+\chi)} .
\]
Assume the integrals and indicated field and scale derivatives are
dominated.  With \(\ell=\log(k/k_0)\),
\begin{equation}
\label{p12:eq:v199-rgas-gaussian}
 \dot C_{\Lambda,k}=-2k^{-2}e^{-P/k^2},\qquad
 \dot V_k=\tfrac12\operatorname{Tr}(\dot C_{\Lambda,k}V_k^{(2)})
              -\tfrac12(V_k^{(1)})^T\dot C_{\Lambda,k}V_k^{(1)} .
\end{equation}
For \(k<m<\Lambda\), Gaussian integration composes with
\(C_{\Lambda,k}=C_{\Lambda,m}+C_{m,k}\).
\end{proposition}
\begin{proof}
Positivity preservation defines \(-\log w_\tau\).
Use \(\partial_\tau w=\Delta w/2\) and
\(\Delta e^{-S}=e^{-S}(|\nabla S|^2-\Delta S)\).
For the Gaussian identity, differentiate the covariance endpoint.
Finite Gaussian convolution satisfies
\(\partial_\ell e^{-V}=\tfrac12\dot C_{ij}\partial_i\partial_j e^{-V}\),
which gives the displayed interaction flow.
Independent Gaussian sums add covariances, proving composition.
\end{proof}
Configuration-space heat time requires a matching before it is identified
with a spacetime shell scale.  The Gaussian formula is an exact
representative when the actual measure admits its split.  The finite
Haar action is blocked by \eqref{p12:eq:v199-rgas-disintegration}.

% BEGIN M8 PREPARED FINITE GEOMETRIC RG CONSEQUENCES
\begin{corollary}[Geometric shell covariance and exact Born symmetry]
\label{p12:cor:m8td-shell-covariance}
On an admitted finite spectral restriction of the kinetic realization
$P=B_R^2$ from Theorem~\ref{p3f:thm:m8td-hodge-dirac}, every eigenvalue
$\lambda\ge0$ has shell covariance
\begin{equation}
 C_{\Lambda,k}(\lambda)=
 \begin{cases}
 (e^{-\lambda/\Lambda^2}-e^{-\lambda/k^2})/\lambda,&\lambda>0,\\
 k^{-2}-\Lambda^{-2},&\lambda=0,
 \end{cases}\qquad 0<k<\Lambda.
 \label{p12:eq:m8td-shell-zero}
\end{equation}
At finite spectral ceiling $\lambda_{\max}$,
\[
 (k^{-2}-\Lambda^{-2})e^{-\lambda_{\max}/k^2}I
 \le C_{\Lambda,k}\le(k^{-2}-\Lambda^{-2})I.
\]
Thus the geometric zero requires no added mass merely to define this
finite-window covariance.  For every finite geometric descendant whose
complete constructor input is Born-unmarked, Paper~II-A,
Theorem~\ref{p2a:thm:generated-born-inheritance}, supplies the inherited
$\mathcal U$-covariance; the functional-calculus calculation below identifies
the concrete target action.  If the actual real sector is invariant under
$\mathcal U$, its Gaussian law is invariant under $\mathcal U$. For a
Born-invariant retained interaction $V_\Lambda$ the exact convolution in
Proposition~\ref{p12:prop:v199-rgas-heat-gaussian} satisfies
$V_k(\mathcal U\phi)=V_k(\phi)$.
\end{corollary}
\begin{proof}
Apply Paper~1B, Theorem~\ref{p1b:thm:finite-flow}
\cite{Kugelmann2026Paper1B}, to the stated finite spectral restriction
of the kinetic realization $P=B_R^2$. Its spectral integral, including
the zero eigenvalue, gives \eqref{p12:eq:m8td-shell-zero} and the two
bounds. The actual real sector is $\mathcal U$-invariant and
$\mathcal UB_R\mathcal U^{-1}=-B_R$, so the same theorem gives
the Gaussian symmetry; applying it to the present Born-invariant
$V_\Lambda$ gives the convolution identity.
\end{proof}

An invariant cutoff can be chosen by spectral functional calculus in
\(D_0^2\): if \(\mathcal UD_0\mathcal U^{-1}=-D_0\), then
\(\mathcal U f(D_0^2)=f(D_0^2)\mathcal U\) for bounded Borel
\(f\), in particular for its finite spectral projectors.  For a separate
unequal carrier window, the level/window mark prevents internal Born
invariance.  Paper~II-A,
Proposition~\ref{p2a:prop:typed-born-branch-transport}, first types transport
to the separately certified reciprocal presentation, and Paper~I-C,
Theorem~\ref{p1b:thm:unequal-exchange}, supplies the concrete exchanged-window
projector/operator identity.  An internal involution exists only on an
exchange-closed sector or the stated invariant comparison.  This geometric
transport is checked before using an invariant Gaussian measure, with
\(m_p=1\) still fixed; it does not by itself imply invariance of the actual
quantum state/measure, regulator, or source domain.

The real field sector must actually have this kinetic realization and
admit the stated Gaussian split; the spinor/form identity does not change
fermionic variables into bosonic ones. The static inverse $P^{-1}$ still
requires separating its kernel. The zero covariance diverges as
$k\downarrow0$, and the finite-ceiling lower bound need not survive
refinement. This is symmetry compatibility of the finite block map; the
geometric scale intertwiner in
Theorem~\ref{p12:thm:v199-rgas-fixedpoint-inheritance} is a separate relation.
% END M8 PREPARED FINITE GEOMETRIC RG CONSEQUENCES

% BEGIN M8 PREPARED AFFINE SUPPORT REDUCTION
\begin{lemma}[Covariance kernel is exactly the affine redundancy]
\label{p12:lem:m8td-affine-support}
Let $\mu$ be a nonzero positive finite measure on $\mathbb R^N$, let
$\mathcal A=a+E$ be its minimal full-measure affine subspace, and suppose
the positive tilted probability
\[
 d\mu_{k,J}=Z_k[J]^{-1}e^{-x^TR_kx/2+J^Tx}\,d\mu(x)
\]
has the indicated finite moments and dominated source derivatives on an
open source domain. Then
\begin{equation}
 \ker W_k^{(2)}=E^\perp,\qquad W_k^{(2)}|_E>0.
 \label{p12:eq:m8td-affine-covariance}
\end{equation}
The rank is independent of finite $k,J$ on this domain. On the minimal
affine presentation, the mean/source map is locally invertible. For
$\dim E>0$, continuous covariance on a compact subset of the domain has
a strictly positive least eigenvalue there.
\end{lemma}
\begin{proof}
The positive measure, minimal affine support and dominated source
domain satisfy Paper~1B, Theorem~\ref{p1b:thm:finite-flow}
\cite{Kugelmann2026Paper1B}, which gives
\eqref{p12:eq:m8td-affine-covariance} and local Legendre invertibility.
The kernel $E^\perp$ is independent of $k,J$. For $\dim E>0$,
continuity of the least covariance eigenvalue on a compact subset
gives a positive minimum; the zero-dimensional case is a constant field
presentation and needs no inverse on a nonzero space.
\end{proof}


\begin{theorem}[Finite effective average action and its exact flow]
\label{p12:thm:v199-rgas-eaa}
On an admitted finite real field presentation, let \(d\nu(x)e^{-S(x)}\)
be the complete convergent scale-independent ultraviolet measure, and
let the symmetric field-independent matrix \(R_k\) be differentiable.
Suppose the integrals and derivatives below exist with the stated
domination on an open source patch. On the nonzero positive real-measure
branch use minimal affine coordinates: Lemma~\ref{p12:lem:m8td-affine-support}
then proves that \(W_k^{(2)}\) is positive definite. For any other admitted
measure/presentation retain the explicit hypothesis \(W_k^{(2)}>0\):
\[
 \begin{gathered}
 Z_k[J]=\int d\nu(x)e^{-S(x)-\frac12x^TR_kx+J^Tx},\\
 W_k=\log Z_k,\quad \phi=\partial_JW_k,\\
 \Gamma_k[\phi]=J^T\phi-W_k[J]-\tfrac12\phi^TR_k\phi .
 \end{gathered}
\]
At fixed \(\phi\) and fixed finite UV regulator,
\begin{equation}
\label{p12:eq:v199-rgas-eaa}
 W_k^{(2)}=(\Gamma_k^{(2)}+R_k)^{-1}=:G_k,\qquad
 \partial_\ell\Gamma_k
  =\tfrac12\operatorname{Tr}(G_k\partial_\ell R_k).
\end{equation}
\end{theorem}
\begin{proof}
On the positive branch in minimal affine coordinates, this is the
Legendre identity and finite trace flow of Paper~1B,
Theorem~\ref{p1b:thm:finite-flow} \cite{Kugelmann2026Paper1B}, applied
to the actual ultraviolet measure $d\nu e^{-S}$.
For any other admitted presentation with the stated domination and
$W_k^{(2)}>0$, local inversion gives
$\partial_\phi J=\Gamma_k^{(2)}+R_k=(W_k^{(2)})^{-1}$.
At fixed source, $\dot W_k=-\langle x^T\dot R_kx\rangle/2$;
source terms cancel in the differentiated Legendre transform and
$\langle xx^T\rangle=\phi\phi^T+W_k^{(2)}$ gives the same trace.
\end{proof}
This is the finite form of the exact flow used in
\cite{Wetterich1993,Reuter1998}.  It neither assumes
\(\Gamma_\Lambda=S\) at arbitrary finite cutoff nor constructs a gravity
measure by writing the trace.  Gauge and fermionic applications require
their actual constraints, ghost or graded measure, multiplicities,
Jacobians and Ward identities in the complete generator before a
gravitational projection or continuum limit.

\begin{corollary}[Exact affine restriction of the finite EAA]
\label{p12:cor:m8td-affine-eaa}
Choose an orthonormal column map $T:\mathbb R^d\to E$ and fixed origin
$a\in\mathcal A$. Let $\mu_a$ be the image of $\mu$ under $x=a+Ty$,
and put
\[
 R_{k,E}=T^TR_kT,\qquad j=T^T(J-R_ka),\qquad
 W_{k,E}(j)=\log\int e^{-y^TR_{k,E}y/2+j^Ty}\,d\mu_a(y).
\]
Then
\begin{equation}
 W_k(J)=J^Ta-\tfrac12a^TR_ka+W_{k,E}(j),\qquad
 \Gamma_{k,\mathcal A}(a+T\varphi)=\Gamma_{k,E}(\varphi),
 \label{p12:eq:m8td-affine-transport}
\end{equation}
where the Legendre transform on the left uses sources modulo $E^\perp$.
In these minimal coordinates its exact finite flow is
\[
 \partial_\ell\Gamma_{k,E}
 =\tfrac12\operatorname{Tr}_E
  [ (\Gamma_{k,E}^{(2)}+R_{k,E})^{-1}\partial_\ell R_{k,E}].
\]
Thus covariance invertibility is a conclusion on this positive-measure
branch, rather than an additional regularity premise. The actual measure,
its positivity and the domination hypotheses remain.
\end{corollary}
\begin{proof}
Paper~1B, Theorem~\ref{p1b:thm:finite-flow}
\cite{Kugelmann2026Paper1B}, gives the shifted-source identity and
finite flow for this positive measure on $a+E$. Substitution in the
modified Legendre transform cancels the constant and linear terms,
giving $\Gamma_{k,\mathcal A}(a+T\varphi)=\Gamma_{k,E}(\varphi)$.
The $k$-dependent shift $-T^TR_ka$ therefore accompanies the regulator.
This is a fixed affine restriction; a subsequent moving chart uses
Theorem~\ref{p12:thm:v199-rgas-moving-eaa} with its measure data.
\end{proof}
The compact-subdomain estimate is finite-dimensional; it supplies no
uniform refinement or source-boundary estimate. A signed, complex or
Grassmann measure requires its own analysis. Paper~II-A's inherited geometric reciprocity does not by itself extend to a
quantum functional integral.  If the actual UV measure is Born invariant and
$R_k$ commutes with the actual involution, change of variables also gives Born
equivariance of $W_k$ and $\Gamma_k$ (and of their affine restrictions with
the induced affine action).  These measure, regulator and source-domain
conditions are independent quantum hypotheses, not descendants erased by the
geometric inheritance theorem.
Minimal affine restriction removes only directions whose linear
observables are constant almost surely; nonlinear gauge reduction retains
its separately defined constraints and measure.
% END M8 PREPARED AFFINE SUPPORT REDUCTION

\begin{corollary}[Born selection of finite effective-action vertices]
\label{p12:cor:born-vertex-selection}
On the positive finite-measure branch just established, let the actual
orthogonal involution \(J\) preserve the complete UV measure, the
minimal affine support and the dominated source domain, and let
\([J,R_k]=0\). At a fixed background \(J\bar\phi=\bar\phi\), every
existing vertex obeys
\[
 D^n\Gamma_k|_{\bar\phi}[v_1,\ldots,v_n]=0
 \quad\hbox{when}\quad Jv_i=\varepsilon_i v_i,\quad
 \prod_i\varepsilon_i=-1 .
\]
Thus the \(+\!-\) Hessian block vanishes and vertices with an odd
number of negative-grade arguments vanish. A \(+\!-\!-\) vertex
is allowed.
\end{corollary}
\begin{proof}
Paper~I-C, Corollary~\ref{p1b:cor:invariant-legendre-vertex-parity},
applies to this actual measure and regulator. Change of variables gives
\(\Gamma_k(J\phi)=\Gamma_k(\phi)\); differentiation at \(\bar\phi\)
gives the sign product. In affine coordinates choose a fixed origin or
retain the induced affine action.
\end{proof}
The selection persists through an equivariant finite flow. An invariant
free covariance alone does not establish it for a noninvariant interaction,
regulator, source domain or background.

\begin{corollary}[Parity blocks of an equivariant RG stability operator]
\label{p12:cor:born-rg-stability-blocks}
Let \(\mathcal R\) be a differentiable normalized RG map on its stated
action space and suppose \(\mathcal R J^*=J^*\mathcal R\).
At a fixed point \(\Gamma_*\) also fixed by \(J^*\),
\[
 D\mathcal R|_{\Gamma_*}\,J^*
   =J^*D\mathcal R|_{\Gamma_*}.
\]
The linearized map consequently preserves even and odd perturbation sectors.
\end{corollary}
\begin{proof}
Differentiate equivariance at \(\Gamma_*\), as in the same I-B corollary.
\end{proof}

\begin{theorem}[Born-reflected family jet parity]
\label{p12:thm:born-reflected-family-jet-parity}
Let $\mathcal B^2=1$ act smoothly on a finite effective-action chart, and let
$\Gamma_s$ be a $C^r$ family for $s\ge0$.  Define the reciprocal comparison
branch for $s<0$ by
\[
 \widetilde\Gamma_s:=\mathcal B^*\Gamma_{-s},
 \qquad
 \widetilde\Gamma_s:=\Gamma_s\quad(s\ge0).
\]
No inverse blocking map is assumed.  The glued family is $C^r$ at $s=0$ if
and only if
\[
 \boxed{
 \mathcal B^*\left(\partial_s^n\Gamma_s\big|_{0^+}\right)
 =(-1)^n\partial_s^n\Gamma_s\big|_{0^+},
 \qquad 0\le n\le r.}
\]
Thus even scale jets are Born-even and odd scale jets are Born-odd.
\end{theorem}

\begin{proof}
For $s<0$, the $n$th left derivative is
$(-1)^n\mathcal B^*(\partial_s^n\Gamma|_{0^+})$, while the right derivative
is $\partial_s^n\Gamma|_{0^+}$.  Equality of the one-sided jets is exactly the
displayed condition.
\end{proof}

\begin{corollary}[Mixed Born scale/field selection in adapted jets]
\label{p12:cor:born-scale-field-parity-selection}
At a Born-fixed background $p=\mathcal Bp$, put
$J=D\mathcal B_p$ and $Jv_i=\varepsilon_i v_i$.
For a smooth reflected family write $f_n=\partial_s^n\Gamma_0$.
The ordinary-derivative zero rule holds when $\mathcal B$ is affine.
For a general smooth involution it holds \emph{in Born-adapted normal
coordinates} and for derivatives associated with a Born-invariant
connection; there
\[
 D^k_{\rm adapted}f_n(p)[v_1,\ldots,v_k]=0
 \quad\hbox{if}\quad\prod_i\varepsilon_i\ne(-1)^n.
\]
In the original nonlinear chart, the exact correction is
\[
 \left((-1)^n-\prod_i\varepsilon_i\right)
 D^kf_n(p)[v_1,\ldots,v_k]
 =\sum_{\substack{\pi\in\Pi_k\\|\pi|<k}}
 D^{|\pi|}f_n(p)
 \left[D^{|C|}\mathcal B_p[v_i:i\in C]\right]_{C\in\pi},
\]
where $\Pi_k$ denotes the set partitions of the $k$ arguments.
Thus the original ordinary-coordinate zero claim is not valid
without linearity or control of nonlinear Born jets.
\end{corollary}
\begin{proof}
In local coordinates centered at $p$, define
$\Psi(x)=\frac12(x+J\mathcal B(x))$. Since
$\mathcal B^2=1$ and $J^2=I$,
$D\Psi_p=I$ and $\Psi\circ\mathcal B=J\circ\Psi$.
The inverse function theorem supplies equivariant local coordinates.
Differentiating $f_n\circ\mathcal B=(-1)^nf_n$ in these coordinates
gives the parity rule. In arbitrary coordinates, the higher chain
rule expressed as a sum over set partitions has top term
$D^kf_n[Jv_1,\ldots,Jv_k]$; moving this term to the other side yields
the displayed exact lower-jet correction.
A connection pulled back from flat connection in the adapted chart
is Born-invariant and gives the equivalent covariant formulation.
\end{proof}


\begin{theorem}[Born-antiequivariant beta law on a reflected scale]
\label{p12:thm:born-antiequivariant-beta}
Let $M$ be a smooth finite-dimensional essential coupling/action chart,
$\mathcal B:M\to M$ a smooth involution, and $\Phi_\ell$ a local $C^1$
autonomous flow with generator $\beta$.  If the physical signed scale is
Born-reflected and
\[
 \boxed{\mathcal B\circ\Phi_\ell=\Phi_{-\ell}\circ\mathcal B,}
\]
then
\[
 \boxed{\beta(\mathcal Bg)=-D\mathcal B_g\,\beta(g).}
\]
At a Born-fixed point $\mathcal Bg=g$, the tangent $\beta(g)$ lies in the
Born-odd eigenspace.  Hence a self-dual seed can leave the fixed locus only in
Born-odd directions while the law remains Born complete.
\end{theorem}

\begin{proof}
Differentiate the reflected-flow identity at $\ell=0$.  The left derivative is
$D\mathcal B_g\beta(g)$ and the right derivative is $-\beta(\mathcal Bg)$.
At a fixed point of $\mathcal B$, rearrangement gives the odd-eigenspace
statement.
\end{proof}

\begin{corollary}[Anticommuting stability at a Born-reflected fixed point]
\label{p12:cor:born-reflected-stability-anticommutation}
Under the hypotheses of Theorem~\ref{p12:thm:born-antiequivariant-beta},
assume additionally $\mathcal Bg_*=g_*$ and $\beta(g_*)=0$.  Put
$A=D\beta|_{g_*}$ and $J=D\mathcal B|_{g_*}$.  Then
\[
 \boxed{AJ=-JA.}
\]
Thus $A$ exchanges the Born-even and Born-odd tangent subspaces.  In a finite
essential chart its spectrum is symmetric with algebraic multiplicity,
\[
 \boxed{\operatorname{Spec}(A)=-\operatorname{Spec}(A)},
\]
and
\[
 \dim\ker A\ge\bigl|\dim V_+-\dim V_-\bigr|.
\]
Consequently nonzero critical eigenvalues occur in $\lambda,-\lambda$ pairs.
\end{corollary}

\begin{proof}
Differentiate the beta anti-equivariance at $g_*$.  The term containing the
second derivative of $\mathcal B$ is multiplied by $\beta(g_*)$ and vanishes,
leaving $AJ=-JA$.  Therefore $JAJ=-A$, so $A$ and $-A$ are similar and have
the same finite-dimensional spectrum with algebraic multiplicity.  The block
form relative to $V_+\oplus V_-$ is off-diagonal, so its rank is at most
$2\min(\dim V_+,\dim V_-)$, proving the kernel bound.
\end{proof}

\begin{remark}[Two distinct Born RG types]
\label{p12:rem:two-born-rg-types}
The preceding result does not replace
Corollary~\ref{p12:cor:born-rg-stability-blocks}.  If Born leaves the scale
orientation fixed, ordinary equivariance gives $[A,J]=0$ and separate parity
blocks.  If Born reverses the signed scale, reflected equivariance gives
$\{A,J\}=0$ and parity exchange.  A nontrivial stability operator cannot in
general satisfy both structures simultaneously.  The physical scale must
therefore be typed as Born-even or Born-reflected before either conclusion is
used.
\end{remark}

This block decomposition constructs neither a quantum fixed point nor its
relevant-direction count. The supplied geometric UV fixed point and the
quantum intertwiner retain their distinct roles in
Theorem~\ref{p12:thm:v199-rgas-fixedpoint-inheritance}.

\begin{theorem}[Exact finite flow in a moving normalized field chart]
\label{p12:thm:v199-rgas-moving-eaa}
Under the preceding theorem, let \(S_k\) be a differentiable invertible
real linear map, \(\chi=S_k\phi\), \(A=\dot S_kS_k^{-1}\),
\(\widehat\Gamma_k(\chi)=\Gamma_k(S_k^{-1}\chi)\) and
\(\widehat R_k=S_k^{-T}R_kS_k^{-1}\).  Then
\begin{equation}
\label{p12:eq:v199-rgas-moving-eaa}
 [\partial_\ell+(A\chi)\cdot\partial_\chi]\widehat\Gamma_k
 =\tfrac12\operatorname{Tr}\left[
 \widehat G_k(\dot{\widehat R}_k+A^T\widehat R_k+\widehat R_kA)
 \right],\quad
 \widehat G_k=(\widehat\Gamma_k^{(2)}+\widehat R_k)^{-1}.
\end{equation}
Implementation in the integral includes the transformed measure,
determinant and sources.
\end{theorem}
\begin{proof}
At fixed \(\chi\),
\(\dot{\widehat\Gamma}=\dot\Gamma-(A\chi)\cdot\partial_\chi\widehat\Gamma\).
Hessian congruence gives \(\widehat G=SGS^T\).
Differentiation yields
\(S^{-T}\dot R S^{-1}=\dot{\widehat R}+A^T\widehat R+\widehat R A\).
Substitution in the preceding trace and cyclicity prove the identity.
\end{proof}
Taking \(S_k=Z_k^{1/2}\) implements
Theorem~\ref{p3f:thm:v199-rgas-kinetic-connection}.
This is the finite linear instance of essential RG field transport
\cite{BaldazziZinatiFalls2022EssentialRG}.
An essential quotient carries all surviving sources and measure data.
Multiplying the whole action at fixed \(\hbar\) is not in general an
invertible change of field coordinates.

\begin{definition}[Finite generated quantum constructor and exact coarse map]
\label{p12:def:finite-generated-quantum-constructor}
On every finite branch for which Paper~IV supplies the reference integration
ray, integrable Euclidean action and action unit, write
\[
 \boxed{
 \mathcal Q_L(X_L)
 :=[\rho_L]
 =\left[e^{-S_{E,L}/\hbar}\mathfrak I_L\right].}
\]
For a generated coarse configuration map $\pi_{M\to L}$, define the exact
quantum coarse map on positive rays by
\[
 \boxed{\mathcal R_{M\to L}[\rho_M]
 :=[(\pi_{M\to L})_*\rho_M].}
\]
On the positive Haar branch this is represented by the exact conditional
blocking action of Theorem~\ref{p12:thm:v199-rgas-disintegration}.  An EAA
regulator is a useful finite representation/readout of the resulting quantum
comparison; it is not part of the definition of $\mathcal Q_L$.
\end{definition}

\begin{theorem}[Quantum inheritance of the supplied geometric fixed point]
\label{p12:thm:v199-rgas-fixedpoint-inheritance}
Let \(X_*\) be the supplied complete UV presentation, whose admitted
geometric scale maps obey \(\mathcal G_bX_*=X_*\).
Let \(\mathcal Q\) construct an admissible dimensionless quantum
action/state class modulo the declared essential field equivalences.
Suppose exact normalized RG intertwines the geometric maps on the
required domain:
\begin{equation}
\label{p12:eq:v199-rgas-intertwiner}
 \mathcal R_b\mathcal Q=\mathcal Q\mathcal G_b .
\end{equation}
Then \(\Gamma_*=\mathcal Q(X_*)\) satisfies
\(\mathcal R_b\Gamma_*=\Gamma_*\).  When its continuous RG generator
exists, \(\beta[\Gamma_*]=0\) in the essential quotient.
\end{theorem}
\begin{proof}
Apply the identity to \(X_*\):
\(\mathcal R_b\mathcal Q(X_*)=\mathcal Q(\mathcal G_bX_*)=\mathcal Q(X_*)\).
Differentiate at the identity for the generator statement.
\end{proof}
The geometric fixed point is already supplied by the framework.
If the complete quantum definition includes
\eqref{p12:eq:v199-rgas-intertwiner}, its fixed-point consequence is
immediate.  On the finite positive action-first branches of
Definition~\ref{p12:def:finite-generated-quantum-constructor}, the candidate
quantum map is no longer abstract: it is the generated state/action ray and
the exact coarse map is pushforward.  Theorem~\ref{p12:thm:v251-inherited-path-coarse-map}
and Corollary~\ref{p12:cor:v251-finite-same-owner-block-square} additionally
close the former finite same-owner coarse-map problem: the reciprocal path
refinement is inherited presentation structure, $\pi_{M\to L}$ is its
holonomy restriction/evaluation dual, and the resulting geometric block,
source/constraint transport, projective measure, reflection-compatible test
pullback and quantum pushforward are descendants of that one incidence.
Thus G33-NATIVE is closed/retired: no independent finite coarse-map or finite same-owner geometric--quantum bridge selector remains on this branch.

This finite statement must not be conflated with the separately supplied
geometric scale maps $\mathcal G_b$ in
Theorem~\ref{p12:thm:v199-rgas-fixedpoint-inheritance}.  Identifying those
complete UV scale maps with the finite block constructor still requires the
actual domain-specific intertwiner
\eqref{p12:eq:v199-rgas-intertwiner}.  The stronger claim-scoped \(\mathrm{G33	ext{-}COMP}\) certificate, when a theorem consumes it, consists of the hard complete-action/source-uniform multiscale control, generated multi-face bounds, stopping/crossover estimates and continuum universality actually consumed; the full interacting continuum state/action witnesses remain under
their separate QFT4/SURG owner.

\begin{proposition}[A convergent autonomous UV trajectory has a fixed limit]
\label{p12:prop:v199-rgas-uv-limit}
Let \(\gamma(t)\), \(t=\log(k/k_0)\), be an essential-action trajectory
in a Hausdorff space.  Suppose continuous coarse maps \(T_s\), \(s\ge0\),
obey \(T_s\gamma(t)=\gamma(t-s)\) for sufficiently large \(t\), and
\(\gamma(t)\to\Gamma_*\) as \(t\to+\infty\).  Then
\(T_s\Gamma_*=\Gamma_*\) for every \(s\).
\end{proposition}
\begin{proof}
Continuity gives
\(T_s\Gamma_*=\lim_{t\to\infty}T_s\gamma(t)
=\lim_{t\to\infty}\gamma(t-s)=\Gamma_*\).
\end{proof}
This implication requires UV convergence.  Removing a finite regulator
at fixed \(k\) is a different limit.

\begin{corollary}[RG trajectory selector removal after beta/boundary closure]
\label{p12:cor:rg-trajectory-unique-descendant}
Suppose the physical essential quantum RG chart carries a generated beta field
\[
 \frac{dg}{d\ell}=\beta(g),
\]
all boundary coordinates consumed by the claim are fixed by generated or
occurring data at $\ell_0$, and the initial-value problem is well posed on the
required UV--IR interval.  Then Paper~III-A
Theorem~\ref{p3a:thm:unique-solution-history-descent} makes the complete RG
trajectory $g(\ell)$ a generated descendant of the beta field and boundary
member; no additional trajectory selector is licensed.  Consequently
RGAS199-16 can contain an irreducible trajectory/member input only in the
unfixed relevant boundary coordinates or in a failure of existence,
uniqueness or domain control---not in a second choice among solutions of a
well-posed fully initialized beta system.
\end{corollary}

\begin{proposition}[Necessary Newton scaling at an interacting fixed point]
\label{p12:prop:v199-rgas-newton}
In a specified gravitational normalization, set
\(Z_{N,k}=(16\pi G_k)^{-1}\), \(\eta_N=-\partial_\ell\log Z_{N,k}\)
and \(g_N=k^{d-2}G_k\).  Then
\begin{equation}
\label{p12:eq:v199-rgas-newton}
 \beta_{g_N}=(d-2+\eta_N)g_N .
\end{equation}
At nonzero fixed \(g_N\) in four dimensions this forces
\(\eta_{N,*}=-2\).
\end{proposition}
\begin{proof}
Logarithmic differentiation gives
\(\partial_\ell\log g_N=d-2+\partial_\ell\log G_k=d-2+\eta_N\).
\end{proof}
This necessary identity does not compute the anomalous dimension or
decide which gravitational parameter combinations are essential.
That requires the complete normalized action
\cite{BaldazziFalls2021EssentialQEG}.

\begin{theorem}[Full Laurent--Legendre fixed-function scaling balance]
\label{p12:thm:v199-rgas-laurent-balance}
In four dimensions write the complete local curvature density as
\(F_k(R)=k^4\varphi_k(r)\), \(r=R/k^2\), and let
\(\mathcal T_k(r)\) denote the \(k^{-4}\) projection of its complete RG
generator in the chosen chart.  At fixed dimensionful \(R\),
\begin{equation}
\label{p12:eq:v199-rgas-fixed-function}
 \partial_\ell\varphi_k+4\varphi_k-2r\varphi_k'=\mathcal T_k(r).
\end{equation}
For the imported restored Laurent--Legendre family, set
\[
 u=U_k/k^2,\quad c=C_{7,k}/k^2=u\alpha^{-7},\quad
 z=Z_{N,k}/k^2,\quad \lambda=\Lambda_{{\rm ren},k}/k^2,
\]
\[
 \widehat f_{\rm L}(r):=k^{-2}f^{\rm act}_{2\mid7,k}(k^2r),\qquad
 p(r):=r\widehat f_{\rm L}'(r)-\widehat f_{\rm L}(r),
 \qquad D=Q^{\,n_{2+}},\quad Q=8\sqrt\alpha .
\]
Here \(f^{\rm act}_{2\mid7,k}\) is the restored fixed-depth
Laurent--Legendre family, equivalently reconstructed from its generated local
susceptibility \(f_{RR}=D_2\chi_7^{\rm jet}\) with terminal data
\(f(C_{7,k})=0\), \(f_R(C_{7,k})=1\).  The older reciprocal rational
profile is only the secant comparison readout and is not used as the physical
Euler-response generator.
Then
\begin{equation}
\label{p12:eq:v199-rgas-laurent-profile}
 \boxed{\varphi(r)=z\left[\widehat f_{\rm L}(r)-2\lambda\right],
 \qquad r\widehat f_{\rm L}'-\widehat f_{\rm L}=p.}
\end{equation}
On a smooth active stratum, if \(u,\alpha,z,\lambda\) and the generated
coefficient character are stationary, the necessary fixed-function balance is
\begin{equation}
\label{p12:eq:v199-rgas-laurent-balance}
 \boxed{\mathcal T_*(r)
 =2z\left[\widehat f_{\rm L}(r)-p(r)-4\lambda\right].}
\end{equation}
The full \(r\)-dependence of the restored Laurent--Legendre action is already
contained in \(\widehat f_{\rm L}\) and its generated local jet.  The
fixed-function balance uses only the inverse-Euler identity
\(r\widehat f_{\rm L}'-\widehat f_{\rm L}=p\); no reciprocal-secant profile
is reinserted as a second physical Hessian and no extra dressing-derivative
insertion occurs.
The coordinate \(\lambda\) is the independent zero-order coordinate of this
enlarged RG atlas.  It is not to be added again on a retained
terminal-normalized representative in which that zero-order quotient has
already been fixed.
\end{theorem}
\begin{proof}
At fixed \(R\), \(\dot r=-2r\); differentiating
\(k^4\varphi_k(r)\) proves
\eqref{p12:eq:v199-rgas-fixed-function}.  Under the stated stationary
homothety data,
\[
4\varphi-2r\varphi'
=z\bigl[4\widehat f_{\rm L}-2r\widehat f_{\rm L}'-8\lambda\bigr].
\]
The Legendre identity
\(r\widehat f_{\rm L}'-\widehat f_{\rm L}=p\) converts the bracket to
\(2(\widehat f_{\rm L}-p)-8\lambda\), proving the displayed balance.
\end{proof}
The homothety representative supplies an equation to test.  It does not
prove preservation of this finite family by quantum blocking or
stationarity of \(U_k/k^2\).  In a moving chart,
\(\mathcal T_k\) includes the connection terms of
\eqref{p12:eq:v199-rgas-moving-eaa} and all retained fluctuation sectors.
The conformal block is \eqref{p3f:eq:v199-rgas-conformal-hessian}.
The smooth-domain conditions retain the hierarchy junctions and endpoint
realization.  Even if this projection vanishes, remaining generated
essential beta functions must vanish too: different curvature invariants
can have the same constant-curvature projection.

\begin{remark}[Precise asymptotic-safety consequence]
\label{p12:rem:v199-rgas-as-scope}
Fixed-point inheritance is proved conditional on an admitted quantum
intertwiner.  A Weinberg-type interacting UV completion additionally
requires an admissible interacting quantum class, a finite-dimensional
UV critical surface of essential directions, and a trajectory in that
surface with the required physical infrared limit.
For \(\dot{\delta g}=B\delta g\), use \(Bv=-\theta v\):
\(\operatorname{Re}\theta>0\) is UV-attractive; marginal directions
require nonlinear analysis.  A finite microscopic carrier does not bound
the number of generated field operators.  A sufficient spectral lemma
and a response-scaling diagnostic are in
Appendix~\ref{p12:app:v199-rgas-diagnostics}.
The finite identities do not assert a complete essential stability
spectrum or a quantum-gravity continuum construction.
\end{remark}




\section[Quantum source identities and inherited charged/kinetic ports]{Quantum source identities and inherited\\charged/kinetic ports}
\label{p12:sec:v237-c6}
Ward, G16, K16 and native gauge-source identities refer to the actual regulated measures, states and action variations. They are not new source selectors and do not assert interacting continuum convergence.

% Source Paper VII-C, original line 4195; inherited proof bodies unchanged.
\subsection{exact graded boundary response and source naturality}
\label{p12:sec:v2141-exact-graded-boundary-response}

\begin{theorem}[Exact finite graded resolution-boundary flow]
\label{p12:thm:v2141-exact-graded-boundary-flow}
On the finite paired boson/fermion/gauge-quotient domain of
Paper~IV, Theorem~\ref{p4:thm:v4100-exact-graded-eaa}, the resolution-boundary
response of Theorem~\ref{p12:thm:v2140-rg-boundary-response} is exact:
\begin{equation}
\label{p12:eq:v2141-exact-graded-boundary-flow}
 \boxed{
 \partial_{\ell}\Gamma_k^{\partial}
 =\frac12\operatorname{STr}\left[
  (\Gamma_k^{\partial(2)}+\mathcal R_k)^{-1}
  \partial_{\ell}\mathcal R_k\right],\qquad \ell=\log k .}
\end{equation}
Its expansion about the microscopic Hessian reproduces the v2.140 graded
one-loop logarithmic tangent; the full right-hand side contains the running
Hessian and is not itself a one-loop approximation.
\end{theorem}

\begin{theorem}[Exact blocking composition and retained-source naturality]
\label{p12:thm:v2141-blocking-source-naturality}
For a compatible exact finite blocking tower, let
$\mathscr R_{NM}$ denote elimination from a finer boundary $N$ to a retained
boundary $M$.  Then
\begin{equation}
\label{p12:eq:v2141-blocking-composition}
 \boxed{\mathscr R_{ML}\mathscr R_{NM}=\mathscr R_{NL}.}
\end{equation}
If a source $A$ is retained rather than integrated out and the finite source
derivatives are dominated, differentiation commutes with exact blocking:
\begin{equation}
\label{p12:eq:v2141-source-naturality}
 \boxed{
 D_A^r\,\mathscr R_{NL}S_N[A]
 =D_A^r S^{\rm eff}_{N\downarrow L}[A]}
 \qquad(r=1,2,\ldots)
\end{equation}
for every derivative order that exists on the stated domain.  Thus source
Hessians, currents and higher vertices are transported by the same blocking
cocycle as the effective action; beta functions are infinitesimal coordinates
of this finite response, not additional dynamics.
\end{theorem}

\begin{proof}
Equation~\eqref{p12:eq:v2141-blocking-composition} is the tower property of
iterated conditional integration on the same finite measure/action.  Retaining
$A$ makes it a parameter of every fibre integral.  Dominated differentiation
therefore passes through each finite integral, and iterating gives
\eqref{p12:eq:v2141-source-naturality}.
\end{proof}


\begin{corollary}[Finite isolated-pole transport is an owned constructor]
\label{p12:cor:v299-finite-pole-source-transport}
On the finite Hermitian analytic branch of Paper~III-F,
Theorem~\ref{p3f:thm:v1140-finite-pole-source-transport}, exact Schur descent
transports isolated poles and their retained-source residues without an
independent matching coefficient.  With the notation of that theorem,
\[
 \operatorname*{Res}_{z=z_*}(j^\dagger D^{-1}j)
 =\frac{|u^\dagger(j_R-B^\dagger C^{-1}j_C)|^2}
        {u^\dagger S'(z_*)u}.
\]
The energy derivative and the complement-source term are part of the same
owner.  On a compatible nested quadratic reduction the result is independent
of the order of elimination by Paper~III-F,
Corollary~\ref{p3f:cor:v1140-pole-transport-composition}.  Thus finite pole/source
transport is closed on this domain; evaluation of the pole and its survival
under refinement or a massless charged infrared limit retain their own
certificates.  No equality between a terminal finite geometric mass and a
pole mass follows from this transport statement.
\end{corollary}
\begin{proof}
Apply the cited Schur theorem to the already owned inverse propagator while
retaining its full source.  Theorem~
\ref{p12:thm:v2141-blocking-source-naturality} retains the source through the
compatible action tower; when a quadratic block reduction is used, the cited
composition corollary identifies its meromorphic source response.  This
composition does not replace an interacting integration by a Gaussian
approximation.
\end{proof}

\begin{proposition}[Primitive heat-shell threshold and co-consumption]
\label{p12:prop:v2141-heat-threshold-coconsumption}
For the primitive heat representative and a massive block
$\Delta_m=P+m^2$, the spectral factor separates as
\[
 e^{-s\Delta_m}=e^{-sm^2}e^{-sP}.
\]
Writing $s=k^{-2}$ gives the native threshold profile
\begin{equation}
\label{p12:eq:v2141-threshold-profile}
 \boxed{w_m(k)=e^{-m^2/k^2}.}
\end{equation}
At Gaussian order a decomposed one-loop gauge coefficient therefore has
\[
 \frac{d\alpha_i^{-1}}{d\log k}
 =-\frac1{2\pi}\sum_r b_i^{(r)}e^{-m_r^2/k^2},
\]
and integration between $k_0$ and $k$ gives
\begin{equation}
\label{p12:eq:v2141-threshold-E1}
 \boxed{
 \alpha_i^{-1}(k)=\alpha_i^{-1}(k_0)
 -\frac1{4\pi}\sum_r b_i^{(r)}
 \left[E_1(m_r^2/k^2)-E_1(m_r^2/k_0^2)\right].}
\end{equation}
The profile is regulator-representative rather than universal.  Its structural
content is co-consumption: below threshold the explicit carrier contribution
is suppressed while the already accumulated effective-action contribution
remains retained.
\end{proposition}

\begin{corollary}[Charged-gap control of the electromagnetic infrared]
\label{p12:cor:v2141-charged-gap-ir}
Assume a strict charged gap $m_r\ge m_{\rm ch}>0$ and uniformly bounded
retained electromagnetic source vertices on the tested branch.  Then for the
primitive heat representative
\[
 |\partial_{\log k}Z_Q(k,0)|\le C e^{-m_{\rm ch}^2/k^2}
\]
for some finite $C$, and
\[
 \int_0^{k_0}\frac{dk}{k}e^{-m_{\rm ch}^2/k^2}
 =\frac12E_1(m_{\rm ch}^2/k_0^2)<\infty.
\]
Hence the $k\to0$ electromagnetic quadratic response has a finite endpoint.
Uniform control in a neighbourhood of $q^2=0$ additionally permits the
$k\to0$ and $q^2\to0$ limits to commute.  The interacting hadronic/QCD
contribution still requires its color-singlet electromagnetic response
control.  The current-weighted criterion below gives a sufficient alternative
to a strict spectral gap on its declared temporal source chart.
\end{corollary}


\paragraph{Finite origin and limiting obligation.}
For the owned finite action/state/current packet, the measure below is already
a spectral readout, as in
Corollary~\ref{p10:cor:finite-qcd-parent-readouts}.
Its weights and contacts must be evaluated, but no independent physical
measure or matching coefficient is constituted here.
The inverse-moment, zero-atom, contact and uniform refinement/IR/UV
requirements are later analytic certificates.  A common endpoint may be used
across a retained compatible family only after its constancy is proved.

\begin{theorem}[Current-weighted infrared endpoint on a temporal source chart]
\label{p12:thm:v299-current-weighted-ir}
Fix an owned finite Hamiltonian $H_L$, normalized vacuum $|0,L\rangle$,
and one Hermitian transverse current test $J_L$ in a reflection-positive
Euclidean source chart.  The current and its normalization are the retained
same-action descendants of Paper~V-C,
Theorem~\ref{p7v2:thm:v2141-em-same-source-hessian}.  Subtract its vacuum mean
and define the connected positive spectral measure
\begin{equation}
\label{p12:eq:v299-current-spectral-measure}
 d\mu_{J,L}(\Delta)
 =\sum_{n:E_{n,L}>E_{0,L}}
 |\langle n,L|(J_L-\langle J_L\rangle_0)|0,L\rangle|^2
 \,\delta_{E_{n,L}-E_{0,L}}(d\Delta).
\end{equation}
Assume the connected current has no zero-energy spectral atom, including
transitions to any other degenerate vacuum.  The sign convention for its
transverse response is the contact-minus-connected-current convention,
\[
 K_L(i\omega)=K_{{\rm loc},L}(i\omega)-\chi_{J,L}(i\omega),
 \qquad
 \chi_{J,L}(i\omega)
 =2\int_{(0,\infty)}\frac{\Delta}{\Delta^2+\omega^2}
                 \,d\mu_{J,L}(\Delta).
\]
Here $K_{{\rm loc},L}$ retains the full owned local kinetic and contact
response.  Suppose its Ward-subtracted difference quotient
\[
 Z_{{\rm loc},L}(\omega^2)
 :=\frac{K_{{\rm loc},L}(i\omega)-K_{{\rm loc},L}(0)}{\omega^2}
\]
has a regular endpoint.  Then, for $\omega\ne0$,
\begin{equation}
\label{p12:eq:v299-current-ir-difference-quotient}
 \boxed{\frac{K_L(i\omega)-K_L(0)}{\omega^2}
 =Z_{{\rm loc},L}(\omega^2)
 +2\int_{(0,\infty)}
       \frac{d\mu_{J,L}(\Delta)}{\Delta(\Delta^2+\omega^2)}.}
\end{equation}
For a limiting source measure with the same spectral representation, finite
$\int\Delta^{-1}d\mu_J$ and
\begin{equation}
\label{p12:eq:v299-current-inverse-cubic-moment}
 \int_{(0,\infty)}\Delta^{-3}\,d\mu_J(\Delta)<\infty
\end{equation}
therefore give the finite temporal endpoint
\begin{equation}
\label{p12:eq:v299-current-ir-endpoint}
 \boxed{Z_J^{\rm tot}(0)=Z_{\rm loc}(0)
                   +2\int_{(0,\infty)}\Delta^{-3}\,d\mu_J(\Delta).}
\end{equation}
A strict positive gap is sufficient for the low-energy part of this
criterion, but is not required by it.
\end{theorem}

\begin{proof}
The finite spectral theorem gives the connected vacuum correlator
\[
\int e^{-\Delta|t|}d\mu_{J,L}(\Delta).
\]
Its Euclidean Fourier transform is $\chi_{J,L}(i\omega)$ because the
transform of $e^{-\Delta|t|}$ is $2\Delta/(\Delta^2+\omega^2)$.  The absence of a zero atom excludes a
nondecaying connected vacuum component.  In the stated sign convention,
\[
 \frac{\chi_{J,L}(0)-\chi_{J,L}(i\omega)}{\omega^2}
 =2\int\frac{d\mu_{J,L}(\Delta)}{\Delta(\Delta^2+\omega^2)},
\]
which proves the finite identity.  The inverse-first moment makes the
unsubtracted connected response finite on the stated limiting chart.  The
integrand in the difference quotient is bounded above by $\Delta^{-3}$;
dominated convergence with \eqref{p12:eq:v299-current-inverse-cubic-moment}
and the regular local/contact endpoint proves
\eqref{p12:eq:v299-current-ir-endpoint}.
\end{proof}

% BEGIN S3 LITERATURE INTEGRATION v3.040 [VII-C: FINITE-CURRENT ENVELOPE]
\begin{corollary}[Two-sided same-source finite current-response envelope]
\label{p12:cor:v340-s3lit-current-envelope}
Under the finite, reflection-positive temporal transverse-current
hypotheses of Theorem~\ref{p12:thm:v299-current-weighted-ir}, fix one
source-owned $H_L$, its vacuum, and its connected Hermitian $J_L$.
Suppose there is no connected zero-gap atom, that the current has
positive variance $V=\operatorname{Var}_0(J_L)>0$, and that the
current-visible excitation spectrum has a positive minimum $m_J$.
Define the same-owner double-commutator moment
\[
 C=\langle0|[J_L,[H_L,J_L]]|0\rangle,
 \qquad \overline\Delta=\frac{C}{2V}>0.
\]
For $s\ge0$ the current \emph{part}, not the total contact-plus-current
photon response, obeys
\[
 \frac{2V}{\overline\Delta(\overline\Delta^2+s)}
 \ \le\ Z_{J,L}(s):=
 2\sum_{\Delta_n>0}\frac{w_n}{\Delta_n(\Delta_n^2+s)}
 \ \le\ \frac{2V}{m_J(m_J^2+s)},
\]
where $w_n=|\langle n|J_L|0\rangle|^2$.  At $s=0$,
\[
 \frac{16V^4}{C^3}\le Z_{J,L}(0)
 \le\min\left\{\frac{2V}{m_J^3},\frac{C}{m_J^4}\right\}.
\]
The Jensen lower bound is saturated precisely when all
current-visible transitions carry one excitation gap (degeneracy
allowed).  A physical Thomson readout still requires the separate
$Z_{\rm loc,L}$, gauge Ward identity, spatial source response and
uniform refinement/infrared certificates.
\end{corollary}
\begin{proof}
The connected weights obey $\sum_nw_n=V$ and
$2\sum_n\Delta_nw_n=C$.  On $(0,\infty)$ the function
$f_s(u)=1/[u(u^2+s)]$ is strictly convex and decreasing for $s\ge0$.
Apply Jensen to weights $w_n/V$ for the lower bound and its equality
criterion.  Apply $\Delta_n\ge m_J$ termwise for the first upper bound.
At $s=0$, $\Delta_n^{-3}\le\Delta_n/m_J^4$, which gives the other.
\end{proof}
% END S3 LITERATURE INTEGRATION v3.040 [VII-C]


\begin{corollary}[Uniform current-moment transport through refinement]
\label{p12:cor:v299-current-ir-uniform-transport}
Assume the same owned source is transported along the finite tower, with
$d\mu_{J,L}$ converging weakly on every compact subinterval of $(0,\infty)$
to $d\mu_J$.  In addition to the regular local/contact term, require the
explicit infrared uniform-integrability and ultraviolet source bounds
\begin{align}
 \lim_{\varepsilon\downarrow0}\sup_L
   \int_{(0,\varepsilon)}\Delta^{-3}d\mu_{J,L}&=0,
   \label{p12:eq:v299-current-ir-ui}\\
 \lim_{M\to\infty}\sup_L
   \int_{(M,\infty)}\Delta^{-1}d\mu_{J,L}&=0.
   \label{p12:eq:v299-current-uv-ui}
\end{align}
Suppose $Z_{{\rm loc},L}$ converges locally uniformly to a function continuous
at $\omega=0$.  Then the temporal difference quotients in
\eqref{p12:eq:v299-current-ir-difference-quotient} converge to
\eqref{p12:eq:v299-current-ir-endpoint} under the joint limit
$L\to\infty$, $\omega\to0$.  A convenient sufficient infrared bound is
\[
 \sup_L\int_{(0,1)}\Delta^{-(3+\epsilon_0)}d\mu_{J,L}<\infty
 \quad\hbox{for some }\epsilon_0>0;
\]
a bound on the inverse-cubic moment alone does not replace uniform
integrability.
\end{corollary}

\begin{proof}
Split the spectral integral into $(0,\varepsilon)$, a compact middle
interval, and $(M,\infty)$.  The first integral is uniformly bounded by the
inverse-cubic tail in \eqref{p12:eq:v299-current-ir-ui}.  For $M\ge1$, the
last is bounded by the inverse-first ultraviolet tail in
\eqref{p12:eq:v299-current-uv-ui}.  On the compact middle interval the
integrand converges uniformly to $\Delta^{-3}$ as $\omega\to0$, and the
stated source-measure convergence passes its integral to the limit.  Smooth
cutoffs give the same argument if an interval endpoint is a spectral atom.
Let the two tails tend to zero and use the locally uniform convergence of
the local/contact quotient.  Finally the stronger low-energy bound gives
$\int_{(0,\varepsilon)}\Delta^{-3}d\mu_{J,L}\le
\varepsilon^{\epsilon_0}\int_{(0,1)}\Delta^{-(3+\epsilon_0)}d\mu_{J,L}$,
proving its sufficiency.
\end{proof}

This contracts the hadronic infrared target to its generated current-weighted
response.  The actual QCD spectral measure and its uniform bounds remain to
be evaluated and certified.  Energy levels and their current residues must
be co-consumed before the response is read out; a terminal mass list alone
omits the required weights.  The result concerns the stated temporal
transverse source chart and its refinement limit.  It does not establish
joint external four-momentum regularity, the full interchange of RG $k\to0$
with $q^2\to0$, the numerical heat-to-Thomson endpoint response, or a vacuum
occurrence certificate.  Paper~V-C Release v2.75 separately proves that this
open endpoint response carries no independent source-normalization matching
coefficient.  Frequency-dependent contact terms remain in $Z_{\rm loc}$;
only a frequency-independent seagull cancels from the difference quotient.

\begin{theorem}[Projective vacuum-cone invariance under undeformed principal flow]
\label{p12:thm:v2141-projective-cone-rg-invariance}
Let a vacuum quadratic descendant have the same principal owner as the
Paper--III-A operational cone and write, in a local isotropic chart,
\[
 \Gamma_{Q,k}^{(2)}
 =Z_t(k)\,\omega^2-Z_s(k)c_{\rm op}^2|\mathbf p|^2+\text{lower order}.
\]
Then
\begin{equation}
\label{p12:eq:v2141-beta-c}
 \boxed{
 \frac{\partial_{\log k}c_Q}{c_Q}
 =\frac12\left(\eta_t-\eta_s\right),
 \qquad
 \eta_{t,s}:=-\partial_{\log k}\log Z_{t,s}.}
\end{equation}
On an undeformed symmetry-preserving same-principal-owner branch,
$Z_t=Z_s$ up to one common positive scalar normalization, so
\begin{equation}
\label{p12:eq:v2141-beta-c-zero}
 \boxed{\partial_{\log k}c_Q=0.}
\end{equation}
A regulator/truncation that breaks the principal symmetry may display
finite-$k$ anisotropic artefacts; a genuine change of the physical front cone
requires a new principal-order tensor/owner or a failure of the stated
symmetry-restoration hypotheses.
\end{theorem}

\begin{remark}[Three distinct scale variables]
\label{p12:rem:v2141-kKq-firewall}
The protected hierarchy/depth label $K$, the RG resolution boundary $k$, and
an external physical probe momentum $q$ are differently typed.  No identity
$K\equiv k$, $k\equiv q$, or $K\equiv q$ is inferred from common monotonicity
or common logarithmic notation.  A relation among them requires an explicit
same-owner intertwiner.
\end{remark}
% END EXACT-GRADED-BOUNDARY-RESPONSE v2.98



% Source Paper VII-C, original line 6019; inherited proof bodies unchanged.
\subsection{Rooted source ancestry and claim-local quantum readouts}
\label{p12:sec:v2109-source-observable-transport}
On the rooted local classical branch, the primitive odd source already
occurs by Paper~V-A's rooted incidence theorem.  Its existence is not a
second \texttt{G0\_FERMION\_SOURCE\_MEMBER} gate.  Member-sensitive
continuum and QCD predictions still require the actual noncentral
same-action family assembly, the finite interacting scheme comparison and
running, and only the state/domain/limit certificates used by that claim.
The primitive source theorem must not be read as a numerical spectrum.

\begin{corollary}[A unique QROOT ground supplies the finite Ward state antecedent]
\label{p12:cor:v2109-qroot-ward-state}
On the same finite positive-energy block, suppose QROOT gives zero
excitation, the ground space is one-dimensional, and the represented
symmetry \(U(g)\) commutes with its Hamiltonian.  Then the root density is
\(P_0\) and is symmetry invariant.  If the same-action evolution
intertwines this symmetry, every evolved state is invariant as well.
The finite Ward identity of
Proposition~\ref{p12:prop:v248-finite-quantum-ward-antecedent} therefore
needs no independently selected invariant state on this common domain.
\end{corollary}
\begin{proof}
Positive support and the simple ground force \(\rho_0=P_0\).
Spectral functional calculus gives \([P_0,U(g)]=0\).  Conjugating by an
intertwining evolution preserves that equality; the stated finite Ward
proposition applies.
\end{proof}
On the native cyclic scalar/tensor perturbation branch QROOT is closed by the
fresh-daughter neutral Fock-unit theorem and positive-family source
compression.  This corollary therefore consumes that generated root state on
that branch.  Generic/noncyclic state occurrence remains separately typed;
ground degeneracy on such an enlarged branch need not force an invariant pure
state or the maximally mixed ground density.  Different cosmological and
QCD/chiral regulator Hamiltonians cannot be identified by this composition;
anomalous/unimplemented symmetries and continuum passage retain their own
conditions.

\begin{proposition}[Interacting scalar vertex as an inverse-propagator derivative]
\label{p12:prop:v2109-scalar-vertex-source-derivative}
Fix a regulated quantum realization with differentiable normalized
source-dependent correlations and fixed microscopic fields.  Let \(S[J]\)
be its full propagator, assumed invertible on the consumed domain, and
\(\mathcal C_S=-\delta S[J]/\delta J\) its connected insertion response.
Then the amputated scalar response is
\[
 \Lambda_S=S^{-1}\mathcal C_SS^{-1}
           =\frac{\delta S^{-1}}{\delta J}.
\]
A uniform scalar mass source gives
\(\Lambda_S(p,p)=\partial_m S^{-1}(p)\); a local derivative retains
nonzero insertion momentum.
\end{proposition}
\begin{proof}
Differentiate \(S^{-1}S=I\):
\(\delta S^{-1}=-S^{-1}(\delta S)S^{-1}\).
Substitute the definition of the connected insertion response.
\end{proof}
The already derived native insertion
\(\mathcal O_m^{\rm BH}=\partial_mD_m
=\tfrac12(I-\Gamma_5\epsilon_H)\) fixes the bare source normalization;
it is not the interacting derivative.  The latter includes the normalized
measure response, including sea/disconnected contributions where present.

\begin{theorem}[Complete finite BH scalar response]
\label{p12:thm:v2123-complete-finite-bh-scalar-response}
On a finite positive vectorlike overlap branch, let
\[
 G_L(U;m):=D_m[U]^{-1},\qquad S_L(m):=\langle G_L\rangle_{L,m},
\]
and suppose the differentiated mass coordinate appears in $n_q$ identical
dynamical determinants while all other source coordinates are held fixed.
With $\mathcal O_m^{\rm BH}=\partial_mD_m$,
\[
 \boxed{
 \partial_mS_L
 =-\langle G_L\mathcal O_m^{\rm BH}G_L\rangle
 +n_q\,\operatorname{Cov}_{L,m}\!\left(
 G_L,\operatorname{Tr}(G_L\mathcal O_m^{\rm BH})\right).}
\]
Therefore the connected scalar response and amputated finite vertex are
\[
 \boxed{
 \mathcal C_{m,L}:=-\partial_mS_L
 =\langle G_L\mathcal O_m^{\rm BH}G_L\rangle
 -n_q\operatorname{Cov}_{L,m}\!\left(
 G_L,\operatorname{Tr}(G_L\mathcal O_m^{\rm BH})\right),}
\]
\[
 \boxed{
 \Lambda_{m,L}
 =S_L^{-1}\mathcal C_{m,L}S_L^{-1}
 =\partial_mS_L^{-1}.}
\]
The covariance term is the normalized determinant/sea response and is not an
optional add-on.
\end{theorem}
\begin{proof}
Differentiate $G_LD_m=I$ to obtain
$\partial_mG_L=-G_L\mathcal O_m^{\rm BH}G_L$.  Differentiating the normalized
expectation contributes the covariance with
$\partial_m\log\det(D_m^{n_q})=n_q\operatorname{Tr}(G_L\mathcal O_m^{\rm BH})$.
Paper~VII-A, Theorem~\ref{p10:thm:v167-overlap-resolvent-accretive}, gives
invertibility of $S_L$ on this positive massive branch, and differentiating
$S_L^{-1}S_L=I$ gives the amputated identity.
\end{proof}

\begin{corollary}[Finite G16 vertex constructor is source-owned]
\label{p12:cor:v2123-finite-g16-constructor}
On the preceding branch, $\Lambda_{m,L}$ is a deterministic functional of the
same finite action, determinant weight and native source derivative.  No
additional finite scalar normalization coefficient or independent external-leg
parameter is inserted in order to define the interacting vertex.
\end{corollary}

\begin{definition}[Leg-normalization-independent projected response]
\label{p12:def:v2123-leg-invariant-scalar-response}
Fix common comparison kinematics $\kappa$ and a nonzero scalar projector
$P_m$.  Let $N_{\psi,L}(\kappa)$ be any nonzero functional of the same full
inverse propagator that is homogeneous of degree one under a common finite
field rescaling $S_L^{-1}\mapsto Z_\psi S_L^{-1}$.  Define
\[
 \widehat\Lambda_{m,L}:=N_{\psi,L}^{-1}\Lambda_{m,L},\qquad
 V_{{\rm BH},L}:=P_m[\widehat\Lambda_{m,L}].
\]
Then $V_{{\rm BH},L}$ is invariant under that common external-leg rescaling.
This is a comparison-coordinate construction, not a new physical datum.
\end{definition}

\begin{theorem}[Finite scheme-comparison constructor on a common noncritical chart]
\label{p12:thm:v2123-finite-scheme-comparison-constructor}
Suppose a named scheme $S$ and the BH presentation describe the same scalar
source subspace on common kinematics $\kappa$, with projected normalized
response matrices $V_S(\mu,\kappa)$ and $V_{\rm BH}(\mu,\kappa)$.  On the
noncritical comparison domain
\[
 \mathcal D_S:=\{(\mu,\kappa):\det V_{\rm BH}(\mu,\kappa)\ne0\},
\]
the finite source conversion is uniquely
\[
 \boxed{
 C_{S\leftarrow{\rm BH}}(\mu,\kappa)
 =V_S(\mu,\kappa)V_{\rm BH}(\mu,\kappa)^{-1}.}
\]
For a one-dimensional scalar source this is
$C_{S\leftarrow{\rm BH}}=v_S/v_{\rm BH}$; if $S$ normalizes $v_S=1$, then
$C_{S\leftarrow{\rm BH}}=v_{\rm BH}^{-1}$.  The source coordinate transforms
contragrediently as in Proposition~\ref{p12:prop:v288-contragredient-mass-source-transport}.
\end{theorem}
\begin{proof}
Linearity of the source insertion gives
$\mathcal O_m^S=C_{S\leftarrow{\rm BH}}\mathcal O_m^{\rm BH}$ and hence the
same linear relation for the normalized projected response.  Invertibility
of the chosen BH response matrix on $\mathcal D_S$ gives the displayed unique
coordinate change.
\end{proof}

\begin{corollary}[Controlled limit transport on certified source charts]
\label{p12:cor:v2123-g16-controlled-limit}
On any fixed bounded positive source chart for which the hypotheses of
Theorems~\ref{p12:thm:v202-bounded-source-inheritance} and
\ref{p12:thm:v283-local-1pi-transfer} hold, the corresponding fixed projected
source responses and local 1PI derivatives converge along each admitted
phase-selected refinement.  If the same-phase canonicality condition QFT4$_C$
holds, the comparison constructor is refinement-path independent.  Critical
points, growing source dimension, unbounded field limits and global
phase-crossing charts remain outside this corollary.
\end{corollary}

\begin{corollary}[Native exact source blocking does not require a named scheme]
\label{p12:cor:v2123-native-source-blocking}
For a source parameter $a$ owned by the finite action, the exact coarse
pushforward theorem gives
\[
 \partial_aS_L^{\rm eff}=\langle\partial_aS_M\rangle_x
\]
whenever its stated differentiation hypotheses hold.  Thus an internal
physical calculation may transport the native source through the exact finite
effective action without first converting it to $\overline{\rm MS}$ or any
other named continuum coordinate.  A named scheme is required only when the
claim itself consumes that comparison coordinate.
\end{corollary}

The structural finite G16 constructor is therefore closed on its declared
positive/noncritical domain.  A numerical named-scheme value remains a
claim-specific evaluation, and a pole certificate remains a separate
claim-scoped continuation/isolation/residue problem.

\begin{theorem}[An observable may converge without a unique cluster state]
\label{p12:thm:v2109-observable-cluster-convergence}
Let \(\omega_\alpha\) be a relatively compact net of states in a topology
in which evaluation of a fixed bounded observable \(A\) is continuous.
Suppose every cluster state \(\omega\) satisfies \(\omega(A)=a\).
Then \(\omega_\alpha(A)\to a\).  The conclusion also applies to
compatible representatives \(A_\alpha\) when their evaluation defects
relative to \(A\) tend to zero uniformly on the retained state family.
\end{theorem}
\begin{proof}
Failure of convergence gives a subnet whose evaluations stay outside some
fixed neighbourhood of \(a\).  Relative compactness gives a convergent
subnet of it.  Continuity makes its limit a cluster-state evaluation still
outside that neighbourhood, contradicting the hypothesis.  A uniformly
vanishing representative defect changes no limit.
\end{proof}
This is a weaker sufficient target for a particular bounded prediction,
not a unique-state or full QFT4 theorem.  Agreement of all surviving cluster
states on the requested readout still has to be proved.  Singular,
unbounded or ordered-limit observables need their own uniformity and
continuity estimates; they are not covered by bounded evaluation alone.

% END INLINE SOURCE: Paper7C-SourceAndObservableTransport-v1_2.tex




% BEGIN INLINE SOURCE: Paper7C-PhysicalChargedPoleHandoff-Main-v1_0.tex
% Framework v2.168: integrated finite charged-source derivations.

% Source Paper VII-C, original line 6235; inherited proof bodies unchanged.
\subsection{Complete finite charged root and its physical identification boundary}
\label{p12:sec:v2168-physical-charged-pole-handoff}
The terminal energies are fixed outputs of V-B's compiler.
Their comparison with a physical charged-particle mass consumes the same
action, source, kinetic and constraint owner. The finite constructors now
make the remaining equation explicit; they do not establish its value on
an unevaluated physical branch.

\begin{proposition}[Ordered physical charged-pole handoff]
\label{p12:prop:v2168-physical-charged-pole-handoff}
On a certified physical stationary rest-comparison chart, use the actual
normalized source response generated in IV and V-C. At a regular
complementary target the finite terminal comparison requires
\[
 [\varepsilon_r I-A-\pi^2 \widehat C K_W(\varepsilon_r)\widehat C^\dagger
 -\pi \widehat C L_{WR}(\varepsilon_r)-\pi L_{RW}(\varepsilon_r)\widehat C^\dagger
 -K_R(\varepsilon_r)]u=0,
 \qquad t_r^\dagger u\ne0.
\]
Here the actual raw source Gram $N$ is restricted to its positive active
support, $\widehat C=N^{-1/2}C$ and $t_r=N^{1/2}e_r$ retain the raw
physical tag. On the canonical CAR chart $N=I$, so $\widehat C=C$ and
$t_r=e_r$. For a different already transported tag use its actual vector.
For a simple normalized null channel its finite residue is
$uu^\dagger/[1+\|R(\varepsilon_r)Bu\|^2]$. At a singular eliminated
complement the full charged matrix test of IV replaces this chart.
If the calculation was made only on the faithful carrier or a containing
support, prove positive expectation of the generated residue effect on the
consumed physical support. A source polynomial root calculated on that
exact actual support already includes this visibility.
The transition then needs the named particle/rest interpretation and each
refinement, resonance or charged infrared limit consumed by the claim.
\end{proposition}
\begin{proof}
The complete marked Schur word is V-C,
Theorem~\ref{p5c:thm:v2168-complete-weak-memory}. IV's spectral source
decomposition and support-family visibility corollary give the source
test; the regular Schur derivative is strictly positive on a simple
channel. Apply the finite same-source pole transport theorem only on
that certified domain. Its finite hypotheses do not imply the physical
identification or limiting hypotheses, which remain separate antecedents.
The source-Gram factors follow from
$\mathcal Q=\mathcal Q_{\rm raw}N^{-1/2}$ and
$B_W=\pi FC^\dagger N^{-1/2}$; raw tag visibility is therefore
$|t_r^\dagger u|^2/[1+\|R(\varepsilon_r)Bu\|^2]$. No projected CAR
or untransported tag is assumed.
\end{proof}

If the physical rest/source/constraint chart has not been certified,
that identification is consumed before the displayed kernel test.
A component band can have nonzero spatial energy; the same internal
current does not put a pole at $\varepsilon_r$ on every external band.
A nonrest target first uses its proved energy/rest dictionary. No primitive
zero-momentum preparation is inserted. If the complete generated root is
shifted, its value is the forward dynamical output; it does not redefine
the earlier terminal coefficient. Charged infrared observables use their
appropriate inclusive, dressed or detector domain where a sharp isolated
pole is not the consumed target. A05 remains
\texttt{CLOSED\_FINITE\_TRANSPORT;OPEN\_PHYSICAL\_IDENTIFICATION\_LIMIT}.

% END INLINE SOURCE: Paper7C-PhysicalChargedPoleHandoff-Main-v1_0.tex



% v2.169: generated action/support constructors and consuming scope.

% BEGIN INLINE SOURCE: Paper7C-PhysicalSupportHandoff-Main-v1_1.tex
% Framework v2.169: Batches 7--13, integrated at the owning source.

% Source Paper VII-C, original line 6303; inherited proof bodies unchanged.
\subsection{Physical-support roots and the remaining mass handoff}
\label{p12:sec:v2169-support-handoff}

\begin{proposition}[Physical finite root with visibility already included]
\label{p12:prop:v2169-physical-support-root}
On a certified actual stationary source/kinetic/constraint/rest-comparison
chart, let $P$ be the exact consumed support. Generate $p_{r,P}$ from the
complete action and the actual raw tag by IV,
Theorem~\ref{p4:v2169:b10:thm:support}. Then first calculate the source-visible root set of $p_{r,P}$ from that
actual word.  If the terminal coefficient itself is claimed as the physical
pole, the finite comparison is $p_{r,P}(\varepsilon_r)=0$, with no
additional source-overlap test on that same support. The full source-visible Schur/rank test in
Proposition~\ref{p12:prop:v2168-physical-charged-pole-handoff} is an
alternative representation. A faithful-full-carrier root or a containing
support does not remove the consumed physical-support test.
After the generated root is identified, compare it with the terminal
coefficient only for the claim that consumes that equality, identify the
transition with the named particle, and prove survival of its pole/residue or
appropriate physical observable in each limit actually consumed.
\end{proposition}
\begin{proof}
Exact-support norm equivalence gives the same invariant null space and
cyclic polynomial for every density of that exact support. Its root is
precisely a nonnull actual-source frequency. Full-resolvent spectral
expansion makes this the same source-visible zero as the Schur/rank test.
The generated positive effect supplies the residue; its numerical weights
remain state dependent. Finite source transport supplies no additional
physical interpretation or limit theorem outside its antecedents.
\end{proof}

If the physical action grammar is balanced, its original fermion-number
conservation is already proved by typing. IV's adjacent-degree construction
then evaluates the complete actual addition/removal blocks. On the licensed
full ground support, retain all its nonzero degree summands and their
least common multiple, without choosing one ground vector. On a claimed
Slater/unit branch not already stationary, first evaluate the complete
creation-coefficient certificate; the available free filled state is not
automatically the interacting vacuum. All weak, occupation, constraint and
mixed contact returns stay in the actual blocks.

On a licensed single-energy global ground the positive terminal root uses
only the addition leg; on a general energy eigen-support both legs remain.
No primitive zero-momentum preparation is introduced. If the nonrest chart's
energy/rest dictionary has not been proved, it is the first consumed step.
If the complete root shifts, that value is the forward dynamical output and
does not redefine the upstream $\varepsilon_r$. A05 remains
\texttt{CLOSED\_FINITE\_TRANSPORT;OPEN\_PHYSICAL\_IDENTIFICATION\_LIMIT}.

% END INLINE SOURCE: Paper7C-PhysicalSupportHandoff-Main-v1_1.tex



% v2.171: conversation proofs, with actual history obligations retained.

% BEGIN INLINE SOURCE: Paper7C-BaryonExportHandoff-Main-v1_0.tex
% Framework v2.171: integrated conversation proof on its stated domain.


% Source Paper VII-C, original line 6360; inherited proof bodies unchanged.
\subsection{Regulated physical baryon-export handoff}
\label{p12:v2171:sec:baryonexporthandoff}
V-A's generated representation coefficients
\eqref{p5a:v2171:eq:exportcoefficients} are finite algebraic data. Transport to a
physical quantum baryon or lepton current consumes the actual regulated
Ward antecedent, including every measure Jacobian, source contact, regulator
residue and retained boundary flux on the same history. The existing
\path{p12:prop:v248-finite-quantum-ward-antecedent} and
\path{p12:cor:v236-ward-dependent} control that transport.
\begin{proposition}[Same-owner vector-current export]
\label{p12:v2171:prop:baryonexport}
Suppose an actually licensed regulated current/history map transports the
displayed V-A representation anomaly to the physical vector currents.
Its gauge contribution to the $B$ and $L$ sources is equal and cancels
from $B-L$. Both occurring weak and hypercharge gauge contributions must
be retained. Additional charged contacts and material-boundary fluxes are
included in their actual balances.
\end{proposition}
\begin{proof}
The weak and hypercharge coefficients for $B$ and $L$ in
\eqref{p5a:v2171:eq:exportcoefficients} coincide. Linearity of the certified
same-owner Ward transport preserves that equality for the gauge source.
Subtract the two transported laws. The conclusion uses the regulated
antecedent; it is not inferred from the representation table alone.
\end{proof}
This current representation and an actual finite generator representation
of the same charge law are alternatives, rather than two additive sources.
The baryon-equivariant finite stationary theorem imposes no conclusion on
an uncomputed anomalous completion outside its hypotheses.
The native regulated current/history map and its nonzero export expectation
remain an A21 construction and execution obligation. A Cartan midpoint,
Euclidean kink, unlicensed filled sea, or conventional plasma rate does not
discharge it. The native direct K8--K16 particle owner remains closed on its
existing domain. The physical solution and charge/entropy continuation are
consumed by VI-B, rather than installed as additional reheating inputs.

% END INLINE SOURCE: Paper7C-BaryonExportHandoff-Main-v1_0.tex


% BEGIN FRAMEWORK v2.228: CHARGED ORIGIN AND CHARACTER TESTS

% Source Paper VII-C, original line 6400; inherited proof bodies unchanged.
\subsection{Charged source origin before continuum transport}
\label{p7c:sec:v228-charged-origin}
Continuum transport of the noncentral charged source starts with the
hierarchy-arrow origin certificate of
Definition~\ref{p5:def:v228-hierarchy-origin} and the conditional handoff
of Proposition~\ref{p5:prop:v228-charged-origin-scope}.
The exact finite character or constant-fibre return identities apply after
identifying the actual source output with their stated domains.
They supply no missing origin arrow, uniform limiting estimate, analytic
continuation or physical pole identification at the continuum interface.

% END FRAMEWORK v2.228: CHARGED ORIGIN AND CHARACTER TESTS


% Source Paper VII-C, original line 7930; inherited proof bodies unchanged.
\subsection{Finite nuclear-readout ownership boundary}
\label{p12:sec:finite-nuclear-readout-ownership-boundary}
The selected-cell closed-shell rank sequence, its twofold occupation readout,
and the split-grammar nuclear diagnostics are finite nuclear-transfer statements
owned by Paper~VII-B.  They have therefore been removed from Paper~VII-C.  Paper~VII-C
imports no nuclear magic-number theorem and uses no nuclear diagnostic in the
OS/Wightman reconstruction, continuum mass-gap transfer, local-field
construction, or subsystem criterion.


% Source Paper VII-C, original line 8571; inherited proof bodies unchanged.
\subsection{charged finite-source continuum synchronization}
\label{p12:sec:v294-charged-source-sync}
The continuum layer imports the charged finite source member as an upstream
theorem output.  Refinement, renormalization and continuum reconstruction may
transport that member but do not select or regenerate its charged masses or
family incidence.  G16 and the claim-indexed continuum certificates retain
their existing scopes.

% END CHARGED-SOURCE-CONTINUUM-SYNC v2.94


% BEGIN CHARGED-COMPARISON-CONTINUUM-SYNC v2.95

% Source Paper VII-C, original line 8583; inherited proof bodies unchanged.
\subsection{charged finite-comparison continuum synchronization}
\label{p12:sec:v295-charged-comparison-continuum-sync}
For charged fermions the native finite source coordinate $T^{\rm BH}$ imported
into the source-dual G16 construction already contains the derived finite
comparison character of Paper~V-B v5.91.  The contragredient map
\[
 M^S(\mu)=C_S(\mu)^{-\dagger}T^{\rm BH}
\]
therefore transports a fixed finite charged source.  Neither its anomalous-
dimension flow nor a choice of reporting scale generates the relative
$24/25$ polar grading.  The statements that no $2\,{\rm GeV}$ or
$E_{32}/\sqrt{240}$ unity-crossing theorem is asserted and that pole extraction
is separately claim-scoped remain unchanged.
% END CHARGED-COMPARISON-CONTINUUM-SYNC v2.95




% BEGIN K16-KINETIC-INTERFACE v2.96

% Source Paper VII-C, original line 8602; inherited proof bodies unchanged.
\subsection{claim-scoped K16 high-temperature kinetic interface}
\label{p12:sec:v296-k16-kinetic-interface}

\begin{theorem}[K16 local kinetic equilibration on an admitted scattering domain]
\label{p12:thm:v296-k16-local-kinetic}
Assume the K16-produced visible ordinary state lies in a high-temperature
quasiparticle/scattering domain for which the generated unitary interaction
action supplies microreversible collision kernels (including the required
infrared-safe effective $2\leftrightarrow2$ and, where used, effective
$1\leftrightarrow2$ processes).  Then the quantum kinetic H-theorem of
Paper~IV applies.  If the reaction graph contains the connected visible
Standard--Model interaction graph of Paper~V-C, the stationary local family is
Bose/Fermi Gibbs modulo exact conserved charges.  If in addition
\[
 \Gamma_{\rm vis}\gg H,
\]
local kinetic/chemical equilibration occurs before appreciable background
expansion.
\end{theorem}

\begin{corollary}[K16 visible stationary-tangent rate benchmark]
\label{p12:cor:v296-k16-rate-certificate}
For the stationary Schur tangent of Paper~VI-A,
\[
 H_{16^+}^{\rm stat}\lesssim35.7732\ {\rm GeV},
 \qquad
 T_{16}^{E}\simeq5.0457\times10^9\ {\rm GeV}.
\]
Thus the stationary-tangent dimensionless benchmark is
\[
 \boxed{\Gamma_{\rm vis}/T_{16}^{E}\gtrsim7.09\times10^{-9}}.
\]
For the exact post-K16 background the scalaron retains its inherited Cauchy
data, so kinetic closure requires the theorem above with the actual dynamical
$H_{16^+}$.  The displayed stationary value is a useful rate benchmark, not a
replacement background or an exact post-K16 Hubble certificate.  This
corollary does not create a scattering certificate where the existing QFT/IR
hypotheses are absent and does not assert global KMS convergence of the entire
cosmology.
\end{corollary}
% END K16-KINETIC-INTERFACE v2.96


% BEGIN G17-ENDPOINT-RESPONSE-SYNC v2.100

% Source Paper VII-C, original line 8687; inherited proof bodies unchanged.
\subsection{native K16 exact gauge-source transport}
\label{p12:sec:v2177-k16-exact-source-transport}

\begin{theorem}[Common-heat boundary as the native resolution boundary]
\label{p12:thm:v2177-common-heat-resolution-boundary}
On the source-native common-heat branch, retain the canonical normalization
\[
 k_i(g_i^{\rm heat})^2=4\pi^2
\]
and the same primitive heat representative used by the exact finite
resolution-boundary flow.  Paper~VI-A fixes the dimensionful heat cutoff by
\[
 \Lambda_g^2=96\pi^2U.
\]
Writing the primitive heat parameter as $s=k^{-2}$, the common-heat datum is
therefore the exact boundary point
\[
 \boxed{k_h=\Lambda_g.}
\]
This identifies the native starting point of the exact blocking flow.  It does
not identify $k$ with cosmological time and does not identify $k_h$ with a
named experimental renormalization scale.
\end{theorem}

\begin{theorem}[K16 resonance generates the physical hard resolution]
\label{p12:thm:v2177-k16-resonance-scale}
On the KIN16 WKB production domain, Paper~VI-B gives the exact leading
stationary-phase condition
\[
 \boxed{2n=a(t_*)m_s(t_*).}
\]
A produced transverse gauge quantum with comoving label $n$ has physical
momentum $p_{\rm phys}=n/a$.  Hence the same produced mode supplies the hard
resolution consumed by its source vertex:
\[
 \boxed{k_*(t_*)=p_{\rm phys}(t_*)=\frac{m_s(t_*)}{2}.}
\]
This is a process-local intertwiner between the exact RG/source coordinate and
the actual K16 production event.  It does not promote the global RG coordinate
$k$ to a cosmic clock.
\end{theorem}

\begin{corollary}[Native K16 scale interval]
\label{p12:cor:v2177-k16-native-scale-interval}
At lower-side K16 entry,
\[
 k_{*,16}=\frac{m_{16}}2,
 \qquad
 k_h=\sqrt{96}\,\pi\sqrt U .
\]
Using the exact lower-branch mass, whose leading value is
$m_{16}=\sqrt{U/6}\,[1+O(\alpha^{-2})]$, gives
\[
 \boxed{\frac{k_{*,16}}{k_h}=\frac1{48\pi}\,[1+O(\alpha^{-2})]},
 \qquad
 \boxed{\log\frac{k_{*,16}}{k_h}=-\log(48\pi)+O(\alpha^{-2}).}
\]
Thus the scale interval required by K16 production is generated internally;
no independent choice of $\mu(t)$ or measured matching point is part of this
step.
\end{corollary}

\begin{theorem}[Exact RG-improved scalaron--gauge vertex]
\label{p12:thm:v2177-exact-scalaron-gauge-vertex}
Let the exact finite effective gauge term at resolution $k$ be written
\[
 \mathcal L_i(k)=-\frac1{4g_i(k)^2}\,\mathcal F_i^2,
 \qquad
 \beta_i(k):=\partial_{\log k}g_i(k),
\]
and retain the Weyl/RG spurion
\[
 \tau=\alpha^{-1/2}\frac{\phi_c}{\bar M_{\rm Pl}},
 \qquad k\mapsto ke^{-\tau}.
\]
Exact source naturality under the finite blocking flow gives, to first order in
$\phi_c$ but without replacing the running source by its Gaussian tangent,
\[
 \Delta\mathcal L_i
 =-\tau\frac{\beta_i(k)}{2g_i(k)}F_{i,c}^2.
\]
Therefore, in the convention
$-\lambda_i^{\rm exact}(k)\phi_cF_{i,c}^2/4$,
\[
 \boxed{\lambda_i^{\rm exact}(k)
 =\frac{2\alpha^{-1/2}}{\bar M_{\rm Pl}}
 \frac{\beta_i(k)}{g_i(k)}.}
\]
Expanding the exact finite flow about the Gaussian branch,
$\beta_i=b_ig_i^3/(16\pi^2)+\cdots$, reproduces exactly the Paper~V-C
coefficient
\[
 \lambda_i=\frac{\alpha^{-1/2}}{\bar M_{\rm Pl}}
 \frac{b_i g_i^2}{8\pi^2}+\cdots .
\]
Hence the printed one-loop source is the Gaussian tangent of the same native
exact source-response functional, not an independent interaction.
\end{theorem}

\begin{corollary}[Exact K16 gauge-pump functional]
\label{p12:cor:v2177-exact-k16-gauge-pump}
Insert Theorem~\ref{p12:thm:v2177-exact-scalaron-gauge-vertex} into the KIN16
production kernel and evaluate the source at the generated resonance
resolution $k_*=m_s/2$.  Then
\[
 \boxed{
 \Gamma_g^{\rm exact}(t)=
 \frac{\alpha^{-1}m_s(t)^3}{16\pi\bar M_{\rm Pl}^{\,2}}
 \sum_i d_i
 \left[\frac{\beta_i(k_*(t))}{g_i(k_*(t))}\right]^2,}
\]
\[
 \boxed{
 Q_g(t)=\Gamma_g^{\rm exact}(t)\rho_{\rm coh}(t)
 \left(1-\frac4{a^2m_s^2}\right)
 [1+O(\varepsilon_{\rm WKB})].}
\]
The threshold factors are generated by the same exact blocking flow; on the
primitive heat representative their Gaussian projection is
$w_r(k)=e^{-m_r^2/k^2}$.  Thus there is no independent gauge coupling,
renormalization-scale selector, threshold vector, decay coefficient or
reheating efficiency in the K16 source law.
\end{corollary}

\begin{remark}[Why long one-loop transport is not the closure theorem]
\label{p12:rem:v2177-no-long-gaussian-extrapolation}
The common-heat boundary is strongly coupled.  The Gaussian beta coefficients
are the local tangent of the exact finite flow, not a license to extrapolate
the common-heat values over the full interval to $k_{*,16}$.  The physical K16
source therefore consumes the exact finite supertrace/source response
\[
 \partial_{\log k}\Gamma_k^\partial
 =\frac12\operatorname{STr}
 [ (\Gamma_k^{\partial(2)}+\mathcal R_k)^{-1}
   \partial_{\log k}\mathcal R_k]
\]
and its gauge projection.  A named-scheme coordinate may be used for
comparison after this native transport is evaluated, but is not a missing
physical owner.
\end{remark}

\begin{corollary}[RG16 scale/matching subproblem retyped]
\label{p12:cor:v2177-rg16-scale-retyped}
The K16 source-scale selection and the scalaron--gauge vertex are closed
analytically:
\[
 \boxed{\mathrm{RG16}_{\rm scale}=\texttt{CLOSED\_EXACT}[k_*=m_s/2]},
\]
\[
 \boxed{\mathrm{RG16}_{\phi FF}=\texttt{CLOSED\_EXACT}
 [\lambda_i=2\alpha^{-1/2}(\beta_i/g_i)/\bar M_{\rm Pl}].}
\]
The first unevaluated object is now finite and explicit: evaluate the exact
finite gauge projection $\beta_i(k_*)/g_i(k_*)$ (including its generated
threshold response) on the actual K16 finite effective action and insert it
into the already closed coupled scalaron/material evolution.  This is an
execution problem, not a missing coupling or scale law.
\end{corollary}


% Framework v2.178: exact-source coordinate retyped below finite production owner.
\begin{corollary}[Continuum exact-source transport is a descendant of the finite K16 owner]
\label{p12:cor:v2178-continuum-source-descendant}
The exact identities
\[
 k_*=m_s/2,
 \qquad
 \lambda_i^{\rm exact}(k)=
 \frac{2\alpha^{-1/2}}{\bar M_{\rm Pl}}\frac{\beta_i(k)}{g_i(k)}
\]
remain valid source-coordinate statements on the admitted finite effective
action and local continuum/WKB reduction.  However, the physical K16
production amount is owned upstream by the action-generated ordered kernel
$K_{\rm ord}(t)$, Paper~IV's finite Nambu/Bogoliubov map, and
$\Gamma_{16}^{\rm ord}$.  Therefore evaluation of $\beta_i/g_i$ is a useful
continuum projection/precision-rate calculation, not the controlling first
open C06 object and not a substitute for the finite production history.
\end{corollary}

\begin{remark}[Framework-native precedence]
\label{p12:rem:v2178-framework-native-precedence}
Named-scheme RG, perturbative Standard-Model running, conventional reheating
widths and textbook fixed-$R^2$ scalaron formulas carry proof weight only
where a current framework theorem derives them as limits or comparison
coordinates of the constituted finite process.  Native finite owners take
precedence.
\end{remark}

% END INLINE SOURCE: Paper7C-K16ExactSourceTransport-Main-v1_1.tex




% BEGIN GEOMETRY-FIRST CONSEQUENCES v2.201
\setlength{\emergencystretch}{2em}


% BEGIN FRAMEWORK v2.323: PETERSEN CONTINUUM NON-PROMOTION CERTIFICATES
\section{Petersen finite channels and the continuum QFT export boundary}
\label{p12:sec:v323-petersen-continuum-boundary}
\noindent\textit{Status: necessary action and refinement conditions, not an
additional continuum field or a derived Yang--Mills mass gap.}
For an occurring embedded $2I$, the rank-fifteen Petersen edge module
has a finite self-adjoint mirror contrast
$D=\Pi_3-\Pi_{3'}$ with spectrum $+1^{(3)},-1^{(3)},0^{(9)}$.
This is an exact finite module over the internal geometry; its
bare spectrum is \emph{not} a spectrum of the continuum QCD Hamiltonian.

\begin{proposition}[Finite Petersen resolution does not select an infrared scale]
\label{p12:prop:v323-petersen-no-continuum-gap}
An $A_5$-invariant real self-adjoint comparison operator on the
fifteen-edge module has the form
\[
 B=b_4\Pi_4+b_5\Pi_5+b_3\Pi_3+b_{3'}\Pi_{3'}.
\]
The geometry fixes its four multiplicities, not the coefficients.
In particular no positive, uniform, physical color-singlet mass gap,
Wilson area law, or asymptotic string tension follows solely from this
finite decomposition or from the signed $\mathbb Z_2$ Petersen loops.
\end{proposition}
\begin{proof}
The real $A_5$ decomposition is multiplicity free, so real Schur's
lemma gives the stated operator family.  Rescaling or varying its
coefficients changes all candidate energy differences without changing
the geometric representation.  The signed loop is a finite internal
central sign; it is not a nontrivial $\mathbb Z_3$ color-center map,
not a family of spatial Wilson contours with increasing area, and
not a uniform refinement estimate.  The actual Paper--VII-A action
and Paper--VII-C limit certificates remain necessary.
\end{proof}

\begin{corollary}[Source-first continuum handoff]
\label{p12:cor:v323-petersen-handoff}
Any finite Petersen-sensitive QCD source must first pass the
Paper--VII-A color/gauge-compatible incidence test and the actual
Paper--VII-B finite pole and residue test.  An ordinary local-field or
scattering interpretation thereafter requires the same state, domain,
refinement and locality certificates as every other Paper--VII-C
observable.  The new geometry neither adds a continuum field by
counting graph vertices nor changes the inherited beta-function
coefficients.
\end{corollary}
% END FRAMEWORK v2.323: PETERSEN CONTINUUM NON-PROMOTION CERTIFICATES


\section{Local composites, source analyticity and claim-indexed fields}
\label{p12:sec:v237-c7}
Bounded reconstruction can precede optional unbounded point fields, LSZ channels, scattering matrices and infrared completion. Each extension is attached only to claims which consume it.

% Source Paper VII-C, original line 1836; inherited proof bodies unchanged.
\subsection{Small-observable and local-field reconstruction}
\label{p12:sec:small-field-reconstruction}


% Source Paper VII-C, original line 1839; inherited proof bodies unchanged.
\subsubsection{Renormalized finite representatives}
\label{p12:sec:renormalized-finite-representatives}

The finite Born--Hopf formulation naturally starts with bounded finite observables.  Continuum fields are typically unbounded distributions.  Therefore the field reconstruction problem must be expressed as a renormalized limit of smeared finite observables.

\begin{definition}[Renormalized small-observable family]
\label{p12:def:small-observable-family}
Let \(\Phi\) be a target local field type and let \(f\in C_c^\infty(\RR^4)\).  A renormalized small-observable representative of \(\Phi(f)\) is a family
\begin{equation}
\label{p12:eq:renormalized-small-observable}
\Phi_L(f)
:=
Z_L^{-1}\sum_{x\in\Gamma_L} a_L^4 f(\chi_L(x))\,\mathcal P_L(x),
\end{equation}
where \(\mathcal P_L(x)\in\A_L\) is a local finite observable and \(Z_L\) is fixed by Born--Hopf normalization, covariance, or an already proved residue theorem, not by comparison with the target continuum value.
\end{definition}


% Source Paper VII-C, original line 1856; inherited proof bodies unchanged.
\subsubsection{Distributional limits and field reconstruction}
\label{p12:sec:distributional-limits-field-reconstruction}

\begin{definition}[Distributional field limit]
\label{p12:def:distributional-field-limit}
The family \(\Phi_L\) reconstructs a local field \(\Phi\) if all mixed moments
\begin{equation}
\label{p12:eq:mixed-field-moments}
\lim_{L}\omega_L\left(A_{0,L}\Phi_L(f_1)\cdots\Phi_L(f_n)B_{0,L}\right)
\end{equation}
exist for local bounded \(A_0,B_0\), are multilinear and continuous in the test functions, and satisfy the same OS positivity and covariance bounds as the bounded-cylinder sector.
\end{definition}

\begin{theorem}[Small-observable field reconstruction]
\label{p12:thm:small-observable-field-reconstruction}
Assume QFT1--QFT7 and QFT10 for a field family \(\Phi_L\).  Then \(\Phi_L\) defines an operator-valued distribution affiliated with the OS-reconstructed local net.
\end{theorem}
\begin{proof}
QFT10 gives existence and test-function continuity of the Euclidean distributional moments.  Reflection positivity and OS regularity pass to the mixed bounded-field moments by the same limiting argument as Lemma~\ref{p12:lem:rp-survives-limit}.  The OS reconstruction theorem then reconstructs the corresponding Wightman distributions.  Locality follows from the local support condition in Definition~\ref{p12:def:small-observable-family} together with asymptotic local commutativity.
\end{proof}


% Source Paper VII-C, original line 8926; inherited proof bodies unchanged.
\subsection{Literature-checked fixed-UV observable and field export certificates}
\label{p12:sec:v2210-lit-export-certificates}
The native branch retains $m_p=1$ and the fixed physical ceiling of
Theorem~\ref{p12:thm:v283-fixed-uv-born-continuum}. The following proofs
identify certificates actually consumed by exact ordinary local-QFT,
source-field or scattering comparisons. They do not reinstate ultraviolet
removal, a universally pure vacuum, an ontic collapse law or an unconditional
continuum Yang--Mills mass-gap target.


% Source Paper VII-C, original line 8935; inherited proof bodies unchanged.
\subsubsection{Representation-wide locality and the bounded-net route}
\begin{lemma}[Locality from native GNS sandwich control]
\label{p12:lem:v2210-lit-locality-sandwich}
Let $A_L,B_L,D_L$ be compatible uniformly bounded cylinder representatives
with all required product moments converging to a positive state $\omega$.
Put $C_L=[A_L,B_L]$. If for every $D$ in a norm-dense cylinder algebra
\[
 \omega_L(D_L^*C_L^*C_LD_L)\longrightarrow0,
\]
then in the limiting GNS representation
$[\pi(A),\pi(B)]=0$. Vacuum control
$\omega_L(C_L^*C_L)\to0$ alone need not imply this operator statement.
\end{lemma}
\begin{proof}
Product compatibility gives
$\|\pi(C)\pi(D)\Omega\|^2=\omega(D^*C^*CD)=0$.
The vectors $\pi(D)\Omega$ are dense by GNS cyclicity, and $\pi(C)$ is
bounded, so it vanishes. For the second assertion take $M_3(\mathbb C)$,
$\Omega=e_1$, $A=E_{23}+E_{32}$ and $B=i(E_{23}-E_{32})$.
Their nonzero commutator kills $\Omega$, while $M_3\Omega=\mathbb C^3$.
Thus its vacuum norm is zero but it does not vanish on the representation.
\end{proof}
This is a sufficient upgrade of Definition~\ref{p12:def:asymptotic-local-commutativity}
when an exact representation-wide commutator is claimed. Norm-locality
control is another sufficient route. The actual product/null-ideal
comparison remains part of the cylinder construction.

\begin{definition}[Bounded Euclidean reconstruction certificate]
\label{p12:def:v2210-lit-bounded-os-certificate}
On a selected native phase, a bounded algebraic OS certificate consists of
the following actual limit data:
\begin{enumerate}
\item an isotonic Euclidean local $C^*$ net generated by the compatible
bounded cylinder questions, commuting for disjoint Euclidean regions;
\item a covariant $E(4)$ action on that net, including time--space rotations,
and a normalized invariant reflection-positive functional;
\item regularity in the form of continuity of
$g\mapsto\omega(a\alpha_g(b)c)$ for each fixed $a,b,c$;
\item a nontrivial time-zero algebra: the reflection-fixed algebra in the
intersection of the positive- and negative-time half-space algebras,
with its actual spatial local subalgebras.
\end{enumerate}
All products, reflections and transports used here must descend through
the same native comparison/null quotient. An independently verified
equivalent reconstruction theorem may replace this sufficient certificate.
A field-based replacement consumes the corrected all-order OS
regularity/growth and locality hypotheses, not temperedness alone
\cite{OsterwalderSchrader1973,OsterwalderSchrader1975}.
\end{definition}

\begin{corollary}[Conditional bounded-first reconstruction on the native phase]
\label{p12:cor:v2210-lit-bounded-os-reconstruction}
If the actual phase-selected bounded native limit has
Definition~\ref{p12:def:v2210-lit-bounded-os-certificate}, its time-zero
observables reconstruct a positive-energy Poincar\'e-covariant Haag--Kastler
net by the algebraic Euclidean reconstruction theorem
\cite{LitQFTSchlingemann1999}. No point-field extension, positive mass
gap or pure-vacuum selection is required by this sufficient route.
Its transported time-zero moments and bounded event probabilities are the
readouts of the same limiting functional. Equality of this reconstructed
subnet with a larger requested observable export additionally consumes
the actual generation/identification of those observables.
\end{corollary}
\begin{proof}
The native common cylinder and quotient maps supply the Euclidean net and
functional in the hypothesis; reflection, locality and $E(4)$ covariance
are certified there rather than inserted as new finite physical owners.
The regular reflection-positive functional gives the reflection quotient
Hilbert space and positive-time semigroup. The cited theorem reconstructs
the positive-energy Poincar\'e action and, from the actual time-zero
algebras, the local bounded net. Its construction retains the functional
on the transported time-zero products, so trace-dual bounded event readout
agrees with the inherited limiting state. That theorem does not assert
that an unproved larger observable family is generated by the time-zero
subnet, explaining the last clause.
\end{proof}
Weak-* compactness, positive transfer and bounded regularity alone do not
discharge the displayed locality, full covariance and time-zero clauses.
Thus Theorem~\ref{p12:thm:main-qft} is an exact ordinary-QFT implication
on its actual reconstruction certificate, while the finite native
probabilities and positive Hamiltonian remain valid on their own domains.

\begin{proposition}[Product native covariance does not transport a mixed rotation]
\label{p12:prop:v2210-lit-covariance-scope}
In the local product charts of this paper, native Euclidean-time translation
and $\Spin(4)$ spatial isometries preserve the covector $d\tau$.
Their geometric flat-chart limits therefore do not themselves represent
an ordinary Euclidean rotation mixing $\tau$ with a spatial coordinate.
An actual limiting state may have that larger covariance only if its
time--space transport/correlation comparison is separately verified.
\end{proposition}
\begin{proof}
Every such product transformation has $\tau' =\tau+a$, hence
pulls $d\tau'$ back to $d\tau$. This relation survives its product-chart
geometric limit. A nontrivial mixed rotation instead pulls $d\tau'$ back
to $\cos\theta\,d\tau+\sin\theta\,dx^i$. It cannot be that product
transformation. Additional invariance of the state is a property of its
correlations, not a conclusion from this geometric action.
\end{proof}
The normalized invariant ground projector still closes native state
covariance by Theorem~\ref{p12:thm:native-bh-covariance}. Its actual
flat-state comparison must include all transformations claimed in an
ordinary $E(4)$ export.


% Source Paper VII-C, original line 9039; inherited proof bodies unchanged.
\subsubsection{A fixed-band comparison diagnostic and its operational scope}
\begin{proposition}[Fixed-band free comparison kernel]
\label{p12:prop:v2210-lit-fixed-band-kernel}
For a canonical free comparison field with spatial Fourier support
$|k|\le\Lambda_F$, the equal-time derivative commutator has the
band-projector kernel
\[
 \delta_{\Lambda_F}(r)=\frac1{(2\pi)^3}
    \int_{|k|\le\Lambda_F}e^{ik\cdot x}\,d^3k
 =\frac{\sin(\Lambda_Fr)-\Lambda_Fr\cos(\Lambda_Fr)}{2\pi^2r^3},
 \qquad r=|x|>0.
\]
It is not identically zero on any open interval of $r>0$ and does not
vanish by reciprocal-position volume growth at fixed $\Lambda_F$.
The short-time free commutator therefore has a nonzero term proportional
to $t\delta_{\Lambda_F}(r)$ outside its nominal cone.
\end{proposition}
\begin{proof}
Angular integration gives $4\pi\sin(kr)/(kr)$, and integrating
$k\sin(kr)$ from zero to $\Lambda_F$ gives the formula. It is an
analytic nonzero function for $r>0$; its regular value at zero is
$\Lambda_F^3/(6\pi^2)$. The canonical comparison commutator vanishes
at $t=0$, while its first derivative is the displayed kernel up to the
fixed canonical factor. Taylor expansion gives the last statement.
\end{proof}
This is a diagnostic for that comparison field, not a newly inserted
native field and not a requirement to remove the protected momentum ceiling.
The actual native observable net may use a different licensed quotient or
support grammar; an exact ordinary local identification requires its map.
The existing approximate-microcausality bounds in Paper~IV remain
appropriate for detector resolutions and separations, with all DLOC,
spectral and occupation errors retained.

\begin{lemma}[Open-region spanning does not create exact spectral support]
\label{p12:lem:v2210-lit-spectral-support}
Let $E$ be a finite retained harmonic space on the round $S^3$, and
$\Pi_E$ its orthogonal projector. For every nonempty open $U\subset S^3$,
$\operatorname{span}\{\Pi_Ef:f\in C_c^\infty(U)\}=E$.
Consequently, labelling every projected smear by its preprojection support
alone need not yield different local field algebras.
\end{lemma}
\begin{proof}
A vector in $E$ orthogonal to all projected smears is orthogonal to all
$f\in C_c^\infty(U)$, so vanishes on $U$. A finite sum of round-sphere
harmonics is real analytic; analytic continuation on the connected sphere
makes it zero everywhere. The orthogonal complement is therefore zero.
If the field constructor uses all those projected mode vectors, it may
generate the same whole retained mode algebra in every such region.
\end{proof}
This preserves Paper~IV's distinction between exact-support subspaces and
approximately localized spectral ranges. Full-spacetime Laplace spectral
cutoffs can also fail reflection positivity
\cite{LitQFTBailleulCutoff2023}. That result does not invalidate the
native spatial-cutoff/positive-Hamiltonian time-transfer construction;
any claimed equality to a spacetime Euclidean cutoff must verify its own
reflected kernel and support comparison.


% Source Paper VII-C, original line 9096; inherited proof bodies unchanged.
\subsubsection{Uniform source analyticity as an executable field certificate}
\begin{theorem}[Fixed-source holomorphic compactness and derivative transfer]
\label{p12:thm:v2210-lit-source-analyticity}
Let $Z_L(J)$ be normalized generating functions of one actually owned,
fixed finite source tuple, holomorphic on a common connected polydisk
$D\subset\mathbb C^q$ centered at zero, with $Z_L(0)=1$.
Assume local uniform boundedness on $D$ and pointwise convergence on a
nonempty real open neighborhood of zero. Then $Z_L$ converges locally
uniformly on $D$ to a unique holomorphic $Z$. After a cofinal restriction, on a smaller common
neighborhood of zero, $W_L=\log Z_L$ has its zero-source branch and
$W_L\to W$ with every fixed derivative.
For concentric polydisks of radii $r<r'<R$ inside $D$, a proved comparison
$\sup_{|J_j|\le r'}|Z_M-Z_L|\le\epsilon_{ML}$ gives
\[
 \sup_{|J_j|\le r}|\partial^\alpha Z_M-\partial^\alpha Z_L|
 \le\alpha!\,(r'-r)^{-|\alpha|}\epsilon_{ML}.
\]
\end{theorem}
\begin{proof}
Local boundedness and Cauchy's formula give a normal family in the compact-open
topology. Every holomorphic cluster limit agrees on the real neighborhood
by the assumed pointwise convergence. The identity theorem, successively
in the $q$ coordinates, makes any two limits equal on connected $D$.
Compactness and uniqueness of cluster limits give convergence of the net.
Since $Z(0)=1$, choose a smaller neighborhood with $|Z-1|<1/4$.
Local uniform convergence makes $|Z_L-1|<1/2$ eventually, giving a common
logarithm there. The multivariable Cauchy formula proves logarithmic
derivative convergence and the displayed error bound for $Z$.
\end{proof}
For the bounded cylinder family, the uniform budget and convergence are
already supplied by Theorem~\ref{p12:thm:v202-bounded-source-inheritance}.
For unbounded or independently normalized residues, common analyticity and
local uniform bounds are additional physical-source estimates; finite
dimensionality does not make them refinement uniform. A derivative of
$\omega(e^{\sum J_jO_j})$ supplies symmetrized moments of noncommuting
$O_j$, not arbitrary ordered products. Each ordered/reflected/contact
family needed for QFT10 must therefore be generated by its actual
ordered source functional or separately certified mixed moments.
For actual bounded observable words,
Theorem~\ref{p12:thm:v256-bounded-ordered-source} now constructs
that ordered readout and its common zero-free chart directly from
the same algebra and positive state.  It does not establish equality
to an independently specified action/contact-source functional.

\begin{corollary}[Source-separated canonicality without a universal gap]
\label{p12:cor:v2210-lit-source-canonicity}
Suppose the preceding actual source certificates are path independent and
their controlled derivatives determine every word of a norm-dense bounded
local cylinder algebra. Then all selected-phase cluster states agree on
that algebra and QFT4$_C$ holds. No positive full-sector gap is required
by this sufficient route. If the sources determine only a proper readout
algebra, only its canonical export follows.
\end{corollary}
\begin{proof}
Derivative convergence fixes each consumed word expectation. For any
bounded local $A$, choose a cylinder polynomial $B$ with
$\|A-B\|<\epsilon$. Two positive normalized cluster states then differ
on $A$ by at most $2\epsilon$ because they agree on $B$. Send
$\epsilon\to0$ and apply the singleton-cluster criterion of
Theorem~\ref{p12:thm:v283-qft4-cluster-split}.
\end{proof}
The existing source-stable polymer and SURG implementations may prove
these analytic hypotheses on their actual activity domain
\cite{KoteckyPreiss1986}. Neither the exact coarse square nor a chosen
optional strong-heat convention proves the actual word comparisons,
growing-volume source bounds or independent action-source identifications.
The fixed bounded-readout analytic chart itself is now supplied by
Theorem~\ref{p12:thm:v256-bounded-ordered-source}.
For field distributions, locality/RP, test-function growth and all required
mixed moment identifications remain the QFT10 hypotheses of
Theorem~\ref{p12:thm:ucs-implies-qft10}.


% Source Paper VII-C, original line 9160; inherited proof bodies unchanged.
\subsubsection{A nonlinear readout is not by itself a dynamical interaction witness}
\begin{proposition}[An exact free-composite cumulant diagnostic]
\label{p12:prop:v2210-lit-composite-cumulant}
Let $X$ be a zero-mean real Gaussian variable of variance $v>0$ and
$Y=X^2-v$. On $|2vJ|<1$ its normalized cumulant generator is
\[
 W_Y(J)=-vJ-\tfrac12\log(1-2vJ),\qquad
 \kappa_n(Y)=2^{n-1}(n-1)!v^n\quad(n\ge2).
\]
In particular $\kappa_3(Y)=8v^3\ne0$, despite Gaussian underlying data.
For Wick squares of a centered free Gaussian field, the separated-point
connected three-point function is $8C_{12}C_{23}C_{31}$.
\end{proposition}
\begin{proof}
The Gaussian integral gives
$\mathbb E e^{J(X^2-v)}=e^{-vJ}(1-2vJ)^{-1/2}$.
Expand its logarithm and differentiate. For the field statement Wick
pairings in the connected three-vertex graph must link each squared vertex
to both other vertices; there are eight such pairings. Subtraction of the
one-vertex contractions is exactly Wick ordering.
\end{proof}
This is a diagnostic, not a new physical source. It agrees with the
Gaussian/Wick-square comparison literature
\cite{LitQFTDangWickSquares2022}. The existing higher-cumulant criterion
therefore establishes non-Gaussianity of its chosen observable family.
For nonlinear Wilson/curvature/composite sources it does not alone exclude
Gaussian underlying dynamics.

\begin{corollary}[Scope of a same-source free-comparator test]
\label{p12:cor:v2210-lit-interaction-comparator}
On an already licensed actual source chart and its licensed noninteracting
comparison family, an executed nonzero difference between the actual
connected/amputated readout and every admissible noninteracting comparator
excludes that comparator family on the stated readout domain. A single
named comparator excludes only that comparator. Nonfactorizing finite
collision probabilities may instead furnish a dynamical witness on their
actual normalized channel domain.
\end{corollary}
\begin{proof}
Equality to any admitted comparator would force equality of the same-source
readout. Its nonzero difference contradicts that equality. The finite
collision version is the same argument applied to the action-owned channel
probability rather than a nonlinear composite source law.
\end{proof}
QFT8 remains optional for QFT existence. This sharper distinction preserves
the finite non-Gaussian observable theorem while keeping a stronger
particle-interaction claim tied to an actual action/source or collision test.


% Source Paper VII-C, original line 9208; inherited proof bodies unchanged.
\subsubsection{Chiral projector and quantum-identity comparison budgets}
\begin{proposition}[Source-gap control of an actual finite chiral projector]
\label{p12:prop:v2210-lit-chiral-projector}
Let $A(s)$ be the complete $C^1$ Hermitian overlap/Wilson kernel on an
admitted finite source path, with $\operatorname{spec}A(s)$ avoiding
$(-g,g)$, $g>0$. Put $P_-(s)=\mathbf1_{(-\infty,0)}(A(s))$.
Then $P_-$ is $C^1$, has constant rank, and
\[
 \|\partial_sP_-\|_{\rm HS}\le\frac{\|\partial_sA\|_{\rm HS}}{2g}.
\]
In particular the native sign-index trace is constant on that path.
\end{proposition}
\begin{proof}
Riesz spectral calculus on the gapped finite path gives differentiability.
Differentiating $P_-^2=P_-$ shows its derivative has only off-diagonal
positive/negative spectral blocks. Differentiating $[A,P_-]=0$ gives
$[A,P_-']=-[A',P_-]$. In an instantaneous spectral basis each retained
off-diagonal entry is divided by an eigenvalue separation at least $2g$.
Sum the squared entries for the Hilbert--Schmidt bound. Rank is a continuous
integer, and the sign-index trace is determined by that rank.
\end{proof}
The gap here belongs to the actual kernel/source family, not a claimed
uniform interacting Hamiltonian mass gap. Uniform gap and derivative
bounds supply projector control but not determinant-line holonomy or
weights of topological sectors. Paper~IV's existing global parallel-line
criterion still requires its flat connection and trivial holonomy.
Abelian GW measure construction, nonabelian/global anomaly tests and
classical overlap-index comparisons
\cite{LitQFTLuscherAbelian1999,LitQFTLuscherNonabelian2000,
LitQFTBarCampos2000,LitQFTAdamsIndex2001} are sufficient routes only on
their actual admitted families. Smooth fine-lattice torus hypotheses are
not an obligation to remove the native fixed UV ceiling. The free-overlap
RP theorem \cite{LitQFTKikukawaUsui2010} does not certify arbitrary
gauge-coupled chiral reflection positivity. These remain the single G39
ordered chiral/topological certificate, not a reopened native G20 owner.

\begin{proposition}[Quantitative transport of a finite quantum identity]
\label{p12:prop:v2210-lit-ward-budget}
Suppose the actual regulated finite identity has
$\sum_{r=1}^N c_{r,L}\omega_L(A_{r,L})=d_L$,
$c_{r,L}\to c_r$, $\omega_L(A_{r,L})\to a_r$,
$|\omega_L(A_{r,L})-a_r|\le\epsilon_{r,L}$ and $|a_r|\le M_r$.
Then
\[
 \left|\sum_rc_ra_r\right|
 \le |d_L|+\sum_r\bigl(|c_{r,L}|\epsilon_{r,L}
              +|c_{r,L}-c_r|M_r\bigr).
\]
If the right side tends to zero, the limiting identity holds, including
every regulator, measure, contact and boundary insertion retained among
the $A_{r,L}$.
\end{proposition}
\begin{proof}
Subtract the finite identity, telescope its coefficients and expectations,
and apply the triangle inequality.
\end{proof}
This provides an executable error budget for the existing finite Ward,
anomaly and Schwinger--Dyson transport theorem. It does not turn classical
Noether invariance into an invariant quantum measure or discard an
unevaluated insertion. Unbounded/contact terms consume the actual mixed
source certificate above; the actual action/measure remains upstream.

\paragraph{Retained claim boundaries.}
The optional QFT9 positive-gap transfer remains conditional on the actual
physical uniform gap and semigroup/vector comparison; a positive gap in
each finite matrix does not prove a positive infimum. The massive
finite-window detector bounds and finite inclusive soft identity are now
\ref{p4:prop:v2210-lit-detector-stability} and
\ref{p4:cor:v2210-lit-inclusive-soft}; neither supplies its remaining
scattering/infrared histories. Finally, exact RG algebra is not radiative
protection: \ref{p4:thm:v2210-lit-determinant-jets} supplies the complete
finite Gaussian source-jet test, while complete interacting effective
response and any stronger M1--M2 crossover remain on their own actual
action/source domains.
% END LITERATURE-CHECKED QFT EXPORT CERTIFICATES v2.210

% BEGIN FRAMEWORK v2.226: KINEMATIC CLOSURES


\section{Main QFT export theorem and explicit limitations}
\label{p12:sec:v237-c8}
The finite theory is already a quantum theory on its finite domains. The phase-selected bounded-continuum export is conditional on QFT certificates; uniform gaps, interacting witnesses and detector scattering have distinct proof obligations.

% Source Paper VII-C, original line 4511; inherited proof bodies unchanged.
% Former logical part: Main Certificate Theorem; integrated into the new thematic argument.

\noindent\textbf{Orientation.}
This part states the final QFT1--QFT10 status census, the continuous-QFT certificate theorem, its downstream use, and the conclusion.


% Source Paper VII-C, original line 4516; inherited proof bodies unchanged.
\subsubsection{Exact finite interaction compilation and its QFT handoff}
\label{p12:sec:v199-finite-tensor-handoff}
Theorems~\ref{p10:thm:v199-quartic-heat-vertex} and
\ref{p10:thm:v199-haar-interaction-compiler} strengthen the formerly G33 finite input, now closed on the native branch: the quartic face vertex is fixed by the existing heat stiffness and
its derivative, and normalized Haar elimination gives every finite blocked
interaction as a convergent representation/intertwiner contraction.
The tensor class is closed under finite gluing.  This supplies the full
finite boundary weight, rather than a product-of-faces closure assumption.

\begin{corollary}[Controlled fixed-block action and bounded readouts]
\label{p12:cor:v199-finite-block-approximation}
Fix a finite block and positive face times.  Let \(F\) be its full boundary
weight, \(F_L\) the total-degree character truncation, and \(\mu\) normalized
Haar measure on its boundary variables.  If \(F\ge m>0\) and the explicit
Paper-VII-A tail bound gives \(\|F-F_L\|_\infty\le\delta<m\), then
\[
 F_L>0,\qquad \|\log F-\log F_L\|_\infty
 \le-\log(1-\delta/m).
\]
Writing \(\langle O\rangle_F=\int OF\,d\mu/\int F\,d\mu\), every
bounded boundary readout satisfies
\[
 |\langle O\rangle_F-\langle O\rangle_{F_L}|
 \le 2\|O\|_\infty\delta/m.
\]
\end{corollary}
\begin{proof}
Positivity and the logarithmic bound follow pointwise from
\(|F_L/F-1|\le\delta/m<1\).  Put \(Z=\int Fd\mu\),
\(Z_L=\int F_Ld\mu\).  Then \(Z\ge m\),
\(|Z-Z_L|\le\delta\), and the difference of normalized expectations is
bounded by \(\|O\|_\infty\delta/Z\) for each of the numerator and
normalization terms.  Sum the two bounds.
\end{proof}

The hypotheses are achievable for each fixed positive-time finite block.
On the isotropic native heat branch, Theorem~\ref{p10:thm:v202-native-heat-bounds}
now supplies the explicit minimum $m=a^N$. Theorems
\ref{p12:thm:v202-bounded-source-domain}--\ref{p12:thm:v202-bounded-source-inheritance}
below control fixed bounded-source derivatives and their state-limit
inheritance. These finite results do not assign the actual continuum state
or provide a size-independent interaction norm as blocks grow and heat times shrink.
In particular the complete-action norm interface of
Definition~\ref{p12:def:v198-complete-action-norm-interface}, mixed
heat/defect polymer domination and strict stopping membership in
Proposition~\ref{p12:prop:v198-conditional-stopping-budget} retain all their
stated premises.  A local quartic coefficient and a finite sup-norm tail do
not supply transfer to that fixed decomposition's interaction norm.

The G20 cover theorem~\ref{p10:thm:v199-native-clock-covers} similarly
provides uniform regulator gaps on its clock-cover family.  Its parameter
is not spatial/ultraviolet refinement, and its gap is not the QFT physical
mass gap.  The locality theorem~\ref{p10:thm:v202-native-overlap-locality}
now also applies on the certified all-link neighborhoods. These results
leave physical continuum reconstruction, reflection, unbounded-source
and growing-domain estimates, and occurrence requirements explicit.

% BEGIN V202 BOUNDED SOURCE CLOSURE

% Source Paper VII-C, original line 4574; inherited proof bodies unchanged.
\subsubsection{Bounded-source domains and connected derivatives}
\label{p12:sec:v202-bounded-sources}
The following statements use the actual positive finite weight and a
fixed bounded source tuple. They supply the source part of the fixed-block
handoff without assuming a growing-block interaction norm.

\begin{theorem}[A common zero-free domain from positivity]
\label{p12:thm:v202-bounded-source-domain}
Let $F\ge0$, $0<Z_F(0)=\int Fd\mu<\infty$, and let real cylinders
$O_j$ satisfy $|O_j|\le B_j$. Define
\[
 Z_F(J)=\int F e^{\sum_jJ_jO_j}d\mu,\qquad
 W_F(J)=\log Z_F(J)-\log Z_F(0).
\]
For every $r<\pi/2$, on $\sum_j|J_j|B_j\le r$,
\begin{equation}
 \Re Z_F(J)\ge\cos r\int F e^{\Re\sum_jJ_jO_j}d\mu
 \ge\cos r\,e^{-r}Z_F(0)>0.
 \label{p12:eq:v202-source-half-plane}
\end{equation}
The partition function is entire and its logarithm has the canonical
right-half-plane branch on this domain. The domain is uniform over any
family of positive weights with the same observable bounds, including
positive interpolating weights $e^{-S_{\rm heat}-s\mathcal D}$,
$0\le s\le1$, whenever these are finite real weights.
\end{theorem}
\begin{proof}
Write $\sum J_jO_j=u+iv$. Pointwise $|u|,|v|\le r$, so
$\Re e^{u+iv}=e^u\cos v\ge e^u\cos r$. Integrate against $F$.
Boundedness gives an entire exponential series and the positive real
part supplies the logarithm.
\end{proof}

\begin{theorem}[Fixed-block source and cumulant errors]
\label{p12:thm:v202-source-errors}
Let $F>0$ and a real approximant obey
$\|F_L/F-1\|_\infty\le\varepsilon<\cos r$, $r<\pi/2$.
Then $F_L>0$ and on the preceding source domain,
\begin{equation}
 |W_{F_L}(J)-W_F(J)|\le C_r(\varepsilon)
 :=-\log(1-\varepsilon/\cos r)-\log(1-\varepsilon).
 \label{p12:eq:v202-source-log-error}
\end{equation}
For $\rho_j>0$, $\sum_j\rho_jB_j\le r$, every multiindex satisfies
\[
 |\partial^\alpha W_{F_L}(0)-\partial^\alpha W_F(0)|
 \le\alpha!\prod_j\rho_j^{-\alpha_j}C_r(\varepsilon).
\]
Independently, for $n$ real bounded variables with respective bounds
$B_1,\ldots,B_n$, the weaker premise $\varepsilon<1$ already gives
\begin{equation}
 |\kappa_L(O_1,\ldots,O_n)-\kappa(O_1,\ldots,O_n)|
 \le2\varepsilon\!\left(\sum_{k=1}^n k!\,{n\brace k}\right)
       \prod_i B_i
 \le\varepsilon n!2^n\prod_iB_i.
 \label{p12:eq:v202-partition-cumulant-error}
\end{equation}
Here ${n\brace k}$ is a Stirling number of the second kind.
\end{theorem}
\begin{proof}
The partition error is bounded by
$\varepsilon\int F e^{\Re\sum J_jO_j}d\mu$. Divide by
\eqref{p12:eq:v202-source-half-plane} to obtain
$|Z_{F_L}/Z_F-1|\le\varepsilon/\cos r$. The series for
$\log(1+z)$ and the zero-source normalization give the logarithmic
bound; the multivariable Cauchy formula gives its derivatives.

For the real comparison, writing a normalized expectation difference
as a weight-difference term over $Z_F$ plus a normalization term bounds
it by $2\varepsilon$ for every unit-bounded observable. Thus
$\operatorname{TV}(\mu_{F_L},\mu_F)\le\varepsilon$. Apply the
cumulant estimate in Paper~1B, Theorem~\ref{p1b:thm:blocking}
\cite{Kugelmann2026Paper1B}, for the first bound in
\eqref{p12:eq:v202-partition-cumulant-error}. Finally an ordered set
partition is encoded by a permutation and composition of $n$, giving
$\sum_k k!{n\brace k}\le n!2^{n-1}$ and the final inequality.
\end{proof}

\begin{corollary}[Certified finite native Wilson cumulants]
\label{p12:cor:v202-native-cumulants}
For the isotropic heat blocks of
Theorem~\ref{p10:thm:v202-native-heat-bounds}, the degree-17 approximation
has, for $N\le96$ and any real $|O|\le1$,
\begin{equation}
 |\kappa_n^{(17)}-\kappa_n|<9.6\,10^{-74},\qquad 1\le n\le4.
 \label{p12:eq:v202-native-cumulant-margin}
\end{equation}
\end{corollary}
\begin{proof}
The ordered partition numbers for $n=1,2,3,4$ are $1,3,13,75$.
The rational bound of \eqref{p10:eq:v202-native-heat-relative-error} at
$N=96$ is less than $6.35\,10^{-76}$. Thus the largest error is less
than $150(6.35\,10^{-76})<9.6\,10^{-74}$.
\end{proof}

Real and imaginary parts of normalized Wilson traces have bound one by
unitarity. If a fine cylinder does not factor through the boundary
variables, insert it before Haar elimination and apply the same face
product estimates there. A different coarse function cannot replace it.
The stated margin concerns approximation error, not a predicted cumulant.

\subsubsection{Sharp bounded ordered-source charts from the same positive state}
\label{p12:sec:v256-bounded-ordered-sources}
The following construction concerns readouts of the actual bounded
observable algebra and its actual normalized positive state.  Positivity
and normalization are already fixed by Paper~II-A,
Definition~\ref{p2:def:standard-finite-state-space} and
Theorem~\ref{p2:thm:canonical-finite-state-test-evaluation}, as imported
in Paper~IV-A.  Thus no additional probability principle is used.  It supplies
ordered words as well as symmetrized exponential moments.  No new action,
preparation, physical source incidence, or refinement selector is introduced.
An independently claimed action-source functional equals this readout
functional only when its own identification has been established.

\begin{lemma}[Accretivity from the integrated operator-norm budget]
\label{p12:lem:v256-ordered-evolution-accretivity}
Let $B\in L^1([0,T];\operatorname{End}\mathcal H)$ on a finite
Hilbert space, and let $U$ be the absolutely continuous evolution
\[
 U'(t)=B(t)U(t),\qquad U(0)=I,\qquad
 \beta=\int_0^T\|B(t)\|dt<\pi/2.
\]
No self-adjointness or commutation assumption on $B(t)$ is needed.  Then
\begin{equation}
 \operatorname{Re}U(T):=\frac{U(T)+U(T)^*}{2}
 \ \ge\ e^{-\beta}\cos\beta\,I,
 \qquad \|U(T)\|\le e^\beta.
 \label{p12:eq:v256-ordered-accretivity}
\end{equation}
Consequently every normalized positive state $\omega$ satisfies
$\operatorname{Re}\omega(U(T))\ge e^{-\beta}\cos\beta>0$.
\end{lemma}
\begin{proof}
The finite integral equation gives $U$ and its inverse, and the usual
norm integral inequality gives $\|U(T)\|\le e^\beta$.  Fix a unit vector
$x$, write $y(t)=U(t)x$, and put $z(t)=y(t)/\|y(t)\|$.
Invertibility ensures that this normalization is defined.  Almost
everywhere, with $a(t)=\operatorname{Re}\langle z(t),B(t)z(t)\rangle$,
\[
 \begin{gathered}
 \frac{d}{dt}\log\|y(t)\|=a(t),\qquad
 z'(t)=B(t)z(t)-a(t)z(t),\\
 \|z'(t)\|^2=\|B(t)z(t)\|^2-a(t)^2\le\|B(t)\|^2.
 \end{gathered}
\]
Thus $\|y(T)\|\ge e^{-\beta}$.  Regard the Hilbert space as a real
Hilbert space with inner product $\operatorname{Re}\langle\cdot,\cdot\rangle$.
For $c(t)=\operatorname{Re}\langle x,z(t)\rangle$, tangency gives
\[
 |c'(t)|
 =|\operatorname{Re}\langle x-c(t)z(t),z'(t)\rangle|
 \le\sqrt{1-c(t)^2}\,\|z'(t)\|.
\]
Equivalently the angular distance $\theta(t)=\arccos c(t)$ increases
at speed at most $\|z'(t)\|$.  To handle the endpoints without division
by zero, apply the displayed inequality to
$\theta_\varepsilon(t)=\arccos(c(t)/(1+\varepsilon))$ and let
$\varepsilon\downarrow0$.  Since $c(0)=1$, it follows that
$\theta(T)\le\int_0^T\|z'(t)\|dt\le\beta$.  Therefore
\[
 \operatorname{Re}\langle x,U(T)x\rangle
 =\|y(T)\|c(T)\ge e^{-\beta}\cos\beta.
\]
The cosine is positive on the stated budget, so both lower bounds may
be multiplied.  This proves the operator inequality.  A positive
normalized state preserves it.
\end{proof}

\begin{theorem}[Exact ordered readouts on a common zero-free chart]
\label{p12:thm:v256-bounded-ordered-source}
Let $O_1,\ldots,O_q$ be the actual finite bounded operators, with
$\|O_j\|\le B_j$, in a normalized positive state $\omega$.
Zero-bound operators may be omitted.  For any fixed ordered slot list
$\mathbf i=(i_1,\ldots,i_m)$, define
\begin{align*}
 Z_{\mathbf i}(z)
 &=\omega\!\left(e^{z_1O_{i_1}}\cdots e^{z_mO_{i_m}}\right),&
 b_{\mathbf i}(z)&=\sum_{a=1}^m|z_a|B_{i_a},\\
 Z_{\rm exp}(J)
 &=\omega\!\left(e^{\sum_{j=1}^qJ_jO_j}\right),&
 b_{\rm exp}(J)&=\sum_{j=1}^q|J_j|B_j.
\end{align*}
Both functions are entire, are one at zero, and satisfy, for either
source convention and any $b\le r<\pi/2$,
\begin{equation}
 |Z|\le e^r,\qquad
 \operatorname{Re}Z\ge c_r:=e^{-r}\cos r>0.
 \label{p12:eq:v256-ordered-source-half-plane}
\end{equation}
Hence $W=\operatorname{Log}Z$, with $W(0)=0$, has the canonical
right-half-plane logarithm on the whole indicated source chart.

The ordered functional generates each actual word without symmetrizing:
\begin{equation}
 \left.\partial_{z_1}\cdots\partial_{z_m}
 Z_{\mathbf i}(z)\right|_{z=0}
 =\omega(O_{i_1}\cdots O_{i_m}).
 \label{p12:eq:v256-ordered-slot-words}
\end{equation}
More precisely, for a block $C\subseteq\{1,\ldots,m\}$ let
$O_C$ denote the subword with indices in their original slot order.
The connected ordered-slot readout is
\begin{equation}
 \left.\partial_{z_1}\cdots\partial_{z_m}
 W_{\mathbf i}(z)\right|_{z=0}
 =\sum_{\pi\in\Pi_m}(-1)^{|\pi|-1}(|\pi|-1)!
       \prod_{C\in\pi}\omega(O_C),
 \label{p12:eq:v256-ordered-slot-connected}
\end{equation}
where $\Pi_m$ is the set of partitions of the slot set.
Repeated operator labels still have independent slot sources.  Actual
bounded adjoint or reflected representatives can be included as slots
with their own inherited norm bounds.
\end{theorem}
\begin{proof}
Absolute norm convergence of the exponential series proves entirety and
the upper bound.  The ordered product is an endpoint evolution with
piecewise constant generators $z_aO_{i_a}$, executed rightmost factor
first.  Its integrated norm budget is at most $b_{\mathbf i}(z)$.
The single exponential is the evolution with constant generator
$\sum_jJ_jO_j$, whose norm is at most $b_{\rm exp}(J)$.
Lemma~\ref{p12:lem:v256-ordered-evolution-accretivity} applies in both
cases.  Since $e^{-s}\cos s$ decreases for $0\le s<\pi/2$, it gives
\eqref{p12:eq:v256-ordered-source-half-plane}.  The right half-plane
has its fixed scalar logarithm, normalized by $Z(0)=1$.

The coefficient of $z_1\cdots z_m$ in the product is the ordered
operator word, proving \eqref{p12:eq:v256-ordered-slot-words}.
Expanding $\log Z=\sum_{k\ge1}(-1)^{k-1}(Z-1)^k/k$ near zero,
the coefficient with each slot used once is a sum over set partitions.
For a partition with $k$ blocks there are $k!$ block orderings and a
factor $(-1)^{k-1}/k$, giving the coefficient in
\eqref{p12:eq:v256-ordered-slot-connected}.  The scalar block
expectations commute; the operators inside each subword retain their
original order.  Thus reversing the slot order generally changes the
functional and its connected readout.
\end{proof}

\begin{corollary}[Ordered-word errors control analytic and connected readouts]
\label{p12:cor:v256-ordered-source-transfer}
Compare two normalized positive presentations $L,M$ of the same finite
operator alphabet, with common bounds $B_j>0$.  They may use different
finite carriers.  Write $\widehat O_{j,L}=O_{j,L}/B_j$ and set
\[
 \delta_n(L,M)=\max_{j_1,\ldots,j_n}
 \left|\omega_L(\widehat O_{j_1,L}\cdots\widehat O_{j_n,L})
       -\omega_M(\widehat O_{j_1,M}\cdots\widehat O_{j_n,M})\right|,
 \qquad 0\le\delta_n\le2.
\]
For $N\ge0$ and $0<r<\pi/2$, define the executable error budget
\begin{equation}
 E_N(r;L,M)=\sum_{n=1}^N\delta_n(L,M)\frac{r^n}{n!}
                +2\sum_{n>N}\frac{r^n}{n!}.
 \label{p12:eq:v256-ordered-word-error}
\end{equation}
For either generating-function convention of the preceding theorem,
on its fixed source budget $b\le r$,
\begin{equation}
 |Z_L-Z_M|\le E_N(r;L,M),\qquad
 |W_L-W_M|\le\frac{E_N(r;L,M)}{c_r}.
 \label{p12:eq:v256-ordered-log-error}
\end{equation}
No small-error side condition is required.  On a source polydisk with
radii $\rho_a>0$ and total budget at most $r$, every multiindex obeys
\begin{equation}
 |\partial^\alpha W_L(0)-\partial^\alpha W_M(0)|
 \le\alpha!\left(\prod_a\rho_a^{-\alpha_a}\right)
          \frac{E_N(r;L,M)}{c_r}.
 \label{p12:eq:v256-ordered-connected-error}
\end{equation}
Consequently convergence of every fixed ordered word implies locally
uniform convergence of $Z,W$ and every fixed source derivative on the
whole budget chart $b<\pi/2$.  This conclusion applies equally when
one presentation is an already established positive limit state.
\end{corollary}
\begin{proof}
At total degree $n$, the sum of absolute exponential-series weights is
at most $r^n/n!$, both for a single exponential and for an ordered
product.  The normalized word errors through degree $N$ give the first
sum in \eqref{p12:eq:v256-ordered-word-error}; the two tails give its
second sum.  This proves the partition-function estimate.

Both values $Z_L,Z_M$ and the straight segment joining them lie in
$\operatorname{Re}z\ge c_r$.  Integration of $d\operatorname{Log}z=dz/z$
along this segment gives
\[
 |\operatorname{Log}Z_L-\operatorname{Log}Z_M|
 \le |Z_L-Z_M|/c_r.
\]
This proves the logarithmic estimate even when the error exceeds
$c_r$.  The multivariable Cauchy formula gives
\eqref{p12:eq:v256-ordered-connected-error}.
For convergence, first make the factorial tail small by fixing $N$,
then use convergence of the finitely many words of lengths at most
$N$.  The same argument works on every compact source budget and on
slightly larger interior polydisks for derivatives at other points.
\end{proof}

\begin{proposition}[The universal bounded-source radius is sharp]
\label{p12:prop:v256-sharp-ordered-budget}
Without extra restrictions on the actual positive state or operators,
the universal radius $\pi/2$ in the preceding results cannot increase.
\end{proposition}
\begin{proof}
Take $O=\operatorname{diag}(1,-1)$ and
$\omega(A)=\operatorname{Tr}(A)/2$.  Then $\|O\|=1$ and the one-slot
source is $Z(z)=\omega(e^{zO})=\cosh z$, with a zero at $z=i\pi/2$.
Any larger universal budget would contain that zero.
\end{proof}

Thus bounded ordered-word generation and its uniform zero-free source
chart are supplied directly by the same finite algebra and positive
state.  Actual ordered-word convergence or path independence is still
an antecedent of the corresponding limit or canonicality claim.
Neither growing-volume source bounds nor unbounded/contact insertions
are supplied by this fixed bounded-source result.  Equality to a
separately specified action-source response retains its actual
same-owner identification requirement.


\begin{theorem}[Bounded-state convergence supplies source convergence]
\label{p12:thm:v202-bounded-source-inheritance}
Suppose the actual positive finite states have convergence of every
fixed polynomial product in a fixed finite tuple $O_{j,L}$, with common
real bounds $B_j$. Their normalized generating functions converge
locally uniformly for complex sources. On
$\sum_j|J_j|B_j<\pi/2$, their normalized logarithms and every fixed
source derivative also converge. Thus this part of bounded QFT4 implies
bounded-source analytic and cumulant convergence.

For noncommuting bounded operator representatives in normalized positive
states, convergence of all fixed ordered words gives the corresponding
conclusion on the sharp universal chart $\sum_j|J_j|\|O_j\|<\pi/2$.
Actual ordered-slot readouts and their connected responses have the same
chart, by Theorem~\ref{p12:thm:v256-bounded-ordered-source} and
Corollary~\ref{p12:cor:v256-ordered-source-transfer}.
\end{theorem}
\begin{proof}
At each degree, the exponential series has finitely many convergent
moments. On a compact source set of budget $r$ the remaining tail is
bounded uniformly by $\sum_{n>N}r^n/n!$. Truncate, pass to the state
limit and remove the truncation. The normalized lower real-part bound
in \eqref{p12:eq:v202-source-half-plane} then gives convergence of the
logarithms; Cauchy's formula gives all fixed derivatives. In the operator
case, $A=\sum J_jO_j$ satisfies the positive real-part bound
of Theorem~\ref{p12:thm:v256-bounded-ordered-source} on every
compact budget $r<\pi/2$.  The ordered-word series gives the entire
partition-function limit.  The common right-half-plane bound and
Corollary~\ref{p12:cor:v256-ordered-source-transfer} give the
logarithm and derivative limits on the whole asserted chart.
\end{proof}

On this fixed bounded-cylinder class, no separate zero-free hypothesis or
full M3 package is needed for source inheritance. Actual QFT4 comparison,
interacting covariance, unbounded-field normalization, growing-block
interaction norms and physical-distance estimates retain their own
premises. In particular an extensive source whose bound grows with
volume is not assigned a uniform budget by this theorem. Appendix
\ref{p12:app:v202-source-comparison} records the sharp domain obstruction
and a supplementary summable-error criterion, without claiming that
criterion for the native interacting tower.
% END V202 BOUNDED SOURCE CLOSURE


% Source Paper VII-C, original line 4712; inherited proof bodies unchanged.
\subsection{Main theorem of Paper VII-C}
\label{p12:sec:main-theorem}


\begin{longtable}{>{\raggedright\arraybackslash}p{0.14\textwidth} >{\raggedright\arraybackslash}p{0.40\textwidth} >{\raggedright\arraybackslash}p{0.35\textwidth}}
\caption[QFT1--QFT10 final status census]{QFT1--QFT10 final status census.}\label{p12:tab:qft1-qft10-final-discharge-census}\\
\toprule
QFT block & Continuum obligation & Status / proof source \\
\midrule
\endfirsthead
\toprule
QFT block & Continuum obligation & Status / proof source \\
\midrule
\endhead
QFT1--QFT3 & spectral-symbol regulator, local cylinder algebra, and bounded representatives & \textbf{Derived for the bounded symbol sector}: Peter--Weyl/Toeplitz arguments; positive exact quadrature is also derived by Proposition~\ref{p12:prop:tier2-quadrature}. \\
QFT4 & bounded local continuum state & \textbf{Split}: QFT4$_E$ cluster-state existence is derived by weak-* compactness; QFT4$_C$ singleton same-phase canonicity remains a stronger analytic theorem when claimed. \\
QFT5 & reflection positivity & \textbf{Derived phase-relatively}: every weak-* cluster state of finite RP presentations is RP. \\
QFT6 & native covariance and full claimed flat-chart comparison & \textbf{Native state covariance derived from the invariant ground projector}; exact ordinary export additionally certifies time--space transport/correlation comparisons as in \ref{p12:prop:v2210-lit-covariance-scope}. \\
QFT7 & regularity on the claimed reconstruction family & \textbf{Bounded test estimates derived}; algebraic group regularity/time-zero locality or an equivalent corrected field-OS certificate remains claim-consumed. Unbounded/composite moment control belongs to QFT10. \\
QFT8 & observable non-Gaussianity; stronger interaction classification & \textbf{Dependent, not a QFT-existence premise}; higher source cumulants establish observable non-Gaussianity. Nonlinear-source dynamical interaction claims additionally use \ref{p12:cor:v2210-lit-interaction-comparator}. \\
QFT9 & positive-gap or controlled gapless branch & Corollary~\ref{p12:cor:paper11-pure-glue-gap-floor}: positive gap is conditional on \(\gamma_*>0\); otherwise gap status remains unresolved or gapless. \\
QFT10 & local-field extension & \textbf{Split}: source-owned fields descend under the common source/UCS package; independently normalized external composites remain claim-scoped mixed-moment certificates.  Neither is required for the bounded local net. \\
\bottomrule
\end{longtable}
\TableNotes{This is the proof-census form of Theorem~\ref{p12:thm:main-qft}.  It records one discharge source for each obligation block before the final aggregation theorem is invoked.}



% Source Paper VII-C, original line 4739; inherited proof bodies unchanged.
\subsubsection{Continuous-QFT certificate theorem}
\label{p12:sec:continuous-qft-reconstruction-theorem}
\begin{theorem}[Phase-selected bounded-local Born--Hopf/QFT compatibility]
\label{p12:thm:main-qft}
Let \(\mathfrak D_{\rm QFT}\) be a finite Born--Hopf spectral/momentum datum satisfying QFT1--QFT3, and let \(\omega_\alpha\) be any weak-* bounded-local cluster state of a supplied/selected physical phase.  Assume finite reflection positivity, the local-flat representative/covariance comparison required on the compact chart exhaustion, the bounded regularity hypotheses of QFT7, and the actual bounded algebraic reconstruction certificate of Definition~\ref{p12:def:v2210-lit-bounded-os-certificate} or an independently verified equivalent corrected OS certificate. Exact ordinary local-net claims include representation-wide locality and full Euclidean comparison in this certificate. Then \(\omega_\alpha\) is OS-admissible on the certified bounded local observable algebra and reconstructs a local-chart bounded OS/Haag--Kastler QFT export
\[
 (\mathcal H_\alpha,\Omega_\alpha,H_\alpha,\mathcal A_\alpha(\mathcal O),U_\alpha).
\]
For every controlled bounded-local observable family whose actual identification with this reconstructed net is certified, its reconstructed moments, connected cumulants, bounded source derivatives and effect/POVM probabilities are the corresponding refinement limits along that phase. Wightman field presentations additionally consume their declared field/reconstruction certificate.  QFT4$_C$ is needed only for the stronger statement that all admissible same-phase refinements yield one canonical export.  QFT8, positive-gap QFT9, QFT10 fields and scattering are claim-scoped extensions.
\end{theorem}
\begin{proof}
Weak-* compactness supplies \(\omega_\alpha\).  Finite reflection-positive quadratic forms are closed under scalar weak-* limits, so the phase state is reflection positive.  Theorem~\ref{p12:thm:native-bh-covariance} gives exact finite state covariance and reduces flat-state covariance to the representative/chart comparison.  Uniform bounded test-function estimates pass to the cluster state and give bounded regularity.  The actual algebraic or equivalent corrected OS certificate then reconstructs the Hilbert space, cyclic vacuum, nonnegative Hamiltonian and the certified local observable net, as made explicit in Corollary~\ref{p12:cor:v2210-lit-bounded-os-reconstruction}. Equality with a larger requested export additionally consumes its actual observable generation/identification.  Moment, cumulant, bounded-source and effect limits follow by applying weak-* convergence to the corresponding bounded local products and by finite cumulant/Cauchy differentiation identities.  None of these steps requires global refinement-path uniqueness, an interaction witness, a positive mass gap, an unbounded point field or an asymptotic scattering construction.
\end{proof}

\begin{corollary}[Canonical same-phase export]
\label{p12:cor:v283-canonical-same-phase-export}
If the selected-phase weak-* cluster set is a singleton, equivalently QFT4$_C$ holds, then the phase-selected QFT export of Theorem~\ref{p12:thm:main-qft} is independent of admissible same-phase refinement path.
\end{corollary}


% Source Paper VII-C, original line 4758; inherited proof bodies unchanged.
\subsubsection{Series-level consequence}
\label{p12:sec:series-level-qft-consequence}

\begin{corollary}[Downstream continuum use]
\label{p12:cor:benefit-for-series}
Once the GP1--GP2 generation conditions of Theorem~\ref{p12:thm:main-qft} or
the sector-import conditions of
Theorem~\ref{p12:thm:sector-import-qft-compatibility} are verified for a sector,
every downstream use of continuum fields, propagators, local operators, Ward
identities, or renormalized couplings may be reclassified from interface
language to derived Born--Hopf continuum output, subject only to the declared
sector scope.  Scattering-level objects remain subject to
Proposition~\ref{p12:prop:scattering-nonclaim} and
Theorem~\ref{p12:thm:scattering-hypothesis-firewall}.
\end{corollary}
\begin{proof}
The continuum objects are reconstructed from the finite regulator tower.  Therefore downstream appearances of those objects cite the reconstruction theorem as framework output.  Classical Noether closure and finite quantum Ward closure remain separately typed: only an invariant finite state/measure identity with all anomaly/contact terms controlled is eligible for continuum transport.
\end{proof}



% Source Paper VII-C, original line 5975; inherited proof bodies unchanged.
\subsection{Continuum certificate status}
\label{p12:sec:v291-continuum-status}
The controlling scope is:
\begin{itemize}
\item \(\mathrm{QFT4}_E\): base phase-selected bounded-local continuum state
existence is available on the stated fixed-UV phase branch;
\item \(\mathrm{QFT4}_C/G39\): required only for a single canonical,
refinement-independent state within the same physical phase;
\item \(\mathrm{G33\text{-}COMP}\): required only for stronger representative,
action-family, or regulator universality claims;
\item C198-04--06 belong to that stronger \(\mathrm{G33\text{-}COMP}\)
comparison route, while C198-07--08 are consumed only by stronger canonicity or
selected-local DLOC claims that explicitly require them;
\item RGAS199-14--16 are a claim-scoped full quantum asymptotic-safety
extension, not prerequisites for the base phase-selected local-QFT export.
\end{itemize}
The remaining ordered chiral/topological comparison, independently normalized
external/unbounded composites, and massless/long-range infrared scattering
certificates remain claim-scoped where asserted.  This typing removes no
hypothesis from a stronger claim that actually consumes it.


% Source Paper VII-C, original line 5996; inherited proof bodies unchanged.
\subsection{Earlier theorem synthesis and scope}
\label{p12:sec:conclusion}
Paper~VII-C now separates the already constituted finite quantum theory from its continuum-QFT export.  The native continuum is fixed-UV and momentum-first: large \(m_x\) densifies the momentum spectrum inside the fixed Born momentum ball while reciprocal-position extent unfolds.  QFT4 is split into QFT4$_E$, existence of a bounded local cluster phase, and QFT4$_C$, the stronger singleton same-phase canonicity theorem.  QFT4$_E$ is obtained by weak-* compactness; finite reflection positivity and bounded regularity pass to every cluster state; exact finite symmetry together with representative/chart convergence supplies state covariance without a separate interacting-vacuum defect.  Consequently every certified cluster phase admits a bounded local OS/Wightman export, and its controlled local moments, cumulants, bounded source responses and effect probabilities are limits of the same finite Born--Hopf data.  QFT8 classifies observable non-Gaussianity or, with its stronger actual comparator, dynamical interaction rather than QFT existence; QFT10 is a field-extension layer, and local positive-measure 1PI vertices are dependent consequences on fixed noncritical source charts.

For the strong sector, exact coarse maps and same-owner finite block squares remain valid, while genuine interacting blocking may generate multi-face terms; this affects canonical action/refinement comparison but does not turn the Born continuum into ultraviolet cutoff removal.  The historical shrinking-position-lattice and naive fixed-geometry nested-cutoff no-go arguments are retained only as comparison diagnostics.  The remaining continuum issues are stronger analytic certificates: canonical same-phase uniqueness if one demands one refinement-independent export, claim-scoped external composite fields, ordered chiral/topological comparison where consumed, and massless/long-range infrared scattering.  These are not new native physical owners and do not determine whether the finite framework is already a quantum theory.

The same dependency discipline now applies to scattering.  Local reconstruction,
finite-window probabilities, light-front operator measures, and isolated
spectral pole/residue limits carry no asymptotic or infrared premise.  An
\(S\)-matrix amplitude, LSZ formula, \(S\)-matrix pole, completeness statement,
or cross section imports only the corresponding clauses of
\(\mathsf{SMAT}_{\mathcal C}\); massless and long-range channels additionally
name and prove an inclusive, dressed, or algebraic-detector route.  Thus the
framework no longer hides a global scattering assumption inside local QFT or
spectral convergence.



% BEGIN INLINE SOURCE: Paper7C-SourceAndObservableTransport-v1_2.tex
% Framework v2.123: closes the finite interacting G16 scalar-vertex constructor and separates native exact source blocking from named-scheme evaluation.
% Sources: RemainingGapTransportAnalysis v1.0 CLS-T06--T07;
% GapReassessment v1.0, observable-level continuum criterion.
% Framework v2.118: native cyclic QROOT is closed; generic/noncyclic state occurrence remains separate.

% Source Paper VII-C, original line 8882; inherited proof bodies unchanged.
\subsection{Geometry-first consequences and physical scope}
\label{p12:sec:v2201-geometry-first-consequences}
\paragraph{Geometry ancestry.}
Rooted matter/color representation $\to$ native gauge indices and common-heat normalization $\to$ licensed local loop and physical threshold transport.
\paragraph{Status.}
These are consequences on the domains stated below.
A mathematical certificate or an executed restriction retains every
physical source, state, clock and readout hypothesis of its inputs.
The native representation sums and one-loop coefficient check
are distinguished from physical coupling matching. A known conditional
coefficient is not a completed low-energy prediction.


% Source Paper VII-C, original line 8894; inherited proof bodies unchanged.
\subsubsection{D20: Executed gauge index transport and its precise loop domain}
\label{p12:sec:v2201-derived-d20}
\noindent\textit{Status: inherited coefficient evaluation.}

The inherited one-family hypercharge square sum is $10/3$; three families
and one complex weak doublet give, in the licensed local one-loop domain,
\[
 b_Y=\frac{41}6,\qquad b_2=-\frac{19}6,\qquad b_3=-7,
 \quad b_1=\frac35b_Y=\frac{41}{10}.
\]
\begin{proof}
Sum $(2/3)T(R)$ for each Weyl fermion and $(1/3)T(R)$ for
each complex scalar, with the vector term $-(11/3)C_A$. Color has
one-family index two and weak has index two; the scalar weak index is
$1/2$. The hypercharge scalar square sum is $1/2$. These reproduce the
conditional coefficients already stated in V-C.
\end{proof}
The common-heat boundary $g_3^2=g_2^2=2\pi^2$ and
$g_Y^2=6\pi^2/5$ then has the explicit same-convention transport
\[
 g_i^{-2}(\mu)=g_i^{-2}(\mu_h)
 -\frac{b_i}{8\pi^2}\log(\mu/\mu_h)
\]
only on its actual perturbative interval and threshold scheme. This is an
executed inherited coefficient check, not a newly selected matching scale.
The representation traces do not create the local loop integral or an
unconditional low-scale coupling.
\par\noindent\textit{Inherited inputs:} \ref{p5:cor:common-heat-absolute-gauge-normalization}, \ref{p7v2:thm:conditional-one-loop-running}.

% END GEOMETRY-FIRST CONSEQUENCES v2.201

% BEGIN LITERATURE-CHECKED QFT EXPORT CERTIFICATES v2.210

\paragraph{Final theorem status and export boundary.}
The fixed-UV bounded-local QFT export is phase-selected and conditional. Canonical uniqueness, uniform spectral gap, full field distributions, scattering and infrared completion remain separate claims.

\paragraph{Technical-reference convention.} References in inherited theorem text to labels formerly described as appendices and now assigned to Paper~VII-D are references to that separate technical companion, not local appendices.



% BEGIN SOURCE-INTEGRATED CROSS-PAPER MATHEMATICS v3_017, 15 November 2026.
%  Derived geometric comparison modules; physical source occurrence not promoted.
%  E3/E4 independent audit of these NEW results remains pending; inherited source statements retained.
\section{Additional exact mathematical consequences from the verified source geometry}

% -------- Research module Paper7C-ThermalZeroPinchObstruction-v1_1 (latest supplied version) --------
% Proposed research insertion for Paper VII-C; conditional analytic sufficiency.
\subsection{A spectral-diameter obstruction to thermal zero pinching}
\label{p12:sec:cc13-zero-pinch}

\begin{corollary}[Uniform finite energy width excludes thermal zero pinching]
\label{p12:cor:cc13-zero-pinch}
Let $\{H_L\}$ be a sequence of finite self-adjoint Hamiltonians,
with dimensions $d_L$ and spectra in intervals
$[c_L-D/2,c_L+D/2]$ for one finite $D>0$ independent of $L$.
Define the shifted, normalized partition sums
\[
 \widehat Z_L(\beta)=d_L^{-1}e^{\beta c_L}
 \operatorname{Tr}e^{-\beta H_L}.
\]
Then:
\begin{enumerate}
\item every $\widehat Z_L$ is entire, equals one at $\beta=0$,
 and has no zeros for $|\operatorname{Im}\beta|<\pi/D$;
\item every subsequence admits a further subsequence converging
 locally uniformly in $\beta\in\mathbb C$ to an entire exponential
 transform of a probability measure supported on $[-D/2,D/2]$;
\item every such limit is zero-free in the same strip, with its
 logarithm analytic there, normalized by $\log \widehat Z(0)=0$.
\end{enumerate}
Consequently, a sequence of thermal partition zeros approaching a finite
real inverse temperature through this normalized spectral-shape family
\emph{requires} unbounded energy spectral diameter along that sequence.
This necessary condition is not sufficient for a phase transition.
\end{corollary}
\begin{proof}
Write $\varepsilon_{j,L}=E_{j,L}-c_L$, so
$|\varepsilon_{j,L}|\le D/2$ and
$\widehat Z_L(\beta)=d_L^{-1}\sum_j e^{-\beta\varepsilon_{j,L}}$.
The common strip follows from the energy-only case of
Theorem~\ref{p11:thm:cc13-zero-free-polytope}.  Each empirical
probability measure
$\nu_L=d_L^{-1}\sum_j\delta_{\varepsilon_{j,L}}$ is supported on
the common compact interval.  Sequential weak compactness of
probability measures gives a weakly convergent subsubsequence
$\nu_{L_k}\rightharpoonup\nu$.  The exponential integrand and all its
$\beta$ derivatives are uniformly bounded on compact subsets of
$\mathbb C\times[-D/2,D/2]$; hence the exponential transforms converge
locally uniformly (including fixed derivatives) to
$\widehat Z(\beta)=\int e^{-\beta\varepsilon}\,d\nu(\varepsilon)$.
The same semicircle argument proves the limit's strict zero-free
bound.  Since the strip is simply connected, the logarithm has a
holomorphic branch fixed at $\beta=0$.  Any sequence of zeros pinching
a real inverse temperature contradicts the common zero-free strip;
therefore a diverging diameter is necessary.
\end{proof}

\begin{remark}[Refinement status]
A fixed finite-dimensional QCD compiler supplies finite $D_L$ but does
\emph{not} supply a uniform bound on $D_L$ as the regulator, volume,
or particle content grows.  The $c_L$ energy shift and dimension
normalization are mathematical comparison choices.  No physical
thermodynamic limit, zero-pinch mechanism, or observed phase transition
is established by this sufficient-condition theorem.
\end{remark}
% END SOURCE-INTEGRATED CROSS-PAPER MATHEMATICS v3_017.

\printbibliography
% BEGIN FRAMEWORK v2.241: ROOTED CONORMAL GRAVITY AUDIT
\section{Whole spectral measures in the boundary cap application}
\label{p12:sec:v229-spectral-measure}
The positive cap representation and its infinite moment rank are proved
in Proposition~\ref{p1b:prop:v229-cap-moment-rank}. They do not replace
the native heat regulator or select an independently weighted mixture
of heat times. Theorem~\ref{p1b:thm:v229-heat-tail} excludes the stated
fixed-cutoff positive compact homothetic heat word, with its full
spectral measure, from exact equality with that cap.
The product determinant calculations of
Propositions~\ref{p3f:prop:v229-determinant-scope} and
\ref{p3f:prop:v229-rooted-product-zeta} retain a fixed spectral unit
and their actual boundary domains. Conditional finite-mode cancellations
do not specify the whole field-theoretic measure. Any physical
determinant cancellation must be derived with the actual field content,
statistics, quotient and boundary integration.

% END FRAMEWORK v2.241: ROOTED CONORMAL GRAVITY AUDIT

% BEGIN FRAMEWORK v2.350: TWO-STAGE FIXED-UV CONTINUUM AND LOCALITY REPAIR
\clearpage
\section{Two-stage fixed-UV reconstruction and the detector-locality certificate}
\label{p12:sec:v350-reconstruction-repair}
The physically nonstandard Born--Hopf limit keeps the ultraviolet
momentum ceiling fixed while increasing position-side volume and
resolution. What requires explanation is the measurable correlation
and detector locality structure, not a priori reproduction of an
unregulated canonical point field. We distinguish a bounded phase
state from an ordinary local Wightman field presentation.

\begin{theorem}[Common-algebra cluster state and bounded identity passage]
\label{p12:thm:v350-common-algebra-cluster}
Let $\mathfrak A$ be one unital $C^*$-algebra of bounded cylinder
observables and let $\pi_L:\mathfrak A\to\mathcal A_L$ be
compatible unital $*$-representations. For normalized finite
states $\omega_L$, let $\widehat\omega_L=\omega_L\circ\pi_L$.
Then some subnet of $\widehat\omega_L$ converges weak-$*$ to a
state $\omega_\alpha$ on $\mathfrak A$.
If one common positive-time test $*$-algebra and reflection
$\Theta$ are supplied and every $\widehat\omega_L$ satisfies
reflection positivity on it, $\omega_\alpha$ is reflection
positive. Every exact bounded-word correlation identity on
this fixed algebra survives along that subnet, provided
all compared words are represented compatibly.
\end{theorem}
\begin{proof}
The state space of a unital $C^*$-algebra is weak-$*$ compact
by Banach--Alaoglu and weak-$*$ closedness of positivity and
normalization. Thus it contains a convergent subnet of the
state net. For every fixed positive-time test $A$, the
reflection-positive scalar
$\widehat\omega_L((\Theta A)^*A)\ge0$ passes to the limit.
The same pointwise convergence sends each fixed bounded
word expectation that vanishes at every finite level to
zero at the cluster state. None of these arguments produces
unbounded point fields or spatial microcausality by itself.
\end{proof}

\begin{theorem}[Ordinary OS/Wightman export is a separate conditional bridge]
\label{p12:thm:v350-os-conditional-bridge}
Suppose the cluster family of
Theorem~\ref{p12:thm:v350-common-algebra-cluster}
additionally admits a common system of Schwinger distributions
with sufficient uniform regularity/growth, full local-flat
Euclidean covariance including time--space mixing, graded
symmetry, reflection positivity, appropriate clustering or
vacuum sector control, and a local-test-function/contact-term
limit satisfying the relevant Osterwalder--Schrader
reconstruction hypotheses. Then the corresponding local
OS/Wightman Hilbert-space and field presentation exists.
Neither mere norm-boundedness, spatial $\Spin(4)$ covariance,
nor weak-$*$ cluster existence separately proves this conclusion.
\end{theorem}
\begin{proof}
The conclusion is precisely the OS reconstruction implication
applied after all its hypotheses have been established on the
same family of limiting distributions. Theorem~\ref{p12:thm:v350-common-algebra-cluster}
provides only a bounded state; it does not supply the required
Schwinger hierarchy, domains, Euclidean mixing symmetry or
locality support. Those hypotheses are explicitly independent
inputs to this implication, not consequences of compactness.
\end{proof}

\begin{proposition}[Canonical fixed-band point-commutator comparison]
\label{p12:prop:v350-fixed-band-nonlocality-bound}
For the canonical spatially bandlimited free comparison field
with Fourier support $|k|\le\Lambda<\infty$ in three flat
spatial dimensions, the equal-time commutator kernel is
\begin{equation}
\label{p12:eq:v350-fixed-band-kernel}
 \delta_\Lambda(r)=
 \frac{\sin(\Lambda r)-\Lambda r\cos(\Lambda r)}{2\pi^2r^3}.
\end{equation}
It is nonzero at some separated points and is not a local
Dirac delta distribution. For $L^1$ detector profiles $f,g$
whose supports are separated by distance at least $d>0$,
\begin{equation}
\label{p12:eq:v350-detector-bound}
 \left|\iint \overline{f(x)}\,\delta_\Lambda(|x-y|)g(y)
                     \,d^3x\,d^3y\right|
 \le\frac{1+\Lambda d}{2\pi^2d^3}\|f\|_1\|g\|_1.
\end{equation}
These statements concern this canonical free cutoff comparison;
they do not prove nonlocality or locality of a different
native interacting Born--Hopf observable net.
\end{proposition}
\begin{proof}
Radial integration gives
$(2\pi)^{-3}\int_{|k|\le\Lambda}e^{ik\cdot x}\,d^3k$
equal to Eq.~\eqref{p12:eq:v350-fixed-band-kernel}.
For $r\ge d$, use
$|\sin(\Lambda r)-\Lambda r\cos(\Lambda r)|\le1+\Lambda r$
and note that $(1+\Lambda r)/r^3$ decreases for positive
$r$. Integrating against $|f(x)g(y)|$ gives the bound.
The Fourier cutoff is fixed, so increasing the reciprocal
position volume alone does not make the kernel a delta.
\end{proof}

\begin{proposition}[Exact continuum Ward identities require ordered moments]
\label{p12:prop:v350-ward-identity-transfer}
Let a finite Ward, gauge, or anomaly-cancellation identity
be an equality of fully source-ordered correlation
\emph{distributions}, with all contact, Jacobian, and
renormalization terms retained. If every term converges
jointly on a common test-function domain, the limiting
correlation identity holds. A limiting operator identity
follows only if all its matrix elements vanish on a
specified dense invariant domain. Cyclicity of a GNS vacuum
alone is not a separating-vector proof for arbitrary operators.
\end{proposition}
\begin{proof}
The finite correlation identity has zero pairing with every
allowed test function. Joint distributional convergence carries
that zero pairing to the limiting correlation functional.
If the limiting difference is an operator-valued distribution,
its vanishing on a dense invariant domain is verified by all
matrix elements, not only vacuum expectations. Such a domain
condition permits the claimed operator equality there.
\end{proof}

\begin{remark}[Native physical target]
\label{p12:rem:v350-operational-locality-frontier}
The remaining native certificate is a concrete family of
same-action physical detector observables $D_L(f)$ and
spacelike-separated profiles whose commutators or response
kernels obey an experimentally adequate leakage bound after
refinement. If ordinary Wightman locality is claimed, an
exact vanishing locality relation is additionally required;
otherwise a quantitatively controlled operationally local
field theory is the appropriate, explicitly distinct output.
The pure-glue continuum gap still requires its independent
uniform spectral lower bound and transfer convergence.
\end{remark}
% END FRAMEWORK v2.350

% BEGIN FRAMEWORK v2.356: gapless current-weighted temporal criteria
\section{Current-weighted Euclidean-time and gapless infrared criteria}
\label{p12:sec:v356-gapless-current-control}

\begin{proposition}[Exact current-correlator integral]
\label{p12:prop:v356-current-time-integral}
On the same positive temporal source chart and zero-ground-atom
domain as Theorem~\ref{p12:thm:v299-current-weighted-ir}, put
$C_{J,L}(t)=\int_{(0,\infty)}e^{-t\Delta}
 d\mu_{J,L}(\Delta)$ for $t\ge0$.  Whenever the inverse-cubic
moment is finite,
\[
 \boxed{Z_{{\rm spec},L}(0)=2\int\Delta^{-3}d\mu_{J,L}
 =\int_0^\infty t^2 C_{J,L}(t)\,dt.}
\]
The complete response still adds the independent same-action
$Z_{{\rm loc},L}(0)$.
\end{proposition}
\begin{proof}
Tonelli's theorem applies to the nonnegative integrand, and
$\int_0^\infty t^2e^{-\Delta t}\,dt=2\Delta^{-3}$.
The source contact is not part of this spectral integral.
\end{proof}

\begin{theorem}[Uniform gapless current-moment sufficient criterion]
\label{p12:thm:v356-uniform-current-moment}
For a family of same-source positive current measures on
$(0,\infty)$ suppose, for some $\varepsilon>0$, $M,W<\infty$,
\[
 \sup_L\int\Delta^{-(3+\varepsilon)}d\mu_{J,L}\le M,
 \qquad \sup_L\mu_{J,L}((0,\infty))\le W.
\]
Then for $\delta,R,T>0$,
\begin{align*}
 \sup_L\int_{(0,\delta)}\Delta^{-3}d\mu_{J,L}
      &\le M\delta^\varepsilon,\\
 \sup_L\int_{(R,\infty)}\Delta^{-1}d\mu_{J,L}
      &\le W/R,\\
 \sup_L\int_T^\infty t^2C_{J,L}(t)\,dt
      &\le\Gamma(3+\varepsilon)M T^{-\varepsilon}.
\end{align*}
If the measures converge against continuous tests on compact positive
energy intervals, and the local/contact temporal coefficients have
their separate source-uniform limit, the spectral zero-frequency
endpoints converge and refinement commutes with the temporal
zero-frequency limit under these hypotheses.  No uniform lower
bound on the \emph{entire} excitation spectrum is required.
\end{theorem}
\begin{proof}
On $\Delta<\delta$, write
$\Delta^{-3}=\Delta^\varepsilon
\Delta^{-(3+\varepsilon)}$.  On $\Delta>R$,
$\Delta^{-1}\le R^{-1}$.  For $t\ge T$,
$t^2\le T^{-\varepsilon}t^{2+\varepsilon}$; Tonelli and
$\int_0^\infty t^{2+\varepsilon}e^{-t\Delta}\,dt
=\Gamma(3+\varepsilon)\Delta^{-(3+\varepsilon)}$
give the time bound.  Split each inverse-moment integral into
$(0,\delta)$, $[\delta,R]$, and $(R,\infty)$.
The two uniform tails and compact-interval convergence give the
required uniform integrability and passage of limits.
For the finite-frequency difference quotient, its integrand
$2/[\Delta(\Delta^2+\omega^2)]$ is bounded by
$2\Delta^{-3}$, so the same uniform-integrability argument
controls the $\omega\to0$ limit.  The local/contact term needs
its independently assumed convergence.
\end{proof}

\begin{remark}[Physical QCD input remains separate]
\label{p12:rem:v356-current-physical-endpoint}
The theorem permits a vanishing current-visible gap when residues
vanish sufficiently rapidly.  It does not establish the bounds for
the actually selected interacting QCD member, instantiate its
Hamiltonian and electromagnetic current, or determine the nonlinear
frequency-dependent local source term.  The static energy-source
Hessian contact cannot automatically replace the temporal
$Z_{\rm loc}(0)$.
\end{remark}
% END FRAMEWORK v2.356


% BEGIN FRAMEWORK v3.023: ACTION-PRESERVING MOBILE 2I GLUING GUARD
\section{Conditional same-owner gluing for mobile $2I$ hadron sources}
\label{p12:sec:v3023-mobile-2i-apbi-gluing}
The five disjoint $2I/2T$ left cosets, the five conjugate $2T$
subgroups and the twenty-five overlapping two-sided transporter
registers are distinct geometric objects. None is a primitive
position lattice, an additional quark species or an extra QCD
regulator. The native Hurwitz stencil uses intra-$2T$ paths; the
600-cell nearest edges join different disjoint cosets. Thus a
mobile $2I$ same-action realisation is not obtained by simply
replicating five fixed QCD stencils.

\begin{proposition}[Conditional mobile-source presentation transport]
\label{p12:prop:v3023-mobile-source-same-owner-gluing}
Suppose the mobile geometric charts admit gauge- and spin-covariant
transition maps $T_{\alpha\beta}$ on their actual overlaps with
$T_{\alpha\beta}T_{\beta\gamma}=T_{\alpha\gamma}$,
and suppose these transitions push forward the \emph{complete}
existing QCD action, Gauss quotient, physical measure and dressed
quark sources without additional counterterms or couplings.
If the resulting common Hilbert-space transport $\mathcal U$ is
unitary and $H_L'=\mathcal U H_L\mathcal U^{-1}$,
$O'=\mathcal U O\mathcal U^{-1}$,
$|\Omega'\rangle=\mathcal U|\Omega\rangle$, then
\[
\boxed{\langle\Omega'|O_1'(z-H_L')^{-1}O_2'|\Omega'\rangle
=\langle\Omega|O_1(z-H_L)^{-1}O_2|\Omega\rangle.}
\]
The same conclusion applies to all retained source derivatives
under the differentiation hypotheses of the existing exact
blocking-source-naturality theorem. Hence an actually certified
same-owner presentation cannot change a hadron correlator merely
by relabeling its geometric $2I$ source chart.
\end{proposition}




\begin{proof}
Unitary conjugation takes the resolvent of $H_L$ to that of
$H_L'$; substitute the source and vacuum transports and cancel
$\mathcal U^\dagger\mathcal U$. The cocycle condition makes
triple-overlap relabeling consistent. The measure/action
pushforward is a separate stated hypothesis for the source
integrals. Exact blocking retains all induced interactions and
source contacts by the previously stated theorem, so no
additional Gaussian or isolated-face factorization is inferred.
\end{proof}

% BEGIN FRAMEWORK v3.025: QCD ACTION/OBSERVABLE OWNER TYPING
\begin{remark}[Three distinct nonlinear action-selection claims]
\label{p12:rem:v3025-qcd-action-selection-ports}
Separate the following logical ports.  \emph{HEAT-REFERENCE} means
that a selected finite compact-group heat member supplies a computable
nonlinear QCD reference action and its inherited quadratic stiffness;
it does not assert unique physical selection among every admissible
magnetic member.  \emph{TEMPORAL-POSITIVE-TRANSFER} concerns a face law
actually used as a central temporal convolution operator; there the
Peter--Weyl coefficients must be nonnegative, and
Proposition~\ref{p10:prop:v3025-sextet-temporal-exclusion}
excludes one named deformation on its stated heat-time interval.
The same test is not imposed on every spatial magnetic potential.
\emph{PHYSICAL-SOURCE-OBSERVABLE} requires a selected same-owner
$H_L$, its complete Gauss/fermion/measure prescription and source map,
or an independently established equality of the complete source-visible
spectral measures on the specified observable family.  Neither equal
quadratic Hessians nor equality of the reference source Gram
$M_0=J^\dagger J$ discharges the physical-source condition.
The exact retained-source blocking theorem
\ref{p12:thm:v2141-blocking-source-naturality} and conditional mobile
presentation in
Proposition~\ref{p12:prop:v3023-mobile-source-same-owner-gluing}
apply only after their same-action and differentiability hypotheses
are satisfied.  No QCD-confinement, hadron-pole, beta-trajectory, or
unique magnetic-law promotion follows from these reference results.
\end{remark}
% END FRAMEWORK v3.025


\begin{remark}[Unresolved global transport and curvature residual]
\label{p12:rem:v3023-mobile-2i-holonomy-frontier}
Geometric matching of each Hurwitz triangle with a 600-cell
triangle does not prove equality of their gauge holonomies.
When both path families belong to a common gauge owner, their
relative loop holonomies $R_i=A_i^{600}(A_i^H)^{-1}$ can be
nontrivial; the source Wilson-determinant discrepancy is bounded
by $\sum_i\|R_i-I\|_{\rm op}$. The exact mobile $2I$ incidence
results do not supply the full chart transitions, prove a bounded-
incidence extensive APBI measure, ensure cross-coset quark kinetic
links, or establish isolated hadronic poles. The physical one-loop
beta response and scale-to-observable map retain their earlier
conditional domain; composite source multiplicities are never
added to the elementary one-loop matter count.
\end{remark}
% END FRAMEWORK v3.023

\end{document}
